\documentclass[11pt,a4paper]{article}
\usepackage[T1]{fontenc}
\usepackage[utf8]{inputenc}
\usepackage{lmodern,microtype}
\usepackage[margin=26mm]{geometry}
\usepackage{amsmath,amssymb,amsthm,mathtools}
\usepackage{booktabs,longtable,array,enumitem}
\usepackage{xcolor}
\usepackage[colorlinks=true,linkcolor=blue!45!black,citecolor=blue!45!black,urlcolor=blue!45!black]{hyperref}
\usepackage{xurl}
\usepackage{fancyhdr}
\pagestyle{fancy}\fancyhf{}\fancyhead[L]{DP\_MULT: formal derivatives and divided powers}\fancyhead[R]{21 November 2026}\fancyfoot[C]{\thepage}
\setlength{\headheight}{14pt}
\setlength{\emergencystretch}{2em}
\newtheorem{theorem}{Theorem}[section]
\newtheorem{lemma}[theorem]{Lemma}
\newtheorem{proposition}[theorem]{Proposition}
\newtheorem{corollary}[theorem]{Corollary}
\theoremstyle{definition}\newtheorem{remark}[theorem]{Remark}
\newcommand{\Z}{\mathbb Z}\newcommand{\Q}{\mathbb Q}\newcommand{\F}{\mathbb F}
\newcommand{\Ga}{\mathbb G_a}\newcommand{\Gm}{\mathbb G_m}
\newcommand{\Spf}{\operatorname{Spf}}\newcommand{\Spec}{\operatorname{Spec}}
\newcommand{\Lie}{\operatorname{Lie}}\newcommand{\Dist}{\operatorname{Dist}}
\newcommand{\ord}{\operatorname{ord}}\newcommand{\gr}{\operatorname{gr}}
\newcommand{\mathscr}[1]{\mathcal{#1}}
\newcommand{\vp}{v_p}\newcommand{\dR}{\mathrm{dR}}\newcommand{\crys}{\mathrm{crys}}
\newcommand{\id}{\operatorname{id}}\newcommand{\Hom}{\operatorname{Hom}}
\newcommand{\RMap}{\operatorname{RMap}}\newcommand{\Int}{\operatorname{Int}}
\newcommand{\ceil}[1]{\left\lceil #1\right\rceil}\newcommand{\floor}[1]{\left\lfloor #1\right\rfloor}
\title{Integral invariant derivatives at every formal height\\
\large Optimal affine error, divided-power distributions,\\Cartier duality, and binomial homotopy types}
\author{DP\_MULT research run\quad\texttt{DP\_MULT\_20261121}}
\date{21 November 2026\quad UTC}
\begin{document}
\maketitle
\begin{abstract}
For a one-dimensional commutative formal group over $W(k)$ of finite height $h$, with normalized integral invariant derivation $D$, we prove coefficientwise
\[
 \vp(D^n f)\ge
 \max\left\{\sum_{\substack{j\ge1\\h\nmid j}}\floor{n/p^j},\;
 \ceil{n\left(\frac1{p-1}-\frac1{p^h-1}\right)-(h-1)},\;0\right\}
 \qquad(f\in W(k)[[t]]).
\]
The affine slope and the uniform constant $h-1$ are optimal. Thus the proposed logarithmic error is unnecessary. The proof does not assume a Lubin--Tate normal form. At height two this gives the coefficient in Alp\"oge's step 2, with uniform affine error $1$. The additional integral theta and period comparisons are identified in recovered author-source modules. Full unramified CM action extends the bounds over ramified bases; residue height alone does not.

We prove the directional higher-height statement and construct an integral model of the actual Honda Cartier dual. Its $h-1$ primitive directions have sharp slopes $p^i/(p^h-1)$, each with optimal uniform affine error $1$; their slopes sum to $c_h$. The generalized scalar kernel has the full $c_h$ slope even though its direct intermediate weight estimate is false. For every integral cyclic CM display the optimal directional affine error is its codimension; the equally weighted sum across a type and its actual dual is maximized exactly by alternating type. For the invariant primitive Poincar\'e kernel, the multivariate weight inflation $p$ is sharp. At odd primes and $h\ge3$, every logarithmic mixed line with primitive projections has exact factorial slope $1/((p-1)(p^k+1))$, where $k$ is its dual restricted-$p$ nilpotency length, even after any integral prefactor with nonzero restriction. Its best value is $1/(p^2-1)$, independent of height. These results separate bare derivative divisibility from the normalized theta kernel. The actual polarized genus-three Jacobian realizes the unit-arrow path theorem. A $p$-divisible tensor exponent has an exact different rate, but the source level denominator outweighs that universal ordinary-path improvement. The full integral polynomial module has sharp slope $1/((p-1)(p^{h-1}+1))$, already attained by a two-column harmonic family; the exact digit-sum error records the number of columns.

Finite sharp jets are globalized at controlled levels. The exact permitted torsion-translate lattice contains sharp individual witnesses whose canonical section norms are independent of level and which remain integral at every good prime. The endpoint period is constructed on an explicit unramified host; prescribed row components and finite selection windows are identified. For every fixed cyclotomic type, the full bare prime budget has the same leading coefficient, with integer rounding and saturation retained.

Grothendieck's boxed zero $p$-operation in cote 54 belongs to the vector kernel of a universal extension, not a supersingularity condition. Oda's theorem supplies the established de Rham/Dieudonn\'e comparison. Identity Frobenius lifts yield binomial structures under flatness; numerical polynomials have a divided-power associated graded, and Horel's theorem gives the precise integral homotopy statement.

\textbf{Status.} The local bounds and the stated algebraic corrections are proved and checked. The literal height-only substitution into the existing logarithm and alternating CM-type interfaces is refuted. A different construction producing a shorter global theta-lift chain remains open at an explicit construction-and-cost gate. That broader part remains under active investigation; no shorter proof is claimed. The final state will be set only when the stopping rule or completion ends the session.
\end{abstract}
\tableofcontents
\newpage

\section{Scope, conventions, and evidence}

The controlling sources are the attached prompt, Alp\"oge's \emph{Remark}, the two supplied Lean files and the supplied map, the Fable suggestions, and the archive transcriptions. Targeted Lean modules were subsequently recovered and compared byte for byte with the author's published archive, as documented below. We follow the instruction to cite Alp\"oge's RH and Dirichlet-GRH results as theorems; this report does not re-audit them. None of the local algebra proved below uses RH or GRH. No Masullo statement or draft is present in this attachment, and no claim conditional on an unspecified Masullo statement is made.

Throughout, $p$ is a prime, $k$ is a perfect field of characteristic $p$, $R=W(k)$, $K=\operatorname{Frac}(R)$, and $\vp(p)=1$. All formal groups in the one-dimensional results are commutative. If
\[
 F(X,Y)\in R[[X,Y]],\qquad v(t)=\frac{\partial F}{\partial Y}(t,0),
\]
then $D=v(t)\,d/dt$ denotes the invariant derivation normalized by $Dt(0)=1$. For a power series $g$, $\vp(g)$ means the infimum of its coefficient valuations; $\vp(0)=+\infty$. Thus our statements about $D^n f$ are stronger than their evaluations at the identity.

Height at least $h$ means
\begin{equation}\label{eq:height}
 [p]_F(t)\in pR[[t]]+t^{p^h}R[[t]].
\end{equation}
At exact finite height $h$, the coefficient of $t^{p^h}$ is a unit. Height infinity means $[p]_F\bmod p=0$. Put
\begin{equation}\label{eq:constants}
 q=p^h,\qquad c_h=\frac1{p-1}-\frac1{q-1},\qquad
 B_h=\vp(q!)-1=\frac{q-1}{p-1}-1=c_h(q-1).
\end{equation}
These are local statements about a formal completion of a scheme over $\Z_p$, and hence about integral lattices, not solely its rational isocrystal.

Stable finding identifiers are printed in brackets. A ``finding'' means a result established or a proposed identification decided in this run; no unqualified priority claim is intended. Existing theorems are cited, not presented as discoveries. In particular the Lubin--Tate case of the affine estimate already follows from Bannai--Kobayashi \cite[\S2, p.~527; Propositions 3.3 and 3.5, pp.~529, 531]{BK}. Our proof treats all formal groups satisfying \eqref{eq:height}.

All archive ``scans'' supplied here are \LaTeX\ transcriptions. Original manuscript pixels were not attached. During this run the official public PDFs for cotes 47 and 54 were recovered, and the three pages underlying the five emendation candidates were visually inspected. Those checks distinguish faithful readings of mathematical slips from actual transcription errors; their exact physical-PDF and manuscript page mappings are recorded in the audit. The supplied complete Lean archive scan is a report of another run; it is not a kernel certificate executed here.

\part{Checked findings and proofs}

\section{The uniform bound, with optimal constant error}

\begin{theorem}[DP1: uniform affine bound]\label{DP1}
For every formal group satisfying \eqref{eq:height}, every $n\ge0$ and every $f\in R[[t]]$,
\begin{equation}\label{eq:affine}
 D^n f\in p^{\max\{0,\ceil{c_hn-(h-1)}\}}R[[t]].
\end{equation}
At slope $c_h$, the constant $h-1$ is the smallest possible real constant in an inequality $\vp(D^nf(0))\ge c_hn-C$ valid for all $n$ and all integral $f$, even for any fixed normalized formal group of height at least $h$.
\end{theorem}

\begin{proof}
We give all steps, including the integral coordinate change. Associativity gives
\[
 v(F(t,u))=F_u(t,u)v(u).
\]
Consequently the normalized logarithm
\[
 L(t)=\int_0^t\frac{ds}{v(s)}=t+\sum_{m\ge2}\ell_m t^m\in K[[t]]
\]
satisfies $L(F(t,u))=L(t)+L(u)$: differentiating in $u$ gives $L'(u)$ on both sides, and evaluation at $u=0$ fixes the constant. In particular $L([p]_F(t))=pL(t)$.

Write $g(t)=[p]_F(t)=pt+\sum_{m\ge2}a_m t^m$. Comparison in degree $m\ge2$ gives
\begin{equation}\label{eq:logrec}
 (p-p^m)\ell_m=a_m+\sum_{i=2}^{m-1}\ell_i[t^m]g(t)^i.
\end{equation}
For $m<q$, all coefficients of $g$ occurring on the right belong to $pR$. Induction therefore gives $\ell_m\in R$: the right side is in $pR$ and $(p-p^m)/p$ is a unit. At $m=q$, the same argument gives $\vp(\ell_q)\ge-1$. Indeed the summands with $i\ge2$ use only coefficients of $g$ in degrees $<q$, so are divisible by $p^i$; $a_q$ is integral.

Use the strict integral coordinate
\[
 z=t+\sum_{2\le m<q}\ell_m t^m.
\]
Its inverse belongs to $R[[z]]$, because its linear coefficient is one. In this coordinate the normalized logarithm has the form
\begin{equation}\label{eq:trunclog}
 L_{\mathrm{new}}(z)=z+bz^q+O(z^{q+1}),\qquad \vp(b)\ge-1.
\end{equation}
Strictness preserves $Dz(0)=1$, and an integral coordinate isomorphism preserves every ideal $p^mR[[t]]$.

Define integral translation operators $\Delta_j$ by
\begin{equation}\label{eq:translation}
 f(F(z,u))=\sum_{j\ge0}(\Delta_j f)(z)u^j.
\end{equation}
Each $\Delta_j$ preserves $R[[z]]$. In logarithmic coordinates translation is ordinary addition, so over $K$
\begin{equation}\label{eq:distidentity}
 \sum_{j\ge0}\Delta_j u^j=\exp(L_{\mathrm{new}}(u)D).
\end{equation}
Every coefficient in $u$ is a finite sum of operators, so this formal identity is legitimate. It is the distribution identity of Katz \cite[Corollary 1.8, p.~669]{Katz81}, in the precise coordinate needed here. Equations \eqref{eq:trunclog}--\eqref{eq:distidentity} give
\[
 \Delta_j=D^j/j!\ (0\le j<q),\qquad
 \Delta_q=D^q/q!+bD.
\]
Set $w=q!/p^{B_h+1}\in R^\times$, $u_0=-wpb\in R$ and $T=w\Delta_q$. Thus $T$ is an integral $R$-linear operator and
\begin{equation}\label{eq:keyop}
 D^q=p^{B_h}u_0D+p^{B_h+1}T.
\end{equation}

For $0\le n<q$, the identity $D^n=n!\Delta_n$ gives $\vp(D^nf)\ge\vp(n!)$. Write $n=\sum_{j=0}^{h-1}d_jp^j$, $0\le d_j\le p-1$. Counting multiples of $p^r$ in $n!$ gives Legendre's formula, hence
\[
 \vp(n!)-c_hn
 =\sum_{j=0}^{h-1}d_j\left(\frac{p^j}{q-1}-\frac1{p-1}\right)\ge1-h.
\]
Every coefficient in parentheses is nonpositive. Its minimum over the allowed digits is attained when every digit is $p-1$, that is $n=q-1$, and then equals $1-h$.

For $n\ge q$, multiply \eqref{eq:keyop} on the right by $D^{n-q}$ and induct on $n$. The first term has coefficient valuation at least
\[
 B_h+c_h(n-q+1)-(h-1)=c_hn-(h-1),
\]
and the second at least
\[
 B_h+1+c_h(n-q)-(h-1)
 =c_hn-(h-1)+(1-c_h)\ge c_hn-(h-1),
\]
since $0\le c_h<1$. Only integrality of $T$ is used; commutation with $D$ is unnecessary. Integer-valued valuations and integrality give \eqref{eq:affine}.

Finally $D^n(t^n)(0)=n!$ in any normalized coordinate. For $n=q-1$ this has valuation
\[
 \vp((q-1)!)=\frac{q-1}{p-1}-h=c_h(q-1)-(h-1).
\]
Thus no smaller real constant $C$ can work.
\end{proof}

\begin{corollary}[DP1r: recursive integer refinement]\label{DP1r}
Set $b_h(n)=\vp(n!)$ for $0\le n<q$ and, for $n\ge q$, set
\[
 b_h(n)=B_h+\min\{b_h(n-q+1),\ 1+b_h(n-q)\}.
\]
Then $D^nR[[t]]\subset p^{b_h(n)}R[[t]]$.
\end{corollary}
\begin{proof}
The base case is $D^n=n!\Delta_n$ in the integral coordinate above. Equation \eqref{eq:keyop}, multiplied on the right by $D^{n-q}$, gives the displayed recurrence by induction. Coordinate invariance returns the assertion to the original ring.
\end{proof}
The recurrence is an explicit proved lower bound, not a claim to compute the exact minimum for every intermediate order and every formal group.

\section{The exponential, floor sums, and optimal slope}

\begin{lemma}[DP2: integral contraction and exact exponential denominator]\label{DP2}
Suppose the special fibre has exact finite height $h$. Adjoin $\pi$ with $\pi^{q-1}=p$, put $R'=R[\pi]$, and define
\[
 F_\pi(X,Y)=\pi^{-1}F(\pi X,\pi Y).
\]
The normalized logarithm $L_\pi$ and its inverse $E_\pi$ belong to $R'[[t]]$. If $E=L^{-1}=\sum e_nt^n$ is the original exponential, then
\begin{equation}\label{eq:expbound}
 \vp(e_n)\ge-\frac{n-1}{q-1},\qquad
 \vp(e_{q^r})=-\frac{q^r-1}{q-1}\quad(r\ge0).
\end{equation}
\end{lemma}
\begin{proof}
The Eisenstein extension $R'$ is a complete discrete valuation ring with $\vp(\pi)=1/(q-1)$; the case $q-1=1$ is the evident trivial extension. Its $p$-series is
\[
 g_\pi(t)=\pi^{-1}g(\pi t)
 =pt+\sum_{m\ge2}a_m\pi^{m-1}t^m\in pR'[[t]].
\]
For $m<q$ use $a_m\in pR$; for $m\ge q$ use $\pi^{m-1}\in pR'$. Coefficient induction \eqref{eq:logrec}, now with every coefficient of $g_\pi$ divisible by $p$, proves $L_\pi\in R'[[t]]$. The inverse of an integral series with linear coefficient one is integral, proving the claim about $E_\pi$.

Uniqueness of the normalized logarithm gives
\[
 L_\pi(t)=\pi^{-1}L(\pi t),\qquad E_\pi(t)=\pi^{-1}E(\pi t).
\]
This proves the first inequality in \eqref{eq:expbound}. For sharpness, put $H=g_\pi/p$. Reduction modulo $\pi$ gives
\[
 \overline H(t)=t+\bar a_q t^q,
\]
since the coefficients below degree $q$ have an extra factor $\pi$, and those above degree $q$ have $\pi^{m-q}$. In the logarithm recursion all summands involving $g_\pi^i/p$ with $i\ge2$ vanish modulo $\pi$, and $(1-p^{m-1})^{-1}\equiv1\pmod\pi$. Therefore
\[
 \overline L_\pi(t)=t+\bar a_qt^q.
\]
Its inverse is $\sum_{r\ge0}\beta_rt^{q^r}$, where $\beta_0=1$ and
\[
 \beta_r=-\bar a_q\beta_{r-1}^{q}.
\]
All $\beta_r$ are nonzero because $\bar a_q\ne0$. The coefficient of degree $q^r$ in $E_\pi$ is thus a unit, yielding the equality in \eqref{eq:expbound}.
\end{proof}

\begin{theorem}[DP3: floor-sum bound and sharpness]\label{DP3}
For finite height $h$ put
\begin{equation}\label{eq:Eh}
 E_h(n)=\sum_{\substack{j\ge1\\h\nmid j}}\floor{n/p^j}
 =\vp(n!)-\sum_{r\ge1}\floor{n/q^r}.
\end{equation}
Then $D^nR[[t]]\subset p^{E_h(n)}R[[t]]$. If $s_b(n)$ denotes the sum of the base-$b$ digits, the error is explicitly
\begin{equation}\label{eq:digiterror}
 E_h(n)=c_hn-\frac{s_p(n)}{p-1}+\frac{s_q(n)}{q-1}.
\end{equation}
The bound is attained at every $n=p^j$: writing $j=hr+s$ with $0\le s<h$, one may use $f=t^{p^s}$. In particular
\begin{equation}\label{eq:slopesharp}
 \vp(D^{q^r}t(0))=c_h(q^r-1).
\end{equation}
Thus $c_h$ is the optimal linear coefficient, even for each fixed height-$h$ formal group.
\end{theorem}
\begin{proof}
Since $D=d/dL$,
\begin{equation}\label{eq:eval}
 D^nf(0)=n![z^n]f(E(z)).
\end{equation}
For $n\ge1$ each term contributing to $[z^n]f(E(z))$ is a product of $m\ge1$ exponential coefficients with total degree $n$. Lemma~\ref{DP2} gives valuation at least $-(n-m)/(q-1)$, hence at least $-(n-1)/(q-1)$. Coefficients belong to $K$, whose valuation is integral, so
\begin{equation}\label{eq:prebound}
 \vp(D^nf(0))\ge A_h(n):=\max\{0,\vp(n!)-\floor{(n-1)/(q-1)}\}.
\end{equation}
This is already the proposed bound with no unspecified error.

It holds coefficientwise for $D^nf(t)$. To see this, apply the evaluation bound to each $t$-coefficient of $f(F(t,u))\in R[[t,u]]$, as a series in $u$. Invariance and associativity give
\[
 \left.D_u^nf(F(t,u))\right|_{u=0}=D_t^nf(t).
\]
All substitutions are valid in the $(t,u)$-adic topology, and each coefficient is a series over $R$.

For $n=p^j$, write $j=hr+s$, $0\le s<h$. Modulo $\pi$, the coefficient of degree $p^j$ in $E_\pi(z)^{p^s}$ is $\beta_r^{p^s}\ne0$. Since
\[
 E_\pi(z)^{p^s}=\pi^{-p^s}E(\pi z)^{p^s},
\]
we have
\[
 \vp([z^{p^j}]E(z)^{p^s})=-\frac{p^j-p^s}{q-1}
 =-\floor{\frac{p^j-1}{q-1}}.
\]
This proves equality in \eqref{eq:prebound} at each $p$-power order.

For general $n=\sum_{j\ge0}d_jp^j$, compose the coefficientwise bounds for $(D^{p^j})^{d_j}$. Legendre's formula and a finite interchange of sums give
\[
 \sum_jd_j\vp((p^j)!)=\vp(n!),\qquad
 \sum_jd_j\floor{\frac{p^j-1}{q-1}}
 =\sum_{r\ge1}\sum_{j\ge hr}d_jp^{j-hr}
 =\sum_{r\ge1}\floor{n/q^r}.
\]
Their difference proves \eqref{eq:Eh}. Summing the base-$q$ expansion gives $\sum_{r\ge1}\floor{n/q^r}=(n-s_q(n))/(q-1)$ and hence \eqref{eq:digiterror}. The previously constructed witnesses show sharpness at $p$-powers. Equation \eqref{eq:slopesharp}, divided by $q^r$ and taken to the limit, rules out every larger slope with bounded or sublinear error.
\end{proof}

Combining DP1 and DP3 gives the bound in the abstract. These are complementary estimates; neither a displayed lower bound nor the maximum is asserted to be pointwise optimal for all $n$. The exact optimality claims concern the affine slope, its uniform error, and the specified witnesses.

\begin{proposition}[DP4: endpoints and normalization]\label{DP4}
At height $1$, the universal floor-sum and affine bounds are zero. For $\widehat\Gm$ with $D=(1+t)d/dt$, the integral function $1+t$ has $D^n(1+t)(0)=1$ for every $n$.

At height infinity, $L,E\in R[[t]]$ and
\[
 D^nR[[t]]\subset p^{\vp(n!)}R[[t]].
\]
Equality at the identity is obtained with $f=L(t)^n$. Its exact error is
\[
 \vp(n!)=\frac{n}{p-1}-\frac{s_p(n)}{p-1}.
\]
There is no finite uniform constant error at slope $1/(p-1)$ for height infinity.

If $D'=uD$ for a constant $u\in K^\times$ and $f\in p^{-M}R[[t]]$, all the bounds for $D^n$ acquire the correction $n\vp(u)-M$. The zero derivation is handled separately and is trivial at positive order.
\end{proposition}
\begin{proof}
The height-one assertions follow directly from $c_1=E_1(n)=0$ and $D(1+t)=1+t$. At infinite height the coefficient recursion \eqref{eq:logrec} works over $R$ itself in every degree, so $L$ and $E$ are integral. Equation \eqref{eq:eval} gives the factorial bound; translation makes it coefficientwise as above. The function $L^n$ gives exactly $n!$ at order $n$. Legendre's formula gives the displayed error; at $n=p^r-1$ its magnitude is $r$, ruling out a finite uniform constant. Finally $(uD)^n=u^nD^n$, and applying a bound to $p^Mf$ proves the stated correction.
\end{proof}

\subsection{Relation to the proposed proof and to the literature}

The height hypothesis does not justify replacing an arbitrary deformation by the particular logarithm $t+t^q/p+t^{q^2}/p^2+\cdots$ in the stated integral coordinate. The proof of Fable suggestion 2 therefore had an unsupported normal-form step. DP1--DP3 remove that step. Even $\widehat\Gm$ in its given coordinate has logarithm $\log(1+t)$; no classification assertion is needed here.

For Lubin--Tate groups over an unramified degree-$h$ extension, Bannai--Kobayashi's period has valuation $c_h$, and their quantity $\overline\rho(0)=|p/p^h|=p^{h-1}$. Their exponential estimate therefore gives $\vp(D^nf(0))\ge nc_h-(h-1)$ \cite[\S2; Propositions 3.3, 3.5]{BK}. This existing case is kept distinct from our general proof. Katz's stronger trace-zero theorem is discussed in Section~\ref{sec:step2}.

\subsection{Finite exact-arithmetic check}

The included script \texttt{research/check\_formal\_bounds.py} computes logarithms from $L(pt+t^q)=pL(t)$ in the Lubin--Tate examples of \cite[initial construction and Lemma 1]{LT}, using exact rational arithmetic, then computes the invariant vector field and its iterates on monomials. It checked 224 derivative orders in six examples $(p,h)=(2,1),(2,2),(2,3),(3,2),(3,3),(5,2)$. All bounds and all tested sharpness witnesses passed. These finite computations are diagnostics, not a formal verification of the theorems. They also illustrate why no pointwise optimality claim should be inferred from a uniform bound: some intermediate orders have strictly larger valuations.

\section{What the divided-power language means}

\begin{proposition}[DP5: distributions, divided powers, and their limit]\label{DP5}
Let $I=(t)\subset R[[t]]$, and filter
\[
 \Dist(F)=\bigcup_{n\ge0}\Hom_R(R[[t]]/I^{n+1},R)
\]
by order, with convolution induced by $F$. Then
\[
 \gr\Dist(F)\simeq\Gamma_R(\Lie F).
\]
Nevertheless, at every finite height the operators $D^n/n!$ do not all preserve $R[[t]]$. Nor does the ideal $(t)$ admit a divided-power structure in $R[[t]]$. After the contraction of DP2 the formal group $F_\pi$ is integrally isomorphic to $\widehat\Ga$ and its normalized invariant derivation does have integral divided powers.
\end{proposition}
\begin{proof}
Let $\epsilon_n$ extract the coefficient of $t^n$. For $m>i+j$, $(\epsilon_i\epsilon_j)(t^m)=0$, because $F(X,Y)$ has order one. For $m=i+j$, only its linear term $X+Y$ contributes, and
\[
 (\epsilon_i\epsilon_j)(t^{i+j})=\binom{i+j}{i}.
\]
Thus $\epsilon_i\epsilon_j\equiv\binom{i+j}{i}\epsilon_{i+j}$ modulo lower order. The associated graded is free of rank one in every degree, with precisely the divided-power multiplication.

For failure of the invariant divided powers, take $n=q$ and $f=t$ in DP2 and \eqref{eq:eval}. Then
\[
 \vp\left(\frac{D^q}{q!}t(0)\right)=\vp(e_q)=-1.
\]
The PD identity $p!\gamma_p(t)=t^p$ would likewise force a coefficient $1/p!$ in $\gamma_p(t)$, which is not in $R$. Thus there is no such PD structure on $(t)$ in the original ring. Finally $L_\pi$ is an integral group isomorphism with integral inverse. In its additive coordinate, $\partial^n/n!$ acts integrally on power series by the binomial coefficient formula. Transport by this isomorphism proves the last assertion.
\end{proof}

The framework is EGA IV$_4$, Theorem 16.11.2, particularly (16.11.2.1)--(16.11.2.2), and SGA 3, Expos\'e VII$_A$, \S\S1--3 \cite{EGA,SGA3}. A PD envelope enlarges the ring to support divided powers; it is not an assertion that division by every factorial was already integral. Likewise a scalar valuation bound does not supply the descent data of a crystal. The precise local crystalline description requires an integrable topologically quasi-nilpotent connection with its compatibilities; see \cite[Proposition 60.17.4, tag 07JH]{Stacks}.

There is also a distinction between restricted nilpotence and zero $p$-curvature. For a connection, the latter compares $\nabla_D^p$ with $\nabla_{D^{[p]}}$. The trivial connection can have zero $p$-curvature when $D^{[p]}\ne0$, as for $\Gm$. It is therefore incorrect to infer a new crystal merely from $D^p\equiv0$, or to regard that congruence as a definition of crystallinity.

\section{Exactly what follows for step 2}\label{sec:step2}

\begin{corollary}[DP6: supersingular local estimate]\label{DP6}
Let $E/R$ be an elliptic scheme with supersingular special fibre. For its normalized invariant derivation and every integral formal function,
\begin{equation}\label{eq:h2}
 \vp(D^nf(0))\ge
 \max\left\{\sum_{r\ge0}\floor{n/p^{2r+1}},\;
 \ceil{\frac{pn}{p^2-1}-1},\;0\right\}.
\end{equation}
Moreover $D^p\equiv0\pmod p$ as an operator. At height one this congruence fails. Its vanishing distinguishes height at least two from height one; it does not distinguish height two from larger height or infinity in general.
\end{corollary}
\begin{proof}
A supersingular elliptic formal group has height two; the ordinary case has height one \cite[\S2.1]{Illusie}. Substitute $h=2$ in DP1 and DP3. The floor sum at $n=p$ equals one when $h\ge2$, so $D^p$ is coefficientwise divisible by $p$. At height one, DP3's sharp witness at $n=p$ makes $D^pt(0)$ a unit. The infinite-height case follows from DP4.
\end{proof}

The leading coefficient is precisely
\[
 c_2=\frac p{p^2-1}=\frac1p+\frac1{p^3}+\frac1{p^5}+\cdots.
\]
This is the coefficient printed in Alp\"oge's \emph{Remark}, p.~3, timeline step 2 \cite{Alpoge}. At a good inert unramified CM prime the elliptic formal group is of this type; Katz's explicit CM setting appears in \cite[\S IV, (4.1)--(4.6), pp.~675--676]{Katz81}. CM is used to select and normalize the relevant arithmetic objects, not in DP6 once the integral height-two group is given.

Here is the precise conditional application to the note. If a local frame and period comparison identify its normalized theta derivatives by
\begin{equation}\label{eq:thetacomparison}
 \frac{D_{\rm an}^n\theta(0)}{\Omega^n}=D^nf(0),\qquad f\in R[[t]],
\end{equation}
then \eqref{eq:h2} applies and improves the logarithmic error to affine error $1$. If $f\in p^{-M}R[[t]]$, subtract $M$; if the normalized analytic derivation is $uD$, add $n\vp(u)$, by DP4. These corrections are not optional. A theta function is a section of a line bundle, and its expression as an integral function and the differentiation of that expression require a chosen frame or an appropriate connection. Supersingularity alone supplies neither the period comparison nor the boundedness of the denominators in \eqref{eq:thetacomparison}.

Thus the local exponent is exactly the height-two mechanism. The entire normalized theta assertion also requires the additional comparisons just specified. The source recovery below identifies these inputs and the later multivariate implementation; this resolves the source-identification question without re-auditing the theorem the user directs us to cite.

\subsection{The exact Katz theorem and the extra hypothesis}

Katz \cite[Main Theorem 3.11(I), p.~672]{Katz81} proves, for the height-two Lubin--Tate group with $[p](X)=pX+\alpha X^{p^2}$ and $\alpha$ a unit, that an integral $f$ satisfying
\[
 \sum_{\lambda\in F[p]}f(X+_F\lambda)=0
\]
has
\[
 \vp(D^nf(0))\ge\floor{np/(p^2-1)}.
\]
The trace-zero hypothesis cannot be discarded. Indeed for $n=p^2-1$ and $f=t^n$, $D^nf(0)=n!$ has valuation $p-1$, whereas the displayed right side is $p$. This is a checked counterexample to that stronger bound for arbitrary integral $f$. Katz's CM application uses modified Weierstrass functions and, in \S IV, assumes $p\ge5$; formulas (4.8)--(4.10) and (4.14) give the integrality and Bernoulli--Hurwitz comparison. The source does not license replacing them by an arbitrary theta expression.

\subsection{Recovered author-source interfaces}\label{sec:recovered}

The two attached Lean files expose the final declarations, not the local lemmas. We recovered targeted modules at public mirror commit \path{7cbf7a0c32fca1badf05931879a37c3f8a2b03b4}, then streamed the author's original fourteen-part archive and compared their SHA-256 hashes. The exact inventory and byte-equality results are included in the source manifests. This verifies the source identity of those modules; it does not execute Lean or independently audit RH. All paths below are relative to \path{lean/src/} in the author archive \cite{AlpogeSource}.

The direct theta normalization appears in \path{Development/AbelianVarietyTheory/BlockTheory/HeightDefect/ThetaFramePullback.lean}: \path{IsGenOver}, \path{IsGenAt}, \path{GenMatch}, \path{ThetaFrame}. Its theta generator generates a section module; it is not automatically a nowhere-vanishing local trivialization. Comparing generic generators uses a nonzero field scalar, which need not be an integral unit. \path{LinePeriodExponential.lean}, in the same directory, specifies the complex uniformization and period identities in \path{IsCurveExp}, \path{IsExp}, \path{IsFamExp}. The integral tangent basis and two quadratic types occur in \path{Development/AbelianVarietyTheory/Definitions/CMLattices/IsDiffOf.lean}, \path{Ell}, \path{IsTen} and \path{HasType}.

For the later multivariate theta bound, the common directory is \path{Formalism/SharpRestriction/Divisibility/}. In \path{ProductZerothFirstParts.lean}, \path{Loc.AdmB} packages \path{B0}--\path{B4}: a normalized theta expression as an integral series times an exponential; unramified coefficients; signed Teichm\"uller invariance on two parity classes; absence of pure-side and bidegree-$(1,1)$ terms; and integral exponential addition coboundaries in both classes. These are extra geometric identities, not consequences of a scalar height bound.

The relevant proved implications, with exact declarations, are as follows. Within \path{JetPolynomialLogs/}, \path{LetterLogIntegrality.integrality_of_letter_log} supplies exponential integrality. \path{UniformBoundAllData.lean}, \path{OneLetterJetKernel.tieB_of} and \path{tieB}, turn \path{B1}--\path{B4} into the unramified field/root-of-unity package \path{BHyp}. \path{OneLetterJetKernel.lean}, \path{kernel_from_twist_holds}, handles the normalized theta kernel; \path{tieFrame} and \path{tieFrameOne} compare coefficients and jets. \path{DworkPowerSeriesWeights.thmA_of} and \path{thmOne_of} assemble the local estimates from these comparisons and the Gauss-jet bounds. The detailed declaration-to-file map is in \path{research/step2_recovered_sources.md}, \S7; a named predicate alone is never counted as a proof of its realization.

The arithmetic bridge is \path{Formalism/LocalToGlobal/Divisibility/JetBounds/JetBoundsProofs/ValuationOfEntryBounds.lean}, \path{Formalism.L2.DIV.valuation_of_entry_bounds}. It assumes the finite path expansion \path{IsExpansion}, integral row/column and final factors, bounded stabilizer denominators, and local entry norms. It then proves the normalized period valuation. These independent geometric path data must not be silently inferred from \path{B0}--\path{B4}. The period requirement is \path{LocalQuadraticParts.lean}, \path{adm2}; its theorem \path{valuation_bound_of_unrounded} passes from the unrounded estimate to the required integer valuation.

There is also a separate controlled-nonvanishing input. In \path{CartierLocalBounds.lean}, \path{vLe(x,b)} means $p^{-b}\le\lVert x\rVert$, hence an \emph{upper} valuation bound. Its \path{ThmP}, \path{CHAt} and \path{ContentAt} find prescribed determinant values with controlled nonzero size. \path{FamilyThetaDenominatorValues.lean}, \path{DepthFromContent}, uses this to control denominator depth. Stronger lower divisibility does not replace it. Moreover the local formulas use both $\gamma=p/(p^2-1)=c_2$ and $\rho=1/(p^2-1)$ in different roles: \path{IsExpansion.expB} has the positive coefficient $\rho$ times a diagram size. Replacing every $p^2$ by $p^h$ would decrease these particular fractions, rather than implement the increasing coefficient $c_h$.

\begin{proposition}[DP21: the unchanged two-step logarithm cannot have higher height]\label{DP21}
The scalar logarithm in \path{Formalism/GaussJetBounds/Statement.lean}, \path{oneLog}, is
\[
 L_2(T)=\sum_{e\ge0}\frac{T^{p^{2e}}}{p^e}
 =T+\frac{T^{p^2}}p+\cdots.
\]
If $F/W(k)$ has height at least three, no integral invertible coordinate change $\phi$ and no unit scalar $a$ can identify $aL_F(\phi(T))$ with $L_2(T)$. The corresponding literal inverse-logarithm substitution in \path{IsSys7.caseB} therefore cannot be retained by replacing its scalar formal group with such an $F$.
\end{proposition}
\begin{proof}
The coefficient induction in DP1 proves $[T^n]L_F\in W(k)$ for all $n<p^h$. Since $p^2<p^h$, the coefficient of degree $p^2$ in $aL_F(\phi(T))$ is a finite sum of integral products. It is integral and hence cannot equal $1/p$. As an independent check of the fixed input, formal inversion gives
\[
 L_2^{-1}(pL_2(T))
 =pT+(1-p^{p^2-1})T^{p^2}+O(T^{p^2+1}).
\]
Thus any integral formal group with this logarithm has height exactly two. This proof concerns an actual scalar logarithm and an integral unit normalization; it does not rule out a different kernel representation for an individual theta function.
\end{proof}

\begin{proposition}[DP24: lower bounds do not give controlled nonvanishing]\label{DP24}
The full set of universal derivative lower bounds does not imply the positive-order determinant nonvanishing in the preceding Cartier interface.
\end{proposition}
\begin{proof}
On any integral product of the formal groups in DP7, take $f=p^N$, $N\ge0$. Every positive-order invariant differential monomial kills it, so every homogeneous determinant differential operator of positive degree also gives zero. All derivative lower bounds hold, while $p^{-b}\le\lVert0\rVert$ is false for every finite real $b$. The actual Cartier theorem assumes a prescribed nonzero exponential main term and unit-minor hypotheses; this example refutes only their inference from the scalar derivative theorem.
\end{proof}

\subsection{The exact exponential gauge and the selected connection}\label{sec:connection}

A further source pass identifies the normalized expression itself. In
\path{ReciprocityTheory/ReductiveGroups/ArithmeticGeometry/BirationalGeometry/ToroidalCompactification/CocycleRightZeroThd.lean},
\path{FB.IsCanonGauge} (lines 174--183), \path{FB.hL} (185), and \path{FB.nabVal} (187--189), it is
\[
 h_L(X)=h(X,[w_0]X),\qquad
 \operatorname{nabVal}(f;P)=\left.P(D)\bigl(f\exp(kh_L)\bigr)\right|_0.
\]
The two addition cocycles, symmetry and CM-adjoint cocycles are integral input; their identities are exponentials of the corresponding coboundaries of the rational gauge $h$. These hypotheses do not assert that $\exp(kh_L)$ is integral. The identity is repeated as \path{GeneratorPowers/ExpTieAlg.lean}, \path{nabVal_eq_const}, lines 413--419; \path{ExpTiePf.tie_holds}, beginning at line 328, identifies the factorial jet of the exponential germ with this value. The exponential-frame comparison in \path{GeneratorPowers/ExpTie.lean}, \path{IsExpGaugeOf}, lines 53--59 is a comparison over a rational coefficient algebra. Its units must not be relabelled integral units.

The integral chart is precise as well. In \path{Formalism/LocalToGlobal/ProjectiveLattices/ChartDirectionsIntegral.lean}, \path{IntegralChartA3}, lines 155--179, supplies integral group, action, biextension-cocycle and tangent data. Its \path{IntData.}\ensuremath{\theta}\path{gen}, lines 74--75, is generation of the global section module. It does not assert that the first nonzero homogeneous Taylor part at the chosen point is primitive, the extra condition in DP35.

There is, however, an actual integral-connection theorem. Set $s_i(X)=-\partial_{Y_i}h(X,0)$ and $S_{ij}=[X_j]s_i$. On $K[[X,Z]]$ the connection and its lattice are
\[
 \nabla_{u,\phi}=D_u+\sum_i(D_us_i+\phi_i)\partial_{Z_i}
       -k\sum_i(Z_i-s_i)u_i,\qquad
 H_h=\{(u,\phi-Su):u,\phi\in R^d\}.
\]
The definitions are in \path{TaylorDiscriminants/ConnectionKillsSeries.lean}, \path{nabla}, lines 31--34, \path{Hset}, lines 41--43, and \path{ConnectionIntegrality}, lines 77--83. Under an integral formal biextension, a canonical gauge, and $w_0+w_0^*=1$, the supplied theorem
\path{TaylorDiscriminants/WickVanishing/ExpConjConnectionIntegral.lean},
\path{exp_conj_connection_integral}, lines 40--82, proves
\[
 e^{-kh_L}\nabla_{u,\phi-Su}(e^{kh_L}G)\in R[[X,Z]]
 \quad(u,\phi\in R^d,\ G\in R[[X,Z]]).
\]
Its \path{connection_first_identity}, lines 84--90, packages that conclusion. Thus full-series nonintegrality does not by itself obstruct the chosen connection operations.

The supplied ordered-column theorem has a denominator budget rather than a height-dependent positive slope. In \path{TaylorDiscriminants/GermSC.lean}, \path{SCSt}, lines 29--41, and \path{GermSCCols.lean}, \path{SCColsSt}, lines 10--22, take a fractional submodule $J\subset K$ containing one, lattice directions in $H_h$, correction entries in $J$, and the displayed vanishing of the three symplectic pairings except for equal colors and opposite sides. If an integral polynomial has columnwise total \path{pairMin} at most $e$, its ordered connection value on $e^{kh_L}f$, for integral $f$, lies in $J^e$. These are supplied by \path{WickVanishing/Closed.lean}, \path{symmetric_canonical_gauge}, lines 25--26, and \path{scaled_connection_columns}, lines 47--63, with \path{lemma_H_columns}, lines 22--23, providing the stated selected-derivative comparison. For $J=p^{-b}R$, the conclusion is a lower valuation $-be$. A new positive higher-height estimate for the actual lattice, corrections and selected operators is additional quantitative work.

Canonical-gauge independence is also supplied, by
\path{Development/AbelianVarietyTheory/Proofs/FormalGroupCharts/CanonGaugeIndependenceAll.lean},
\path{GlobalTheory.canonGauge_independence_all}, lines 23--26. The difference of two canonical gauges is balanced and biadditive; its diagonal in logarithm coordinates is quadratic. The selected value is unchanged when the associated second-order operator annihilates $P$. This removes the specified quadratic ambiguity. It does not remove an arbitrary higher-degree primitive kernel: the series \path{be} in \path{ProductZerothFirstParts.lean}, lines 186--201, is an arbitrary multivariable series. DP38--DP39 below give exact tests of that distinction. The complete source-read scope, including which of the recovered modules were actually read, is recorded in \path{research/theta_normalization_harmonic_trace.md}.

\subsection{A genuine higher-height scalar kernel test}\label{sec:kerneltest}

The unchanged-logarithm obstruction DP21 leaves a concrete alternative: replace the logarithm itself and test the actual analytic hypotheses. Fix an odd prime, $h\ge2$, $q=p^h$, and put
\[
 L_h(T)=\sum_{e\ge0}T^{q^e}/p^e,\qquad E_h=L_h^{-1},\qquad
 \rho_h=(q-1)^{-1},\quad s=(p-1)^{-1}.
\]
The Honda formal group $F_h(S,T)=E_h(L_h(S)+L_h(T))$ is integral and has height $h$: integrality is the functional-equation lemma of Hazewinkel \cite[Lemma 2.3 and (2.5), p.~74]{Hazewinkel}, with coefficient ring $\Z_p$, $\sigma=\id$, and first functional-equation coefficient $1/p$. Its $p$-series begins $pT+(1-p^{q-1})T^q$. DP2, or direct scaling of this logarithm, gives
\begin{equation}\label{eq:honda-scale}
 E_{h,\pi}(Z)=\pi^{-1}E_h(\pi Z)\in R'[[Z]],\qquad
 \overline E_{h,\pi}(Z)=\sum_{r\ge0}(-1)^rZ^{q^r},\quad \pi^{q-1}=p.
\end{equation}
Here and below $R'$ denotes the valuation ring after adjoining $\pi$. In particular $v([T^n]E_h)\ge-\rho_h(n-1)$, with equality for $n=q^r$.

For this test, let $K$ be a valued extension of $\Q_p$, with $v(p)=1$, containing the $(q-1)$-st roots of unity, and let $R$ be its valuation ring, and choose a primitive root $\zeta$. For $\beta\in K[[X,Y]]$ impose the following explicitly modified scalar hypotheses:
\begin{enumerate}[label=(H\arabic*),leftmargin=2.8em]
\item its coefficients lie in an unramified subfield;
\item $\beta(\zeta X,\zeta^{-p}Y)=\beta(X,Y)$;
\item every nonzero monomial has positive bidegrees, and $\beta_{1,1}=0$;
\item the two addition coboundaries
\[
 H_X=\exp\bigl(\delta_X\beta(L_h(S),L_h(S'),L_h(T))\bigr),\qquad
 H_Y=\exp\bigl(\delta_Y\beta(L_h(S),L_h(T),L_h(T'))\bigr)
\]
are integral units with constant coefficient one. Here $\delta_X\beta=\beta(X+X',Y)-\beta(X,Y)-\beta(X',Y)$, and similarly for $Y$.
\end{enumerate}
This keeps the characters $1,-p$, replaces the root order by $p^h-1$, and replaces both scalar logarithms. It is a well-defined analytic proposal, not the assertion that a one-dimensional height-$h$ group is its own Cartier dual or that it realizes an $h$-letter theta correspondence. At $h=2$ these are the scalar conditions in \path{GaussJetBounds/Statement.lean}, \path{BHyp}, lines 39--94.

\begin{proposition}[DP28: the direct higher-height intermediate estimate is false]\label{DP28}
For every odd $p$ and $h\ge3$, hypotheses (H1)--(H4) do not imply
\[
 v(\beta_{a,b})+\rho_h(a+b)\ge s-1.
\]
This refutes the stated substitution in the intermediate \path{LemmaS}, not the final derivative estimate.
\end{proposition}
\begin{proof}
Take $I(S,T)=1+S^pT$ and
\[
 \beta(X,Y)=\log I(E_h(X),E_h(Y)).
\]
Its coefficients lie in $\Q_p$. Since $E_h(\eta X)=\eta E_h(X)$ for $\eta^{q-1}=1$, the torus action cancels in $E_h(X)^pE_h(Y)$, proving (H2). Positive bidegrees and order at least $p+1$ give (H3). Formal logarithm and exponential identities give
\[
 H_X=\frac{I(F_h(S,S'),T)}{I(S,T)I(S',T)},\qquad
 H_Y=\frac{I(S,F_h(T,T'))}{I(S,T)I(S,T')}.
\]
All numerators are integral and all denominators are integral units, proving (H4).

Set $k=p^{h-1}<q$. In
\[
 \beta=\sum_{m\ge1}\frac{(-1)^{m-1}}m E_h(X)^{pm}E_h(Y)^m
\]
the coefficient of $Y^k$ can occur only for $m=k$, because $E_h(Y)=Y+O(Y^q)$. The coefficient of $X^{pk}=X^q$ in the remaining factor is one. Thus
\[
 \beta_{q,p^{h-1}}=\frac{(-1)^{p^{h-1}-1}}{p^{h-1}},\qquad
 v(\beta_{q,p^{h-1}})+\rho_h(q+p^{h-1})
 =2-h+\frac{p^{h-1}+1}{p^h-1}<s-1.
\]
The last inequality follows from $2-h\le-1$ and $(p-1)(p^{h-1}+1)<p^h-1$, equivalently $p<p^{h-1}$.
\end{proof}

The first specific height-two input in this source dependency is \path{OneLetterRadiusIdentity.lean}, \path{weightQuantization}, lines 48--99, especially lines 73--76: $(p+1)\rho_2=s$. It is called in \path{ConstantCoeffEvenDerivativeEqZero.lean}, \path{asmS_diff}, lines 98--163, before \path{lemmaS_of}, lines 165--180. The invariant monomial $X^pY$ already disproves that numerical quantization with $\rho_h$ when $h>2$. The logarithm inverse and its radius estimate themselves do extend. The exact source and author-byte verification are in the accompanying kernel memo and manifest.

\begin{proposition}[DP29: a valid scalar estimate and a sharp nonzero example]\label{DP29}
Under (H2)--(H4), let $U\in R[[S_j:j\in J]]$, with $J$ finite, let $\ell_j(\mathbf t)$ be integral linear forms in $N$ variables, let $\epsilon\in R^\times$, and let $\sigma\in\{1,-1\}$. Set
\[
 \Theta(\mathbf t)=U(E_h(\ell_j(\mathbf t)):j\in J)
 \exp\left(\frac\sigma2\sum_j\beta(\ell_j(\mathbf t),\epsilon\ell_j(\mathbf t))\right).
\]
For $n=|\kappa|$,
\begin{equation}\label{eq:scalar-old-residual}
 v\bigl(\kappa![\mathbf t^\kappa]\Theta\bigr)
 \ge \rho_2n-\frac{\sum_i s_p(\kappa_i)}{p-1}
 \ge\rho_2n-N(\log_p n+1)\quad(n\ge1).
\end{equation}
For the nonzero correction $\beta$ in DP28, the particular function
\[
 \Theta_h(t)=E_h(t)\left(1+E_h(t)^pE_h(\epsilon t)\right)^{1/2}
\]
has, for every $n=q^r$,
\begin{equation}\label{eq:sharp-corrected-scalar}
 v(n![t^n]\Theta_h)=v_p(n!)-\frac{n-1}{q-1}=c_h(n-1).
\end{equation}
Thus the failed intermediate estimate is compatible with exact nonvanishing at the larger derivative slope.
\end{proposition}
\begin{proof}
Write $w_r(A)=\inf_\nu(v(A_\nu)+r|\nu|)$. Products add lower weight bounds, and substitution of series of weight at least $r$ preserves them; these assertions hold coefficient by coefficient. Formula \eqref{eq:honda-scale} gives $w_r(E_h)\ge r$ for $r\ge\rho_h$.

Torus invariance and positive bidegrees force $\ord\beta\ge p+1$: if $a,b>0$ and $a+b\le p$, then $a-pb\ne0$ and $|a-pb|<p^2-1\le q-1$. Put $r_0=p/(p^2-1)$. Then $r_0\ge\rho_h$ and $(p+1)r_0=s+1$. Both integral series $H_X-1,H_Y-1$ have order at least $p+1$. Substituting $E_h$ gives weight at least $s+1$, so their formal logarithms have the same lower bound: the $m$-th term has weight at least $m(s+1)-v_p(m)\ge s+1$. Hence
\[
 w_{r_0}(\delta_X\beta),\ w_{r_0}(\delta_Y\beta)\ge s+1.
\]
For every integer $a\ge2$ there is $0<i<a$ with $v_p\binom ai\le1$. If $a$ is not a power of $p$, expanding $(1+T)^a=\prod_r(1+T^{p^r})^{a_r}$ modulo $p$ supplies an interior $i=p^r$ with unit coefficient. If $a=p^r$, take $i=p^{r-1}$ and use $\binom{p^r}{p^{r-1}}=p\binom{p^r-1}{p^{r-1}-1}$, the second factor being a unit. Every monomial of $\beta$ has one exponent at least two. Its corresponding coefficient in one coboundary is exactly this binomial coefficient times $\beta_{a,b}$. Thus $w_{r_0}(\beta)\ge s$.

Integral linear substitutions and the factor $\sigma/2$ preserve this bound, since $p$ is odd. The exponential has weight at least zero by $v_p(m!)\le ms$; the $U$ factor does also. Therefore $v([\mathbf t^\kappa]\Theta)\ge-r_0n$. Adding Legendre's exact formula for $v_p(\kappa!)$ gives \eqref{eq:scalar-old-residual}, since $s-r_0=\rho_2$. The digit estimate gives its final inequality; zero coordinates contribute zero. At $n=0$ only integrality of $U(0)$ is required.

For the explicit example, $\binom{1/2}{m}\in\Z_p$ for odd $p$: approximate $1/2$ by nonnegative integers in $\Z_p$ and use continuity of each binomial polynomial. Hence
\[
 \pi^{-1}\Theta_h(\pi z)
 =E_{h,\pi}(z)\left(1+\pi^{p+1}E_{h,\pi}(z)^pE_{h,\pi}(\epsilon z)\right)^{1/2}
 \equiv E_{h,\pi}(z)\pmod\pi.
\]
Its coefficient at $z^{q^r}$ is a unit by \eqref{eq:honda-scale}. Rescaling and $s_p(q^r)=1$ prove \eqref{eq:sharp-corrected-scalar}.
\end{proof}

The explicit correction is an integral unit after substitution of the formal logarithm: $\exp\beta(L_h(S),L_h(T))=I(S,T)$. The following classification shows that this behavior holds throughout this scalar class over an unramified base. It supplies neither the prescribed Cartier minors nor a theta-lift comparison.

\begin{theorem}[DP31: the scalar higher-height correction is integral]\label{DP31}
Let $h>2$, let $k$ contain $\F_q$, and now take $R=W(k)$ and $K=R[1/p]$. Under (H2)--(H4),
\begin{equation}\label{eq:scalar-gauge}
 B(S,T):=\exp\beta(L_h(S),L_h(T))\in1+ST\,R[[S,T]].
\end{equation}
Consequently the scalar function of DP29 satisfies the stronger estimate
\begin{equation}\label{eq:scalar-full-slope}
 v_p(\kappa![\mathbf t^\kappa]\Theta)
 \ge c_h|\kappa|-\frac{\sum_i s_p(\kappa_i)}{p-1}.
\end{equation}
The slope $c_h$ is optimal within this class, by \eqref{eq:sharp-corrected-scalar}. This theorem concerns two scalar copies of the height-$h$ group, not the genuinely dual pair.
\end{theorem}
\begin{proof}
Let $G$ have law $F_h$, and write $R_n=R/p^nR$. We use the established formal biextension theorem: Zink \cite[\S4.4, Proposition 121 and (243), pp.~241, 246]{Zink} gives an injection
\[
 \operatorname{Biext}^1(G_n\times G_n,\widehat{\Gm})
 \hookrightarrow\Hom(G_n,G_n^t).
\]
The formal groups here come from displays; the existence/equivalence applies over the Artinian local rings $R_n$ by \cite[Theorem 9, p.~132, or Theorem 103, p.~221]{Zink}. The dual display gives the Cartier-dual formal group; its dual operators are specified in Definition~19 and (22)--(24), pp.~143--144, and the formal-dual identification in Proposition~121. Both $G$ and its Cartier dual are connected for $h>1$, with special-fibre slopes $1/h$ and $(h-1)/h$ respectively. The hypotheses of the injection are therefore satisfied; we do not need the additional surjectivity in Zink's Theorem 122.

First, the integral units $C_X=H_X$ and $C_Y=H_Y$ define a biextension without presupposing \eqref{eq:scalar-gauge}. Over $K$ they are the two multiplicative coboundaries of $B$. Therefore they satisfy symmetry, associativity in each variable, and the interchange identity. These identities hold over $R$, since $R\hookrightarrow K$. On the trivial formal multiplicative torsor over $G\times G$, define the first partial product by
\[
 (S,T,u)\star_X(S',T,u')
  =(F_h(S,S'),T,uu'C_X(S,S',T)),
\]
and define the second analogously with $C_Y$. The identities just stated make a biextension. A normalized splitting is exactly a unit $I$ with
\begin{equation}\label{eq:gauge-split}
 \frac{I(F_h(S,S'),T)}{I(S,T)I(S',T)}=C_X,\qquad
 \frac{I(S,F_h(T,T'))}{I(S,T)I(S,T')}=C_Y,
 \qquad I(S,0)=I(0,T)=1.
\end{equation}

Second, $\Hom(G_n,G_n^t)=0$. Over $k$, distinct slopes give zero rational Hom, by the semisimple isocrystal classification after faithful extension to an algebraic closure \cite[\S1.3, first paragraph of p.~120]{Katz79}, hence zero integral Hom. The integral Hom embeds in its rationalization: a homomorphism killed by $p^a$ is zero since multiplication by $p^a$ on a $p$-divisible group is surjective as a sheaf. To lift the vanishing, let $f:G_n\to G_n^t$ reduce to zero, and suppose its formal-coordinate components are in $p^jR_n[[T]]$, $1\le j<n$. Write $f=p^jg$ componentwise. The target $p$-series has linear part $p$ times the identity and higher terms of order at least two, so
\[
 f([p]_G(T))=[p]_{G^t}(f(T))\equiv0\pmod {p^{j+1}}.
\]
After division by $p^j$ and reduction, each component satisfies $\bar g([p]_{G_k}(T))=0$. Substitution of the nonzero series $[p]_{G_k}(T)$, whose initial degree is $q$, is injective on $k[[T]]$. Thus $f$ is divisible by $p^{j+1}$, and induction proves $f=0$.

The biextension injection now yields a splitting $I_n\in1+ST\,R_n[[S,T]]$ at every $n$. We verify compatibility rather than assume that independently chosen splittings lift. Two normalized splittings differ by a bihomomorphism $J$, which obeys
\[
 J([p]_G(S),T)=J(S,T)^p.
\]
Modulo $p$, let $b>0$ be the smallest $T$-degree in $J-1$, if it is nonzero. Its coefficient on the left remains nonzero under substitution of $[p]_{G_k}(S)$; on the right every positive $T$-degree is at least $pb>b$. This contradiction proves $J=1$ modulo $p$. If $J-1=p^jV$ for $j\ge1$, the right side minus one is divisible by $p^{j+1}$. Dividing and reducing gives $\bar V([p]_{G_k}(S),T)=0$, and injectivity of this substitution on $k[[S,T]]$ proves the next divisibility. Hence $J=1$ over every $R_n$.

Thus the normalized splittings are unique and automatically compatible. Their coefficientwise inverse limit is $I\in1+ST\,R[[S,T]]$ satisfying \eqref{eq:gauge-split}. In additive coordinates,
\[
 A(X,Y)=\beta(X,Y)-\log I(E_h(X),E_h(Y))
\]
is separately additive, by equality of the two coboundaries. In characteristic zero, interior binomial coefficients force a separately additive series to be $aXY$. The logarithm and exponential are Teichm\"uller-equivariant. Consequently the assumed torus symmetry of $\beta$ makes $I_n(\zeta S,\zeta^{-p}T)$ another normalized splitting. Uniqueness makes $I$, and hence $A$, invariant. But $XY$ has character $\zeta^{1-p}\ne1$, so $a=0$. This proves $B=I$ and \eqref{eq:scalar-gauge}.

Finally, write the correction factor in $\Theta$ as
\[
 \prod_j B(E_h(\ell_j),E_h(\epsilon\ell_j))^{\sigma/2}.
\]
Each binomial power is an integral series because $p$ is odd. The whole expression is an integral function of finitely many inverse-logarithm coordinates of weight at least $\rho_h$. Therefore $v_p([\mathbf t^\kappa]\Theta)\ge-\rho_h|\kappa|$. Legendre's factorial formula gives \eqref{eq:scalar-full-slope}. The explicit nonzero example in DP29 attains the slope along $q$-power orders.
\end{proof}

This closes the scalar analytic replacement question under its stated hypotheses despite DP28's failed intermediate estimate. It also identifies the geometric limitation precisely. For the genuine pair $G\times G^t$, the biextension target is $\Hom(G,(G^t)^t)=\operatorname{End}(G)$, which is nonzero; the second factor has dimension $h-1$. The scalar proof therefore does not trivialize or bound that dual kernel. Its construction and arithmetic comparison remain separate tasks.

\section{Higher-height directions in abelian schemes}

\begin{theorem}[DP7: integral directional and mixed derivatives]\label{DP7}
Let $A/R$ be an abelian scheme and $i:F\to\widehat A$ an integral homomorphism from a one-dimensional formal group of finite height $h$. Let $D_i$ be the invariant tangent direction induced by $D_F$. For every integral formal function $f$ on $\widehat A$,
\[
 (D_i^nf)(0)=D_F^n(i^*f)(0),\qquad
 \vp((D_i^nf)(0))\ge C_h(n),
\]
where
\[
 C_h(n)=\max\{E_h(n),\ceil{nc_h-(h-1)},0\}.
\]
If an integral formal product has one-dimensional factors $F_1,\ldots,F_r$ of heights $h_i$, its mixed invariant derivatives satisfy
\[
 \vp(D_1^{n_1}\cdots D_r^{n_r}f(0))\ge\sum_i C_{h_i}(n_i).
\]
Use $C_\infty(n)=\vp(n!)$ for an infinite-height factor.
\end{theorem}
\begin{proof}
A formal-group homomorphism intertwines translations; differentiation at the identity therefore intertwines the indicated invariant derivations. Iterating and evaluating gives the identity, and $i^*f$ is integral. Apply DP1, DP3 and DP4. For a product, apply the coefficientwise operator bound one variable at a time, with the other formal variables as coefficient parameters. Each operator preserves the ideals already obtained from the other factors, so their exponents add. All substitutions are finite at a fixed multidegree.
\end{proof}

An integral subgroup is enough for evaluation; a direct factor is not essential. Conversely a summand only on the special fibre, or a factor only up to isogeny after inverting $p$, does not supply the integral map and tangent lattice assumed in DP7. These are hypotheses that a CM construction must verify.

\subsection{Ramified base rings and the integral CM action}

\begin{theorem}[DP19: full unramified CM action over a ramified base]\label{DP19}
Let $A$ be a complete discrete valuation ring of mixed characteristic $(0,p)$, with $v(p)=1$. Let $\mathcal O\subset A$ be the integers of an unramified degree-$h$ extension of $\Q_p$, and let $F/A$ be a one-dimensional formal group with an integral $\mathcal O$-action satisfying $[a]'(0)=a$. Suppose its special fibre has $p$-height $h$. Then both coefficientwise bounds of DP1 and DP3 hold over $A$:
\[
 D^nA[[t]]\subset p^{E_h(n)}A[[t]],\qquad
 D^nA[[t]]\subset p^{\max\{0,\ceil{nc_h-(h-1)}\}}A[[t]].
\]
Thus DP7 applies to an integral lifted local CM factor of this type even when the base is ramified. An action only on the special fibre is insufficient.
\end{theorem}
\begin{proof}
Put $q=p^h$. Average over the Teichm\"uller group to obtain the strict integral coordinate
\[
 x(t)=\frac1{q-1}\sum_{\zeta\in\mu_{q-1}}\zeta^{-1}[\zeta](t).
\]
Its derivative is one, and reindexing gives $x([\eta](t))=\eta x(t)$. In this coordinate $[\zeta](x)=\zeta x$. Commutation with $[p]$, and uniqueness of the normalized logarithm, imply that $[p](x)$, $L(x)$ and its inverse $E(x)$ have support only in degrees congruent to $1$ modulo $q-1$. For any other degree a suitable $\zeta^{m-1}-1$ is a unit, forcing that coefficient to vanish. In particular
\[
 [p](x)=px+a_qx^q+\sum_{m>q}a_mx^m,
 \qquad a_q\in A^\times.
\]

Adjoin $\pi^{q-1}=p$ and work in the valuation ring of the finite extension. All coefficients of $\pi^{-1}[p](\pi x)$ belong to $pA'$, so the coefficient induction of DP2 makes $L_\pi$ and $E_\pi$ integral. Scaling gives
\[
 v([z^n]E(z)^m)\ge-\frac{n-m}{q-1}.
\]
For a nonzero coefficient, the support condition makes $n-m$ divisible by $q-1$. Thus, for $n,m\ge1$, the right side is an integer at least $-\floor{(n-1)/(q-1)}$. This supplies the integer bound without assuming that the valuation group of $A$ is $\Z$. The identity $D^nf(0)=n![z^n]f(E(z))$ now gives the preliminary bound of DP3. Translation makes it coefficientwise, and composition according to the base-$p$ digits gives exactly $E_h(n)$ by the finite identities already proved there.

For the affine assertion, $L(x)=x+bx^q+O(x^{q+1})$, with $b\in p^{-1}A$. Integral translation operators give $D^j=j!\Delta_j$ for $j<q$ and
\[
 D^q=p^{B_h}uD+p^{B_h+1}T,
 \qquad u\in A,\quad T\text{ integral}.
\]
Use the \emph{integer} recurrence $b_h(n)$ from DP1r. It proves containment in $p^{b_h(n)}A[[x]]$ even over this ramified base. The same numerical induction gives $b_h(n)\ge nc_h-(h-1)$, so its integrality gives the asserted ceiling. One must not instead round a possibly fractional valuation in $A$. The averaging coordinate preserves the normalized tangent and the ideals in question. If $h=1$ both bounds are zero and follow already from integrality; when $q=2$ the averaging group is trivial. This also covers that endpoint.
\end{proof}

\begin{proposition}[DP20: residue height alone fails over a ramified base]\label{DP20}
For every $h\ge2$ and every integer $e>p$, there is an integral one-dimensional formal group over a complete discrete valuation ring with $v(p)=1$, $v(\varpi)=1/e$, whose special fibre has height $h$ but
\[
 v(D^pt(0))=\frac1e<1,
 \qquad
 v(D^{p^s}t(0))=\frac{p^s-1}{e(p-1)}\quad(s\ge0).
\]
In particular the proposed slope $c_h$ with an $O(\log n)$ error fails for this fixed ramified deformation. The special fibre can have full unramified CM action while that action fails to lift integrally to the deformation.
\end{proposition}
\begin{proof}
Choose a finite field $k$ and a height-$h$ group $\Phi/k$ normalized to be additive below degree $p^h$. Set $A=W(k)[\varpi]/(\varpi^e-p)$. Lubin--Tate \cite[Proposition 1.1, pp.~50--51]{LT66} gives an actual integral deformation over $W(k)[[t_1,\ldots,t_{h-1}]]$ whose degree-$p$ term is $t_1C_p$, where
\[
 C_p(X,Y)=\frac{(X+Y)^p-X^p-Y^p}{p}.
\]
Specialize $t_1=\varpi$ and the other parameters to zero. A sign convention for $C_p$ is absorbed by changing the sign of $t_1$. The result is an integral formal group, still of residue height $h$, with
\[
 F(X,Y)=X+Y+\varpi C_p(X,Y)+O((X,Y)^{p+1}).
\]
Thus $D=(1+\varpi t^{p-1}+O(t^p))\partial_t$, and $D^pt(0)=(p-1)!\varpi$. Summing the degree-$p$ cocycle in the $p$-fold group sum telescopes and gives
\[
 [p]_F(t)=pt+a t^p+O(t^{p+1}),\qquad
 a=(p^{p-1}-1)\varpi.
\]

Choose $\beta^{p-1}=p/\varpi$ in a finite extension and put $r=v(\beta)=(1-1/e)/(p-1)>1/p$. The contracted $p$-series belongs to $pA'[[x]]$, and after division by $p$ reduces modulo the maximal ideal to $x+cx^p$, with $c\ne0$. Indeed its degree-$p$ coefficient is $a/\varpi$, a unit, while each degree $m\ge p+1$ has valuation at least $(m-1)r-1>0$ after this division. Its logarithm and exponential are integral, and the logarithm recursion shows that the reduced logarithm is exactly $x+cx^p$. Its inverse has nonzero coefficients at every $p^s$, by the recurrence $b_0=1$, $b_s=-cb_{s-1}^p$. Scaling back therefore gives
\[
 v([z^{p^s}]E(z))=-(p^s-1)r,
\]
and multiplication by $(p^s)!$ gives the claimed equality. Since $1/(e(p-1))<c_h$, the deficit from the proposed slope is linear on this subsequence and cannot be absorbed in a logarithmic error.

One may choose $\Phi$ as the reduction of an unramified Lubin--Tate $\mathcal O$-module. If that full action lifted to the deformation just constructed, DP19 would give $v(D^pt(0))\ge1$, a contradiction. This proves the final distinction.
\end{proof}

\begin{proposition}[DP8: dimension matters, and the CM pairing constraint]\label{DP8}
The one-dimensional formula cannot use the total height of an arbitrary multidimensional formal summand. Under full integral CM by a degree-$2g$ CM field $E$, good reduction, and $p$ unramified in $E$, let $G_v$ denote a local $p$-divisible factor, of height $h_v$ and dimension $d_v$. Then
\[
 h_v=[E_v:\Q_p],\qquad d_v=\#\Phi_v,\qquad
 d_v+d_{\bar v}=h_v.
\]
In particular a self-conjugate one-dimensional factor has height two. A one-dimensional factor of height $h>2$ has a distinct conjugate factor of dimension $h-1$, and the ambient dimension satisfies $g\ge h$.
\end{proposition}
\begin{proof}
For the counterexample take $\widehat\Gm^{\,r}$, $r\ge2$, whose dimension and total height are both $r$. In the first coordinate set $D=(1+t_1)\partial/\partial t_1$ and $f=1+t_1$. Then $D^nf=f$ and $D^nf(0)=1$ for every $n$. Substituting $h=r$ into a positive one-dimensional slope contradicts this equality for large $n$. Such ordinary factors also arise in products of ordinary CM elliptic curves; the formal counterexample does not depend on that realization.

For the CM assertion, the integral idempotents of $\mathcal O_E\otimes\Z_p$ decompose $A[p^\infty]$. The height and dimension counts are the standard CM decomposition, expressed by the local CM type \cite[\S5, definition p.~106, (10*) and Theorem 3 p.~107]{Tate352}. One can see the counts directly from the rank-one $E\otimes\Q_p$ Tate module in characteristic zero and the eigenspace decomposition of $\Lie A$ after a splitting base extension. Each pair of conjugate embeddings contains exactly one member of the CM type. Counting those over $v$ and $\bar v$ gives $d_v+d_{\bar v}=h_v$. If $v=\bar v$ and $d_v=1$, this says $h_v=2$. Otherwise the two distinct factors contribute $1+(h_v-1)=h_v$ to the dimension of $A$, proving the last assertion.
\end{proof}

\begin{proposition}[DP9: size of the gain and a refuted implication]\label{DP9}
For $h>2$, $c_h>c_2$, but
\[
 \frac{c_h}{c_2}=\frac{(p+1)(p^{h-1}-1)}{p^h-1}<1+\frac1p.
\]
The derivative-divisibility assertions alone do not imply positive generic-fibre vanishing multiplicity of an integral formal function.
\end{proposition}
\begin{proof}
Subtracting consecutive terms in $c_h=1/(p-1)-1/(p^h-1)$ proves strict increase; their limit is $1/(p-1)$. Division by $c_2=p/(p^2-1)$ gives the bound, or cross-multiplication gives the displayed exact ratio. For the second assertion, $f=1+t$ is a unit in $R[[t]]$, so has order zero and empty local zero divisor, while satisfying DP1 and DP3. Also $p^Nf$ has valuations increased by $N$ without changing its generic-fibre divisor, because $p^N$ is invertible there. These are counterexamples to the stated implication.
\end{proof}

This does not refute the existence of a specially constructed shorter theta chain. Step 8 of the note is an automorphic construction between totally definite unitary groups; steps 9--10, 12, 14--15 and 18 add integral structures, periods, norm and level estimates, cost and covolumes. No operation in DP7 changes that chain. The missing comparison is specified in Part II. In particular ``$m$ lifts achieve multiplicity $m$'' is not an upper bound saying that fewer lifts could never achieve $m$.

\begin{proposition}[DP22: a concrete height-three CM factor]\label{DP22}
Let $J$ be the Jacobian of $C:y^2=x^7-1$. Over $R=W(\F_{11^3})$ it has good reduction and integral CM by $\Z[\zeta_7]$. Its $11$-divisible group has two integral factors of height three and dimensions two and one. The dimension-one factor is an integral formal direct factor of $\widehat J$, so DP7 applies with
\[
 c_3(11)=\frac{66}{665},\qquad c_2(11)=\frac{11}{120},
 \qquad\frac{c_3(11)}{c_2(11)}=\frac{144}{133}.
\]
\end{proposition}
\begin{proof}
Glue the charts $y^2=x^7-1$ and $v^2=u-u^8$ by $u=x^{-1}$ and $v=y/x^4$. They give a finite double cover of $\mathbf P^1$, hence a projective curve. At $11$ the affine chart is smooth by its partial derivatives; at infinity the derivative in $u$ is a unit. Infinity is a section. The two generators $y$ and $v=x^{-4}y$ give $\pi_*\mathcal O_C=\mathcal O_{\mathbf P^1}\oplus\mathcal O_{\mathbf P^1}(-4)$. On each geometric fibre its two-chart \v Cech quotient for the second summand is
\[
 \bar k[x,x^{-1}]/\bigl(\bar k[x]+x^{-4}\bar k[x^{-1}]\bigr),
\]
with basis $x^{-1},x^{-2},x^{-3}$; the first summand has zero $H^1$. Thus the genus is three. The fibres are geometrically integral because $x^7-1$ has simple zeros and is not a square in $\bar k(x)$. The affine two-chart cover has affine intersection, so $H^2(C_s,\mathcal O)=0$. The Picard representability, smoothness and properness theorems therefore make its relative Jacobian an abelian scheme \cite[expos\'e 232, Theorem 3.1; expos\'e 236, Corollaries 1.7 and 2.3, Proposition 2.10(ii)]{FGAPic}.

The ring $R$ contains $\zeta_7$, since $7\mid11^3-1$. The automorphism $\sigma(x,y)=(\zeta_7x,y)$ extends to the model and fixes infinity. Let it act covariantly on divisor classes of $J$. The regular forms $x^{i-1}dx/y$, $i=1,2,3$, have eigenvalues $\zeta_7^i$; their orders at infinity are $4,2,0$, respectively. Abel--Jacobi pullback identifies these with the forms of $J$. Thus $\Phi_7(\sigma)$ has zero differential in characteristic zero and is the zero endomorphism; a characteristic-zero abelian-variety homomorphism with zero differential has zero-dimensional connected image. Irreducibility of $\Phi_7$ gives the integral action $\Z[\zeta_7]\hookrightarrow\operatorname{End}(J)$, of degree $6=2\dim J$. The tangent type is $\{1,2,3\}$.

Frobenius at $11\equiv4\pmod7$ has embedding orbits $\{1,2,4\}$ and $\{3,5,6\}$. The integral local CM order is a product of two unramified cubic orders; its idempotents split $J[11^\infty]$ integrally. The CM height/type formula \cite[\S5, (10*) and Theorem 3, pp.~106--108]{Tate352} gives heights three and dimensions two and one. Their slopes are $2/3$ and $1/3$. The latter factor therefore has no \'etale part and is connected. The connected formal-group equivalence \cite[\S2.2, Proposition 1, p.~162; \S2.3, Example (b), p.~166]{Tate67} transports the idempotent product to the formal completion. It identifies this factor with an integral one-dimensional formal group embedded as a direct factor in $\widehat J$; its tangent lattice is split, so no nonunit tangent normalization is hidden. The numerical formulas follow by substituting $p=11,h=3$.
\end{proof}

This example realizes the local hypothesis. Its length-three orbits do not satisfy the unchanged even alternating type or the logarithm of DP21. No global theta cost or shorter chain is asserted for this Jacobian.

\begin{proposition}[DP23: the dual-product attempt does not restore the old blocks]\label{DP23}
Let $G_h$ be a one-dimensional unramified Lubin--Tate $p$-divisible group of height $h>2$. The balanced product $G_h\times G_h^D$ cannot be a product of rank-one alternating CM blocks of the existing two-type model, even up to isogeny. Direct products and finite unramified restrictions of scalars do not remove this obstruction.
\end{proposition}
\begin{proof}
The two factors have height $h$ and dimensions $1$ and $h-1$. Their CM slopes are $1/h$ and $(h-1)/h$, by the same height/type formula. A rank-one alternating orbit of even length $d$ has $d/2$ selected embeddings, hence dimension $d/2$ and slope $1/2$. If $h$ is odd, an alternating length-$h$ orbit is impossible; if $h$ is even and greater than two, the counts $1,h-1$ disagree with $h/2$. Combining the factors does not interlace their Frobenius orbits, since their central CM idempotents remain Frobenius-stable. More invariantly, isogeny preserves Newton slopes. Direct products concatenate the slope multisets, and an unramified restriction of scalars becomes a product of conjugates after an unramified splitting extension, repeating the same slopes. Since neither original slope equals $1/2$, these operations cannot produce the required alternating blocks.
\end{proof}

The general higher-rank source object can have rank defects and additional blocks; DP23 is explicitly about retaining the stated rank-one alternating inputs. A different signature would require a different kernel, denominator and cost analysis.

\begin{proposition}[DP25: the induced CM type blocks a larger-field substitution]\label{DP25}
Let a CM field $E$ contain an imaginary quadratic field $K$, and let
\[
 \Phi=\{\tau:E\hookrightarrow\mathbf C:\tau|_K=\kappa\}
\]
be the full CM type induced from one embedding $\kappa$ of $K$. Assume good reduction and full integral CM at $p$, with $p$ unramified in $E$. If $p$ is inert in $K$, every rank-one local CM factor has slope $1/2$; a one-dimensional factor has height two. If $p$ splits in $K$, each factor has slope zero or one; a one-dimensional connected factor has height one. Enlarging $E$ while retaining this type cannot produce a one-dimensional factor of height greater than two.

The basic CM blocks of the inspected theta construction have precisely these induced types. The Jacobian of DP22 cannot replace such a block; balancing it with its conjugate also fails the source's stated surjective comparison.
\end{proposition}
\begin{proof}
Restriction of embeddings from $E$ to $K$ commutes with unramified Frobenius. At an inert prime, Frobenius exchanges the two embeddings of $K$, so membership in $\Phi$ alternates around each orbit. Every orbit has even length $d$, with exactly $d/2$ selected embeddings. At a split prime that membership is constant, so the count is zero or $d$. The full-CM height/type formula \cite[\S5, (10*) and Theorem 3]{Tate352} gives height $d$, dimension equal to the count, and slope their ratio. The assertions follow, without requiring an orbit to be stable under complex conjugation.

The source premise is explicit: \path{Formalism/SharpRestriction/RamificationFrobenius/TotientGaloisClasses/ConjEMemHostFIff.lean}, \path{HostRaw.K} and \path{HostRaw.cmType}, define $K=\Q(\sqrt{-d})$ and select exactly one of its embeddings. In \path{Development/AbelianVarietyTheory/Definitions/CMLattices/IsDiffOf.lean}, \path{HasType}, \path{IsTen}, and \path{CMObj.typeP/typeM} impose the two conjugate induced types and their Serre-tensor construction from a quadratic CM elliptic curve. This is a definition of the blocks, not an inference from definiteness of the unitary groups.

For DP22, write $S=\zeta_7+\zeta_7^2+\zeta_7^4$ and $T=\zeta_7^3+\zeta_7^5+\zeta_7^6$. Multiplication gives $S+T=-1$, $ST=2$, and $(S-T)^2=-7$. Thus the unique quadratic subfield has embedding classes $\{1,2,4\}$ and $\{3,5,6\}$. The type $\{1,2,3\}$ is neither class. Relabeling by a Galois automorphism preserves or exchanges its counts $(2,1)$; an $E$-linear isogeny preserves the type. At $11$, $-7\equiv4$ is a square, so that prime splits in $K$. It is outside the source's \path{InertK} gain class, by \path{LocalRingCardinality/DataHost.lean}, \path{primeClass}, and \path{SharpRestriction/GlobalTheoremInputs.lean}, \path{InertK}.

For $J^a\times\overline J^{b}$ the multiplicities on the two quadratic classes are $(a,a,b)$ and $(a,b,b)$, hence become constant only if $a=b$. For $a=b>0$ the target has type $(a,a)$, but \path{CMObj.onto} also requires a surjection from the corresponding copies of the two induced-type blocks. Equal dimensions would make this an isogeny. Those source blocks are ordinary at $11$, whereas this target has only slopes $1/3$ and $2/3$. Isogeny invariance of Newton slopes rules out that comparison. These statements concern the specified CM action and the existing source data, not every possible mixed-type construction.
\end{proof}

\subsection{The prime-sum gain under a pure height substitution}

\begin{proposition}[DP17: a summable change in the ramp certificate]\label{DP17}
For $A\ge0$, a prime $p$, an integer $h\ge2$ and weights $w_i\ge0$, set
\[
 R_h(A,p;w)=\sum_{\substack{i\ge1\\h\nmid i}}\max\{0,A/p^i-w_i\}.
\]
Then
\begin{equation}\label{eq:rampcomparison}
 |R_h(A,p;w)-R_2(A,p;w)|\le\frac{A}{p^2-1}.
\end{equation}
Consequently, with $C=\sum_p\log(p)/(p^2-1)<\infty$, any common prime set and cutoff satisfy
\[
 \left|\sum_p\log(p)\bigl(R_{h(p)}(A,p;w)-R_2(A,p;w)\bigr)\right|\le CA.
\]
The height may vary with the prime. If $w_i\ge1$, the sums already terminate when $p^i>A$. In the zero-weight case the slopes satisfy
\[
 0\le c_h-c_2=\frac1{p^2-1}-\frac1{p^h-1}\le\frac1{p^2-1}.
\]
\end{proposition}
\begin{proof}
Put $b_i=\max\{0,A/p^i-w_i\}$, so $0\le b_i\le A/p^i$. The terms added to the odd-index sum are even indices not divisible by $h$; the terms removed are odd indices divisible by $h$. The added sum is at most $A\sum_{r\ge1}p^{-2r}=A/(p^2-1)$. If $h$ is even there are no removed terms. If $h$ is odd, then $h\ge3$, and their sum is at most
\[
 A\sum_{r\ge0}p^{-(2r+1)h}=\frac{Ap^h}{p^{2h}-1}\le\frac{A}{p^2-1}.
\]
The last inequality follows on multiplying the positive denominators. The difference of two nonnegative numbers bounded by the same constant has absolute value at most that constant, proving \eqref{eq:rampcomparison}. Summing it with weights $\log p$ proves the second assertion; convergence follows by comparison with $\sum_{n\ge2}\log(n)/n^2$. The cutoff assertion is immediate. For zero weights the geometric sums give $R_h=Ac_h$ and the final formula.
\end{proof}

This is the ramp shape in the recovered \path{Formalism/LocalRingCardinality/DataParams}\textsubscript{2}\path{.lean}, \path{SizeData}\textsubscript{2}\path{.vonMangoldtClassSum}, with \path{wPart} at least one. The place-weight identity \path{FactEtaMean} makes the corresponding total change at most $2C\sum_{j,t}A_t$ for fixed remaining inputs. It compares a \emph{lower-bound certificate}: \path{GlobalTheoremInputs.LinkG_at} does not assert that this certificate is the exact or largest possible gain. Thus pure height substitution leaves unchanged any established leading $A\log X$ prime-sum coefficient, but the proposition is not an upper bound on the gain of every new geometric construction. In particular it does not exclude a saving of one lift when the number of lifts grows, or an improvement after leading terms cancel. See the exact source record in Section~\ref{sec:recovered}.

\begin{proposition}[DP26: a finite numerical saving is not excluded]\label{DP26}
For the ramps of DP17, let $A=p^3$ and $w_i=1$. Then
\[
 R_2=p^2-1=(p-1)(p+1),\qquad
 R_3=p^2+p-2=(p-1)(p+2),
\]
so $(p+1)R_3=(p+2)R_2$. At $p=11$ this gives
\[
 12\cdot130=13\cdot120,
 \qquad 24\cdot130=26\cdot120=3120.
\]
The second equality is compatible with even lift counts. It does not construct a theta chain or assert that its other costs are unchanged.
\end{proposition}
\begin{proof}
Only $i=1,2$ can contribute positively: their terms are $p^2-1$ and $p-1$. The $i=3$ term is zero, and all $i\ge4$ give $p^{3-i}<1$. Height two retains only the first positive term; height three retains both. This proves the formulas. The source-shaped finite cutoff is \path{range}$(\lfloor\log_p A\rfloor+2)=\{0,1,2,3,4\}$, giving the same values. The term $i=0$ is excluded by both divisibility-index conventions, and every weight satisfies $w_i\ge1$.
\end{proof}
This example shows why boundedness and summability of the local improvement cannot themselves rule out every finite saving. For the particular Jacobian at $11$, DP25 independently prevents its insertion into the old inert-prime gain data.

\begin{proposition}[DP18: limitation of the unchanged error budget]\label{DP18}
Let $x\ge1$, $\epsilon\ge0$ and $b,m,n,N_1,N_2,k>0$, with $A\ge xb$ and $N_1,N_2\le n$. Define
\begin{align*}
 \mathcal B&=N_1A(1+m^2/k)
 +m\left(2nmbx^\epsilon+\frac{N_2kbx^\epsilon}{n(nx^{-\epsilon})}\right),\\
 \mathcal S&=\frac{N_1mbN_2}{2n}.
\end{align*}
Then
\[
 \frac{\mathcal B}{\mathcal S}\ge\frac{2x}{m}+4m.
\]
For $m\le x^{1/2-\delta}$, $\delta>0$, the first term alone is at least $2x^{1/2+\delta}$. Multiplying the signal by a positive uniformly bounded factor cannot change that exponent.
\end{proposition}
\begin{proof}
Set $d_1=N_1/n$ and $d_2=N_2/n$, both in $(0,1]$. Direct division gives
\[
 \frac{\mathcal B}{\mathcal S}
 =\frac{2A}{mbd_2}(1+m^2/k)
  +\frac{4mx^\epsilon}{d_1d_2}
  +\frac{2kx^{2\epsilon}}{d_1n^2}.
\]
Use $A\ge xb$, $x^\epsilon\ge1$ and $d_1,d_2\le1$, and discard only nonnegative terms. Substitution of the upper bound on $m$ proves the last assertion.
\end{proof}

The displayed inner allowance and signal occur in \path{GrowthBoundGivesRiemannHypothesis.lean}, \path{SizeWindows}; the additional outer factor $x^{C\epsilon}\ge1$ only enlarges the allowance, and the scale $A\ge xb$ is present in the actual growth data, not in the bare generic window definition. DP18 concerns this \emph{unchanged upper-bound allowance}. A lower bound on an error allowance is not a lower bound on the actual error. The proposition therefore obstructs certification by a pure bounded-coefficient replacement, and does not prove that a new construction or sharper estimate cannot use a shorter chain.

\begin{proposition}[DP27: parity and endpoint constraints of the existing chain]\label{DP27}
In \path{Formalism/LocalRingCardinality/Params.lean}, \path{Ranks} requires an even lift count $M\ge2$. With its $k,t$ fixed, its endpoint ranks satisfy
\[
 N(M)-N(L)=(M-L)(t+1)>0\qquad(0\le L<M).
\]
Thus a replacement by $M-1$ violates the same parameter specification, and literal truncation cannot preserve the old hermitian endpoint by an isometric isomorphism. The auxiliary signature at $1\le L<M$ is $(k+L-1,M-L)$, whereas at $M$ it is $(k+M-1,0)$. The column spaces $X_j$ themselves are definite at every stage. A truncated column does not, by definition, supply the period identification at the old endpoint.
\end{proposition}
\begin{proof}
The parity statement is the field \path{Ranks.m_even}; the difference between two allowed even counts is even. The definitions are $N(j)=k+j(t+1)$ and $\operatorname{sig}W(j)=(k+j-1,M-j)$. Since $t+1>0$, the displayed rank difference is positive. The source field $k\ge M+3$ and $1\le L<M$ make both entries of the intermediate auxiliary signature positive. These numerical ranks and signatures are imposed on the actual spaces by \path{DataSpaces.lean}, \path{Spaces.X_rank}, \path{W_sigval}; its separate \path{X_def} requires all column spaces to be definite. Rank is invariant under an isometric isomorphism, proving the endpoint obstruction.

Finally \path{DataKernels.lean}, \path{ColumnChain.col}, is \path{chainOpM} at exactly $M$, with output \path{C.cls M}\,$\to$\,\path{C.Y M}\,$\to\mathbf C$. The field \path{PeriodMap.evz} consumes that terminal space. Truncation gives stage $L$ instead. A numerical lower bound on values supplies none of the new representation or period-comparison data. This is a limitation of literal truncation, not a prohibition on separately constructing admissible data with a smaller even count.
\end{proof}

The source also permits formal regrouping of two kernels: \path{Formalism/ThetaLiftIndependence/Statement.lean}, \path{kerMul}, expands the composition as the sum over the intermediate variable. Its \path{ChainIntM} retains the product of the denominator factors. That identity shortens the written composition only; it does not identify the composite with one geometrically realized theta kernel having the cost of a single old lift. A new such identification would be an additional theorem.

\subsection{The actual Cartier dual has different directional divisibility}

\begin{proposition}[DP30: a dual direction with a unit $p$-th derivative]\label{DP30}
Let $A$ be a complete mixed-characteristic discrete valuation ring with perfect residue field $k$. Let $G/A$ be a one-dimensional strict formal $\mathcal O$-module, where $\mathcal O\subset A$ is the ring of integers of the unramified degree-$h$ extension of $\Q_p$, and suppose its special fibre has height exactly $h\ge2$. Let $H=G[p^\infty]^D$ be its actual Cartier dual and $\widehat H$ its connected formal group. Then $\Lie(H_k)$ has a basis $e_0,\ldots,e_{h-2}$ with
\[
 e_i^{[p]}=e_{i+1}\ (i<h-2),\qquad e_{h-2}^{[p]}=0.
\]
Thus the restricted $p$-operation has kernel dimension one and image dimension $h-2$. If $h>2$, there are a primitive integral invariant derivation $D$ on $\widehat H$ and an integral formal coordinate $f$ such that $D^pf(0)\in A^\times$.
\end{proposition}
\begin{proof}
Teichm\"uller averaging as in DP19 gives an integral coordinate $x$ with $[\zeta](x)=\zeta x$ for every $\zeta^{p^h-1}=1$. The group law then satisfies $F(\zeta X,\zeta Y)=\zeta F(X,Y)$, so its nonzero monomials have total degree congruent to one modulo $q-1$, where $q=p^h$. In particular $F=X+Y+$ terms of degree at least $q$, integrally and after Frobenius twisting.

For the next calculation work over $k$. Relative Frobenius $F_G:G\to G^{(p)}$ and Verschiebung $V_G:G^{(p)}\to G$ compose in either order to multiplication by $p$; their kernel orders are $p$ and $p^{h-1}$. These facts and the duality of $F,V$ are Tate \cite[\S2.2, diagram (3), pp.~163--164; \S2.3, Proposition 3, p.~166]{Tate67}. Restrict $V_G$ to $G^{(p)}[p]$. Its image has order $p$ and lies in $G[F_G]$, so it is exactly $G[F_G]$. Exact finite Cartier duality \cite[\S1.2--1.3, pp.~159--160]{Tate67} gives
\[
 H[F_H]^D\simeq G[p]/G[F_G].
\]
Height $h$ implies $[p](x)=u(x)x^q$ with $u(0)\ne0$, so $G[p]=G[F_G^h]$ as subgroup schemes, by equality of their defining ideals. If $r=p^{h-1}$, relative Frobenius induces
\[
 G[p]/G[F_G]\simeq G^{(p)}[F^{h-1}].
\]
On coordinates this is $k[z]/(z^r)\hookrightarrow k[x]/(x^q)$, $z\mapsto x^p$. The target is free of rank $p$ with basis $1,x,\ldots,x^{p-1}$, proving faithful flatness and the asserted quotient. In $k[X,Y]/(X^r,Y^r)$ every surviving monomial has degree at most $2r-2<q$. The nonlinear part of the twisted group law vanishes. Hence
\begin{equation}\label{eq:dual-F-kernel}
 H[F_H]^D\simeq\alpha_{p^{h-1}}.
\end{equation}

Put $P=H[F_H]$. Its distribution Hopf algebra is, by duality,
\[
 \Dist(P)=k[t]/(t^r),\qquad\Delta(t)=t\otimes1+1\otimes t.
\]
The Lie algebra is its primitive subspace: the augmentation Leibniz rule for a tangent functional is precisely the primitive-element identity in the dual Hopf algebra. This subspace has basis $t,t^p,\ldots,t^{p^{h-2}}$. Indeed these powers are primitive, while for any other exponent $1\le n<r$ a nonzero base-$p$ digit gives an interior binomial coefficient $\binom n{p^j}\not\equiv0\pmod p$. The resulting mixed tensor term survives; different total degrees cannot cancel it. Constants are not primitive.

The restricted operation is the $p$-th power in this algebra: convolution powers of a tangent functional correspond to powers of its invariant derivation. It shifts the displayed basis and kills its last member. Since $dF_H=0$, the inclusion $P\hookrightarrow H$ identifies their Lie algebras and their restricted operations. Perfectness of $k$ gives the stated image and kernel dimensions.

The connected formal group exists by Tate \cite[\S2.2, Proposition 1, p.~162 and (4), p.~164]{Tate67}; its tangent lattice is free of rank $h-1$, since dual dimensions sum to height by Proposition 3. For $h>2$, lift a tangent vector $v$ with $v^{[p]}\ne0$ to this lattice and let $D$ be its invariant derivation. The lift is primitive. Some integral coordinate $f$ has differential detecting $v^{[p]}$. Reduction commutes with $D$ and its iterates, so
\[
 \overline{D^pf(0)}=df_0(v^{[p]})\ne0.
\]
Thus $D^pf(0)$ is a unit. No assertion that the full $H$ is connected is needed for this last argument.
\end{proof}

The original one-dimensional height-$h$ line has divisibility at least $p$ in order $p$. DP30 proves that this cannot be assigned to every primitive direction of its actual dual. It does not prove that a prescribed theta entry detects the bad direction, or quantify a global denominator cost. For the genus-three case the dual restricted Lie algebra is already explicit: $e_0^{[p]}=e_1$, $e_1^{[p]}=0$.

\begin{theorem}[DP32: sharp asymptotic slopes on the actual Honda dual]\label{DP32}
Let $p$ be any prime, $h\ge2$, $q=p^h$, $d=h-1$, and $\mathcal O=W(\F_q)$. The Cartier dual $H=G_h^D$ of the Honda group with logarithm $L_h(T)=\sum_{n\ge0}T^{q^n}/p^n$ has integral coordinates $X_1,\ldots,X_d$ and normalized primitive invariant derivations $D_i$ for which, setting
\[
 c_i=\frac{p^i}{q-1},\qquad r_i=\frac1{p-1}-c_i\quad(1\le i\le d),
\]
every integral formal function satisfies coefficientwise
\begin{equation}\label{eq:dual-mixed}
 v_p(D_1^{a_1}\cdots D_d^{a_d}f)
 \ge\max\left\{0,\ \sum_i c_i a_i-\frac{\sum_i s_p(a_i)}{p-1}\right\}.
\end{equation}
For every $i$ and $n=q^r$,
\begin{equation}\label{eq:dual-sharp}
 v_p(D_i^nX_i(0))=c_i(n-1).
\end{equation}
Thus every displayed directional slope is optimal, and $\sum_i c_i=c_h$. These assertions persist over $W(k)\supset\mathcal O$. In particular the order-$p$ obstruction DP30 does not imply absence of positive asymptotic gain on the dual.
\end{theorem}
\begin{proof}
We construct an integral group and prove its dual identification before estimating it. Over $\Z_p$, define a vector of monomials
\[
 M_1(\mathbf X)=X_d^{p^2},\qquad M_i(\mathbf X)=X_{i-1}^p\ (i\ge2),
 \qquad \mathcal L(\mathbf X)=\sum_{n\ge0}\frac{M^{\circ n}(\mathbf X)}{p^n}.
\]
For $d=1$, this means $M_1=X_1^{p^2}$. The degrees tend to infinity, so the series is well defined and has identity linear term. Its functional equations are
\[
 \mathcal L_1=X_1+p^{-1}\mathcal L_d(\mathbf X^{p^2}),\qquad
 \mathcal L_i=X_i+p^{-1}\mathcal L_{i-1}(\mathbf X^p)\ (i\ge2).
\]
Apply Hazewinkel \cite[Lemma 2.3, (2.2)--(2.6), p.~74]{Hazewinkel} with $B=\Z_p$, ideal $pB$, $\sigma=\id$, the lemma's $q=p$, matrix $s_1$ having $1/p$ just below the diagonal, and $s_2$ having $1/p$ in entry $(1,d)$. When $d=1$, $s_1=0$. It gives an integral commutative law $F_H=\mathcal L^{-1}(\mathcal L(\mathbf X)+\mathcal L(\mathbf Y))$ and an integral $p$-series. The exact identity
\[
 \mathcal L(M(\mathbf X))=p(\mathcal L(\mathbf X)-\mathbf X)
\]
and the congruence clause \cite[(2.8), p.~75]{Hazewinkel} imply $[p]_H\equiv M\pmod p$.

Here the finite-flat condition can be checked directly. Put $C=\Z_p[[U_1,\ldots,U_d]]$, $B=\Z_p[[X_1,\ldots,X_d]]$ and send $U_i$ to $[p]_{H,i}(\mathbf X)$. Modulo $p$, $B/p$ is free over $C/p$ on the $p^h$ monomials
\[
 X_1^{b_1}\cdots X_d^{b_d},\qquad
 0\le b_i<p\ (i<d),\quad 0\le b_d<p^2.
\]
Their lifts define $C^{p^h}\to B$. Successively lifting the coefficient expansion modulo $p$ and dividing the remainder by $p$ proves surjectivity by $p$-adic completeness. A relation has all coefficients divisible by $p$; dividing by $p$ is valid since $B$ is $p$-torsion-free. Repetition and $p$-adic separatedness prove injectivity. Thus $[p]_H$ is finite free of degree $p^h$. Its iterated kernels form a connected $p$-divisible group of height $h$ and dimension $d$, as in Tate \cite[\S2.2, Proposition 1, pp.~162--164]{Tate67}.

Everything is defined over $\Z_p$. On the special fibre over $\F_p$, a full cycle of $M$ has total Frobenius exponent $2+(d-1)=h$, so $[p]_H^d=M^d=F_H^h$. For $G=H^D$, dual dimensions give $\dim G=1$ \cite[\S2.3, Proposition 3, p.~166]{Tate67}, and duality gives $V_G^h=[p]_G^d$. Since $F_G^hV_G^h=[p]_G^h$ and $d=h-1$, cancellation gives $F_G^h=[p]_G$. Cancellation is valid because multiplication by $p^d$ is surjective as an fppf sheaf. All these maps are over $\F_p$, so they are actual endomorphisms with no unidentified twists. Their Frobenius kernels show that $G$ is connected. It therefore has an integral one-dimensional formal group over $\Z_p$ by Tate's equivalence. In an integral coordinate $T$ its $p$-series $f$ has linear term $pT$ and reduction $T^q$.

Over $\mathcal O=W(\F_q)$, this is precisely a Lubin--Tate series for uniformizer $p$. The known uniqueness and strict-isomorphism theorem \cite[\S5, pp.~76--77, (5.1)--(5.4)]{Hazewinkel} identifies the group with this $p$-series with the normalized Lubin--Tate group $F_f$, and identifies all such groups for the same uniformizer integrally and strictly. The Honda group $G_h$ has a $p$-series in this class. Thus $G\simeq G_h$ integrally over $\mathcal O$, and dualizing identifies the constructed $H$ with the actual Cartier dual claimed in the statement. No guessed dual logarithm or unproved lifting of an action is involved.

We now compute its radius in each coordinate. Put $S_h=(q-1)/(p-1)$, choose $\pi^{q-1}=p$, and set $\alpha_i=\pi^{S_h-p^i}$, so $v_p(\alpha_i)=r_i>0$. Scale coordinate $i$ in input and output by $\alpha_i$. For every arrow $X_j^{p^e}$ in $M_i$ one has $p^er_j-r_i=1$. Consequently the $n=1$ term of
\[
 \mathcal L'_i(\mathbf Z)=\alpha_i^{-1}\mathcal L_i(\alpha_1Z_1,\ldots,\alpha_dZ_d)
\]
has coefficient exactly one. If an $n$-fold cyclic iterate has successive exponents $e_1,\ldots,e_n$, its coefficient valuation after scaling is
\[
 1+p^{e_1}+p^{e_1+e_2}+\cdots+p^{e_1+\cdots+e_{n-1}}-n.
\]
This is positive for $n\ge2$. Hence $\mathcal L'$ is integral and reduces to $\mathbf Z+M(\mathbf Z)$. Its inverse $\mathcal E'$ is integral. In characteristic $p$, the monomial map $M$ is additive, and therefore
\[
 \overline{\mathcal E'}(\mathbf Z)=\sum_{n\ge0}(-1)^nM^{\circ n}(\mathbf Z).
\]
Composition with $\id+M$ verifies the cancellation, also at $p=2$. In component $i$, after setting other variables to zero, terms survive exactly at cyclic step numbers $dr$, with degrees $q^r$ and nonzero coefficients $(-1)^{dr}$. Undoing the input and output scales gives
\[
 v_p([t^{q^r}]\mathcal E_i(t e_i))=-(q^r-1)r_i.
\]

The normalized invariant derivations are integral and commute. In additive logarithm coordinates they are partial derivatives, so
\[
 D^{\mathbf a}f(0)=\mathbf a![\mathbf t^{\mathbf a}]f(\mathcal E(\mathbf t)).
\]
Integrality of $\mathcal E'$ says that every coefficient of $\mathcal E_j$ has weighted valuation at least $r_j$ when variable $t_i$ is given weight $r_i$. Integral substitution therefore bounds $[\mathbf t^{\mathbf a}]f(\mathcal E)$ by $-\sum_i r_i a_i$. Adding Legendre's exact factorial valuation proves \eqref{eq:dual-mixed} at zero; integrality also supplies the bound zero. Apply the same computation in $\mathbf Y$ to the integral translated function $f(F_H(\mathbf X,\mathbf Y))$, treating $\mathbf X$ as coefficient variables. Translation invariance gives the coefficientwise bound. At $n=q^r$, the exact exponential coefficient above and $s_p(n)=1$ prove \eqref{eq:dual-sharp}. Finally
\[
 \sum_{i=1}^{h-1}c_i=\frac{p+\cdots+p^{h-1}}{p^h-1}=c_h.
\]
This proves all assertions, including sharpness against a larger slope with sublinear error.
\end{proof}

For the equally weighted list of these directions on $G_h\times G_h^D$, the mean slope is $2c_h/h$. If $h>2$, DP9 gives
\[
 \frac{2c_h}{hc_2}<\frac{2(1+1/p)}h\le1.
\]
Thus this particular mean is smaller than $c_2$. This is a precise comparison of the displayed coordinate directions, not a lower bound on an automorphic cost; a theta construction need not weight them equally. The actual dual's integral model and directional asymptotics are now determined. Its prescribed biextension/theta entries and the global cost comparison remain to be supplied.

\begin{theorem}[DP33: optimal affine errors for the height-three dual]\label{DP33}
In the two-coordinate model of DP32 with $h=3$, write $X=X_1$, $Y=X_2$, $q=p^3$, $c_X=p/(q-1)$ and $c_Y=p^2/(q-1)$. For every integral $f$ the coefficientwise bounds improve to
\[
 v_p(D_X^nf)\ge\max\{0,\ceil{c_Xn-1}\},\qquad
 v_p(D_Y^nf)\ge\max\{0,\ceil{c_Yn-1}\}.
\]
Both slopes and both constant errors $1$ are optimal. Mixed derivatives have the sum of the two displayed integer exponents as a lower bound. In addition to the exact diagonal sequences in DP32,
\[
 v_p(D_X^{pq^r}Y(0))=\frac{p^2(q^r-1)}{q-1},\qquad
 v_p(D_Y^{p^2q^r}X(0))=\frac{p(q^{r+1}-1)}{q-1}.
\]
These statements include $p=2$.
\end{theorem}
\begin{proof}
The actual-dual model and its integral coordinates are already established by DP32. Its logarithm begins
\[
 \mathcal L_1=X+Y^{p^2}/p+X^q/p^2+\cdots,\qquad
 \mathcal L_2=Y+X^p/p+Y^q/p^2+\cdots.
\]
Let $\Delta_{a,b}$ be the integral coefficient operators of translation by $(S,T)$. They commute with one another and with $D_X,D_Y$, by commutativity of the group. In characteristic zero the translation operator is $\exp(\mathcal L_1(S,T)D_X+\mathcal L_2(S,T)D_Y)$. Comparing coefficients gives
\[
 \Delta_{p,0}=D_X^p/p!+D_Y/p,\qquad
 \Delta_{0,p^2}=D_Y^{p^2}/(p^2)!+D_X/p,\qquad
 \Delta_{0,p}=D_Y^p/p!.
\]
Consequently, with integral commuting operators $T,S$ and units $a,b$,
\begin{equation}\label{eq:dual-low-ops}
 D_X^p=aD_Y+pT,\quad a=-(p-1)!,\qquad
 D_Y^{p^2}=p^p bD_X+p^{p+1}S,\quad b=-\frac{(p^2)!}{p^{p+1}}.
\end{equation}
Also $D_Y^p$ is divisible by $p$ as an operator. Raising the first identity to $p$ shows $D_X^{p^2}$ is divisible by $p$: the leading term uses $D_Y^p$, and every remaining binomial term has a factor $p$. Thus for all $n$ the elementary bounds are $\floor{n/p^2}$ in direction $X$ and $\floor{n/p}$ in direction $Y$.

Raise the first identity in \eqref{eq:dual-low-ops} to $p^2$. Its nonleading term indexed by $1\le j\le p^2$ has valuation at least
\[
 v_p\binom{p^2}{j}+j+\floor{(p^2-j)/p}\ge p+1.
\]
For $j=p^2$ this is $p^2\ge p+1$. Otherwise $v_p\binom{p^2}{j}=2-v_p(j)$. If $j=ap+b$ with $1\le b<p$, the expression is $p+1+a(p-1)+b$; if $j=ap$, $1\le a<p$, it is $p+1+a(p-1)$. Applying the second identity to the leading term proves
\[
 D_X^q=p^p u_XD_X+p^{p+1}U_X
\]
with $u_X$ a unit and $U_X$ integral. Raise the second identity of \eqref{eq:dual-low-ops} to $p$. Its leading term uses the first identity; every mixed term $1\le j<p$ has valuation at least $1+p(p-j)+(p+1)j=p^2+j+1$, and its last term has valuation $p(p+1)\ge p^2+1$. Hence
\[
 D_Y^q=p^{p^2}u_YD_Y+p^{p^2+1}U_Y.
\]

For either direction put $(r,B,c)=(p^2,p,c_X)$ or $(p,p^2,c_Y)$. Then $c=B/(q-1)<1$ and $cr=q/(q-1)$. For $0\le n<q$, the maximum of $cn-\floor{n/r}$ is one, attained at $n=q-1$: within a block it increases, and successive block maxima increase by $cr-1>0$. This proves the affine inequality at the initial orders. For $n\ge q$, apply the common recurrence $D^q=p^BuD+p^{B+1}U$ to $D^{n-q}f$. Induction bounds the first term by $cn-1$ and the second by $cn-1+(1-c)$. Integrality and integer valuations give the ceiling. Since the bounds are on the whole formal function ring, composing the two operator containments proves the mixed assertion.

The affine constants are attained by the same monomial $f=X^{p-1}Y^{p^2-1}$ at $n=q-1$. Along the $X$ logarithmic axis the leading terms of the exponential are $X=t+\cdots$, $Y=-t^p/p+\cdots$, so $f$ has first degree $q-1$ with coefficient $(-1/p)^{p^2-1}$. Since $v_p((q-1)!)=p^2+p-2$, the derivative has valuation $p-1=c_X(q-1)-1$. Along the $Y$ axis the leading terms are $X=-t^{p^2}/p+\cdots$, $Y=t+\cdots$; the coefficient is $(-1/p)^{p-1}$ and the derivative valuation is $p^2-1=c_Y(q-1)-1$.

For the additional exact sequences, DP32's scaled inverse here reduces to
\[
 \overline{\mathcal E'}_1(u,v)=\sum_{r\ge0}u^{q^r}-\sum_{r\ge0}v^{p^2q^r},\qquad
 \overline{\mathcal E'}_2(u,v)=\sum_{r\ge0}v^{q^r}-\sum_{r\ge0}u^{pq^r}.
\]
The nonzero off-diagonal coefficients, the respective input/output weights $r_X=(p^2+1)/(q-1)$ and $r_Y=(p+1)/(q-1)$, and Legendre's formula give exactly the two stated valuations. The diagonal sequences from DP32 prove slope optimality. In particular $D_X^pY(0)=-(p-1)!$ is a unit, while $c_X>0$.
\end{proof}

\begin{theorem}[DP34: optimal constant-one affine bounds on every Honda dual]\label{DP34}
In the actual-dual coordinates of DP32, for every $h\ge2$, every $i\in\{1,\ldots,h-1\}$ and every integral $f$,
\[
 D_i^nf\in p^{b_i(n)}R[[\mathbf X]],\qquad
 b_i(n)=\max\{0,\ceil{c_i n-1}\},\qquad c_i=\frac{p^i}{p^h-1}.
\]
Both the slope and the constant $1$ are optimal. With $q=p^h$ and $d=h-1$, the same monomial
\[
 f_*(\mathbf X)=\left(\prod_{j<d}X_j^{p-1}\right)X_d^{p^2-1}
\]
satisfies $v_p(D_i^{q-1}f_*(0))=p^i-1$ for every $i$. Mixed derivatives have the sum $\sum_i b_i(a_i)$ as a lower bound. The displayed ideal containments persist under integral base change.
\end{theorem}
\begin{proof}
Work first over $\Z_p$ in the model already identified by DP32. Integral translation coefficients and the first nonlinear terms of its logarithm give, exactly as in DP33,
\begin{equation}\label{eq:all-dual-low}
 D_i^p=aD_{i+1}+pT_i\quad(i<d),\qquad
 D_d^{p^2}=p^pbD_1+p^{p+1}T_d,\qquad D_d^p\in p\operatorname{End}_{\rm int}.
\end{equation}
Here $a=-(p-1)!$, $b=-(p^2)!/p^{p+1}$ are units; all remainder operators preserve the integral function ring and commute. The last two identities also hold when $d=1$. Descending induction, raising the first identity to $p^{h-i-1}$, shows that $D_i^{R_i}\in p\operatorname{End}_{\rm int}$ for $R_i=p^{h-i}$. Hence $D_i^n$ has the preliminary divisibility $\floor{n/R_i}$.

Fix $i<d$, put $B=p^i$, and start with $D_i^q$. Successively substitute the first identity of \eqref{eq:all-dual-low} to move its leading term through directions $j=i,i+1,\ldots,d-1$. At the step $j=i+t$, this identity is raised to $N=p^{h-t-1}$, and $N/R_{j+1}=B$. Its $m$-th nonleading term has valuation at least
\[
 v_p\binom Nm+m+\floor{(N-m)/R_{j+1}}
 =B+v_p\binom Nm+m-\ceil{m/R_{j+1}}.
\]
For $1\le m<N$, the binomial coefficient is divisible by $p$ because $N$ is a $p$-power, and $m-\ceil{m/R_{j+1}}\ge0$. For $m=N$ the bound is $N\ge p^{i+1}\ge B+1$. Thus every discarded term belongs to $p^{B+1}\operatorname{End}_{\rm int}$. The leading term at the end is a unit times $D_d^{p^{i+1}}$. Raise the second identity of \eqref{eq:all-dual-low} to $p^{i-1}$ and factor out $p^B$. Modulo $p$, what remains is a unit times $D_1^{p^{i-1}}$, which the first identities modulo $p$ identify with a unit times $D_i$. Therefore
\begin{equation}\label{eq:all-dual-q}
 D_i^q=p^{p^i}u_iD_i+p^{p^i+1}U_i
\end{equation}
with a unit $u_i$ and an integral operator $U_i$. For $i=d$, use the same last step directly by raising the second identity to $p^{d-1}$. This includes $d=1$.

Now $c_iR_i=q/(q-1)$ and $c_i(q-1)=p^i$. For $0\le n<q$, the increasing block maxima give $\floor{n/R_i}\ge c_i n-1$, with equality in the real difference at $n=q-1$. For $n\ge q$, applying \eqref{eq:all-dual-q} to $D_i^{n-q}f$ bounds its first term by $c_i n-1$ and its second by $c_i n-1+(1-c_i)\ge c_i n-1$. This proves the affine bound by induction. All coefficients and operators were over $\Z_p$, so integer valuations give the asserted ideal containment there. Base change preserves that containment; no fractional valuation is rounded on a ramified base. Composition proves the mixed bound.

It remains to prove the sharp common witness, including exclusion of possible lower-degree terms. The explicit logarithm and its inverse satisfy the Teichm\"uller equivariance
\[
 \mathcal L_j(\zeta^pX_1,\zeta^{p^2}X_2,\ldots,\zeta^{p^d}X_d)
 =\zeta^{p^j}\mathcal L_j(\mathbf X),\qquad\zeta^{q-1}=1.
\]
For the ordinary arrows this is $p\cdot p^{j-1}=p^j$; for the exceptional arrow it is $p^2p^d\equiv p\pmod{q-1}$. Thus a nonzero degree $n$ in $\mathcal E_j(te_i)$ requires $np^i\equiv p^j\pmod{q-1}$. Its smallest possible positive degree is
\[
 m_{ij}=p^{j-i}\ (j\ge i),\qquad m_{ij}=p^{h+j-i}\ (j<i).
\]
DP32's scaled inverse gives a nonzero coefficient at this first path degree, with exact valuation $r_j-m_{ij}r_i$. Equivariance excludes all smaller degrees. If $e_j=p-1$ for $j<d$ and $e_d=p^2-1$, then
\[
 \sum_j e_jm_{ij}=q-1,\qquad \sum_j e_jr_j=d.
\]
The first equality telescopes: the ranges $i\le j<d$, $j=d$, and $j<i$ contribute $p^{d-i}-1$, $p^{h-i+1}-p^{d-i}$, and $q-p^{h-i+1}$ respectively. For the second, use $\sum e_j=(p-1)(h+p-1)$ and $\sum e_jp^j=p(q-1)$. The unique lowest-degree term of $f_*(\mathcal E(te_i))$ therefore has degree $q-1$ and valuation $d-(q-1)r_i$. Finally
\[
 v_p((q-1)!)=\frac{q-1}{p-1}-h
\]
makes the derivative valuation $c_i(q-1)-h+d=p^i-1$. This proves the optimal constant in every direction. Slope optimality is already the exact sequence in DP32. Every step applies at $p=2$ and $h=2$.
\end{proof}

\subsection{The denominators of a primitive Poincar\'e splitting}

\begin{theorem}[DP35: denominator rigidity and linear growth]\label{DP35}
Let $R$ be a complete discrete valuation ring extending $\Z_p$, with fraction field $K$ and $v(p)=1$. Let $F$ be an integral formal group in variables $\mathbf S$, and let $\mathbf T$ be a further finite list of variables. Suppose
\[
 B\in K[[\mathbf S,\mathbf T]],\qquad B(0,\mathbf T)=B(\mathbf S,0)=1,
 \qquad C_F=\frac{B(F(\mathbf S,\mathbf S'),\mathbf T)}{B(\mathbf S,\mathbf T)B(\mathbf S',\mathbf T)}
       \in R[[\mathbf S,\mathbf S',\mathbf T]]^\times.
\]
Then the following statements hold.
\begin{enumerate}[leftmargin=*,label=(\alph*)]
\item If the coefficients of $B$ have bounded denominators, $B$ is an integral unit.
\item For every rational $t>0$, the uniform weighted infimum
\[
 w_t(B)=\inf_\alpha\bigl(v(B_\alpha)+t|\alpha|\bigr)
\]
is either $0$ or $-\infty$. If a degree-$N$ coefficient has valuation $-m<0$, then
\[
 \liminf_{|\alpha|\longrightarrow\infty}\frac{v(B_\alpha)}{|\alpha|}\le-\frac mN.
\]
\item An integral unit cannot clear unbounded denominators. More generally, if $f\in R[[\mathbf S,\mathbf T]]$ has primitive first nonzero homogeneous part, then, for every integer $M\ge0$,
\[
 p^M fB\in R[[\mathbf S,\mathbf T]]\quad\Longrightarrow\quad
 p^M B\in R[[\mathbf S,\mathbf T]].
\]
\item For the actual pair $G_h,H=G_h^D$ of DP32 over $W(\F_q)$, the primitive formal Poincar\'e class has integral normalized cocycle representatives and rational compatible splittings $B$. Every such splitting is nonintegral, hence has the unbounded and linearly growing denominators in (a)--(b). No fixed power of $p$, integral unit frame, or multiplier as in (c) repairs it.
\end{enumerate}
The last conclusion concerns a compatible splitting of both biextension laws. It does not assert that every rational gauge of a trivial class is integral, or determine the selected harmonic jets of a geometric theta section.
\end{theorem}
\begin{proof}
For a nonzero bounded-denominator series $A$, discreteness gives a minimum coefficient valuation $\nu(A)$. It is multiplicative:
\[
 \nu(AA')=\nu(A)+\nu(A').
\]
Indeed, scale each factor by a power of a uniformizer to have least valuation zero. Its reduction is a nonzero formal series over the residue field; the residue series ring is a domain. Integral substitutions with zero constant term cannot lower $\nu$. Iterating the displayed coboundary through $p$ additions gives
\[
 B([p]_F(\mathbf S),\mathbf T)=B(\mathbf S,\mathbf T)^p U(\mathbf S,\mathbf T),\qquad
 U=\prod_{j=1}^{p-1}C_F([j]_F(\mathbf S),\mathbf S,\mathbf T).
\]
Here $U$ is integral and has constant term one. If $B$ has bounded denominators,
\[
 \nu(B)\le\nu(B([p]_F(\mathbf S),\mathbf T))=p\nu(B).
\]
Thus $\nu(B)\ge0$, and the constant term one makes it zero. This proves (a), over a ramified ring as well.

For (b), adjoin an element $\lambda$ of valuation $t$ in a finite ramified extension. The contracted law
\[
 F_\lambda(\mathbf S,\mathbf S')=\lambda^{-1}F(\lambda\mathbf S,\lambda\mathbf S')
\]
is integral: each coefficient of degree $e\ge2$ is multiplied by $\lambda^{e-1}$. The series $B_\lambda=B(\lambda\mathbf S,\lambda\mathbf T)$ has integral coboundary $C_F(\lambda\mathbf S,\lambda\mathbf S',\lambda\mathbf T)$. If $w_t(B)>-\infty$, part (a) makes $B_\lambda$ integral, so $w_t(B)=0$. A coefficient with valuation $-m$ and degree $N$ has negative weight whenever $0<t<m/N$. Thus $w_t(B)=-\infty$ for each such rational $t$. There are only finitely many coefficients of bounded total degree, so coefficients realizing arbitrarily negative weighted valuations have degrees tending to infinity. They give a liminf at most $-t$. Let $t$ increase rationally to $m/N$.

An integral unit and its inverse preserve boundedness of denominators, proving the first assertion of (c). For the stronger assertion write $A=p^MB$, and let $f_m$ be the first nonzero homogeneous part of $f$. Induct on the degree $n$ of $A_n$. The homogeneous part of $fA$ of degree $m+n$ is
\[
 f_m A_n+\sum_{j>m}f_j A_{m+n-j}.
\]
The latter sum is integral by induction. Hence $f_mA_n$ is integral. The polynomial Gauss valuation is multiplicative by the same residue-domain argument, and $\nu(f_m)=0$. Consequently $A_n$ is integral. This includes the degree-zero start and proves the assertion.

We supply the rational-splitting argument needed for (d), rather than assume that a frame is already a bilinear splitting. An integral normalized frame gives symmetric normalized unit cocycles $C_F(S,S',\mathbf T)$ and $C_H(S,\mathbf T,\mathbf T')$ satisfying both cocycle identities and the interchange identity. Use the explicit characteristic-zero logarithms of $G_h$ and $H$ to write additive coordinates $x,\mathbf y$. Put $A=\log C_F$ and $C=\log C_H$ in those coordinates. With
\[
 \delta_x\beta=\beta(x+x',\mathbf y)-\beta(x,\mathbf y)-\beta(x',\mathbf y)
\]
and the analogous vector coboundary, interchange says $\delta_{\mathbf y}A=\delta_x C$.
Set
\[
 a(x,\mathbf y)=\left.\partial_{x'}A(x,x',\mathbf y)\right|_{x'=0},\qquad
 \beta_0(x,\mathbf y)=\int_0^x a(u,\mathbf y)\,du.
\]
Differentiating the normalized two-cocycle identity gives
\[
 \partial_{x'}A(x,x',\mathbf y)=a(x+x',\mathbf y)-a(x',\mathbf y).
\]
Since $a(0,\mathbf y)=0$, integration and normalization prove $\delta_x\beta_0=A$. Interchange now makes $\delta_{\mathbf y}\beta_0-C$ additive in $x$. In characteristic zero an additive scalar series is linear, so this difference is $x\,d(\mathbf y,\mathbf y')$, where $d$ is a symmetric normalized vector two-cocycle. Put
\[
 b_i(\mathbf y)=\left.\partial_{y_i'}d(\mathbf y,\mathbf y')\right|_{\mathbf y'=0}.
\]
Its cocycle identity gives
\[
 \partial_{y_i'}d(\mathbf y,\mathbf y')=b_i(\mathbf y+\mathbf y')-b_i(\mathbf y').
\]
Commuting the derivatives shows that $\partial_jb_i-\partial_i b_j$ is translation invariant and therefore constant. Symmetry of $d$, evaluated by mixed derivatives at the origin, makes that constant zero. The formal one-form $\sum_i b_i\,dy_i$ is closed. Its radial primitive
\[
 \gamma(\mathbf y)=\int_0^1\sum_i y_i b_i(t\mathbf y)\,dt
\]
has gradient $(b_i)$, zero constant term, and $\delta_{\mathbf y}\gamma=d$ by the differentiated identity and normalization. Thus
\[
 \beta=\beta_0-x\gamma(\mathbf y),\qquad B=\exp\beta
\]
solve both rational coboundary equations in logarithm coordinates. Returning to the integral group coordinates means $B(S,\mathbf T)=\exp\beta(L_G(S),L_H(\mathbf T))$; parts (a)--(c) apply in these group coordinates. The integrations are coefficientwise formal integrations over $K$. The axes are normalized. The difference of two logarithmic solutions is additive separately in $x$ and $\mathbf y$, hence is a bilinear form $x\sum_i a_i y_i$.

For existence of the integral primitive class and its frame, use the display evaluation pairing. Zink's construction \cite[\S4.3, (216)--(218), pp.~238--239]{Zink}, followed by Proposition~121, (243), and Theorem~122, identifies this class with $\id_G\in\operatorname{End}(G)$. Both dual displays are connected; the nilradical hypothesis of Theorem~122 holds over each $R/p^n$. The pairings and classes are compatible under reduction. Zink's Lemma~111, \S4.2, p.~230 supplies an underlying frame over the prorepresentable formal product. Axial rigidifications can be restored by integral units. The same argument can be made compatibly modulo $p^n$: after choosing a frame at one level, lift a frame to the next and multiply by a lift of the unit measuring its discrepancy with the previous frame. The maps of formal unit groups under reduction are surjective. The inverse limit is an integral normalized frame over $R$.

The identity class is nonzero already on the special fibre and is primitive: an equality $pf=\id_G$ with an integral endomorphism would give $p\,df=1$ on the tangent line. An integral compatible splitting $B$, having constant term one, would be a unit splitting both laws and would make that class zero. Therefore every rational compatible splitting just constructed is nonintegral. Parts (a)--(c) give all the asserted denominator conclusions. The same conclusions hold for every nonzero integer multiple $B^k$: its class is $[k]$, which is nonzero since its characteristic-zero differential is $k$. For odd $p$ they also hold for $B^{k/2}$, $k\ne0$, since the normalized square roots of its integral cocycles are integral and the class has differential $k/2$. These assertions do not estimate selected harmonic combinations.
\end{proof}

The primitive-first-part hypothesis cannot be silently replaced by arbitrary nonvanishing: $f=p-z$ and $B=(1-z/p)^{-1}$ satisfy $fB=p$, although $B$ has unbounded denominators. This example is a warning about multiplication, not an asserted biextension example. Conversely, a trivial biextension can acquire a nonintegral rational bilinear gauge. DP35 is consequently a one-way obstruction, not a numerical invariant or a test characterizing the class. In particular, it does not bypass the source's selected harmonic-jet and connection identities.

\subsection{A conditional kernel estimate for the genuine dual pair}

\begin{proposition}[DP36: unequal weights restore the kernel quantization]\label{DP36}
Let $p$ be odd, $h\ge2$, $q=p^h$, $d=h-1$, and $R=W(k)$ with $\F_q\subset k$. Use the actual pair $G_h,H=G_h^D$ of DP32, with logarithms $L_0,\mathcal L$ and inverses $E_0,\mathcal E$. Put
\[
 s=(p-1)^{-1},\quad\rho=(q-1)^{-1},\quad r_0=\rho,\quad r_i=s-p^i\rho\ (1\le i\le d).
\]
Suppose $\beta(x,\mathbf y)\in R[1/p][[x,\mathbf y]]$ vanishes on either axis, has no bidegree-$(1,1)$ monomial, and satisfies
\[
 \beta(\zeta x,\zeta^{-p}y_1,\ldots,\zeta^{-p^d}y_d)=\beta(x,\mathbf y)
\]
for a primitive $(q-1)$-st root $\zeta$. Suppose its two exponential addition coboundaries
\[
 C_X=\exp\bigl(\delta_x\beta(L_0(S),L_0(S'),\mathcal L(\mathbf T))\bigr),\quad
 C_Y=\exp\bigl(\delta_{\mathbf y}\beta(L_0(S),\mathcal L(\mathbf T),\mathcal L(\mathbf T'))\bigr)
\]
are integral units with constant coefficient one. Here $\delta_{\mathbf y}$ uses addition of the whole vector in logarithm coordinates. For any integral series $U$ in the group coordinates and $\sigma\in\{1,-1\}$, set
\[
 \Theta(x,\mathbf y)=U(E_0(x),\mathcal E(\mathbf y))\exp\bigl((\sigma/2)\beta(x,\mathbf y)\bigr).
\]
Then
\begin{equation}\label{eq:true-pair-weight}
 v_p(\mathbf a![x^{a_0}\mathbf y^{\mathbf a'}]\Theta)
 \ge (s-p\rho)a_0+\sum_{i=1}^d\left(\frac{p^{i+1}}{q-1}-1\right)a_i
       -\frac{\sum_{i=0}^d s_p(a_i)}{p-1}.
\end{equation}
On the specific path $x=t$, $y_d=\epsilon t$, $y_i=0$ for $i<d$, with $\epsilon\in R^\times$, this gives
\begin{equation}\label{eq:true-pair-path}
 v_p(n![t^n]\Theta(t,0,\ldots,0,\epsilon t))
 \ge\rho n-\frac{s_p(n)}{p-1}.
\end{equation}
These are sufficient estimates, with no asserted optimality or identification with a prescribed geometric theta expression or path. DP37 below realizes the hypotheses for a normalized representative of the primitive formal Poincar\'e class.
\end{proposition}
\begin{proof}
If $x^a\mathbf y^{\mathbf b}$ is invariant, $a-\sum_i p^ib_i$ is divisible by $q-1$. Its weighted degree is
\[
 r_0a+\sum_i r_ib_i
 =\rho\left(a-\sum_i p^ib_i\right)+s|\mathbf b|\in s\Z.
\]
All $r_i$ are positive. Every nonconstant invariant monomial therefore has weighted degree at least $s$. The same calculation applies with duplicated $x$-variables or duplicated $\mathbf y$-vectors. Integral coefficient valuations are nonnegative, so the lower bound remains valid. The actual logarithms, exponentials and group laws are equivariant for these characters by their explicit formulas in DP32. Thus the integral series $C_X-1,C_Y-1$ have weighted valuation at least $s$.

The inverse-logarithm estimates give weight at least $r_0$ to $E_0$ and at least $r_i$ to $\mathcal E_i$; substituting them preserves this bound. Taking the formal logarithm consequently gives
\[
 w_{\mathbf r}(\delta_x\beta),\ w_{\mathbf r}(\delta_{\mathbf y}\beta)\ge s.
\]
Indeed its $m$-th term has weight at least $ms-v_p(m)\ge s$, since $v_p(m)\le(m-1)/(p-1)$.

For a monomial $x^a\mathbf y^{\mathbf b}$ of $\beta$, both side degrees are positive and not both one. If $a\ge2$, choose an interior binomial coefficient of valuation at most one as in DP29 and read the corresponding coefficient of $\delta_x\beta$. If $a=1$ and $|\mathbf b|\ge2$, either two components of $\mathbf b$ are nonzero, in which case a split sending one whole component to the first vector has coefficient one, or a single component has exponent at least two, in which case use the same binomial choice in that component. Recombining the split exponents recovers the original monomial uniquely, so there is no contamination by another coefficient. It follows that $w_{\mathbf r}(\beta)\ge s-1$.

Each monomial of $\beta$ has positive weighted degree at least $s$. Increasing all weights from $r_i$ to $pr_i$ therefore adds at least $(p-1)s=1$ to every coefficient weight. Thus $w_{p\mathbf r}(\beta)\ge s$, and its exponential with multiplier $\sigma/2$ has nonnegative weight by $v_p(m!)\le ms$. The integral $U$ factor has nonnegative weight as well; increasing the positive weights preserves the established inverse bounds. Hence the ordinary coefficient is at least $-p\sum_i r_ia_i$ in valuation. Adding Legendre's exact factorial identity proves \eqref{eq:true-pair-weight}.

For the stated path assign $t$ weight $pr_d$. Since $r_d\ge r_0$, both nonzero substitutions have at least their required weights. Setting the other coordinates to zero preserves the bound. The resulting factorial coefficient is bounded by $(s-pr_d)n-s_p(n)/(p-1)$. The identity $s-pr_d=\rho$ proves \eqref{eq:true-pair-path}.
\end{proof}

At $h=2$ all weights in this pair are equal, and the final residual is the source's $\rho_2$. At larger height, the actual dual's unequal weights restore the quantization that failed in DP28's two-identical-scalar model. On the displayed path the sufficient residual is instead $\rho_h$, which decreases with $h$. This is not an upper bound on actual gain and does not rule out a sharper primitive-class estimate or a different path. It does show why replacing the increasing formal-function slope $c_h$ directly in a normalized-kernel formula is not the analytic implication just proved.

\begin{proposition}[DP37: the true dual torus and a primitive normalized kernel]\label{DP37}
In the actual model of DP32 over $\mathcal O=W(\F_q)$, the automorphism
\[
 A_\zeta(\mathbf X)=(\zeta^pX_1,\ldots,\zeta^{p^{h-1}}X_{h-1}),\qquad\zeta^{q-1}=1,
\]
satisfies $A_\zeta^D=[\zeta]_{G_h}$. Consequently the primitive formal Poincar\'e biextension has an integral normalized frame invariant under
\[
 ([\zeta],A_\zeta^{-1}).
\]
In that frame there is a unique invariant rational logarithmic splitting $\beta$. It vanishes on both axes, has no bidegree-$(1,1)$ term, and has the integral exponential coboundaries required in DP36. Thus for odd $p$ DP36 applies to an actual primitive formal Poincar\'e representative. Its compatible splitting still has the denominators proved in DP35. This does not identify the particular section, gauge, path or harmonic operators of a global theta correspondence.
\end{proposition}
\begin{proof}
The explicit monomial map $M$ of DP32 commutes with $A_\zeta$: the ordinary arrows use $p\cdot p^{i-1}=p^i$, and the exceptional arrow uses $p^{h+1}\equiv p\pmod{q-1}$. Hence its logarithm and integral group law commute as well. This proves that $A_\zeta$ is an integral automorphism. It suffices first to take $\zeta$ primitive. Put $a=A_\zeta^D$ and write $\lambda=da$ on the one-dimensional tangent space of $G_h$.

Tate's functorial Hodge--Tate sequence \cite[\S4, Corollary 2 following Theorem 3, p.~180]{Tate67}, over the completed algebraic closure $C$ of $\operatorname{Frac}(\mathcal O)$, is
\[
 0\longrightarrow \Lie(H)_C
 \longrightarrow \Hom_{\Z_p}(T_pG_h,C)
 \longrightarrow \Lie(G_h)_C^*\otimes_C\Hom_{\Z_p}(\Z_p(1),C)
 \longrightarrow0.
\]
On its middle term $a$ acts by precomposition, on the first term by $d(a^D)=dA_\zeta$, and on the cotangent factor by pullback, with scalar $\lambda$. These are adjoint actions, so neither Hom nor the cotangent introduces an inverse eigenvalue. This also follows directly by differentiating the natural Cartier identity $\langle aP,Q\rangle=\langle P,a^DQ\rangle$. The cyclotomic factor is unchanged.

The characteristic polynomial $P_a(Z)$ of $a$ on the rank-$h$ $\Z_p$ Tate module is consequently
\[
 P_a(Z)=(Z-\lambda)\prod_{i=1}^{h-1}(Z-\zeta^{p^i}).
\]
It is monic of degree $h$ over $\Q_p$. The primitive root $\zeta^p$ has degree exactly $h$: its Frobenius conjugates are $\zeta^{p^j}$, and $p^m\not\equiv1\pmod{p^h-1}$ for $0<m<h$. Therefore its degree-$h$ minimal polynomial divides and equals $P_a$. Comparing the one missing root proves $\lambda=\zeta$. In characteristic zero a one-dimensional formal endomorphism is determined by its differential: conjugate by the normalized logarithm and compare the coefficient of $X^{n-1}Y$ in additivity to kill all degrees $n\ge2$. Thus $a=[\zeta]$. Taking powers proves the equality for every root, and its integral identity persists under base change.

The evaluation pairing giving the primitive Poincar\'e biextension is preserved by $([\zeta],A_\zeta^{-1})$. Functoriality of the pairing construction, not just equality of isomorphism classes, gives a compatible action of the finite group $\Gamma=\mu_{q-1}$ on the biextension and its axial rigidifications. Choose an integral normalized frame $e$ as in DP35 and write $g(e)=c_g e$. The units $c_g$ are one on both axes and satisfy $c_{\sigma g}=\sigma(c_g)c_\sigma$. Set $N=q-1$ and
\[
 u=\left(\prod_{g\in\Gamma}c_g\right)^{1/N}.
\]
The unique root with constant term one is integral: coefficientwise recursion divides only by the unit $N$. It is one on both axes. The cocycle identity gives
\[
 \sigma(u)=u/c_\sigma,
\]
so $e'=ue$ is an invariant integral normalized frame. Its two addition cocycles are invariant integral units.

Use DP35 to obtain a rational logarithmic splitting $\beta_0$ of these cocycles in the additive logarithm coordinates, and average
\[
 \beta=\frac1N\sum_{g\in\Gamma}g\beta_0.
\]
Both additive coboundaries remain the logarithms of the same invariant integral cocycles. The axes remain zero. The characters are exactly $(\zeta,\zeta^{-p},\ldots,\zeta^{-p^{h-1}})$ by the first part of the proof. No bilinear monomial $xy_i$ is invariant, because $1-p^i$ is not divisible by $q-1$ for $1\le i<h$. Thus $\beta$ has no bidegree-$(1,1)$ term. Any two rational splittings differ by a bilinear form by DP35, so the same character observation proves uniqueness among invariant splittings. All hypotheses of DP36 have now been obtained for the primitive class. Its nontriviality and denominator conclusion are unchanged.
\end{proof}

\begin{theorem}[DP38: the first primitive Poincar\'e coefficients]\label{DP38}
For every $p$ and $h\ge2$, over $W(\F_q)$ there are integral Teichm\"uller coordinates $S,Y_1,\ldots,Y_d$ on the Honda line and its actual dual, $d=h-1$, and a normalized integral Poincar\'e frame in which every rational compatible splitting satisfies
\[
 [S^pY_1]B=-\frac1p,\qquad [SY_d^p]B=\frac1p.
\]
These coefficients are unchanged by a rational bilinear gauge and change only by integral amounts under an axes-normalized integral unit frame. The frame is invariant for the true action $(S,Y_i)\mapsto(\zeta S,\zeta^{-p^i}Y_i)$. Its unique invariant logarithmic splitting has no positive term below degree $p+1$, and
\begin{equation}\label{eq:primitive-first}
 (\log B)_{p+1}=B_{p+1}=\frac{-S^pY_1+SY_d^p}{p}.
\end{equation}
The same first homogeneous term holds in the normalized logarithm coordinates. Consequently
\[
 D_G^pD_1 B(0)=-(p-1)!,\qquad D_GD_d^p B(0)=(p-1)!.
\]
Both values are units, although both operators are divisible by $p$ on integral formal functions. In the group coordinates,
\[
 w_t(B)=-\infty\quad\left(0<t<\frac1{p+1},\ t\in\Q\right),\qquad
 \liminf_{|\alpha|\to\infty}\frac{v_p(B_\alpha)}{|\alpha|}\le-\frac1{p+1}.
\]
For this invariant splitting, odd $p$, $h>2$, and any integral $U$ with unit constant term, the last-dual logarithmic path $x=t$, $y_d=\epsilon t$, $y_i=0$ for $i<d$, gives
\begin{equation}\label{eq:primitive-path-unit}
 v_p\left((p+1)![t^{p+1}]\,
 U(E_G(t),E_H(\epsilon te_d))
 \exp\left(\tfrac12\log B(E_G(t),E_H(\epsilon te_d))\right)\right)=0
\end{equation}
for every unit $\epsilon$. This is an exact finite-order assertion, not an asymptotic upper bound on every selected theta jet or a shorter-chain impossibility theorem.
\end{theorem}
\begin{proof}
We give an explicit display calculation, working compatibly over $\mathcal O/p^n$, where $\mathcal O=W(\F_q)$, and then passing to the inverse limit. Let $P$ be free on $e_0,\ldots,e_d$ over the Witt ring, with normal decomposition
\[
 T=We_0,\qquad L=\bigoplus_{j=1}^dWe_j,\qquad Q=L\oplus IT.
\]
With $e_h=e_0$, specify the Frobenius-semilinear operators by
\[
 Fe_0=e_1,\qquad V^{-1}e_j=e_{j+1},\qquad Fe_j=pe_{j+1}\quad(1\le j\le d).
\]
On $IT$ the relation is $V^{-1}(V\xi\,e_0)=\xi Fe_0$. The linearization of $V^{-1}\oplus F$ on $L\oplus T$ is a cyclic permutation and hence an isomorphism. The actual dual display has
\[
 L'=We_0^*,\quad T'=\bigoplus_{j=1}^dWe_j^*,\quad
 V'^{-1}e_0^*=e_1^*,\quad F'e_0^*=pe_1^*,\quad
 F'e_j^*=e_{j+1}^*\ (j\ge1).
\]
These formulas satisfy the pairing identities of Zink, Definition~19 and (22), p.~143. The $h$-fold Frobenius cycles send a basis vector to $p^{h-1}$ times itself in the first display and to $p$ times itself in the dual, with the corresponding Frobenius on coefficients. Both nilpotence conditions hold for $h\ge2$.

By Zink's Theorem~81 and (145)--(146), p.~204, the formal groups are the cokernels of $V^{-1}-\id$. The maps represented by
\[
 \sigma(S)=[S]e_0,\qquad \tau(\mathbf Y)=\sum_{j=1}^d[Y_j]e_j^*
\]
give integral formal coordinates: their tangent maps to $P/Q$ and $P'/Q'$ are identities, so the formal inverse function theorem, proved coefficientwise by its unit linear term, applies. On the special fibre $p[S]=V[S^p]$, and the display relations carry this around the cycle to $[S^q]e_0$. Thus $[p]_G(S)\bmod p=S^q$. The linear coefficient is $p$, so Hazewinkel's integral Lubin--Tate uniqueness theorem, \S5, p.~77 identifies this line with $G_h$. The second display is its actual Cartier dual. No equality of its entire coordinate logarithm with DP32's chosen formula is assumed in the calculation that follows.

Use the evaluation pairing of the displays. The scalar-one section represented by $(1,\sigma(S),\tau(\mathbf Y))$ in Zink's Proposition~106, p.~226 gives a normalized integral frame. If $\kappa_G,\kappa_H$ are the defects of these lifts, defined by
\[
 (V^{-1}-\id)\kappa_G=\sigma(S)+\sigma(S')-\sigma(F_G(S,S'))
\]
and its dual analogue, its cocycles from (218), p.~239 are
\begin{equation}\label{eq:display-cocycles}
 C_G=\operatorname{hex}\langle V^{-1}\kappa_G,\tau(\mathbf Y)\rangle,
 \qquad C_H=\operatorname{hex}\langle\sigma(S),\kappa_H\rangle^{-1}.
\end{equation}
Here $\operatorname{hex}$ is Zink's Artin--Hasse map (196), p.~232, and (210)--(211), p.~236. His convention sends $[X]$ to $\exp(-\sum_{n\ge0}X^{p^n}/p^n)$, with linear coefficient $-1$. The inverse is the minus sign in Zink's additively written multiplicative group. Indeed, the quotient relation in Proposition~106 replaces a lift differing by a defect with its scalar multiplied by the corresponding pairing. Thus these signs agree with $C_G=B(F_G)/(BB)$, and similarly for $C_H$. DP35 supplies rational compatible splittings.

Put
\[
 C_p(S,S')=\frac{(S+S')^p-S^p-S'^p}{p},\qquad
 [S]+[S']-[S+S']=V\xi.
\]
The first two ghost polynomials give $\xi_0=-C_p(S,S')$. The integral Witt addition polynomials are weighted homogeneous, so $\xi_i$ has ordinary degree $p^{i+1}$, and every component of $F^r\xi$ has degree at least $p^{r+1}$. This assertion uses this homogeneous carry vector; it is not an assertion about Frobenius and arbitrary coefficient-degree filtrations. The explicit element
\[
 \kappa_0=-V\xi\,e_0-\sum_{j=1}^dF^{j-1}\xi\,e_j
\]
satisfies the exact telescoping identity
\[
 (V^{-1}-\id)\kappa_0=V\xi\,e_0-F^d\xi\,e_0.
\]
The last term has degree at least $p^{d+1}=q$. Modulo the degree-$q$ ideal the group law is therefore ordinary addition. Moreover the true defect equals $\kappa_0$ to that precision: reduce Zink's exact sequence (146) to this nilpotent test algebra and use its injectivity. No inverse of an uncomputed operator is assumed to preserve degree.

The $e_1$ coefficient of $V^{-1}\kappa_0$ is the Witt vector $w=-\xi$. Multiplication by $[Y_1]$ sends its $i$th Witt component to $w_iY_1^{p^i}$; no componentwise-negation formula is assumed. In the Artin--Hasse product only the zeroth component contributes linearly in $Y_1$, with coefficient $-1$ in this convention, and $w_0=-\xi_0$. Thus
\[
 [Y_1]\log C_G=-C_p(S,S')+O((S,S')^q).
\]
Writing $b=\log B$, extraction of $S^{p-1}S'Y_1$ from $\delta_G b$ uses ordinary addition and gives $p[S^pY_1]b=-1$. Since $b$ has positive degree on both sides, exponentiation cannot alter its $Y$-linear coefficient. This proves $[S^pY_1]B=-1/p$.

For the second coefficient restrict both dual inputs to the last coordinate. Again write $[Y]+[Y']-[Y+Y']=V\xi$. This time
\[
 \kappa'_0=-V\xi\,e_d^*-\xi\,e_0^*,\qquad
 (V'^{-1}-\id)\kappa'_0=V\xi\,e_d^*-F\xi\,e_1^*.
\]
The endpoint has degree at least $p^2$. Addition on this source axis is ordinary to that precision, and the true defect has the same reduction by injectivity. Pair with $[S]e_0$ and use the inverse in \eqref{eq:display-cocycles}. It gives
\[
 [S]\log C_H=C_p(Y,Y')+O((Y,Y')^{p^2}).
\]
Hence $p[SY_d^p]b=1$, and exponentiation leaves this $S$-linear coefficient unchanged.

These same addition calculations imply $\log_G(S)=S+O(S^q)$ and
\[
 \log_H(Ye_d)=Ye_d+O(Y^{p^2}):
\]
compare the low-degree terms of the logarithm's additive identity on the indicated inputs. Therefore a rational bilinear change $\log_G(S)\sum_j a_j\log_{H,j}(\mathbf Y)$ contributes to neither of the two coefficients. An axes-normalized unit frame change adds or subtracts $\log I$. Its $Y$-linear and $S$-linear coefficients are integral, since higher powers of $I-1$ cannot contribute there. This proves the claimed invariance of their valuations.

For the full first homogeneous part, the diagonal display map $e_j\mapsto[\zeta^{p^j}]e_j$ commutes with the operators, including Frobenius on its coefficients. It sends $\sigma(S)$ to $\sigma(\zeta S)$, so gives the true $[\zeta]$ on the line. Its inverse dual sends $\tau(\mathbf Y)$ to $\tau(\zeta^{-p^j}Y_j)$. The evaluation pairing and the scalar-one quotient frame are invariant. As in DP37, averaging gives the unique invariant rational $b$.

Suppose an invariant mixed monomial has side degrees $a,b\ge1$ and total degree at most $p+1$. Then
\[
 a-\sum_jp^jb_j\equiv0\pmod{q-1},\qquad
 pb\le\sum_jp^jb_j\le p^db\le q.
\]
If the sum equals $a$, the degree bound forces $b=1$, $a=p$, and $j=1$. If the sum exceeds $a$, their difference can only be $q-1$; the upper bound forces $a=1$, $b=p$, and all dual factors to be $Y_d$. These are precisely the two monomials already computed, both of degree $p+1$. This proves \eqref{eq:primitive-first} and absence of smaller terms. Substitution of the inverse logarithms, whose linear parts are identities, preserves that first homogeneous part. The two stated factorial mixed derivatives follow.

For an integral function, the translation coefficient at source degree $p$ on the line is $D_G^p/p!$, since its logarithm has no nonlinear term of that degree. The corresponding last-dual coefficient is $D_d^p/p!$ by the last-axis calculation above. Both translation operators are integral, so $D_G^p$ and $D_d^p$ are divisible by $p$ on the entire integral function ring. This proves the asserted contrast with the two unit mixed jets. DP35 applied to a degree-$p+1$ coefficient of valuation $-1$ gives the numerical denominator rate.

Finally assume $h>2$. On the stated logarithmic path the first monomial in \eqref{eq:primitive-first} vanishes because $y_1=0$, leaving $\epsilon^pt^{p+1}/p$. Put $D=D_G+\epsilon D_d$. Its two pure $p$th powers and every interior binomial term are divisible by $p$, so $D^p$ takes integral functions into $p\mathcal O[[S,\mathbf Y]]$. Consequently the integral $U$ factor contributes a multiple of $p$ to its $(p+1)$st derivative. There is no earlier positive term in the exponential gauge; its contribution at that order is
\[
 U(0)\epsilon^p\frac{(p+1)!}{2p},
\]
a unit for odd $p$. It cannot cancel the first contribution. This proves \eqref{eq:primitive-path-unit}.
\end{proof}

At $h=2$ the path coefficient is $(\epsilon^p-\epsilon)/p$ and can vanish, which explains the restriction in \eqref{eq:primitive-path-unit}. For any prime, still with $h>2$, the same unit-jet proof works with $\exp(k\log B)$ for a unit scalar $k$ in place of the half-power. The asymptotic denominator rate also survives on a formal generic group-coordinate diagonal after an unramified extension with transcendental residue unit $\epsilon$: first restrict to $S,Y_1$, which retains the integral first coboundary and the coefficient $-1/p$, then put $S=t,Y_1=\epsilon t$. In each nonzero homogeneous part the least-valuation residue is a nonzero polynomial in $\bar\epsilon$, so no cancellation of that minimum occurs. This proves the same rate there, but makes no claim for a prescribed algebraic CM parameter or for an arbitrary logarithmic diagonal.

\begin{proposition}[DP39: what harmonicity removes, and what it does not]\label{DP39}
The following two assertions distinguish full-series denominators from selected theta-type jets.
\begin{enumerate}[leftmargin=*,label=(\alph*)]
\item There are integral formal-biextension data, two canonical gauges, and a primitive germ $f=1$ for which one normalized series has unbounded denominators, while all the specified harmonic values agree. In fact the agreement holds for every germ $f$ whenever the polynomial $P$ is killed by the specified second-order contraction.
\item For odd $p$ and every $h\ge2$, two copies of the actual primitive pair have an integral polynomial satisfying the summed harmonic, row-invariance and fixed-row-degree conditions defining \path{HaO}, whose normalized kernel value is a unit. The same invariant polynomial takes values in $pR$ on every integral function on that product. Therefore those algebraic harmonic conditions alone do not transfer the bare integral-function divisibility through the primitive exponential.
\end{enumerate}
Neither assertion constructs a new global mixed-type theta correspondence or asserts that the particular germ $1$ is a global theta section.
\end{proposition}
\begin{proof}
For a polynomial $P$ and a formal series $f$ over a characteristic-zero field, write $\langle P,f\rangle=P(\partial)f(0)$. For every polynomial $Q$,
\begin{equation}\label{eq:fischer-ideal}
 \langle P,Qf\rangle=\langle Q(\partial)P,f\rangle.
\end{equation}
For monomials $P=Z^\alpha,Q=Z^\beta,f=Z^\gamma$, both sides vanish unless $\alpha=\beta+\gamma$, when both are $\alpha!$. Linearity and the finite jet read by $P$ prove the identity for formal series. If $Q_j(\partial)P=0$ for every $j$, the pairing kills the formal ideal $(Q_j)$. If $\eta$ has zero constant term and belongs to that ideal, then $\exp\eta-1=\eta\sum_{n\ge1}\eta^{n-1}/n!$ also belongs to it. Hence
\[
 \langle P,f\exp\eta\rangle=\langle P,f\rangle.
\]

For (a), let $F=\widehat\Ga^2$ over $\Z_p$, let $O=\Z\times\Z$ act diagonally, and let its involution exchange the two factors. Take trivial addition, symmetry and adjoint cocycles. Both $h_0=0$ and
\[
 h_1(X,Y)=X_1Y_2+X_2Y_1
\]
are canonical gauges: their axes vanish, their addition coboundaries vanish by biadditivity, and symmetry is immediate. For $(a,b)\in O$ the adjoint identity is explicitly
\[
 h_1((aX_1,bX_2),Y)=aX_1Y_2+bX_2Y_1
 =h_1(X,(bY_1,aY_2)).
\]
Take $w_0=(1,0)$, so $w_0+w_0^*=1$ and $h_{1,L}=X_1X_2$. For $f=1$ the normalized series is $\exp(X_1X_2)$, whose coefficients $1/n!$ have unbounded denominators. If $\partial_U\partial_VP=0$, its selected value agrees with that in gauge zero for every $f$, by \eqref{eq:fischer-ideal} with $Q=X_1X_2$. Equivalently $P$ is a sum of pure polynomials in $U$ and in $V$, and the exponential is one on each axis. The second-order operator in the source's symmetric-matrix convention is $2\partial_U\partial_V$, with the same kernel over $\Q_p$, including when $p=2$. For $f=1$ all positive-degree selected values vanish. This proves (a) and shows why primitive initial terms and full-series nonintegrality do not by themselves decide harmonic integrality.

For (b), use DP38's invariant splitting in logarithm coordinates, whose coefficient is $[x^py_1]B=-1/p$. On two copies put
\[
 \Theta=B(x_1,\mathbf y^{(1)})^{1/2}B(x_2,\mathbf y^{(2)})^{1/2},\qquad
 P=X_1^pY_{1,1}-pX_1^{p-1}X_2Y_{1,2}.
\]
Give the positive type one row and the negative type $d$ rows, each with two columns. Give the positive row degree $p$, the first negative row degree one, and all other rows degree zero. Every coefficient is integral and every monomial has those degrees. The positive row has no earlier row. Adding an earlier negative row to a later one leaves $P$ unchanged because it uses only the first negative row. Thus the prescribed row invariances hold. For the identity-form contractions
\[
 \Delta_i=\partial_{X_1}\partial_{Y_{i,1}}+
          \partial_{X_2}\partial_{Y_{i,2}},
\]
one has $\Delta_iP=0$ for $i>1$ and
\[
 \Delta_1P=pX_1^{p-1}-pX_1^{p-1}=0.
\]
These are exactly the algebraic definitions \path{lap}, \path{harmK}, \path{RowInv}, \path{RowDeg}, and \path{HaO} in
\path{Formalism/LocalToGlobal/ProjectiveLattices/RowDeg.lean},
lines 21--30 and 37--48; \path{IsThRule.lean}, \path{harmA}, lines 86--87, specializes this definition in the source. They are sums over columns, not individual contractions.

Every monomial of $B-1$ has positive degree on both sides. Therefore the $y_1$-linear coefficient of $B^{1/2}$ is half that of $B$, and the first monomial of $P$ contributes
\[
 p!\,[x^py_1]B^{1/2}=-\frac{(p-1)!}{2},
\]
a unit. The second monomial contributes zero: it differentiates the first copy $p-1>0$ times in $x_1$, with its entire dual variable set zero, where that factor is identically one. Hence $P(\partial)\Theta(0)=-(p-1)!/2$.

For an integral function $f$ on $(G_h\times H)^2$, the first term of $P(D)f$ is divisible by $p$ because $D_{x_1}^p$ is coefficientwise divisible by $p$ at height at least two. Its other invariant derivation preserves that ideal. The second term has its explicit factor $p$ and otherwise only integral invariant derivations. Thus $P(D)f(0)\in pR$. This proves (b), with no assumption that the formal row labels already define a new global source datum.
\end{proof}

The associated ideal in \eqref{eq:fischer-ideal} is generated by column sums $\sum_bX_{\tau,a,b}X_{\bar\tau,a',b}$, not every separate product. In DP22's pure CM type $\{1,2,3\}\subset(\Z/7)^\times$, the conjugate type is $\{6,5,4\}$ and is disjoint. If this formula is applied with one row on the type and zero outside it, each conjugate-row contraction has an empty index set. It does not impose $\partial_x\partial_{y_i}$ on the height-three line and its dual directions. This observation is about the stated harmonic formula; DP25 still prevents treating that row list as the old induced-type global host.

\subsection{General cyclic CM types and the sum of sharp directional slopes}

\begin{theorem}[DP40: cyclic displays and the alternating maximum]\label{DP40}
Let $h\ge2$, $q=p^h$, $s=1/(p-1)$, and let $A$ be a nonempty proper subset of $\Z/h\Z$, of cardinality $m$. Over $R=W(k)$ take a Witt module with cyclic basis $e_i$, normal decomposition
\[
 T=\bigoplus_{i\in A}W(R)e_i,\qquad
 L=\bigoplus_{i\notin A}W(R)e_i,\qquad Q=L\oplus IT,
\]
and arbitrary units $u_i\in W(R)^\times$. Define
\[
 Fe_i=u_ip^{\epsilon_i}e_{i+1},\qquad
 \epsilon_i=1-1_A(i),\qquad
 V^{-1}e_i=u_ie_{i+1}\quad(i\notin A),
\]
with the defining display rule on $IT$. This gives a formal $p$-divisible group $H_{A,u}$ of height $h$ and dimension $m$. Its actual Cartier dual is the complementary display $H_{A^c,u^{-1}}$.

For $j\in A$, put
\[
 r_j=\frac{\sum_{i\in A}p^{[j-i]_h}}{q-1},\qquad
 c_j=s-r_j=\frac{\sum_{i\notin A}p^{[j-i]_h}}{q-1},
\]
where $[a]_h$ is the least nonnegative residue. In the Teichm\"uller coordinates $X_j$, the normalized integral invariant directions satisfy, coefficientwise,
\begin{equation}\label{eq:cyclicmixed}
 \vp\left(\prod_{j\in A}D_j^{n_j}f\right)
 \ge\max\left\{0,\sum_{j\in A}c_jn_j-
                   \frac{\sum_{j\in A}s_p(n_j)}{p-1}\right\}.
\end{equation}
Each slope is sharp, even with the arbitrary unit arrows:
\[
 \vp(D_j^{q^r}X_j(0))=c_j(q^r-1)\qquad(r\ge0).
\]
For unit arrows equal to one the exact logarithm is
\[
 L_A=\sum_{n\ge0}\frac{M^{\circ n}}{p^n},\qquad
 M_j=X_{\operatorname{prev}(j)}^{p^{g_j}},
\]
where $g_j$ is the positive cyclic gap from the preceding selected index to $j$.

Across the group and its actual dual define
\[
 \gamma_i(A)=\frac{\sum_{j:\,1_A(j)\ne1_A(i)}p^{[i-j]_h}}{q-1},
 \qquad N_k=|A\mathbin\triangle(A+k)|,\qquad
 S_h(A)=\sum_i\gamma_i(A).
\]
Then
\begin{equation}\label{eq:cyclicmax}
 S_h(A)=\frac{\sum_{k=1}^{h-1}N_kp^k}{q-1}
 \le h c_2,\qquad
 \sum_{i\in A}\gamma_i=\sum_{i\notin A}\gamma_i=\frac{S_h(A)}2.
\end{equation}
Equality in the first inequality holds exactly when $h$ is even and $A$ is alternating. More precisely,
\[
 S_h(A)\le \frac{4m(h-m)}h\,c_2.
\]
For nonconstant types equality in this last inequality again requires an even alternating type. These are equally weighted sums, not componentwise bounds or global theta costs.
\end{theorem}

\begin{proof}
The linearized normal map $V^{-1}\oplus F:L\oplus T\to P$ is a monomial matrix of units, hence invertible. Complete cycles accumulate $h-m>0$ powers of $p$ under $F$, and the complementary cycle accumulates $m>0$. The corresponding nilpotence conditions hold over each $R/p^n$ and are compatible under reduction. Zink's normal-decomposition construction, p.~129, and Theorem~5, p.~131, give a formal $p$-divisible group of height $\operatorname{rank}P=h$; its tangent module $P/Q$ has rank $m$. On the dual basis the evaluation pairing forces
\[
 F'e_i^*=u_i^{-1}p^{1-\epsilon_i}e_{i+1}^*,\qquad
 T'=L^*,\quad L'=T^*.
\]
Indeed $\langle Fe_i,F'e_j^*\rangle=p\delta_{ij}$, and the other identities follow on the normal generators. Both displays are connected. Zink's Definition~19/(22), Proposition~121 and Theorem~122 identify this complementary display as the actual Cartier dual \cite{Zink}.

The section $\sum_{i\in A}[X_i]e_i$ has identity tangent map to $P/Q$, so gives integral formal coordinates. Write $z_n=w_n(z)/p^n$ for divided Witt ghosts. With unit arrows one, let $B$ denote the Frobenius matrix. The resolution operator on these ghosts is $B\operatorname{shift}-1$, and its logarithm is
\[
 \operatorname{pr}_T\sum_{n\ge0}B^nz_n.
\]
It vanishes on the relation module by telescoping and has identity differential. For arbitrary units, let $B_n$ have nonzero entries $w_n(u_i)p^{\epsilon_i}$. Each $w_n(u_i)$ is a unit, and the operator and logarithm become
\[
 z_n\longmapsto B_nz_{n+1}-z_n,\qquad
 \operatorname{pr}_T\left(z_0+\sum_{n\ge1}B_0\cdots B_{n-1}z_n\right).
\]
Here is an explicit completion justification. On a degree truncation with nilpotent ideal $N^D=0$, put $C=\floor{\log_p(D-1)}$ if $D>1$. Integral Witt vectors satisfy
\[
 z_n=\sum_{r=0}^{\min(n,C)}
       \frac{a_{n-r}^{p^r}}{p^r};
\]
thus their divided ghosts have uniformly bounded denominators. Every complete matrix cycle contributes $h-m$ powers of $p$ times units. The sums therefore converge in the completed integral resolution, and their telescoping boundaries tend to zero. The particular Teichm\"uller sections are finite on every degree truncation. This is a rational calculation with the already constructed integral quotient, not an application of a display theorem over a field with $p$ inverted.

From one selected index to the next, a path of length $g_j$ has exactly one Frobenius arrow without a factor $p$. Its Teichm\"uller ghost denominator is $p^{g_j}$, leaving precisely a unit divided by $p$. A path through $a$ selected gaps has coefficient $\eta/p^a$, with $\eta\in R^\times$. Each starting index and ordinary degree specifies a unique path, so there is no cancellation. With all units one, these paths give exactly the asserted formula for $L_A$.

The weights satisfy
\begin{equation}\label{eq:cyclicedge}
 p^{g_j}r_{\operatorname{prev}(j)}-r_j=1.
\end{equation}
Indeed multiplication of the defining numerator by $p^{g_j}$ advances each cyclic distance to $j$; only the term belonging to $j$ wraps, changing $1$ to $q$. Adjoin $\pi^{q-1}=p$ and scale coordinate $j$ by $\alpha_j=\pi^{(q-1)r_j}$. A path of one selected gap then has valuation zero. A path of $a\ge2$ selected gaps has scaled valuation
\[
 1+p^{g_1}+p^{g_1+g_2}+\cdots+
       p^{g_1+\cdots+g_{a-1}}-a>0,
\]
with gaps ordered backwards from the output. Thus the contracted logarithm is integral, has identity differential, and reduces to
\[
 \overline{L'}=\id+M_\eta,\qquad
 (M_\eta)_j=\bar\eta_jX_{\operatorname{prev}(j)}^{p^{g_j}},
 \quad\bar\eta_j\ne0.
\]
Formal inversion is integral. Since $M_\eta$ is additive in characteristic $p$ and raises ordinary degree, its inverse reduction is
\[
 \overline{E'}=\sum_{a\ge0}(-1)^aM_\eta^{\circ a}.
\]
In component $j$, with only input coordinate $j$ retained, degree $q^r$ comes from exactly $mr$ arrows and has a nonzero coefficient. Hence
\[
 \vp([t^{q^r}]E_j(te_j))=-(q^r-1)r_j.
\]
For any integral $f$, contraction bounds the coefficient of $\mathbf t^{\mathbf n}$ in $f(E(\mathbf t))$ below by $-\sum r_jn_j$. Multiplication by $\prod n_j!$ and Legendre's formula give \eqref{eq:cyclicmixed} at zero. Apply it to the integral translated function to obtain the coefficientwise statement. For $f=X_j$, $n=q^r$, the exact coefficient gives
\[
 \vp(D_j^nX_j(0))=\vp(n!)-(n-1)r_j=(n-1)c_j.
\]
This proves sharpness and works also at $p=2$.

It remains to prove the sum comparison. At displacement $k$ there are $N_k$ pairs of opposite type. The number with output in $A$ is
$|A|-|A\cap(A+k)|$; the same formula counts those with output in $A^c$. Each is $N_k/2$, proving both sum identities. Also $N_k=N_{h-k}$.

Set $x_i=1$ on $A$ and $x_i=-1$ otherwise. For the circulant operator
\[
 (Wx)_i=(q-1)^{-1}\sum_{k=1}^{h-1}p^kx_{i-k}
\]
we have
\[
 S_h(A)=\frac12(hc_h-x^TWx).
\]
At the discrete frequencies $\theta_\ell=2\pi\ell/h$, the symmetric part of $W$ has eigenvalues
\[
 \lambda_\ell=\frac{\sum_{k=1}^{h-1}p^k\cos(k\theta_\ell)}{q-1}
 =\frac{p\cos\theta_\ell-1}{p^2-2p\cos\theta_\ell+1}
       -\frac1{q-1}.
\]
The rational expression follows from the finite geometric sum using
$e^{ih\theta_\ell}=1$; it is not asserted for arbitrary real frequencies.
As a function of $t=\cos\theta_\ell$ its derivative is
$p(p^2-1)/(p^2-2pt+1)^2>0$. Consequently
\[
 \lambda_\ell\ge\lambda_*=-\frac1{p+1}-\frac1{q-1},
\]
with equality only at $\theta_\ell=\pi$. For normalized Fourier coefficients
$\widehat x_\ell=h^{-1/2}\sum_jx_je^{-2\pi i\ell j/h}$, Parseval gives
$\sum|\widehat x_\ell|^2=h$. Hence
\[
 S_h(A)\le\frac h2(c_h-\lambda_*)=hc_2.
\]
Equality forces Fourier support only at $h/2$, so $h$ is even and
$x_j=\pm(-1)^j$. Conversely those two words give equality. Explicitly,
\[
 hc_2-S_h(A)=\frac{p(p-1)}{2(p+1)}
 \sum_{\ell=0}^{h-1}
 \frac{1+\cos\theta_\ell}{p^2-2p\cos\theta_\ell+1}
 |\widehat x_\ell|^2.
\]
Finally retain the zeroth eigenvalue $\lambda_0=c_h$ and
$|\widehat x_0|^2=(2m-h)^2/h$ instead of bounding it by $\lambda_*$. This gives
$S_h(A)\le hc_2-c_2(2m-h)^2/h=4m(h-m)c_2/h$.
Equality permits only modes $0,h/2$; a nonconstant $\{\pm1\}$ word of that form is alternating. Constant types, if included, have sum zero and equality in this cardinality bound.
\end{proof}

The display hypotheses have a useful exact CM interface. Suppose an integral action of an unramified degree-$h$ order splits the display into free rank-one eigenspaces, Frobenius cycles those eigenspaces, and the Hodge filtration has $Q_i=IP_i$ on $A$ and $Q_i=P_i$ off $A$. Choosing generators, commutation with the action forces the arrow positions above, and invertibility of the normal map forces their coefficients to be units. Thus DP40 applies. This argument does not infer those integral hypotheses from height alone.

For $h=6$ and $A=\{0,1,2\}$, direct counting gives
\[
 (N_1,\ldots,N_5)=(2,4,6,4,2),\qquad
 (\gamma_0,\gamma_1,\gamma_2)
 =(\gamma_3,\gamma_4,\gamma_5)
 =\frac{(p,p^2,p^3)}{(p-1)(p^3+1)}.
\]
Therefore
\[
 S_6(A)=6c_2-\frac{4p(p-1)}{p^3+1},\qquad
 \frac{S_6(A)}{6c_2}=\frac{p^2+p+1}{3(p^2-p+1)}<1.
\]
Nevertheless $\gamma_2-c_2=p/(p^3+1)>0$, which explicitly refutes a componentwise interpretation. DP22's CM type $\{1,2,3\}$ for $y^2=x^7-1$ gives this word on the Frobenius orbits
$(1,3,2,6,4,5)$ when $p\equiv3\bmod7$ and
$(2,3,1,5,4,6)$ when $p\equiv5\bmod7$.
The finite check of all binary words through $h=10$ at $p=2,3,5,11$ verifies 8,184 cases of the sum identities, inequalities and equality classifications; it supplements the proof above.


\subsection{The full primitive logarithm and its exact exponential threshold}

\begin{theorem}[DP41: optimal multivariate inflation for the primitive class]\label{DP41}
Let $p$ be odd, $h\ge2$, $d=h-1$, $q=p^h$, and $\mathcal O=W(\F_q)$. For the actual Honda line and its Cartier dual use the Teichm\"uller coordinates and invariant primitive frame of DP38. Put
\[
 s=\frac1{p-1},\qquad \rho=\frac1{q-1},\qquad
 r_0=\rho,\qquad r_j=s-p^j\rho\quad(1\le j\le d).
\]
Let $b(S,\mathbf Y)$ be the invariant rational logarithmic splitting and
$\beta(x,\mathbf y)=b(E_0(x),E_H(\mathbf y))$ its expression in additive coordinates. For $\lambda\in\mathcal O^\times$ and every rational $a\ge0$,
\[
 w_{a\mathbf r}\bigl(\exp(\lambda\beta)\bigr)=
 \begin{cases}-\infty,&a<p,\\0,&a\ge p.\end{cases}
\]
The same formula holds for $\exp(\lambda b)$ in group coordinates. Multiplication by any nonzero integral $U(E_0,E_H)$ preserves this exact boundedness threshold. Above the threshold the product is integral after contraction, with finite coefficient infimum; that infimum is zero if $U(0)$ is a unit. Thus even a nonunit integral prefactor cannot improve the threshold. Here $w_{\mathbf v}(\sum a_\alpha Z^\alpha)=\inf_\alpha(\vp(a_\alpha)+\mathbf v\cdot\alpha)$.

The entire primitive logarithm, with Zink's negative $\operatorname{hex}$ convention, is
\begin{equation}\label{eq:fullprimitive}
\begin{split}
 b(S,\mathbf Y)
 ={}&-\sum_{j=1}^d\sum_{l,n\ge0}
       \frac{S^{p^{n+lh+j}}Y_j^{p^n}}{p^{n+l+1}}\\
 &+\left(\sum_{u\ge0}\frac{S^{q^u}}{p^u}\right)
      \sum_{j=1}^d\sum_{l\ge0}
       \frac{Y_j^{p^{lh+h-j}}}{p^{dl+h-j}}.
\end{split}
\end{equation}
The logarithms in these exact coordinates are $L_0=\sum S^{q^u}/p^u$ and the $L_H$ of DP32. The threshold concerns the full multivariate series; it does not by itself determine a restriction or a selected harmonic value.
\end{theorem}

\begin{proof}
\emph{Divided ghosts and logarithms.}
Use DP38's display matrices
\[
 Ae_0=e_1,\quad Ae_j=pe_{j+1}\ (j\ge1),\qquad
 A'e_0^*=pe_1^*,\quad A'e_j^*=e_{j+1}^*\ (j\ge1),
\]
with cyclic indices. Thus $\langle Au,A'v\rangle=p\langle u,v\rangle$,
$A^h=p^d$, and $(A')^h=p$. Write $u_n(z)=w_n(z)/p^n$ for divided Witt ghosts. The rational resolution operator is $A\operatorname{shift}-1$, with inverse
$-\sum_{r\ge0}(A\operatorname{shift})^r$ on finite-support nilpotent ghosts. The additive logarithm of a lift is
\begin{equation}\label{eq:ghostlog}
 \operatorname{pr}_T\sum_{n\ge0}A^nu_n(z).
\end{equation}
On the relation module this telescopes to $-\operatorname{pr}_Tu_0(Q)=0$, and on the section with only a zeroth ghost in $T$ it is the identity.

All these operations are justified on the integral construction before inverting $p$. On a truncation $N^D=0$ with $D\ge2$, the formula
\[
 u_n(z)=\sum_{r=0}^{\min(n,C)}
              z_{n-r}^{p^r}/p^r,\qquad C=\floor{\log_p(D-1)},
\]
uniformly bounds the denominators of divided ghosts of integral Witt vectors. Since the matrices contract after each complete cycle, the sums and telescoping identities converge on the completed integral resolution. The particular Teichm\"uller sections and rational zeroth-ghost sections below have finite ghost support on each nilpotent truncation; passage to formal series is coefficientwise. This proves both completion and rationalization steps. The source context is Zink (147)--(156), pp.~204--207 \cite{Zink}.

For $\sigma(S)=[S]e_0$, the ghosts are $S^{p^n}e_0/p^n$, and cycles returning to $e_0$ in \eqref{eq:ghostlog} give exactly $L_0(S)$. For $\tau(\mathbf Y)=\sum_{j=1}^d[Y_j]e_j^*$, put
$v_n=\sum_jY_j^{p^n}e_j^*/p^n$ and
$Z=\sum_{n\ge0}(A')^nv_n$. It satisfies
$Z=\sum_jY_je_j^*+(A'/p)Z(\mathbf Y^p)$.
Write $C(\mathbf Y)$ for its $e_0^*$ component and $L_j$ for the others. Then
\[
 C=p^{-1}L_d(\mathbf Y^p),\qquad
 L_1=Y_1+p^{-1}L_d(\mathbf Y^{p^2}),\qquad
 L_j=Y_j+p^{-1}L_{j-1}(\mathbf Y^p)\ (j\ge2).
\]
These are exactly DP32's logarithm equations. The path from $e_j^*$ to $e_0^*$ has length $h-j$, and each further cycle contributes one numerator factor $p$, so
\[
 C(\mathbf Y)=\sum_{j=1}^d\sum_{l\ge0}
                \frac{Y_j^{p^{lh+h-j}}}{p^{dl+h-j}}.
\]

\emph{Comparison with an additive bilinear section.}
Let $x=L_0(S)$ and $\mathbf y=L_H(\mathbf Y)$. Choose additive rational lifts $g_0(x)$, $g'_0(\mathbf y)$ with only zeroth ghosts $xe_0$, $\sum y_je_j^*$. The scalar-one section in that additive frame is bilinear. For the comparison defects put
\[
 k=(A\operatorname{shift}-1)^{-1}(\sigma-g_0),\qquad
 l=(A'\operatorname{shift}-1)^{-1}(\tau-g'_0).
\]
Their zeroth $T$ and $T'$ projections vanish by \eqref{eq:ghostlog}, so they belong to the rationalized normal submodules. Explicitly,
\[
 k_n=-\sum_{r\ge0}A^ru_{n+r}+\mathbf1_{n=0}xe_0,\qquad
 l_n=-\sum_{r\ge0}(A')^rv_{n+r}
                 +\mathbf1_{n=0}\sum_jy_je_j^*.
\]
Zink's Proposition~106, p.~226, and (218), p.~239, express the Teichm\"uller frame as the additive frame multiplied by
\[
 \operatorname{hex}\langle V^{-1}k,\tau\rangle\,
 \operatorname{hex}\langle g_0,l\rangle^{-1}.
\]
The bilinear splitting scalar in that frame is the inverse product. Zink's (196), p.~232, and (210)--(211), p.~236, have the convention
\[
 \log\operatorname{hex}(z)=-\sum_{n\ge0}w_n(z)/p^n.
\]
For two divided-ghost sequences the pairing therefore has logarithm
$-\sum_np^n\langle u_n,v_n\rangle$. Substitution gives
\[
 b=-\sum_{n,r\ge0}p^n\langle A^{r+1}u_{n+r+1},v_n\rangle
          +L_0(S)C(\mathbf Y).
\]
A path of length $lh+j$ from $e_0$ ends at $e_j$ with numerator $p^{dl+j-1}$. Combining both ghost denominators and the factor $p^n$ yields $p^{-n-l-1}$, proving \eqref{eq:fullprimitive}. Every fixed ordinary degree has finitely many summands. Some terms may cancel, but no disjoint-support assumption was made. Every monomial is invariant under
$(\zeta S,\zeta^{-p}Y_1,\ldots,\zeta^{-p^d}Y_d)$; both axes vanish. The additive lifts are equivariant, so the uniqueness argument of DP37--DP38 for this fixed frame identifies this with its invariant primitive splitting. No bilinear correction has been inserted.

\emph{Exact initial form after contraction.}
Fix rational $a\ge p-1$ and choose powers $\alpha_i$ of a common ramified uniformizer with $\vp(\alpha_i)=ar_i$. Set
$b_a=b(\alpha_0S,\alpha_1Y_1,\ldots,\alpha_dY_d)$ and $t=as-1\ge0$.
All monomials of unscaled weight $s$ have the same scalar after division by $p$; call it $\mu$, with $\vp(\mu)=t$. We claim
\begin{equation}\label{eq:primitiveinitial}
 \vp(b_a)=t,\qquad
 \overline{\mu^{-1}b_a}
       =-\sum_{j=1}^dS^{p^j}Y_j+SY_d^p=:Q.
\end{equation}
For a term in the first family of \eqref{eq:fullprimitive}, its weight is
$p^n(s+p^j\rho(q^l-1))$. Its scaled valuation minus $t$ is
\[
 as(p^n-1)-n+a p^{n+j}\rho(q^l-1)-l.
\]
The first part is at least $p^n-1-n$, strictly positive for $n\ge1$ since $p$ is odd. The second is at least $((p-1)p^{n+j}-1)l$, strictly positive for $l\ge1$. Only $n=l=0$ survives.

For the product family put $k=h-j$. Its coefficient valuation is $-u-dl-k$ and weight
$W=\rho q^u+q^l(sp^k-q\rho)$. The baseline $u=l=0,k=1$ has weight $s$, and
\[
 W-s=\rho(q^u-1)+(q^l-1)(sp^k-q\rho)+s(p^k-p).
\]
After multiplication by $p-1$, the three terms are at least
$(p-1)u$, $(q-p)l$, and $p^k-p$, respectively. Here
$sp^k-q\rho\ge sp-q\rho=s-\rho>0$.
Consequently the scaled valuation minus $t$ is at least
\[
 (p-2)u+(q-p-d)l+(p^k-p-k+1).
\]
Every term is nonnegative and equality requires $u=l=0,k=1$.
The surviving monomials in \eqref{eq:primitiveinitial} are distinct, and
$\partial Q/\partial Y_1=-S^p\ne0$, including when $d=1$.

\emph{Exponential obstruction, including arbitrary nonzero bounded multipliers.}
We prove the needed algebraic lemma in full. Let $R'$ be a complete discrete valuation ring over $\Z_p$, with uniformizer $\pi$, and let $z$ have zero constant term and bounded coefficients. Suppose its coefficient infimum is $t\in[0,s)$ and its initial form has a nonzero formal partial derivative. Put
$B=R'[[Z_1,\ldots,Z_m]][1/\pi]$.
Its coefficient infimum is multiplicative: remove each attained minimum power of $\pi$ and reduce the product in the domain $k'[[\mathbf Z]]$. It extends to a valuation on $\operatorname{Frac}(B)$, with residue field $\operatorname{Frac}(k'[[\mathbf Z]])$.

The series $\exp(pz)$ is integral and
\[
 \vp(\exp(pz)-1)=1+t
\]
with the initial form of $pz$. Indeed every exponential order $n\ge2$ has valuation at least $n(1+t)-\vp(n!)>1+t$, since
$\vp(n!)\le(n-1)s$ and $1+t-s>0$.
Suppose a nonzero $f\in B$ made $fE$ bounded, where $E=\exp z$.
Then $E=(fE)/f$ belongs to $\operatorname{Frac}(B)$. The identity
$E^p=\exp(pz)$ gives $\vp(E)=0$ and $\bar E^p=1$, hence $\bar E=1$.
Put $\delta=\vp(E-1)>0$. If $\delta\ge s$, every term in
$(1+(E-1))^p-1$ has valuation at least $1+s>1+t$, a contradiction.
If $\delta<s$, its unique least term is $(E-1)^p$. Its initial form is a scalar multiple of a $p$th power in the residue fraction field. Formal partial derivatives extend to that field by the quotient rule and kill such powers, whereas the initial form of $pz$ has a nonzero derivative. This is again a contradiction. Thus every nonzero bounded $f$ gives unbounded coefficients for $f\exp z$.

Apply the lemma to $z=\lambda b_a$. For $p-1\le a<p$, \eqref{eq:primitiveinitial} has $0\le t<s$, so every nonzero bounded multiplier leaves coefficient infimum $-\infty$. For $a\ge p$, $t\ge s$ and
\[
 \vp((\lambda b_a)^n/n!)\ge ns-\vp(n!)
                    =s_p(n)/(p-1)\ge0.
\]
Its exponential is integral with constant one. For smaller $a$, monotonicity of weighted infima supplies the unbounded conclusion.

\emph{Coordinate changes and prefactors.}
For $a\ge1$, both contracted logarithms and their inverses are integral with identity differential. For the line this follows from
$a(q^u-1)/(q-1)-u\ge0$; for the dual it follows from DP32's path-weight increments, each at least its number of arrows. Integral substitution with an integral inverse preserves the extended coefficient infimum: the two substitution inequalities prove equality in the bounded case, and an allegedly bounded image of an unbounded series would become bounded on applying the inverse.
This transfers the result to $\beta$ for $a\ge1$; monotonicity handles $a<1$.
For a nonzero integral $U$, its expression under these contracted inverses is bounded and nonzero for $a\ge1$. The lemma handles $p-1\le a<p$, monotonicity all smaller $a$, and integrality handles $a\ge p$. If $U(0)$ is a unit, its contraction is an integral unit and the finite infimum is zero. This proves every assertion.
\end{proof}

DP41 proves that DP36's inflation factor $p$ is exact in this family of weights. The arbitrary-multiplier conclusion uses the nonzero-derivative initial form and the identity $E^p=\exp(pz)$; it does not generalize to every series with unbounded denominators, as DP35's example $f=p-z$ shows. A path can still annihilate a prefactor or cancel an initial form. The next theorem treats the primitive mixed paths explicitly.


\subsection{Every mixed line with primitive projections}

\begin{theorem}[DP42: sharp primitive mixed-line slopes]\label{DP42}
Assume $p$ is odd and $h\ge3$, and use exactly DP41's group, dual, coordinates and invariant primitive frame over $\mathcal O=W(\F_q)$. Let $u\in\mathcal O^\times$ and let $\mathbf v=(v_1,\ldots,v_d)\in\mathcal O^d$ be primitive. Put
\[
 j=\min\{i:v_i\in\mathcal O^\times\},\qquad k=h-j,\qquad
 N_k=p^k+1,\qquad S_k=\frac{p^k-1}{p-1},
\]
and
\[
 w_k=\frac{p^k}{(p-1)(p^k+1)},\qquad
 c_k^{\mathrm{path}}=\frac1{(p-1)(p^k+1)}.
\]
For an integral $U\in\mathcal O[[S,\mathbf Y]]$ assume only that
\[
 V(t)=U(E_0(ut),E_H(\mathbf vt))
\]
is not identically zero. It may vanish at the origin. For $\lambda\in\mathcal O^\times$ set
$\Theta(t)=V(t)\exp(\lambda\beta(ut,\mathbf vt))=\sum a_nt^n$.
For every rational $w\ge0$,
\[
 \nu_w(\Theta):=\inf_n(\vp(a_n)+wn)=
 \begin{cases}
 -\infty,&w<w_k,\\
 \nu_w(V)\in[0,\infty),&w\ge w_k.
 \end{cases}
\]
The finite value is zero if $U(0)$ is a unit. In particular
\begin{equation}\label{eq:pathsharp}
 \vp(n!a_n)\ge c_k^{\mathrm{path}}n-\frac{s_p(n)}{p-1},
 \qquad
 \liminf_{n\to\infty}\frac{\vp(n!a_n)}n=c_k^{\mathrm{path}}.
\end{equation}
The logarithmic restriction also has the exact congruence
\begin{equation}\label{eq:pathcoefficient}
 p^{S_k}[t^{N_k}]\beta(ut,\mathbf vt)
       \equiv(-1)^{k-1}u v_j^{p^k}\pmod p.
\end{equation}
The integer $k$ is the restricted-$p$ nilpotency length of the reduced dual direction. Thus the largest slope among all these primitive mixed lines is
$c_1^{\mathrm{path}}=1/(p^2-1)$, independent of $h$; longer restricted-$p$ chains give strictly smaller slopes.
\end{theorem}

\begin{proof}
Put $s=1/(p-1)$. We isolate the full weighted initial form, including its behavior under inverse logarithms; a first pole alone would not determine the exponential radius.
For rational
\[
 w\in I_k:=\left[\frac{S_k}{p^k+1},\,w_k\right],
 \qquad \delta=s-w,
\]
define
\[
 \omega_0=w,\qquad \omega_i=s-p^{i-j}\delta\quad(1\le i\le d).
\]
Then $\omega_j=w$,
\[
 \frac{s}{p^k+1}\le\delta\le\frac{2s}{p^k+1}<s,\qquad
 0\le\tau:=2s-(p^k+1)\delta=(p^k+1)w-S_k\le s.
\]
All weights are positive and exceed DP41's original contraction weights:
since $p\ge3$,
\[
 \delta\le\frac{2s}{p^k+1}\le\frac1{p^k+1}
 <\frac1{p^k-p^{-j}}=p^j\rho,
 \qquad \omega_i-r_i=p^i(\rho-p^{-j}\delta)>0.
\]
Also $w\ge1/(p+1)>\rho$. These imply
\begin{equation}\label{eq:pathstrict}
 w(q-1)>1,\qquad p^{h-i}\omega_i(q-1)>d.
\end{equation}
For the second inequality it suffices to replace $\omega_i$ by $r_i$:
the resulting expression is at least
$(q-p)/(p-1)=p+\cdots+p^{h-1}>d$.

\emph{Integral logarithms and their inverse reductions.}
Choose $\pi^e=p$ clearing all weight denominators and put
$\alpha_i=\pi^{e\omega_i}$, $\gamma=\alpha_0=\alpha_j$.
The contracted scalar logarithm reduces to the identity, since each nonlinear coefficient has valuation $(q^a-1)w-a>0$.
For the dual arrows,
\[
 p\omega_{i-1}-\omega_i=1\quad(i\ge2),\qquad
 p^2\omega_d-\omega_1
       =p+1-p^{1-j}(q-1)\delta>1.
\]
A composition of at least two arrows has weight increment at least
$1+p+\cdots+p^{a-1}>a$, so all higher iterates divided by $p^a$ vanish in reduction. The two contracted logarithms are integral with identity differential and integral inverses, and the dual reduction is exactly
\[
 \overline{\widetilde L}_1=Y_1,\qquad
 \overline{\widetilde L}_i=Y_i+Y_{i-1}^p\quad(i\ge2).
\]
The scaled additive input in coordinate $i$ is
$z_i(T)=v_i\gamma\alpha_i^{-1}T$.
For $i>j$ its valuation is positive because $\omega_i<w$.
For $i<j$ we have $v_i\in p\mathcal O$, whereas
$0<\omega_i-w=\delta(1-p^{i-j})<s<1$, so its valuation is again positive. The unramified integral lattice is used here. Only $z_j=v_jT$ survives. Successive inversion therefore gives
\begin{equation}\label{eq:pathinversereduction}
 \overline{\widetilde E}_i(z(T))=
 \begin{cases}
 0,&i<j,\\
 (-1)^{i-j}\bar v_j^{p^{i-j}}T^{p^{i-j}},&i\ge j,
 \end{cases}
 \qquad
 \overline{\widetilde E}_0(uT)=\bar uT.
\end{equation}
Oddness of $p$ preserves these signs under $p$th powers.

\emph{The unique initial term of the primitive logarithm.}
Apply the weights $\boldsymbol\omega$ to \eqref{eq:fullprimitive}, and put
$\mu=\alpha_0\alpha_d^p/p$, of valuation $\tau$.
For its first family set
\[
 A_i=p^iw+\omega_i=s+p^i\bigl(s-\delta(1+p^{-j})\bigr).
\]
The bracket is positive since $p^k+1>2(1+p^{-j})$; thus $A_i$ increases with $i$. Its smallest baseline has excess
\[
 A_1-1-\tau
       =\delta\bigl(p^k-p+1-p^{1-j}\bigr)>0.
\]
For $k=1$ the bracket is $1-p^{1-j}>0$ because $h>2$; for $k>1$ it is positive even at $j=1$. A general first-family term has excess over its baseline
\[
 (p^n-1)A_i-n+p^{n+i}w(q^l-1)-l.
\]
Both parts are nonnegative, strictly positive for a positive corresponding index, using $A_i>1$ and \eqref{eq:pathstrict}. Thus the first family lies strictly above $\tau$.

In the product family put $e_i=h-i$. Its baseline excess is
\[
 w+p^{e_i}\omega_i-e_i-\tau
       =\frac{p^{e_i}-1}{p-1}-e_i.
\]
This is zero precisely for $i=d$, and positive otherwise. Its further excess at indices $a,l$ is
\[
 w(q^a-1)-a+p^{e_i}\omega_i(q^l-1)-dl,
\]
strictly positive if either index is positive, by \eqref{eq:pathstrict}. Consequently the unique surviving term is $SY_d^p/p$, and
\begin{equation}\label{eq:pathinitial}
 \overline{\mu^{-1}b(\alpha_0S,\alpha_1Y_1,\ldots,\alpha_dY_d)}
       =SY_d^p.
\end{equation}
No cancellation can occur at this weight.

Composition of \eqref{eq:pathinitial} with the integral inverse maps commutes with reduction. Equation \eqref{eq:pathinversereduction} now gives the entire restricted initial form
\[
 \overline{\mu^{-1}\beta(u\gamma T,\mathbf v\gamma T)}
       =(-1)^{k-1}\bar u\bar v_j^{p^k}T^{p^k+1}.
\]
Its derivative is nonzero in characteristic $p$. With the common chosen powers of $\pi$ one has exactly
$\mu=\gamma^{p^k+1}/p^{S_k}$.
Comparing that coefficient proves \eqref{eq:pathcoefficient}: its original coefficients lie in the unramified field, so positive valuation of the difference implies divisibility by $p$.

\emph{Exponential threshold and the prefactor.}
For $w\in I_k$ below $w_k$, the logarithmic coefficient infimum is
$\tau\in[0,s)$ with the displayed nonzero-derivative initial form.
DP41's fraction-field multiplier lemma therefore shows that every nonzero bounded factor leaves $\exp(\lambda\beta(u\gamma T,\mathbf v\gamma T))$ unbounded.
The factor $V(\gamma T)$ is integral by the integral contracted inputs and inverse maps, and is nonzero by hypothesis. Thus $\nu_w(\Theta)=-\infty$ on this interval; monotonicity gives the same conclusion for every smaller nonnegative $w$.

At $w=w_k$ the logarithm has coefficient infimum $s$, so its exponential is integral with coefficient infimum zero. Gauss multiplicativity gives
$\nu_{w_k}(\Theta)=\nu_{w_k}(V)\in[0,\infty)$.
For larger $w$ both factors remain integral by monotonicity and the exponential retains constant one, proving the asserted finite equality. A unit $U(0)$ makes this infimum zero.
In particular $\vp(a_n)\ge-w_kn$. If $\liminf\vp(a_n)/n$ were greater than $-w_k$, some rational $w<w_k$ would make the weighted infimum finite, a contradiction. Therefore $\liminf\vp(a_n)/n=-w_k$.
Legendre's formula and $s_p(n)/n\to0$ give \eqref{eq:pathsharp}. Zero coefficients have valuation $+\infty$; unboundedness already forces infinitely many nonzero ones.

\emph{Restricted-$p$ interpretation.}
Along dual axis $i<d$, the logarithm through degree $p$ is
$te_i+t^pe_{i+1}/p+O(t^{p+1})$.
The translation coefficient is
$D_i^p/p!+D_{i+1}/p$, an integral operator. Thus modulo $p$,
$\bar D_i^{[p]}$ is a nonzero scalar multiple of $\bar D_{i+1}$.
Along the last axis the first nonlinear logarithm term has degree $p^2$, so $D_d^p/p!$ is integral and $\bar D_d^{[p]}=0$.
Invariant directions commute, and their restricted operation is additive and Frobenius-semilinear. The first nonzero coordinate $j$ of the reduced vector therefore survives exactly $k-1=d-j$ iterations and dies on the next. Later coordinates cannot cancel its leading surviving component. This proves the intrinsic interpretation and, since $s/(p^k+1)\le s/(p+1)=1/(p^2-1)$, the maximum assertion.
\end{proof}

The result covers every mixed logarithmic line whose two projections are primitive in these unramified lattices, with any integral prefactor having nonzero restriction. It strengthens DP36's sufficient last-axis residual $\rho_h$ to the exact $\rho_2$. It does not cover a zero or nonprimitive projection, a nonunit exponential parameter, another prescribed gauge, or arbitrary selected harmonic differential polynomials. In particular it is a local obstruction to a proposed improvement on these paths, not a nonexistence theorem for every shorter global correspondence.


\subsection{Optimal affine error for every cyclic type}

\begin{theorem}[DP43: the optimal affine loss is the codimension]\label{DP43}
For the arbitrary-unit cyclic display of DP40, let $m=|A|$, $\ell=h-m$, and define at every cyclic position $a$, selected or not,
\[
 c_a=\frac{\sum_{b\notin A}p^{[a-b]_h}}{q-1},\qquad
 r_a=\frac1{p-1}-c_a.
\]
For $i\in A$ put
\[
 E_i(n)=\sum_{\substack{u\ge1\\i+u\notin A}}
                  \floor{n/p^u},
\]
where membership is read modulo $h$. Then, coefficientwise for every integral $f$,
\begin{equation}\label{eq:cyclicaffine}
 \vp(D_i^nf)\ge
 \max\{E_i(n),\ceil{c_in-(h-m)}\}.
\end{equation}
The slope $c_i$ and the constant $h-m$ are both optimal for every fixed normalized eigendirection. If $g_a$ is the forward gap from $a\in A$ to its next selected index, the single monomial
\[
 f_A=\prod_{a\in A}X_a^{p^{g_a}-1}
\]
attains the affine bound in every direction:
\[
 \vp(D_i^{q-1}f_A(0))=c_i(q-1)-(h-m).
\]
All primes, including $p=2$, are covered. The floor sum and affine bound are two lower bounds; no exact optimum at every order is asserted.
\end{theorem}

\begin{proof}
Write $s=1/(p-1)$ and $\epsilon_a=1$ off $A$, zero on $A$. The weights at all cyclic positions, including the exponent-zero contribution off $A$, satisfy
\[
 0<c_a,r_a<s,\qquad
 c_a=\sum_{u\ge1}\epsilon_{a+u}p^{-u}.
\]
Consequently
\[
 p^tc_i=b_i(t)+c_{i+t},\qquad
 b_i(t)=\sum_{u=1}^t\epsilon_{i+u}p^{t-u}\in\Z_{\ge0}.
\]
DP40's contraction at $n=p^t$ gives the coefficientwise bound
\[
 \vp(D_i^{p^t}f)\ge p^tc_i-s=b_i(t)+(c_{i+t}-s).
\]
All finite coefficient valuations are integers because $R=W(k)$ is unramified, and $-1<c_{i+t}-s<0$, also at $p=2$.
Thus $D_i^{p^t}$ preserves $p^{b_i(t)}R[[\mathbf X]]$ in the sense that it sends every integral function into that ideal. Compose these operator inclusions for the base-$p$ digits $n=\sum_tn_tp^t$. Interchanging finite sums gives
\[
 \vp(D_i^nf)\ge\sum_tn_tb_i(t)
       =\sum_{u\ge1}\epsilon_{i+u}\floor{n/p^u}=E_i(n).
\]

\emph{A primitive reduction of the normalized $q$th derivative.}
Fix $i\in A$ and put
\[
 C_i=c_i(q-1)=
       \sum_{\substack{1\le u<h\\i+u\notin A}}p^{h-u}\in\Z_{>0}.
\]
Because $i+h\in A$, $b_i(h)=E_i(q)=C_i$. Therefore
$\mathscr T_i=p^{-C_i}D_i^q$ is an integral continuous operator commuting with translations.
For $0<k<q$, put $t=\vp(k)$.
The equality
\[
 \vp\binom qk=h-t
\]
follows from $\binom qk=(q/k)\binom{q-1}{k-1}$, since every coefficient of
$(1+Z)^{q-1}$ is a unit modulo $p$.
For $1\le u<h$ one has
\[
 \floor{k/p^u}+\floor{(q-k)/p^u}
       =q/p^u-\mathbf1_{u>t}.
\]
It follows that
\[
 E_i(k)+E_i(q-k)
       =C_i-\#\{u:t<u<h,\ i+u\notin A\}.
\]
Every cross term in the Leibniz expansion of $\mathscr T_i(fg)$ is therefore divisible by $p$, since its valuation is at least
\[
 h-t-\#\{u:t<u<h,\ i+u\notin A\}\ge1.
\]
Hence $\overline{\mathscr T}_i$ is an invariant derivation of the special-fibre formal group. An invariant derivation is determined by its tangent value at the identity, by differentiating the translation identity. It is therefore a residue-field linear combination of the $\bar D_a$ for $a\in A$.

\emph{The torus isolates the same direction.}
Make a finite unramified coefficient extension containing $\F_q$ and a primitive $(q-1)$st root $\zeta$. The display automorphism
\[
 e_a\longmapsto[\zeta^{p^a}]e_a
\]
commutes with its arrows, even for arbitrary $u_a$, because Witt Frobenius takes $[\zeta^{p^a}]$ to $[\zeta^{p^{a+1}}]$. It induces exactly $X_a\mapsto\zeta^{p^a}X_a$ on the Teichm\"uller coordinates.
The character of $D_a$ in
$D_aa_\zeta^*=\zeta^{p^a}a_\zeta^*D_a$ is $\zeta^{p^a}$.
The character of $\mathscr T_i$ is $\zeta^{qp^i}=\zeta^{p^i}$.
These $h$ characters are distinct also in the residue field. Consequently
$\overline{\mathscr T}_i=\bar a_i\bar D_i$.
DP40's exact $q$th coordinate coefficient makes
$a_i:=\mathscr T_iX_i(0)$ a unit already in the original ring.
The equality descends from the unramified extension, so
\begin{equation}\label{eq:cyclicqoperator}
 D_i^q=p^{C_i}a_iD_i+p^{C_i+1}\mathscr R_i,
\end{equation}
where $\mathscr R_i$ is an integral continuous $R$-linear operator. The entire operator identity, including the extra $p$, has now been proved; one sharp coordinate value alone would not suffice.

\emph{Affine induction.}
For $0\le n<q$ with $n=\sum_{t=0}^{h-1}n_tp^t$,
\[
 c_in-E_i(n)=\sum_{t=0}^{h-1}n_tc_{i+t}
       \le(p-1)\sum_{a\in\Z/h\Z}c_a=\ell.
\]
The last equality counts each of the $\ell$ complementary indices against
$\sum_{a=0}^{h-1}p^a=(q-1)/(p-1)$.
For $n\ge q$, multiply \eqref{eq:cyclicqoperator} by $D_i^{n-q}$ and induct. The valuations of its two terms are at least
\[
 C_i+c_i(n-q+1)-\ell=c_in-\ell,\qquad
 C_i+1+c_i(n-q)-\ell=c_in-\ell+1-c_i.
\]
Since $c_i<s\le1$ and $\mathscr R_i$ preserves integral coefficient ideals, the induction holds. Integer rounding and the independently proved floor sum give \eqref{eq:cyclicaffine}.

\emph{A common sharp witness.}
Let $e_a=p^{g_a}-1$. The selected-edge weight identities give
$p^{g_a}r_a-r_{\operatorname{next}(a)}=1$; summing yields
\[
 \sum_{a\in A}e_ar_a=m.
\]
Fix an input index $i\in A$. If $d(i,a)$ is the forward distance from $i$ to $a$, set $N_a=p^{d(i,a)}$.
The first nonzero pure-$t$ term of $E_a(te_i)$ has degree exactly $N_a$ and coefficient valuation $r_a-N_ar_i$.
Indeed torus equivariance forces $Np^i\equiv p^a\bmod(q-1)$, whose least positive solution is $N_a$.
In DP40's contracted inverse $\sum_{b\ge0}(-1)^bM_\eta^{\circ b}$, the unique path before a complete cycle supplies a nonzero coefficient at that degree.
These two observations prove both minimality and exact valuation, with no cancellation from unit arrows.

List the selected indices cyclically from $i$. Multiplication of $N_a$ by $p^{g_a}$ gives the next degree, with the last giving $q$; thus
$\sum_{a\in A}e_aN_a=q-1$.
The first coefficient of $f_A(E(te_i))$ occurs at that degree and is the product of the corresponding first coefficients, with valuation
\[
 \sum_{a\in A}e_a(r_a-N_ar_i)=m-(q-1)r_i.
\]
Finally $\vp((q-1)!)=s(q-1)-h$, so the factorial derivative has valuation
$c_i(q-1)-(h-m)$, as asserted.
A smaller real affine constant would contradict the ceiling bound at this order, because $c_i(q-1)$ is an integer. The slope itself is sharp by DP40's $q^r$ sequence.
\end{proof}

The singleton type gives the affine loss $h-1$ of DP1, while its complementary type gives loss one as in DP34. Mixed derivative estimates follow by composing the coefficientwise operator inclusions. This codimension statement belongs to the integral display and distribution lattice; it does not follow merely from the rational Newton slopes of an arbitrary deformation.


\subsection{What can actually be rebuilt with fewer even lift parameters}

\begin{proposition}[DP44: smaller rank records and the exact remaining comparison]\label{DP44}
Let $R$ be a recovered \path{Ranks} record with lift count $M$, fixed frame, and parameters $k,t$. For every even $M'$ with $2\le M'\le M$:
\begin{enumerate}[label=(\alph*),leftmargin=*]
\item Replacing $M$ by $M'$ gives a valid rank record. Every \path{HCParam R} restricts by keeping $u,v$ and the first $M'$ entries of $b,e$.
\item Every \path{Host R} transports with the same fields and a weaker size proof. Given the same prime $T$ and norm element $z\bar z=T$, the source's actual generic constructor supplies the new Hermitian \path{Spaces}.
\item At fixed $k,t$, retained positive spaces $X_j$ keep their standard forms and ranks $k+j(t+1)$, but the auxiliary signatures and terminal period change. This reconstruction is not literal truncation with the old endpoint.
\item For the source's gain and main-term definitions, if the other data are held fixed, put
\[
 g_j=\sum_w\operatorname{sz}_w(\operatorname{gainD}_{j,w}),\quad
 e_j=\sum_w\operatorname{sz}_w(E_{u,j,w}),\quad
 T_\tau=\sum_w\operatorname{sz}_w(\tau_w),
\]
\[
 H=\sum_t(\operatorname{phF}(A_t)-A_t\log\Omega),\qquad
 K_u=\sum_t A_t\log T_u.
\]
Then
\[
 \operatorname{Ga}_m=\sum_{j=1}^m g_j-T_\tau,\qquad
 \operatorname{Ee}_m=\sum_{j=1}^m e_j,\qquad
 \operatorname{MAIN}_m=mH-K_u.
\]
In particular, with $Q_m=\operatorname{Ga}_m-\operatorname{Ee}_m-\operatorname{MAIN}_m$,
\[
 Q_{M'}-Q_M=-\sum_{j=M'+1}^{M}(g_j-e_j-H).
\]
If an allowance $B_m$ is also subtracted, the right side acquires $B_M-B_{M'}$. For genuinely rebuilt families, every changed input must instead be included in the comparison.
\end{enumerate}
The proposition neither constructs a nonzero shorter theta family with its required estimates nor asserts minimality of the old lift count.
\end{proposition}

\begin{proof}
The full \path{Ranks} fields are in \path{Formalism/LocalRingCardinality/Params.lean}, lines 6--22, and its \path{HCParam} fields at 50--69. The inequalities
\[
 M'+3\le k,\quad k+M'-1+5\le t,\quad 3M'<t,\quad
 M'\le2\operatorname{ord}\le k
\]
follow from their old versions. The parity assumptions on $k,t$ and the frame are unchanged. The remaining nonlinear field follows from $M'\le t+1$:
$(M')^2\le M'(t+1)\le k+M'(t+1)$.
This proves every rank requirement.
The interlacing conditions on $u,v$ and the first $b$ are unchanged; every retained inequality
$e_j-b_j\ge(t+3)/2$ and $b_{j+1}-e_j\ge(t+3)/2$ is an old one. Integral and half-integral entries remain so. The same restriction preserves the word-family conditions in \path{SpanPowerMembership/NZRowAsm.lean}, lines 15--29.

The full \path{Host} structure in \path{DataHost.lean}, lines 39--109, depends on the changed $m$ only through its size condition
$4t(k+m-1)+1<q,d$ at line 54. That condition weakens when $M$ is replaced by $M'$. All other fields use the same frame, $k,t$, or independent data.
The generic \path{hostSpacesOfNorm} in
\path{SpanPowerMembership/RankRecordRowDimensionZero.lean}, lines 137--142, calls \path{spacesOfRanks}, using the same host's \path{cmData}, integral involution and \path{exists_adelic_neg_one}, together with the assumed norm equation.
The actual constructor is \path{SpanPowerMembership/BaseFormAnisotropicOrthogonalSingle.lean}, hypotheses at 200--205 and \path{spacesOfRanks} at 325--386. It constructs $X_j$ as the positive standard space of rank $k+j(t+1)$, $Y_j$ of rank $k+M'-1+jt$ with negative index $\floor{(M'-j)/2}$, and $W_j$ of rank $k+M'-1$ with signature $(k+j-1,M'-j)$; its distinguished vectors, complement and adelic isometry fields are filled in the cited constructor. Thus the new spaces exist under the stated inputs.

Their retained $X_j$ are unchanged, but
\[
 N_M-N_{M'}=(M-M')(t+1)>0\quad(M'<M).
\]
The old $W_{M'}$ has signature $(k+M'-1,M-M')$; the new terminal $W'_{M'}$ has signature $(k+M'-1,0)$. This proves the endpoint qualification.
Even equality of terminal positive dimensions alone does not force the lift count if $t$ changes: given even $2\le M'<M$, even $k\ge M+3$ and a frame with $M\le2\operatorname{ord}\le k$, choose a sufficiently large integer $c$ and put
$t+1=2M'c$, $t'+1=2Mc$.
Both rank records then satisfy all their inequalities and
$k+M(t+1)=k+M'(t'+1)$.
This numerical construction transports no host bounds, gaps, signatures or period estimates.

Finally \path{DataParams}\textsubscript{2}\path{.lean} defines \path{sz} at 203--206 as a scalar multiple, the lift set at 208, \path{gainD} at 212--213, \path{Fo/Ga/Ee} at 218--222, and \path{MAIN} at 232--234. Distributing the finite sums gives the three displayed identities. Subtracting them for $M,M'$ gives the identity for $Q$ and the optional allowance. No sign assumption on $H$ was used. This proves all parts.
\end{proof}

The next actual input is now precise. The generic nonzero-row constructor
\path{SpanPowerMembership/NZRowE.lean}, \path{hostLevelRecord} at line 194, uses the generic spaces at 198, but its hypotheses at 165--189 include class lattices, a level window, norm congruences, signature weights, local gap and $G$ data, and at 185 a period
\[
 \operatorname{per}_u:\operatorname{FK}(\operatorname{cls}_u(M))
        \longrightarrow\operatorname{En}(d)\longrightarrow\mathbf C.
\]
At $M'$ that period has a different domain. Composing the old period with discarded kernels still uses those kernels; no cheaper geometric operation has thereby been exhibited. The similarly named \path{rankRecord} in the constructor source has empty class sets at lines 25--44 and zero period at 77--83, so it is not a nonzero shorter-chain witness.

The quantitative inputs remain explicit as well.
\path{NZRowAsm.lean}, \path{ExtraNZIn}, lines 61--86, requires the level-error and lattice comparisons, actual local-to-global kernels, entry and entropy ties, the $\operatorname{Kap}$ data, orbit-filling data, normalized dimensions and gauges. Its \path{extraNZ}, lines 91--123, assembles these inputs; it does not eliminate them. \path{GlobalTheoremInputs.lean}, lines 158--160, then requires \path{LinkG}, \path{LinkS}, \path{LinkIRR}, \path{LinkDIV} and \path{LinkMu} for the new family.

In particular \path{LinkG_at}, lines 122--136, is the one-way lower certificate
\[
 \sum_{j=1}^m\sum_{t\in\mathrm{TT}_w}
       (R_{2,j,t,p}-\operatorname{loG})-m\operatorname{flG}
 \le\sum_{j=1}^m\operatorname{gain}_{j,w}-\tau_w.
\]
The odd-level ramp $R_2$ is \path{vonMangoldtClassSum},
\path{DataParams}\textsubscript{2}\path{.lean}, lines 240--242, with cutoff
$\operatorname{Nat.log}_p(a_T^{\mathrm{toNat}})+2$ and weights at least one. This lower certificate cannot be reversed into an upper bound on actual gain.
The upper comparison \path{LinkDIV_at}, lines 77--89, contains
$\operatorname{MAIN}$, $m(B_0+N_c)$ and further costs; the local
\path{LocalQuadraticParts.lean}, \path{Conclusion4}, lines 677--680, uses
$\operatorname{Ga}+\operatorname{Fo}-\operatorname{Ee}$.
Other quantities also depend on $m$: \path{Gam_eq} at
\path{DataParams}\textsubscript{2}\path{.lean}, line 176, is
$C_\Gamma m^5fL$, and \path{TuGe} at 252 requires $b^{m/2}\le T$.
\path{flG}, lines 100--102, and \path{loG}, lines 109--120 of
\path{GlobalTheoremInputs.lean}, contain terminal ranks and further orbit/scale costs.

Thus rank/HC/host compatibility and generic Hermitian spaces are not the remaining obstruction for a smaller even parameter in the existing host. The next gate is a new endpoint period with its ties, followed by the actual selected local inputs and net cost comparison. A mixed CM host additionally has to replace DP25's induced-type input. DP40's equal-direction sum and DP42's primitive-line rates settle their stated local families, but neither identifies the full selected determinants with those rates or reverses the source's gain certificate.


\subsection{The height-six CM Jacobian has the required integral normal form}

\begin{proposition}[DP45: an actual three-direction CM example]\label{DP45}
Let $p\equiv3$ or $5\bmod7$, $R=W(\F_{p^6})$, and
$J=\operatorname{Jac}(y^2=x^7-1)$ with the good integral model of DP22.
Then $\widehat J$ is the actual connected height-six, dimension-three
$p$-divisible group with a unit-arrow cyclic display of type
$A=\{0,1,2\}\subset\Z/6\Z$.
In primitive integral eigendirections its sharp slopes are
\[
 (\gamma_0,\gamma_1,\gamma_2)
       =\frac{(p,p^2,p^3)}{(p-1)(p^3+1)}.
\]
Its actual dual has the complementary type and the same triple.
The mixed bound and exact $q^r$ coordinate witnesses of DP40 apply to this Jacobian, and DP43 gives optimal affine loss three in every direction.
The mean of its three slopes is
\[
 c_2\frac{p^2+p+1}{3(p^2-p+1)}<c_2,
\]
although the largest slope exceeds $c_2$ by $p/(p^3+1)$.
\end{proposition}

\begin{proof}
Since $p\ne2,7$, the smooth projective two-chart model
$y^2=x^7-1$, $v^2=u-u^8$ of DP22 has a genus-three Jacobian abelian scheme; the same FGA inputs cited there apply.
Let $a$ be the \emph{covariant} Jacobian automorphism
$a[P-\infty]=[(\zeta_7x(P),y(P))-\infty]$.
The integral basis $x^{i-1}dx/y$, $i=1,2,3$, has eigenvalues
$\zeta_7^i$ and orders $4,2,0$ at infinity. The induced cotangent action, and hence its ordinary transpose on $\Lie J$, has those same three characters. As in DP22, $\Phi_7(a)=0$: its differential on the characteristic-zero generic fibre vanishes, so the homomorphism vanishes there and hence on the abelian scheme.
The action extends to $\Z_p[\zeta_7]$ on $J[p^\infty]$.
Because $p$ has order six modulo seven, its fraction field is the unramified degree-six field $E_p$.

Write $G=J[p^\infty]$. Tate's connected--\'etale sequence is functorial.
If the \'etale quotient were nonzero, the unital $E_p$ action on its isogeny endomorphisms would be injective. Its rational Tate space would have dimension a positive multiple of six. Since $G$ has height six, the quotient would have full height and $G$ would be \'etale, contradicting its dimension three.
Thus $G$ is connected. The dual has the same height and, by the dimension formula, dimension three; the same field-action argument proves it connected. Tate's formal-group equivalence and abelian-scheme example identify $G$ with $\widehat J$.
These facts use \cite[\S1.4, p.~160; \S2.1 Example (a), p.~161; \S2.2 Proposition 1 and (4), pp.~162,164; \S2.3 Proposition 3 and Example (b), p.~166]{Tate67}.

Work first over $R_n=R/p^n$. Zink's Theorem~9, p.~132, supplies the display of this connected group over the excellent local $p$-nilpotent base; Theorem~5 gives its Witt rank six. Full faithfulness transports the cyclotomic action.
Let $\widehat\zeta=[\zeta_7]\in W(R_n)$. Since seven is invertible, the operators
\[
 e_b=\frac17\sum_{r=0}^6\widehat\zeta^{-br}a^r
       \quad(0\le b\le6)
\]
are orthogonal idempotents, with $e_0=0$ by $\Phi_7(a)=0$.
They split $P_n$ into six finite projective summands $P_{b,n}$.
The ring $W(R_n)$ is local (a Witt vector is a unit exactly when its zeroth component is a unit), so these summands are free.
Semilinear Frobenius sends $P_{b,n}$ to $P_{pb,n}$: it sends
$\widehat\zeta^b$ to $\widehat\zeta^{pb}$.
Its linearization is invertible after $p$ is inverted, since relative to any normal decomposition it differs from the invertible normal map by multiplying each $L$-column by $p$.
The localized Witt ring is nonzero, as it maps to
$W(\F_{p^6})[1/p]$.
The six free ranks are therefore equal along the single Frobenius cycle and sum to six. Each rank is one.

The Hodge sequence
\[
 0\longrightarrow Q_n/I_nP_n\longrightarrow P_n/I_nP_n
       \longrightarrow P_n/Q_n=\Lie G_{R_n}\longrightarrow0
\]
is equivariant. On characters $1,2,3$ its middle and right terms are free of rank one; their surjection is an isomorphism. On the other three characters the right term is zero. Hence
\[
 Q_n\cap P_{b,n}=
 \begin{cases}I_nP_{b,n},&b=1,2,3,\\P_{b,n},&b=4,5,6.\end{cases}
\]
Order the characters as
\[
 (1,3,2,6,4,5)\quad(p\equiv3\bmod7),\qquad
 (2,3,1,5,4,6)\quad(p\equiv5\bmod7).
\]
The first three positions are the tangent type. In chosen free generators, $F$ and $V^{-1}$ on the normal summands move to the next line. Their normal matrix is monomial and invertible, so its coefficients are units. This is exactly DP40's arbitrary-unit display.

The choices can be made compatibly in $n$: display base change preserves the idempotents, and a free-line generator lifts through the surjective Witt reductions to a generator. The resulting arrow units lie in
$\varprojlim_n W(R_n)=W(R)$.
Thus the formal isomorphism is integral at all levels, with primitive tangent basis; no isogeny rescaling is used.
DP40 identifies its actual dual with the complementary reciprocal-unit display. For clarity, $a^D$ acts on a covector by precomposition, retaining the eigenvalue $\widehat\zeta^b$, not inverting it. Its tangent characters are therefore the complementary three, consistent with the dual Hodge pairing in Zink Definition~19/(22).

The three numerator sums for DP40 are
$p+p^2+p^3$, $p^2+p^3+p^4$, and $p^3+p^4+p^5$, divided by $p^6-1$.
Cancelling gives the displayed triple; all sharp sequences follow because DP40 allowed arbitrary units. DP43 has codimension $6-3=3$ and supplies the exact affine loss. The mean and last-direction excess follow by the DP40 height-six algebra.
\end{proof}

This removes the conditional rank-one-display hypothesis for this concrete Jacobian at these primes. It leaves its nonalternating CM type unchanged and does not identify its equal-direction mean with a global theta cost.


\subsection{Primitive path slopes with arbitrary unit arrows}

\begin{theorem}[DP46: unit-arrow and invariant-frame independence]\label{DP46}
Let $p$ be odd, $h\ge3$, $R=W(k)$ with perfect $k\supset\F_{p^h}$.
Take DP40's singleton cyclic display with arbitrary units $u_i$:
\[
 Fe_0=u_0e_1,\qquad Fe_i=pu_ie_{i+1}\quad(1\le i\le d=h-1),
\]
and its actual complementary dual. Put $a_i=w_0(u_i)\in R^\times$.
Use its Teichm\"uller coordinates and the invariant primitive Poincar\'e frame.
Every conclusion of DP42 about the exact primitive mixed-line threshold, finite norm, factorial bound and liminf remains valid, with
$k_*=h-j$ the reduced dual direction's restricted-$p$ nilpotency length.
The exact coefficient congruence is now
\[
 p^{S_{k_*}}[t^{p^{k_*}+1}]\beta(ut,\mathbf vt)
 \equiv(-1)^{k_*-1}\theta_j u v_j^{p^{k_*}}\pmod p,\qquad
 \theta_j=\prod_{r=j}^{d}a_r^{-p^{d-r}}\in R^\times.
\]
The same assertions hold for every normalized invariant integral frame with its invariant logarithmic splitting, for every unit exponential parameter $\lambda\in R^\times$.
They also transport under supplied integral CM isomorphisms over unramified extensions. No finite unramified trivialization of arbitrary twists is assumed.
\end{theorem}

\begin{proof}
Let $B_n$ be the $n$th ghost matrix of $F$, and $B'_n$ its complementary reciprocal-unit counterpart. Thus
$B_n^{\mathsf t}B'_n=pI$, and every ghost of each $u_i$ is a unit.
Set $P_{n,0}=I$ and
$P_{n,r}=B_nB_{n+1}\cdots B_{n+r-1}$, with analogous $P'$.
On divided ghosts the resolution is $z_n\mapsto B_nz_{n+1}-z_n$, whose inverse on the finite rational sections is
$-\sum_{r\ge0}P_{n,r}z_{n+r}$.
The logarithm is the selected projection of $\sum_nP_{0,n}z_n$.
DP40's uniform ghost-denominator bound and positive $p$-power per complete cycle justify these sums on the completed integral quotient.

For the Teichm\"uller sections put
$\xi_n=S^{p^n}e_0/p^n$ and
$\upsilon_n=\sum_iY_i^{p^n}e_i^*/p^n$.
Let $C_H$ be the $e_0^*$ component of $\sum_nP'_{0,n}\upsilon_n$.
The additive rational sections with only zeroth ghosts
$L_G(S)e_0$ and $\sum_i(L_H)_i e_i^*$ have relation-module defects
\[
 -\sum_rP_{n,r}\xi_{n+r}+\mathbf1_{n=0}L_G(S)e_0,\qquad
 -\sum_rP'_{n,r}\upsilon_{n+r}
       +\mathbf1_{n=0}\sum_i(L_H)_ie_i^*.
\]
Their selected zeroth components vanish. Substitution in the exact quotient-frame comparison of DP41, using the same negative $\operatorname{hex}$ convention, gives
\begin{equation}\label{eq:unitprimitive}
 b(S,\mathbf Y)=
 -\sum_{n,r\ge0}p^n
       \langle P_{n,r+1}\xi_{n+r+1},\upsilon_n\rangle
       +L_G(S)C_H(\mathbf Y).
\end{equation}
This uses no commutation between distinct ghost matrices.

Unique cyclic paths show that the first sum, the line logarithm, and the complementary component have the exact forms
\[
 -\sum_{i=1}^d\sum_{l,n\ge0}
   \eta_{i,l,n}\frac{S^{p^{n+lh+i}}Y_i^{p^n}}{p^{n+l+1}},\qquad
 L_G=\sum_{a\ge0}A_aS^{q^a}/p^a,
\]
\[
 C_H=\sum_{i=1}^d\sum_{l\ge0}
        C_{i,l}Y_i^{p^{lh+h-i}}/p^{dl+h-i},
\]
where every $\eta,A_a,C_{i,l}$ is a unit and $A_0=1$.
A line path of length $lh+i$ contributes $dl+i-1$ numerator powers of $p$; a dual path to index zero first has length $h-i$ and gains one numerator $p$ per additional cycle. These counts prove the denominators, not just lower bounds for them.
The unique coefficient of $SY_d^p$ is $a_d^{-1}/p$.
The ordinary first dual-logarithm arrows are
$a_{i-1}^{-1}Y_{i-1}^p/p$ into coordinate $i\ge2$.

Use exactly DP42's interval and weights
$\omega_0=w$, $\omega_i=s-p^{i-j}(s-w)$.
Choose all scales $\alpha_i=\pi^{e\omega_i}$ using a common $\pi^e=p$, with $e$ clearing the weights, and put $\gamma=\alpha_0=\alpha_j$.
The supports and denominator valuations just proved are precisely those used in its strict initial-form inequalities. Its first-family baseline gap
$(s-w)(p^{k_*}-p+1-p^{1-j})$ is positive; its product-family gap
$(p^{h-i}-1)/(p-1)-(h-i)$ vanishes only for $i=d$; additional cycles strictly increase valuation. Thus
\[
 \overline{\mu^{-1}b(\alpha_0S,\alpha_1Y_1,\ldots,\alpha_dY_d)}
       =\bar a_d^{-1}SY_d^p,\qquad \mu=\alpha_0\alpha_d^p/p.
\]
The reduced contracted logarithm has ordinary arrows
$Y_i+\bar a_{i-1}^{-1}Y_{i-1}^p$, with the exceptional arrow absent.
Only the first unit input $j$ survives, exactly as in DP42.
Recursive inversion gives its last coordinate
\[
 (-1)^{d-j}
 \left(\prod_{r=j}^{d-1}\bar a_r^{-p^{d-r-1}}\right)
 \bar v_j^{p^{d-j}}T^{p^{d-j}}.
\]
Hence the entire restricted initial form is
$(-1)^{k_*-1}\bar\theta_j\bar u\bar v_j^{p^{k_*}}T^{p^{k_*}+1}$.
It has nonzero derivative. The identity
$\mu=\gamma^{p^{k_*}+1}/p^{S_{k_*}}$ gives the coefficient congruence.
DP41's nonzero-multiplier lemma below the threshold, the factorial estimate at the threshold, monotonicity, and Gauss multiplicativity now prove every asserted norm and factorial conclusion.
The first dual arrows also show that the restricted operation takes each $D_i$ to a unit times $D_{i+1}$ and kills the last. Its length on the primitive vector is $h-j$, as asserted.

The display torus $e_i\mapsto[\zeta^{p^i}]e_i$ commutes with all unit arrows and acts on the paired coordinates as
$(S,Y_i)\mapsto(\zeta S,\zeta^{-p^i}Y_i)$.
The pairing frame and additive sections are equivariant. No bilinear additive monomial $xy_i$ is invariant, so the invariant logarithmic splitting is unique in a fixed invariant frame.
If another such integral frame differs by an invariant normalized unit $J$, its splitting changes by $\pm\log J$.
For an invariant nonconstant monomial $S^a\mathbf Y^{\mathbf b}$ its original contraction weight is
\[
 \rho\left(a-\sum_ip^ib_i\right)+s|\mathbf b|.
\]
The parenthesis is divisible by $q-1$; this positive weight lies in $s\Z$ and is at least $s$.
Each DP42 profile weight strictly exceeds its original contraction weight. Uniform strictness follows from the minimum of these finitely many positive increments. Therefore the contracted $J-1$, and then $\log J$, have coefficient infimum strictly greater than $s$.
For every integral $\lambda$, $\exp(\pm\lambda\log J)$ is an integral unit there. It leaves both the threshold and the initial coefficient congruence unchanged; monotonicity handles larger and smaller scalar path weights.
This argument does not assert unrestricted binomial integrality for arbitrary $\lambda\in R$. For arbitrary, possibly noninvariant integral frames, integral binomial powers give the corresponding frame conclusion when $\lambda\in\Z_p$.

Finally a supplied integral $O_q$-linear isomorphism $\phi:G\to G_0$ and its inverse dual $\psi:H\to H_0$ transport the primitive pairing. Their tangent maps are a unit scalar $a$ and an invertible integral matrix $B$. Normalized logarithms satisfy
$L_{G_0}\phi=aL_G$ and $L_{H_0}\psi=BL_H$.
Primitive projections remain primitive, restricted-$p$ length is unchanged by functoriality, and integral prefactors and nonzero restrictions are preserved. Invariant frames compare by the unit just treated. Unramified extension does not change coefficient valuations, so numerical conclusions descend whenever the original expression is defined over the smaller ring.
\end{proof}


\subsection{The actual polarized block and the trace-one theta diagonal}

\begin{proposition}[DP47: the hyperbolic Poincar\'e pullback]\label{DP47}
Let $G$ be the actual Honda line of height $h\ge2$, or an arbitrary-unit line model under the hypotheses of DP46. Put $H=G^D$, and let
$B=\exp b$ be its invariant primitive Poincar\'e splitting in the integral frame of DP37--DP41 or DP46. Set $A=G\times H$. Suppose the integral symmetric biextension on $A\times A$ is, with axial rigidifications, the actual hyperbolic evaluation biextension
\[
 \mathcal P_{\mathrm{hyp}}((x,y),(x',y'))
       =\mathcal P(x,y')\otimes\mathcal P(x',y).
\]
Let the commutative CM action preserve this decomposition with its adjoint involution, and let $w_0+w_0^*=1$. In the block coordinates write
$w_0=(a,(1-a)^D)$, where $a$ is an integral endomorphism of $G$ commuting with its Teichm\"uller action.

Under the source convention $\mathrm{ph}=\exp(-\delta h)$, the product frame has canonical gauge
$h_0((x,y),(x',y'))=-b(x,y')-b(x',y)$.
There is an integral normalized unit $U_0$ such that
\[
 \exp(h_0(z,w_0z))=U_0(x,y)B(x,y)^{-1}.
\]
For any other source canonical gauge in any normalized integral frame,
\begin{equation}\label{eq:hyperbolicpullback}
 \exp(h(z,w_0z))=U(x,y)B(x,y)^{-1}
                       \exp Q(\log_Gx,\log_Hy),
\end{equation}
where $U$ is an integral normalized unit and $Q$ is a homogeneous quadratic form over the fraction field.
The inverse biextension convention reverses both signs coherently.
Thus for an integral section germ $f$ and integer tensor power $\kappa\ge0$,
\[
 f\exp(\kappa h_L)=\widetilde f\,B^{-\kappa}\exp(\kappa Q),
 \qquad \widetilde f=fU^\kappa\ \hbox{integral}.
\]
The primitive coefficient is exactly $-\kappa$; for DP42's unit case one requires $p\nmid\kappa$.
\end{proposition}

\begin{proof}
Write the integral addition factors as
\[
 C_G(x,x';y)=\frac{B(x+_Gx',y)}{B(x,y)B(x',y)},\qquad
 C_H(x;y,y')=\frac{B(x,y+_Hy')}{B(x,y)B(x,y')}.
\]
The actual Poincar\'e adjunction
$(a\times1)^*\mathcal P\simeq(1\times a^D)^*\mathcal P$
is integral and preserves the laws and rigidifications. In the integral frame its scalar is a normalized integral unit.
The ratio of the two rational bilinear sections after this adjunction is a rational bihomomorphism. Its logarithm is bilinear in additive logarithm coordinates: additivity in each variable forces degree one in each by comparing binomial coefficients in characteristic zero.

That bilinear ambiguity vanishes here. The adjunction is natural, $a$ commutes with the Teichm\"uller action, and the chosen frame and splitting are invariant under
$(x,y_i)\mapsto(\zeta x,\zeta^{-p^i}y_i)$.
No additive monomial $XY_i$ has trivial character, since
$0<p^i-1<p^h-1$ for $1\le i<h$.
Therefore
\[
 J_a(x,y):=\frac{B(ax,y)}{B(x,a^Dy)}
\]
is exactly an integral normalized unit, with no additional bilinear exponential.
This is the scalar of an actual integral adjunction, not an appeal to arbitrary integral-unit powers.

Put $c=1-a$. Cartier duality is additive, so
$a^Dy+_Hc^Dy=y$. Consequently
\[
 B(x,c^Dy)B(ax,y)
 =J_a(x,y)\,C_H(x;a^Dy,c^Dy)^{-1}B(x,y).
\]
All substituted addition factors are integral units. Taking inverses proves the first formula with $U_0=C_H/J_a$.
The trace-one relation therefore gives one primitive factor, with no difference of eigenvalues.
Every integral endomorphism of the one-dimensional $G$ commutes with its Teichm\"uller action: normalized logarithms turn both into scalar multiplications over the fraction field, and equality there is equality of the integral series.
The same adjunction units verify the source's adjoint identities for $h_0$; its symmetry is literal, its axes vanish, and its addition identities are the negative coboundaries displayed above. This verifies that it is a canonical gauge in the stated convention.

In a fixed integral frame, divide the exponential addition identities for $h$ and $h_0$. Injectivity of formal exponential over the characteristic-zero field makes their difference additive in both arguments, hence bilinear in logarithm coordinates.
Symmetry and adjoint balance follow from the other gauge identities.
Every formal endomorphism becomes its tangent linear map under logarithms. The diagonal difference is therefore $L(Z,dw_0Z)$, a homogeneous quadratic form.
A change of normalized integral frame contributes the logarithm of an integral normalized unit; after integral pullback its exponential is another such unit. Adjust the frames first, then apply the bilinear-difference argument. This proves \eqref{eq:hyperbolicpullback}. Taking the ordinary integer power $\kappa$ preserves integrality of the unit and proves the last assertion.
\end{proof}

The hyperbolic hypothesis itself is geometric under a precise integral criterion. Suppose a polarized CM formal block has conjugate integral idempotents $e,e^*$ with
$A=eA\oplus e^*A=G\oplus H$, its polarization is prime to $p$ on the block, and its adjoint is CM conjugation.
For $z,z'\in eA$, balance gives
$\mathcal B(z,z')\simeq\mathcal B(z,e^*z')=\mathcal B(z,0)$,
canonically trivial; the same holds on $H\times H$.
The perfect polarization identifies the cross factor $H$ with the actual $G^D$.
Symmetry identifies the opposite cross restriction with its transpose, and biadditivity identifies the entire block with the product of these two cross restrictions. This proves the hyperbolic hypothesis integrally. It does not follow just from a slope count. The alternating finite-torsion Weil pairing does not change this symmetric-biextension transposition convention.

The exact source sign and diagonal are
\path{CocycleRightZeroThd.lean}, \path{IsCanonGauge}, lines 174--183,
\path{hL}, line 185, and \path{nabVal}, lines 187--189.
\path{WickVanishing/ExpConjConnectionIntegral.lean}, line 40, explicitly assumes $w_0+\operatorname{star}(w_0)=1$.
\path{Development/AbelianVarietyTheory/Proofs/FormalGroupCharts/CanonGaugeSubIsBiadd.lean}, lines 16--60, supplies the gauge-difference identities; \path{CanonGaugeIndependenceAll.lean}, lines 19--26, supplies the specified selected-value comparison.
The actual pairing and its functoriality are Zink (216)--(220), Proposition~121, (243) and Theorem~122 \cite{Zink}.

On a primitive mixed logarithmic line, at odd $p$, $h>2$ and $p\nmid\kappa$, the resulting expression is
$F(t)\exp(-\kappa\beta(ut,\mathbf vt)+ct^2)$.
The necessary condition is $F\ne0$ as a restricted series; global nonvanishing alone does not prove this.
DP48 treats the quadratic term. These results identify the actual local kernel under the polarized-block hypotheses, while the prescribed global path, selected harmonic module, nonzero minors and cost comparison remain separate inputs.


\subsection{Quadratic canonical-gauge changes on a primitive mixed line}

\begin{proposition}[DP48: a quadratic gauge cannot improve the path slope]\label{DP48}
Use all hypotheses and notation of DP42 over the unramified coefficient field $K$, including the nonzero integral-function restriction $V(t)$. For $\lambda\in\mathcal O^\times$ and $c\in K$ put
\[
 \Theta_c(t)=V(t)\exp\bigl(\lambda\beta(ut,\mathbf vt)+ct^2\bigr)
        =\sum_{n\ge0}a_nt^n.
\]
If $c=0$ set $r_c=-\infty$; otherwise set
$r_c=(s-\vp(c))/2$. Write $R_c=\max\{w_k,r_c\}$.
Then
\[
 \nu_w(\Theta_c)=
 \begin{cases}
 -\infty,&0\le w<R_c,\\
 \nu_w(V)\in[0,\infty),&w\ge R_c,
 \end{cases}
\]
for rational $w$, and
\[
 \vp(n!a_n)\ge(s-R_c)n-\frac{s_p(n)}{p-1},\qquad
 \liminf_{n\to\infty}\frac{\vp(n!a_n)}n=s-R_c.
\]
If $c$ is integral, the sharp slope of DP42 is unchanged. If
$\vp(c)=-M\le-1$, the new slope is $(s-M)/2<0$, strictly worse.
\end{proposition}

\begin{proof}
The case $c=0$ is DP42. Otherwise $m=\vp(c)\in\Z$, by unramifiedness.
If $m\ge0$, then $r_c\le s/2<w_k$.
At every $w\ge w_k$, DP42 makes the primitive exponential integral, and $m+2w\ge s$ makes both $\exp(c\gamma^2T^2)$ and its inverse integral for $\vp(\gamma)=w$. Both have coefficient infimum zero. The prefactor is integral and nonzero, so Gauss multiplicativity gives the claimed finite value.
For $w<w_k$ sufficiently close to $w_k$, the quadratic exponential and its inverse remain integral, while DP42 makes $V\exp(\lambda\beta)$ unbounded. Their product cannot become bounded, since multiplication by the integral inverse would then contradict DP42. Monotonicity gives unboundedness at every smaller nonnegative weight. Thus $R_c=w_k$ in this case.

If $m=-M\le-1$, then, since $s\le1/2$,
$r_c=(s+M)/2>s>w_k$.
At every $w\ge r_c$ both exponential factors are integral with coefficient infimum zero, proving the finite equality again.
For $w<r_c$ sufficiently close to $r_c$, we have $w>w_k$ and
$0\le m+2w<s$.
The factor $V(\gamma T)\exp(\lambda\beta(u\gamma T,\mathbf v\gamma T))$ is nonzero and bounded by DP42.
The exponent $c\gamma^2T^2$ has initial form a nonzero scalar multiple of $T^2$, with nonzero derivative since $p$ is odd.
DP41's fraction-field multiplier lemma therefore proves unboundedness of the full product. Monotonicity extends it to every $w<r_c$.

At $w=R_c$ integrality gives $\vp(a_n)\ge-R_cn$.
If the coefficient liminf were greater than $-R_c$, a rational smaller weight sufficiently close to $R_c$ would have finite coefficient infimum, contradicting the established threshold. Hence $\liminf\vp(a_n)/n=-R_c$. Legendre's exact formula gives the two factorial conclusions.
\end{proof}

The source's difference of canonical logarithmic gauges is balanced, symmetric and biadditive; after logarithms it is bilinear, and after a linear logarithmic path and the $w_0$ diagonal it becomes $ct^2$. Thus DP48 applies once the prescribed expression is identified with the primitive mixed kernel up to an integral prefactor and this ambiguity. It does not supply that geometric identification. Nor does it replace the exact annihilation condition in the source's selected harmonic gauge-independence theorem. The integer valuation of $c$ separates the two cases; no statement for arbitrary ramified quadratic coefficients is asserted here.


\subsection{A concrete polarized theta germ at the height-three prime}

\begin{proposition}[DP49: actual Jacobian theta germs and nonzero primitive paths]\label{DP49}
Let $R=W(\mathbf F_{11^3})$, $K=\operatorname{Frac}R$, and let
$J=\operatorname{Jac}(y^2=x^7-1)$ with DP22's good integral model.
Let $a$ be the covariant Jacobian automorphism induced by
$(x,y)\mapsto(\zeta_7x,y)$, and let $\lambda$ be its canonical principal polarization.
Put
\[
 w_0=-(a+a^2+a^4).
\]
Then $a^\dagger=a^{-1}$, $w_0+w_0^\dagger=1$, and $w_0w_0^\dagger=2$.
The height-three factor $G$ on the orbit $(3,5,6)$ has dimension one; the other factor $H$ has dimension two and is integrally its actual Cartier dual.
Their integral displays are the arbitrary-unit singleton/complement models of DP46.
The polarized formal completion is the actual hyperbolic block $G\times H$ of DP47.

Let $\mathcal P_\lambda$ be the symmetric polarization biextension and
$\mathcal L_{w_0}=(1,w_0)^*\mathcal P_\lambda$.
This trace-one line has associated biextension $\mathcal P_\lambda$.
Fix an integer $\kappa>0$ prime to $11$ and a nonzero integral section germ $f$ of
$\mathcal L_{w_0}^{\otimes\kappa}$ in an integral frame.
Use the unit-arrow logarithm coordinates $X,Y_1,Y_2$.
The corresponding tangent directions have original Jacobian CM characters $3,2,1$.
For a logarithmic line
\[
 X=ut,\qquad (Y_1,Y_2)=(v_1t,v_2t),
 \qquad u\in R^\times,\quad (v_1,v_2)\ \hbox{primitive},
\]
suppose the integral algebraic prefactor has nonzero restriction.
In the invariant product canonical gauge, writing the normalized germ as
$\sum_n A_nt^n$, one has
\[
 v_{11}(n!A_n)\ge c_kn-\frac{s_{11}(n)}{10},
 \qquad
 \liminf_{n\to\infty}\frac{v_{11}(n!A_n)}n=c_k,
\]
where $s_{11}(n)$ is the base-$11$ digit sum and
\[
\begin{array}{c|c|c|c}
\text{dual direction}&k&w_k&c_k=1/10-w_k\\ \hline
v_1\in R^\times&2&121/1220&1/1220\\
v_1\in11R,\ v_2\in R^\times&1&11/120&1/120 .
\end{array}
\]
The integer $k$ is the reduced dual tangent vector's restricted-$11$ nilpotency length.

For any rational canonical gauge let $ct^2$ be its remaining quadratic correction after the integral-frame adjustment of DP47.
With $r_c=(1/10-v_{11}(c))/2$, or $r_0=-\infty$, the exact coefficient threshold is
$R_c=\max(w_k,r_c)$ and the exact factorial-slope limit is $1/10-R_c$.
An integral $c$ leaves the displayed slope unchanged; if $v_{11}(c)=-M<0$, the limit is
$(1/10-M)/2<0$.

Finally, for every fixed $\bar u\in(R/11R)^\times$ and every nonzero
$\bar v\in(R/11R)^2$, some lifts of these exact residues give a nonzero restriction of $f$.
Thus both displayed rates are attained in the product gauge for a nonzero germ.
This is existence in every primitive residue class, not nonvanishing of a path prescribed in advance.
\end{proposition}

\begin{proof}
DP22 supplies the actual smooth genus-three curve and Jacobian abelian scheme, with the integral cyclotomic action.
The differentials $dx/y,x\,dx/y,x^2dx/y$ have characters $1,2,3$.
The covariant Abel--Jacobi convention
$a[P-\infty]=[\sigma(P)-\infty]$ gives the same scalar characters on the tangent space: ordinary transposition retains each scalar.
The curve automorphism fixes infinity and preserves effective divisors, so it preserves the canonical theta polarization.
Thus $a^\vee\lambda a=\lambda$, giving $a^\dagger=a^{-1}$.

Since $11\equiv4\pmod7$ has order three, the integral unramified cyclotomic algebra has the two conjugate factors
$\{1,2,4\}$ and $\{3,5,6\}$.
Their heights and dimensions are those of DP22.
A nonzero etale quotient of either height-three factor would have a rational Tate space over the unramified cubic field, hence positive height divisible by three.
It would exhaust the factor, contradicting its positive dimension.
Both factors are therefore connected; the same applies to their duals.
These are the functorial connected--etale and dimension statements of
\cite[\S2.2, (4); \S2.3, Proposition 3]{Tate67}.

Here the integral unit-arrow normal form is also explicit.
Over $R_n=R/11^n$, use the actual display from
\cite[Theorems 5 and 9, pp.~131--132]{Zink}.
Its rank-six Witt module splits by
\[
 e_b=\frac17\sum_{r=0}^6[\zeta_7]^{-br}a^r,\qquad 1\le b\le6,
\]
into free eigenspaces; the zeroth summand is zero by $\Phi_7(a)=0$.
Frobenius sends character $b$ to $11b$.
Its linearization becomes invertible after inverting $11$:
in a normal decomposition the isomorphism $V^{-1}\oplus F$ becomes that linearization after its $L$ columns are multiplied by $11$.
The coefficient ring after this inversion is nonzero, as it maps to
$W(\mathbf F_{11^3})[1/11]$.
Consequently ranks are equal along each orbit.
Both orbit ranks are positive because the tangent quotient contains characters $1$ and $3$.
Their common ranks $r_1,r_2$ satisfy $3r_1+3r_2=6$, so both are one.

The exact Hodge quotient $P_n/Q_n=\operatorname{Lie}J$ now gives
$Q_n\cap P_{b,n}=I_nP_{b,n}$ for $b=1,2,3$ and
$Q_n\cap P_{b,n}=P_{b,n}$ for $b=4,5,6$.
In the orbit $(3,5,6)$ there is exactly one selected line.
Choose generators in this order.
The normal-decomposition isomorphism is monomial on them, so every arrow coefficient is a Witt unit:
\[
 Fe_0=u_0e_1,\qquad Fe_1=11u_1e_2,\qquad
 Fe_2=11u_2e_0 .
\]
Generators lift compatibly through the surjective Witt reductions, and units remain units.
Their inverse limits give the actual integral unit-arrow model over $R$, exactly as in DP45; no isogeny denominator or nonunit tangent map has been introduced.
The dual display is the complementary reciprocal-unit display
\cite[Definition 19 and (22), p.~143]{Zink}.

Rosati conjugation exchanges the two integral idempotents.
For two points in the same factor, balance by its idempotent identifies the polarization biextension with its zero-axis restriction.
The diagonal restrictions are thus trivial as biextensions.
The principal polarization is perfect on the $11$-divisible group and identifies $H$ with the actual $G^D$.
Symmetry identifies the opposite cross restriction with its transpose; biadditivity gives the hyperbolic product.
This verifies DP47's integral geometric hypotheses.
The ordinary dual action on its two tangent positions retains characters $5,6$.
The original Jacobian action on $H$ is $(a_G^{-1})^D$, so those become $-5=2$ and $-6=1$ modulo seven.
The unit arrows give the nonzero restricted-power chain from character $2$ to character $1$, followed by zero.
This proves the intrinsic description of $k$.

For the trace-one calculation put $\eta=a+a^2+a^4$.
Then
\[
 \eta+\eta^\dagger=\sum_{r=1}^6a^r=-1,\qquad
 \eta\eta^\dagger=3+\sum_{r=1}^6a^r=2.
\]
Thus $w_0=-\eta$ has trace one and norm two.
Also $\eta^2+\eta+2=0$, so $\delta=1+2\eta$ satisfies
$\delta^2=-7$ and $w_0=(1-\delta)/2$.
Its polynomial expression proves integrality directly.

The associated biextension of $\mathcal L_{w_0}$ is
\[
 \mathcal P_\lambda(z,w_0z')\otimes
 \mathcal P_\lambda(z',w_0z).
\]
Symmetry and adjunction turn its second factor into
$\mathcal P_\lambda(z,w_0^\dagger z')$.
The trace-one relation and addition in the second argument give exactly
$\mathcal P_\lambda(z,z')$.
Hence the trace-one line induces the actual principal polarization.

DP47 now identifies the normalized germ, in an integral frame, with
\[
 \widetilde f(S,Y)B(S,Y)^{-\kappa}
       \exp Q_\kappa(\log_GS,\log_HY),
\]
where $\widetilde f$ is integral and $Q_\kappa$ is quadratic.
The coefficient $-\kappa$ is a unit.
On the prescribed logarithmic line the quadratic is $ct^2$.
DP46 gives the exact primitive threshold and factorial limit for the actual unit-arrow model.
DP48's proof applies without change, since it uses precisely that threshold, bounded nonzero restricted prefactors and the bounded-multiplier lemma.
It yields $R_c=\max(w_k,r_c)$ and the stated limit.
In the invariant product gauge $Q_\kappa=0$.
Substituting $p=11$ and $k=1,2$ gives all entries of the table and both factorial inequalities.

It remains to prove that the nonzero restrictions exist.
Let $f_1$ be the integral germ including its unit frame factor.
Inverse logarithms define an automorphism of $K[[X,Y_1,Y_2]]$, so
$F=f_1(E_G(X),E_H(Y))$ is nonzero.
Let $P_d$ be its first nonzero homogeneous polynomial.
The inverse logarithms have identity tangent, so this is the corresponding leading polynomial of $f_1$ as well.
The lifts of any fixed primitive residues $(\bar u,\bar v)$ form a nonempty full $11$-adic open box.
A nonzero polynomial over $K$ cannot vanish on it:
in one variable it has only finitely many roots, whereas an open ball is infinite; induction on the number of variables applies this fact to its coefficient polynomials.
Choose lifts with $P_d(u,v)\ne0$.
Then the coefficient of $t^d$ in $F(ut,vt)$ is nonzero.
The residue class fixes the same restricted-power length $k$, proving the last assertion.

For a nonzero algebraic section on the integral Jacobian model, its local germ is nonzero because vanishing on a nonempty open set would force vanishing at the generic point and everywhere.
Its formal completion remains nonzero by the local Krull intersection theorem: for a proper ideal in a Noetherian local ring, the intersection of all its powers on a finite module is zero
(Stacks Project, Lemma 10.51.4, Tag 00IP,
\url{https://stacks.math.columbia.edu/tag/00IP}).
This requires neither a nonzero value at the origin nor a primitive initial homogeneous part.
\end{proof}

DP49 supplies actual normalized theta-germ rates on a concrete higher-height polarized Jacobian block.
The best ordinary primitive mixed-line rate is $1/(11^2-1)$.
It does not identify the global construction's prescribed paths, selected harmonic polynomials or minors with these ordinary lines, and the polynomial argument does not establish nonvanishing on a previously fixed path.
The Jacobian still fails the old induced-type host condition of DP25.
The local primitive Poincar\'e pullback, however, is an actual proved identity in this example.


\subsection{Nonunit tensor powers and the actual level normalization}

\begin{theorem}[DP50: the exact effect of a $p$-divisible tensor exponent]\label{DP50}
Use the arbitrary-unit singleton/complement display and its actual primitive Poincar\'e kernel from DP46. Assume $p$ odd, $h\ge3$, $q=p^h$, $d=h-1$, and $R=W(k)$ with perfect $k\supset\F_q$. Put
\[
 s=\frac1{p-1},\qquad \rho=\frac1{q-1},\qquad
 r_i=s-p^i\rho\quad(1\le i\le d).
\]
Let $u\in R^\times$, let $\mathbf v\in R^d$ be primitive, and set
\[
 j=\min\{i:v_i\in R^\times\},\qquad
 k_*=h-j,\qquad S_{k_*}=\frac{p^{k_*}-1}{p-1}.
\]
For $\kappa=p^a\lambda$ with $a\ge1$ and $\lambda\in R^\times$, write
\[
 K_\kappa(t)=\exp\bigl(\kappa\beta(ut,\mathbf vt)\bigr)
       =\sum_{n\ge0}A_nt^n.
\]
Then, for every rational $w\ge0$,
\begin{equation}\label{eq:nonunitradius}
 \nu_w(K_\kappa)=
 \begin{cases}-\infty,&w<r_j,\\0,&w\ge r_j.\end{cases}
\end{equation}
More precisely, for every $m\ge0$,
\begin{equation}\label{eq:nonunitsequence}
 N_m=p^{k_*}q^m+1,\qquad
 \vp(A_{N_m})=a-S_{k_*}-r_jp^{k_*}(q^m-1).
\end{equation}
Consequently
\begin{equation}\label{eq:nonunitslope}
 \liminf_{n\to\infty}\frac{\vp(n!A_n)}n
       =s-r_j=\frac{p^j}{q-1}.
\end{equation}
These statements, including the exact valuations in
\eqref{eq:nonunitsequence}, hold in every normalized invariant integral frame of the same primitive biextension.

For every integral formal function $U$, put
$V(t)=U(E_G(ut),E_H(\mathbf vt))$ and
$V(t)K_\kappa(t)=\sum_n C_nt^n$. Then
\begin{equation}\label{eq:nonunitprefactor}
 \vp(n!C_n)\ge\frac{p^j}{q-1}n-\frac{s_p(n)}{p-1}.
\end{equation}
This is a sharp slope as a bound valid for all integral $U$, because $U=1$ attains \eqref{eq:nonunitslope}. Equality of asymptotic slopes is not asserted for each individual prefactor.
\end{theorem}

\begin{proof}
\emph{The boundary contraction.}
Set
\[
 \Delta=r_j-\rho>0,\qquad
 \omega_0=r_j,\quad\omega_i=r_i\ (1\le i\le d).
\]
The smallest possible numerator of
\[
 \Delta=\frac{1+p+\cdots+p^{h-1}-p^j-1}{q-1}
\]
occurs at $j=d$ and is $p+\cdots+p^{d-1}>0$; this uses $h\ge3$.
Choose $e$ clearing the weight denominators, adjoin $\pi$ with $\pi^e=p$, and put
$\alpha_i=\pi^{e\omega_i}$ and $\gamma=\alpha_0=\alpha_j$.
The contracted scalar logarithm reduces to the identity, since its nonlinear coefficient at degree $q^n$ has valuation
$(q^n-1)r_j-n>0$.

For the dual, every selected first arrow has weight defect one. Its ordinary arrows have degree $p$ and its exceptional arrow degree $p^2$.
A path through $b\ge2$ selected arrows has weight increment at least
$1+p+\cdots+p^{b-1}>b$; its denominator $p^b$ therefore leaves positive valuation.
DP40's exact path logarithm, including arbitrary unit coefficients, gives
\begin{equation}\label{eq:nonunitreducedlog}
 \overline{\widetilde L}_G=S,\qquad
 \overline{\widetilde L}_H=\id+M_\eta,
\end{equation}
where
\[
 (M_\eta)_1=\eta_1Y_d^{p^2},\qquad
 (M_\eta)_i=\eta_iY_{i-1}^p\ (2\le i\le d),
 \qquad\eta_i\in k^\times.
\]
The scaled logarithms and their formal inverses are integral with identity linear part. Since $M_\eta$ is additive over the residue field,
\[
 \overline{\widetilde E}_H
       =\sum_{b\ge0}(-1)^bM_\eta^{\circ b}
\]
in the formal degree topology. Its coefficients need not be Frobenius-fixed.

The contracted input to dual coordinate $i$ is $v_i\gamma/\alpha_i$.
For $i>j$ its valuation is positive because $r_i<r_j$.
For $i<j$, use $v_i\in pR$ and
$0<r_i-r_j=(p^j-p^i)\rho<1$ to reach the same conclusion.
Exactly the $j$th unit input survives reduction.
The unique path from $j$ to $d$, followed by any number of complete cycles, therefore gives
\begin{equation}\label{eq:nonunitinversecycle}
 \overline{\widetilde E}_{H,d}(z(T))
       =\sum_{m\ge0}\xi_mT^{p^{k_*-1}q^m},
       \qquad\xi_m\in k^\times.
\end{equation}
The exponents are distinct and there is one path for each, so none of these coefficients cancels. The scalar inverse reduces to $\bar uT$.

\emph{The entire primitive logarithm at this boundary.}
DP46's variable-ghost formula is
\begin{align*}
 b(S,\mathbf Y)
 ={}&-\sum_{i=1}^d\sum_{\ell,n\ge0}
 \epsilon_{i,\ell,n}
 \frac{S^{p^{n+\ell h+i}}Y_i^{p^n}}{p^{n+\ell+1}}\\
 &+\left(\sum_{b\ge0}U_b\frac{S^{q^b}}{p^b}\right)
   \left(\sum_{i=1}^d\sum_{\ell\ge0}
       V_{i,\ell}\frac{Y_i^{p^{\ell h+h-i}}}{p^{d\ell+h-i}}\right),
\end{align*}
where the displayed coefficients are units and $U_0=1$.
Its $SY_d^p$ coefficient is $\upsilon_d^{-1}/p$, where
$\upsilon_d=w_0(u_d)\in R^\times$ is the zeroth ghost of the last display arrow.
Put
\[
 \tau=r_j+pr_d-1=s-1+\Delta,\qquad
 \mu=\alpha_0\alpha_d^p/p,
 \quad\vp(\mu)=\tau.
\]
We verify that every other displayed term has strictly larger contracted valuation, even though $\tau$ may be negative.

For the first family set $A_i^*=p^ir_j+r_i=s+p^i\Delta>s$.
Its baseline excess is
\[
 A_i^*-1-\tau=(p^i-1)\Delta>0.
\]
The additional excess at general $n,\ell$ is
\[
 (p^n-1)A_i^*-n+p^{n+i}r_j(q^\ell-1)-\ell.
\]
Both parts are nonnegative, strictly positive for a positive corresponding index: use $(p^n-1)s\ge n$, $A_i^*>s$, and $r_j(q-1)>1$.
For the product family put $e_i=h-i$. The baseline excess is
\[
 r_j+p^{e_i}r_i-e_i-\tau
       =(p^{e_i}-p)s-e_i+1\ge0,
\]
with equality exactly at $e_i=1$.
The extra excess at positive line index $b$ or dual cycle index $\ell$ is
\[
 r_j(q^b-1)-b+p^{e_i}r_i(q^\ell-1)-d\ell>0.
\]
For the last assertion use
\[
 p^{e_i}r_i(q-1)\ge\frac{q-p}{p-1}
       =p+p^2+\cdots+p^{h-1}>d.
\]
Thus the unique minimum is the $+\upsilon_d^{-1}SY_d^p/p$ term:
\begin{equation}\label{eq:nonunitloginitial}
 \overline{\mu^{-1}b(\alpha_0S,\alpha_1Y_1,\ldots,\alpha_dY_d)}
       =\bar\upsilon_d^{-1}SY_d^p.
\end{equation}
This checks the full ghost series; it does not extrapolate the narrower interval from the unit-exponent proof.

\emph{An infinite initial form survives the exponential.}
Compose \eqref{eq:nonunitloginitial} with the integral inverse maps.
By \eqref{eq:nonunitinversecycle},
\begin{equation}\label{eq:nonunitinfiniteinitial}
 \overline{\mu^{-1}\beta(u\gamma T,\mathbf v\gamma T)}
 =\bar\upsilon_d^{-1}\bar uT
   \left(\sum_{m\ge0}\xi_mT^{p^{k_*-1}q^m}\right)^p
 =\sum_{m\ge0}\chi_mT^{N_m},
 \qquad\chi_m\in k^\times.
\end{equation}
The restricted logarithm consequently has infimum $\tau$, attained at infinitely many coefficients after scaling.
For $z(T)=\kappa\beta(u\gamma T,\mathbf v\gamma T)$ its infimum is
\[
 L=a+\tau=a+s-1+\Delta>s.
\]
For every exponential order $b\ge2$,
\[
 \nu(z^b/b!)-L
       \ge(b-1)L-\vp(b!)
       \ge(b-1)(L-s)>0.
\]
It follows that $\exp z$ and its inverse are integral and that the entire initial form of $\exp z-1$, divided by $\kappa\mu$, is the same nonzero series as in \eqref{eq:nonunitinfiniteinitial}. Therefore
\[
 \vp(A_{N_m})+r_jN_m=a+\tau
       \qquad(m\ge0).
\]
The identity $\tau-r_j(p^{k_*}+1)=-S_{k_*}$ proves
\eqref{eq:nonunitsequence}. Its right side is integral, since
$r_j(q-1)=1+p+\cdots+p^{h-1}-p^j$ is an integer.

At weight $r_j$ the kernel has infimum zero, and the same holds at every larger weight. At a smaller weight $w$, its displayed coefficients satisfy
\[
 \vp(A_{N_m})+wN_m=L-(r_j-w)N_m\longrightarrow-\infty.
\]
This proves \eqref{eq:nonunitradius} and the coefficient slope $-r_j$ directly. Legendre's formula gives \eqref{eq:nonunitslope}; no radius assertion has been inferred from the logarithm alone.

\emph{Frames and prefactors.}
Let $J$ compare two normalized invariant integral frames.
Its axes are one. Every positive supported monomial has positive $S$ degree and, under the original weights $(\rho,r_1,\ldots,r_d)$, has weight in $s\Z_{>0}$ by the torus quantization of DP46.
Increasing the scalar weight by $\Delta$ gives
\[
 \nu(J-1)\ge s+\Delta>s,\qquad
 \nu(\log J)\ge s+\Delta.
\]
The logarithm estimate follows termwise from $\vp(b)\le(b-1)s$.
The invariant splitting changes by $\pm\log J$, with no invariant bilinear remainder. Its contribution after multiplication by $\kappa$ has valuation at least
$a+s+\Delta=L+1$, and cannot change the infinite initial form or its coefficient valuations.

Finally $V(\gamma T)$ is integral for every integral $U$, by the integral inverse maps and inputs already checked. Hence $V K_\kappa$ is integral at weight $r_j$ and
$\vp(C_n)\ge-r_jn$. This proves \eqref{eq:nonunitprefactor}.
The choice $U=1$ proves universal sharpness. An arbitrary prefactor is not multiplied by $\kappa$ as a frame logarithm is, so it can enter at the first nonconstant boundary weight. Equality for each such prefactor does not follow from this argument.
\end{proof}

Thus one factor of $p$ changes the pure-kernel factorial slope from
$1/((p-1)(p^{h-j}+1))$ to $p^j/(p^h-1)$; further factors only change the constant valuation in \eqref{eq:nonunitsequence}.
For the actual height-three block at $p=11$ of DP49, these values are
\[
\begin{array}{c|c|c}
\text{dual residue class}&p\nmid\kappa&p\mid\kappa\\\hline
v_1\in R^\times&1/1220&11/1330\\
v_1\in11R,\ v_2\in R^\times&1/120&121/1330
\end{array}
\]
The last column is exact for the pure kernel and the sharp universal lower-bound slope with integral prefactors. These are genuine local improvements, not the result of adding $\vp(\kappa)$ to an old asymptotic slope.

\begin{corollary}[DP50: the recovered level denominator outweighs this improvement]
The recovered definition \path{thExpr} in
\path{Formalism/LocalToGlobal/ProjectiveLattices/IsThRule.lean}, lines 66--70, uses tensor exponent $M^2$ at lines 68 and 70 and the factor $M^{-P.\mathrm{totalDegree}}$ at line 69.
For a homogeneous degree-$n$ ordinary-path specialization, let $b=\vp(M)\ge1$.
Then the normalized pure kernel $K_{M^2}(t/M)$ has exact factorial slope
\[
 \liminf_n\frac{\vp(n!A_n/M^n)}n
       =\frac{p^j}{q-1}-b<0.
\]
Relative to a unit-level pure kernel, the change is
\[
 \frac{p^j}{q-1}-b-\frac{s}{p^{h-j}+1}
       <s-b\le-\frac b2.
\]
Thus forcing $p\mid M$ cannot improve this universal ordinary-path slope after the actual normalization.
\end{corollary}

\begin{proof}
The tensor exponent $M^2$ has valuation $2b\ge1$, so DP50 applies.
Every coefficient of $K_{M^2}(t/M)$ is $A_n/M^n$, whose valuation is decreased by $bn$.
This proves the exact slope.
Since $p^j/(q-1)<s\le1/2$ and $b\ge1$, subtracting the positive unit-level slope gives the displayed comparison.
For an arbitrary integral prefactor the corresponding sharp universal lower bound is
\[
 \vp(n!C_n/M^n)
       \ge\left(\frac{p^j}{q-1}-b\right)n
           -\frac{s_p(n)}{p-1}.
\]
Its universal sharpness again uses the pure kernel; no equality for every individual section has been asserted.
\end{proof}

This corollary concerns the actual displayed tensor/degree normalization and a homogeneous ordinary-path specialization. It allows additional divisibility from a particular algebraic section, a special selected harmonic value, or a compensating factor elsewhere in the construction; none is supplied merely by $p\mid M$.
For a fixed nonzero integer $\kappa$, its exceptional odd primes also satisfy the elementary bound
\[
 0\le\sum_{\substack{p\mid\kappa\\p\ \mathrm{odd}}}\log p
 \left(\frac{p^{j_p}}{p^{h_p}-1}
       -\frac1{(p-1)(p^{h_p-j_p}+1)}\right)
 \le\sum_{\substack{p\mid\kappa\\p\ \mathrm{odd}}}\frac{\log p}{p-1}
 \le\frac12\log|\kappa|.
\]
The difference in parentheses is positive: its numerator after a common positive denominator is
$(p-2)p^{h_p}+(p-1)p^{j_p}+1$.
At derivative order $n$, these finitely many primes contribute only $O_\kappa(n)$ to a prime-sum comparison, without changing an $n\log X$ coefficient as the cutoff $X$ grows. This does not dismiss a finite improvement in a finite inequality, or supply a cost theorem when $\kappa$ varies with the construction.

A rational quadratic gauge remains a separate input. The new radius $r_j$ can be below $s/2$, so an arbitrary integral quadratic coefficient is not automatically harmless as it was in the unit-exponent situation of DP48. If the quadratic is $\kappa Q$ with integral $Q$, its extra factor of $p$ makes its exponential and inverse integral at every nonnegative weight. No selected harmonic determinant or shorter global chain is identified by the ordinary-path theorem alone.


\subsection{A sharp asymptotic family of harmonic selected values}

\begin{theorem}[DP51: the two-column harmonic rate]\label{DP51}
Use DP46's arbitrary-unit singleton/complement pair over the unramified ring
$R=W(k)$, with $p$ odd and $h\ge3$.
Let $d=h-1$, $s=1/(p-1)$, and write $\beta$ for the primitive logarithm in additive coordinates.
Fix $1\le j\le d$ and $\lambda\in R^\times$, and put
\[
 N=p^{h-j},\qquad S=\frac{N-1}{p-1},\qquad
 w=\frac{sN}{N+1},\qquad c=\frac{s}{N+1},
 \qquad F(x,y)=\exp\bigl(\lambda\beta(x,e_jy)\bigr).
\]
For $a,b\ge0$ define the integral polynomial
\[
 P_{a,b}=\sum_{i=0}^{\min(a,b)}(-1)^i
 \binom ai\binom bi
 X_1^{a-i}Y_1^{b-i}X_2^iY_2^i
\]
and its actual primitive-kernel value
\[
 H_{a,b}=P_{a,b}(\partial)
       \bigl(F(x_1,y_1)F(x_2,y_2)\bigr)\big|_0.
\]
Then $P_{a,b}$ has fixed row degrees $(a,b)$ and is killed by
$\Delta=\partial_{X_1}\partial_{Y_1}+\partial_{X_2}\partial_{Y_2}$.
Moreover
\begin{equation}\label{eq:harmonicsharp}
 \vp(H_{a,b})\ge c(a+b)-\frac{s_p(a)+s_p(b)}{p-1},
 \qquad
 \liminf_{n\to\infty}\min_{a+b=n}\frac{\vp(H_{a,b})}{n}=c.
\end{equation}
Zero values have valuation $+\infty$; infinitely many nonzero values attain the limiting slope.
Any sequence attaining it necessarily satisfies
\[
 \frac{a}{a+b}\longrightarrow\frac1{N+1},\qquad
 \frac{b}{a+b}\longrightarrow\frac{N}{N+1}.
\]

For $j=1$, give the positive type one row, the negative type $d$ rows, and use only its first row.
Then these polynomials satisfy every algebraic condition in the recovered \path{HaO}: summed harmonicity, row invariance, fixed row degrees and integral coefficients.
The sharp slope on that family is $1/((p-1)(p^{h-1}+1))$.
For general $j$ the harmonic statement holds, but a full \path{HaO} interpretation requires that the chosen direction be a permitted first row; no reordering of a prescribed geometric datum is assumed.
The values are unchanged by balanced quadratic gauge factors lying in the formal ideal generated by the column-sum contractions that annihilate $P_{a,b}$.
\end{theorem}

\begin{proof}
Write $z(x,y)=\lambda\beta(x,e_jy)$.
We first recover the entire two-variable initial form from DP42 and DP46.
For every rational
\[
 t\in I=\left[\frac{S}{N+1},w\right],
\]
their contraction profile gives
\begin{equation}\label{eq:harmonicinitial}
 \nu_{(t,t)}(z)=(N+1)t-S,
\end{equation}
with a unique least monomial, a unit multiple of $xy^N/p^S$.
Indeed replace the scalar input $uT$ in those proofs by an independent variable $X$ and the surviving dual input $v_jT$ by $Y$.
The scalar inverse reduces to $X$.
The reduced dual inverse is zero before $j$ and then has successive coordinates equal to unit multiples of $Y,Y^p,\ldots,Y^{p^{d-j}}$.
The full primitive logarithm has the unique least group-coordinate term $S_{\rm group}Y_d^p/p$.
Composition therefore gives a unit multiple of $XY^N$, while all other terms remain strictly higher by the same integral-substitution reduction.
The common contraction powers obey $\mu=\gamma^{N+1}/p^S$.
This proves \eqref{eq:harmonicinitial}, including the coefficient valuation $-S$, with the arbitrary unit arrows and $\lambda$ included.
Invariant integral-frame corrections are strictly higher by DP46.

Choose a rational $t_0$ in the interior of $I$ below $w$.
Boundedness at $(t_0,t_0)$ gives, for every supported exponent,
\[
 \vp([x^ay^b]z)+w(a+b)
 \ge (N+1)t_0-S+(w-t_0)(a+b).
\]
These quantities tend to $+\infty$ with total degree.
At $(w,w)$ the unique least monomial $xy^N$ has weight $s$.
There is a small positive rational $\varepsilon$ such that both vectors
\begin{equation}\label{eq:harmonicperturb}
 (w-N\varepsilon,w+\varepsilon),\qquad
 (w+N\varepsilon,w-\varepsilon)
\end{equation}
have positive entries and give $z$ infimum $s$, with that same unique least monomial.
For rigor, first require $N\varepsilon<(w-t_0)/2$.
Every perturbed weight is then at least $t_0+(w-t_0)/2$, so all sufficiently large degrees have weight greater than $s+1$, uniformly over these perturbations.
Only finitely many other monomials remain; shrink $\varepsilon$ to preserve their strict inequalities.
The weight of $xy^N$ is unchanged because $-N\varepsilon+N\varepsilon=0$.
Thus no unjustified uniform-gap assumption about an infinite series is needed.

Since $\vp(m!)\le ms$, exponentiating a zero-constant series of coefficient infimum at least $s$ gives an integral series after contraction.
The inverse exponential is integral as well.
Therefore $F$ is integral with constant one at both vectors \eqref{eq:harmonicperturb} and at $(w,w)$.

Write $f_{a,b}=[x^ay^b]F$ and set
\[
 D(T)=\sum_{i\ge0}f_{i,i}T^i,\qquad
 C(x,y)=F(x,y)D(-xy)=\sum_{a,b\ge0}C_{a,b}x^ay^b.
\]
These formal series are well defined and $D(0)=1$.
The first vector in \eqref{eq:harmonicperturb} gives
\[
 \vp(f_{i,i})+i\bigl(2w-(N-1)\varepsilon\bigr)\ge0.
\]
Consequently $D(-T^2)$ is integral at the strictly smaller weight
\[
 w_D=w-\frac{N-1}{2}\varepsilon<w.
\]
Also $D(-xy)$ is integral at both vectors \eqref{eq:harmonicperturb}, whose coordinate sums are
$2w-(N-1)\varepsilon$ and $2w+(N-1)\varepsilon$.
Thus $C$ is integral at both vectors.

The first contraction of the $i$th summand of $P_{a,b}$ cancels the second contraction of its $(i+1)$st summand, because
\[
 \binom ai\binom bi(a-i)(b-i)
 =\binom a{i+1}\binom b{i+1}(i+1)^2.
\]
The end terms vanish when an exponent is zero.
This proves $\Delta P_{a,b}=0$.
Factorial coefficient extraction and
\[
 \binom ai\binom bi(a-i)!(b-i)!i!^2=a!b!
\]
give the exact generating identity
\begin{equation}\label{eq:harmonicextract}
 H_{a,b}=a!b!\sum_i(-1)^if_{a-i,b-i}f_{i,i}
         =a!b!C_{a,b}.
\end{equation}
Every harmonic correction term has been retained in this identity.

On the diagonal,
\[
 C(T,T)=F(T,T)D(-T^2).
\]
At rational weights sufficiently close to $w$ from below,
\eqref{eq:harmonicinitial} puts the logarithm of $F(T,T)$ in $[0,s)$, with initial form a unit times $T^{N+1}$.
Its derivative in the residue field is nonzero because $N+1$ is prime to $p$.
The factor $D(-T^2)$ is nonzero and bounded at every weight at least $w_D$.
DP41's bounded-multiplier lemma therefore makes $C(T,T)$ unbounded at all those weights below $w$.
Monotonicity gives unboundedness at every smaller nonnegative weight.
At and above $w$ both factors are integral with constant one.
Thus the diagonal has exact coefficient threshold $w$.

If the two-variable coefficients $C_{a,b}$ were bounded at $(t,t)$ for some $t<w$, their finite sums at each total degree would make the diagonal bounded there, a contradiction.
They are integral at $(w,w)$, so
\[
 \liminf_{n\to\infty}\min_{a+b=n}\frac{\vp(C_{a,b})}{n}=-w.
\]
A strictly larger lower limit would give boundedness at a sufficiently close smaller rational weight, which proves the equality explicitly.
Legendre's identity
\[
 \vp(a!b!)=s(a+b)-s\bigl(s_p(a)+s_p(b)\bigr)
\]
and the uniform $O(\log(n+1))$ digit-sum bound for $a+b=n$ prove
\eqref{eq:harmonicsharp} from \eqref{eq:harmonicextract}.

Both inequalities from \eqref{eq:harmonicperturb} further give
\[
 \vp(C_{a,b})\ge-w(a+b)+\varepsilon|Na-b|.
\]
Hence any sequence attaining the slope $c$ has
$|Na-b|/(a+b)\to0$, proving the asserted degree ray.

For the source conditions, use \path{RowDeg.lean}: \path{lap} lines 21--25,
\path{harmK} 27--30, \path{RowInv} 37--38, \path{RowDeg} 40--41 and
\path{HaO} 46--48; \path{IsThRule.lean}, \path{harmA} 86--87, specializes them.
For $j=1$ the first negative-row contraction is exactly $\Delta$, and all later-row contractions vanish.
Adding an earlier negative row to a later one leaves $P_{a,b}$ unchanged because it involves no later row.
The sole positive row has no row operation.
Each monomial has row degrees $a,b,0,\ldots,0$, and all coefficients are integers.
This proves precisely the stated algebraic \path{HaO} membership.

Finally DP39's Fischer identity makes the pairing with $P_{a,b}$ vanish on the formal ideal generated by the column-sum quadratics whose differential operators annihilate it.
If a zero-constant $Q$ lies in that ideal, so does $\exp(Q)-1$.
Multiplication by this balanced gauge factor therefore does not change the selected value.
This proves the gauge assertion with its exact scope.
\end{proof}

\begin{corollary}[DP51: the sharp universal polynomial bound]
Put $N_*=p^{h-1}$, $w_*=sN_* /(N_*+1)$ and $c_*=s/(N_*+1)$.
For $\ell$ independent copies, integral formal functions $U_b$, and units $\lambda_b$, set
\[
 \Theta(\mathbf Z)=\prod_{b=1}^{\ell}
 U_b(E_G(x_b),E_H(\mathbf y_b))
 \exp\bigl(\lambda_b\beta(x_b,\mathbf y_b)\bigr),
\]
with $h$ additive variables in each copy.
For every nonzero integral homogeneous differential polynomial $P$ of degree $n$,
\begin{equation}\label{eq:harmonicuniversal}
 \vp\bigl(P(\partial)\Theta(0)\bigr)
 \ge c_*n-s\max_{\alpha\in\operatorname{supp}P}
                  \sum_{i=1}^{\ell h}s_p(\alpha_i).
\end{equation}
For $n\ge1$ this implies the bound
$c_*n-\ell h(1+\floor{\log_p n})$.
For fixed $h$ and $\ell\ge1$, the slope $c_*$ is optimal uniformly over these integral polynomials and prefactors.
For $\ell\ge2$ it is already optimal on the source-shaped \path{HaO} submodule with one positive row and $d$ negative rows.
If the number of columns varies with $n$, the exact digit-sum error must be kept; no sublinear error uniform in a growing $\ell$ is asserted.
\end{corollary}

\begin{proof}
Use the profile of DP42/DP46 with $j=1$ at $w=w_*$.
Its scalar and first dual weights are $w_*$, and for $i>1$ its dual weights are
\[
 \omega_i=s-p^{i-1}(s-w_*)<w_*.
\]
All weights are positive and exceed the bare contraction weights.
The scaled inverse logarithms are integral, and the full group-coordinate primitive logarithm has coefficient infimum $s$ at this profile.
Give every independent additive variable the common weight $w_*$.
Its input into each contracted variable has nonnegative extra weight: zero on the scalar and first dual coordinates, and $w_*-\omega_i$ on later dual coordinates.
The inverse substitutions therefore remain integral and cannot decrease the primitive logarithm's infimum below $s$.
Its exponential with unit parameter is integral on the full polydisc.
Every integral prefactor remains integral after these substitutions as well.
Multiplying the copies gives
\[
 \vp([\mathbf Z^\alpha]\Theta)+w_*|\alpha|\ge0.
\]
The invariant derivative of multiorder $\alpha$ is $\alpha!$ times this additive coefficient.
Legendre's formula gives
$\vp(\alpha!)=s|\alpha|-s\sum_i s_p(\alpha_i)$.
The coefficients of $P$ are integral and all supported multiindices have degree $n$, so the nonarchimedean triangle inequality proves \eqref{eq:harmonicuniversal}.
For $0\le\alpha_i\le n$ and $n\ge1$ one has
$s_p(\alpha_i)\le(p-1)(1+\floor{\log_p n})$, giving the stated simpler error.

For the full polynomial module with one copy, take prefactor one and $P=(X+Y_1)^n$.
Its value is the ordinary derivative along $x=t,y_1=t$, whose exact slope is $c_*$ by DP42/DP46.
Extra copies with prefactor one do not affect that value.
For sharpness on the $\ell\ge2$ \path{HaO} submodule take two equal kernels with prefactors one and the polynomials $P_{a,b}$ above with $j=1$.
All unused variables are set to zero in their selected value; extra copies with prefactor one contribute their constant value one.
Their values are precisely $H_{a,b}$ and attain $c_*$ by \eqref{eq:harmonicsharp}.
Their full algebraic \path{HaO} membership was proved above.
A larger uniform slope with sublinear error is therefore impossible on either module.
This does not assert that every prescribed geometric prefactor or weight sequence attains that minimum.
\end{proof}

The full-series bound uses the stated invariant product normalization.
A balanced quadratic gauge factor killed by the harmonic contractions leaves the selected value unchanged, as proved above; an arbitrary rational quadratic before selection is not covered by that assertion.

DP51 closes an asymptotic local shortcut left open by a single-jet example:
the stated harmonic conditions alone do not give a larger uniform slope on this family.
It does not identify the pure kernel germ with a global theta section, force the prescribed global row degrees to approach the displayed ray, or decide the complete gain/cost inequality for a shorter chain.
Those questions still require the actual section, selected weights, path, period and global links.


\subsection{Shorter endpoint parameters and a mixed-type period construction}

\begin{proposition}[DP52: the exact endpoint and period interface]\label{DP52}
The following statements hold for the recovered source interfaces.
\begin{enumerate}[label=(\alph*),leftmargin=*]
\item Let $D,D'$ have the same $k$ and identified complex places, with even lengths
$2\le M'\le M$. Rebuild the smaller rank record and spaces as in DP44.
Put $r=k+M-1$, $r'=k+M'-1$.
For $P:\operatorname{TPar}(D)$, retain its $t,u,v,\eta$ and restrict $b,e$ to
their first $M'$ entries. This defines $P':\operatorname{TPar}(D')$.
Interlacing, parity and centredness pass to $P'$, and for $j\le M'$,
\[
 X_{P'}(j,\tau)=X_P(j,\tau),\qquad
 D_{P'}(j,\tau)\subseteq D_P(j,\tau),
\]
including the Boolean signs in $D_P$.
The noncompact-gap condition passes as well.
For any source theta packages $X,X'$ on these records,
\[
 \operatorname{Hyp3}(X,P),\qquad 1801\le k+M'-1+t
 \quad\Longrightarrow\quad \operatorname{Hyp3}(X',P').
\]
The remaining inequality follows, for example, from
$r+t\ge1801+(M-M')$.
\item The six statements forming $\operatorname{AdmS}(D,X)$, together with
$\operatorname{Hyp3}(X,P)$, imply the source's $\operatorname{PointsLemma}$
at $P$. Their finite downward induction works for every permitted even $M$.
Thus a rebuilt shorter family satisfying these same six statements has the
same nonzero-period conclusion; no additional minimum induction length is
needed.
\item The actual endpoint period is uniform in the rank record.
More strongly, its row and point goodness certificates and its inner-product
formula can be constructed without $\operatorname{Fields.Ok.quad}$.
It suffices to assume the exact conjugation compatibilities
\[
 \phi(\bar x)=\overline{\phi(x)}
       \quad(\phi:E\hookrightarrow\mathbf C),\qquad
 \sigma_0(\bar y)=\overline{\sigma_0(y)}
       \quad(y\in K),
\]
for one coefficient-field embedding $\sigma_0$, and
\[
 \operatorname{emb}_\tau(\bar x)
       =\overline{\operatorname{emb}_\tau(x)},
\]
together with the specified weight and endpoint-level inputs.
These assumptions hold for $E=K=\mathbf Q(\zeta_7)$,
$F=E^+$ and the actual type $\{1,2,3\}$.
Consequently that mixed type has the actual finite-class row/point spaces,
endpoint period and inner-product formula at every admissible rebuilt rank.
This assertion does not say that the period is nonzero on a prescribed
iterated theta image.
\item The unchanged field wrapper nevertheless fails for that type:
\[
 E_{\{1,2,3\}}
 :=\{z\in E:\sigma_1(z)=\sigma_2(z)=\sigma_3(z)\}
 =\mathbf Q,
 \qquad \sigma_a(\zeta_7)=\zeta_7^a.
\]
Thus its field $\operatorname{Fields.Kq}$ has degree one over $\mathbf Q$,
whereas $\operatorname{Fields.Ok.quad}$ requires degree two.
This failure is independent of $M$ and persists under coefficient-field
enlargement, simultaneous relabeling and conjugation of the type.
For a Galois CM field, the degree-two condition is equivalent to the type
being induced from one embedding of an imaginary quadratic subfield.
\end{enumerate}
Part (c) explicitly bypasses the failed wrapper at the endpoint period-map
layer. It does not construct the full mixed-type theta correspondence or
its selected quantitative certificates.
\end{proposition}

\begin{proof}
\textit{Parameter restriction.}
The exact definitions are
\path{Formalism/LocalRingCardinality/DataTheta.lean},
$\operatorname{TPar}$ at 501--511, interlacing at 529--533, blocks at 535--536,
$D_P$ at 538--543, $X_P$ at 545--548 and centredness at 550--554.
The parity, gap and $\operatorname{Hyp3}$ definitions are in
\path{Formalism/SharpRestriction/CuspSpanNormBounds/CayleySignature.lean},
106--118.
The first package requires $t$ odd, $r+5\le t$, the indicated arrays and
$\eta_\tau\in\{1/2,-1/2\}$. Since $r'\le r$, its inequality weakens.
The interlacing conditions involving $u,v,b_0$ are unchanged; those involving
$b_j,e_j,b_{j+1}$ are merely restricted. Parity removes entries from integral
arrays. Every centredness condition retains its old block length $t$ and
either keeps its old indices or restricts them. Hence all these conditions
remain true.

Before its final translation, the source's signed set $D_P(j,\tau)$ is
\[
 \begin{aligned}
 &\{(v_i,\mathrm{true})\}_{i<k-1}\\
 &{}\cup\{(-b_i,[\,i\le j\text{ or }i-j\text{ even}\,])\}_{i<M}\\
 &{}\cup\bigcup_{i<j}\{(x,\mathrm{true}):x\in\operatorname{block}(-e_i,t)\}.
 \end{aligned}
\]
Natural subtraction in the Boolean test gives the identical answer on
each retained pair of indices. For $j\le M'$, the first set is unchanged,
the second loses precisely the indices $i\ge M'$, and the third is unchanged
because $i<j$ implies $i<M'$. Both records translate by the same $\eta_\tau$
when $j$ is odd. This proves the signed inclusion.
The set $X_P$ consists of the unchanged $u$-set and the unchanged blocks
of length $t+1$ with indices $i<j$, proving its asserted equality.

Every pair of opposite-sign entries in the smaller set therefore occurs in
the old set. Its old separation exceeds $(r+Mt)^3$, and
\[
 (r'+M't)^3\le(r+Mt)^3.
\]
Thus the strict gap inequality survives.
The remaining infinity-type identity transports automatically:
every source $\operatorname{ThetaX}$ satisfies $\operatorname{etaT\_pos}$
at \path{DataTheta.lean}, line 614, requiring infinity type $1/2$.
The old identity consequently says $\eta_\tau=-1/2$, which is also the
new package's required identity.
The only unaccounted $\operatorname{Hyp3}$ conjunct is
$1801\le r'+t$, precisely the displayed condition.
Since $r'=r-(M-M')$, the sufficient margin follows. This proves (a).
It does not identify the old indefinite $Y_{M'}$ with the rebuilt endpoint.

\textit{Finite nonzero-period induction.}
Here are the exact six inputs in \path{CayleySignature.lean}:
\path{RallisSiegelWeilFormula}, 154--161;
\path{LiftParameterConditions}, 163--166;
\path{ThetaUnfoldingFreeData}, 168--175;
\path{ThetaStablePeriodVanishing}, 177--184;
\path{RealThetaLiftDiscreteSeries}, 186--194; and
\path{RealConstantNonvanishing}, 196--199.
Their conjunction is \path{AdmS}, 201--202.
The conclusion \path{PointsLemma} is at 138--148 and the auxiliary
\path{PerNZ} at 150--152.
The name \path{TowerRallisLemma} at line 204 denotes the proposition of
the implication; it is not itself a theorem declaration.
We prove that implication from the six inputs here.

The first input supplies a character sequence $\mu$ and irreducible nonzero
representations $\pi_j$ of the definite groups, with their prescribed types.
For every $j<M$, every vector in $\pi_{j+1}$ is a finite sum of upward theta
lifts of vectors in $\pi_j$, using factorizable finite-type Schwartz data.
The second input supplies terminal $\phi_M\in\pi_M$ and a cuspidal vector
$F_M$ of the required type, with nonzero top period.

Suppose $\phi_{j+1},F_{j+1}$ have been constructed with these properties.
Expand $\phi_{j+1}$ as the finite sum supplied by the first input.
The period is additive in this vector; the conjugation in the lower
period is additive as well. Since the sum of period values is nonzero,
some summand has nonzero period. Choose its vector $\phi_j\in\pi_j$ and
upward Schwartz datum $\Phi'_j$.
The third input supplies a factorizable finite-type downward datum $\Phi_j$
whose lift will have nonzero \path{PerNZ}, once its cuspidality is supplied.

For that last requirement, let $W_{j+1}$ be the subspace of automorphic
vectors which are cuspidal and have the specified \path{IsDv} type.
It is a complex subspace: the latter condition says that the conjugate
vector lies in a fixed complex type submodule, and inverse images of
subspaces under a conjugate-linear additive map are subspaces.
It is stable under the finite and real group actions by the sixth input
and the cusp-stability fields of \path{ThetaX}.
For every vector in this subspace, the second clause of the fifth input
gives all lower-tower vanishings needed by the fourth input.
That fourth input therefore makes
\[
 F_j=v\theta_j(\Phi_j,F_{j+1})
\]
cuspidal. The first clause of the fifth input gives its \path{IsDv} type,
and the third input now gives its nonzero period against $\phi_j$.
This completes one step, hence the finite descent from $M$ to zero.
Indices outside the constrained finite range may be filled by zero.
The required character and $\pi_0$ came from the first input.

At zero, the terminal alternative in \path{PerNZ} is unavailable since
$M\ge2$. Thus a lower period is nonzero for some position and
Fourier--Jacobi data. Its rank difference is
\[
 \operatorname{rank}Y_0-\operatorname{rank}X_0
       =(k+M-1)-k=M-1,
\]
which is odd. The field \path{ThetaX.Bper_odd} at
\path{DataTheta.lean}, line 663, replaces the Fourier--Jacobi datum by zero
without changing the value. This is the exact zero-datum conclusion of
\path{PointsLemma}. The required additivity, stability and conjugation
identities are in the full \path{ThetaX} interface at 573--694.
No additional lower bound on the number of steps entered this proof.
This proves (b).

\textit{Actual endpoint period and its weaker field hypotheses.}
The source \path{Formalism/LocalToGlobal/Columns/Record.lean} takes arbitrary
rank, weight and endpoint-level inputs at 1434--1438, defines its row and
point data at 1440--1442 and proves point goodness at 1444--1448.
For a basis $\{h_a\}$ of the row invariant space and a point-class
representative $c$, coefficient pairing with $h_a(c)$ gives a linear
functional on the endpoint coefficient representation.
The evaluation map has linearly independent images of its basis.
The source's extension functional, at 154--205, consequently extends
this coefficient functional to the ambient complex function space and
agrees with it on evaluated vectors.

The definition \path{periodC}, 1568--1571, is the finite sum of these
functionals over point classes and row coordinates, with reciprocal
stabilizer-order weights.
The coefficient Gram matrix is invertible by
\path{Columns/RecordInst.lean}, \path{ptEv_gram_unit}, 47--58.
Thus each coefficient functional is represented by a vector under its
classwise inner product. The theorem
\path{Record.periodC_eq_ipCl}, 1580--1590, and its version in
\path{RecordInst.lean}, 60--64, identify the entire period with the
classwise inner product against \path{jMCL}.
All definitions use the general rank record, with endpoint dimension
$N_R(m)-1$.
The original \path{recPer}, 1603--1610, combines this map with evaluation
and the branching map \path{lowC}, 1558--1564, from $N_R(m)$ to $N_R(m)-1$.
It has no separate numerical existence assumption for a particular $m$.

More precisely, the argument $h_X:X.\operatorname{Ok}$ in the displayed
definition of \path{periodC} is used only through its point-goodness
proof \path{ptOk W \Lambda_M hX}.
Replace that proof by an explicit
\[
 h_{\rm pt}:(\operatorname{ptGSD}W\Lambda_M).\operatorname{Ok}
\]
and supply likewise
$h_{\rm row}:(\operatorname{rowGSD}W\Lambda_M).\operatorname{Ok}$.
Leave the same classwise functionals and row basis in place.
This defines a weakened endpoint map.
The same finite-sum proof gives its inner-product identity.
The Gram theorem just cited has no $X.\operatorname{Ok}$ argument,
so that input has not been reintroduced.

The two goodness proofs can actually be supplied under the weaker
conjugation assumptions in (c).
Read the entire proof \path{formGSD_ok} in
\path{Columns/ColFormsLemmas.lean}, 83--119.
Its five \path{GSD.Ok} fields are obtained as follows.
First, compatibility with complex conjugation at all embeddings of $E$
makes the standard hermitian form positive at all infinite places.
The source's decorated-lattice finiteness theorem supplies finitely many
representatives. The representative move is an isometry paired with the
identity level element; it already lies in \path{FormG}.
Thus this is a finite-orbit proof for the required group itself, not an
inference from finiteness for a larger group.
Second, positivity makes each decorated-lattice stabilizer finite.
The inclusion \path{FormG.incl} is injective, so its stabilizers are finite.
These are exactly the arguments at 84--101.

Third, every value of the finite level character is a root of unity,
by \path{LevFam.chr_root}, 77--81.
Such a scalar and its inverse are algebraic integers, so multiplication
by it preserves an integral lattice.
Combining this with the two supplied coefficient-lattice equivariance
identities proves the \path{nab_smul} and \path{del_smul} fields
(102--109).
Fourth, compatibility of one embedding $\sigma_0$ of the coefficient
field with conjugation gives $\bar c c=1$ for every such root of unity:
after the injective embedding its complex norm is one.
The invariant hermitian coefficient pairing remains invariant after
the character twist, giving the final \path{inv} field (110--119).
These steps cover all five fields; no quadratic-subfield hypothesis is used.

For point coefficients, \path{Record.ptOk}, 1444--1448, supplies the
needed coefficient equivariance from the round coefficient-lattice
identities and \path{roundPk_inv}; its only field input there is
\path{emb_star}.
For row coefficients,
\path{Columns/RowFormsLemmas.lean}, \path{rowVal_eqv}, 20--23,
uses the row lattice identities and \path{squarePk_inv} with precisely
the same input.
The asserted selected-embedding conjugation identity supplies it.
Hence both explicit goodness certificates exist.
For $E=K=\mathbf Q(\zeta_7)$, canonical complex conjugation commutes
with all Galois embeddings, so every one of these weaker hypotheses
holds for the selected embeddings $\sigma_1,\sigma_2,\sigma_3$.
This proves (c), for every specified admissible weight package and level.
The exact weight interfaces are \path{Columns/AbelianHostFieldConditions.lean},
\path{Wts} at 443--449, its goodness at 455--460 and endpoint weights
at 472--476.
The full original \path{recPer} is still attached to a
\path{recChain hX}; the preceding proof constructs the entire endpoint
period-map layer and does not silently supply the remaining global kernels.

\textit{The unchanged field wrapper.}
The carrier of \path{Fields.Kq} is displayed at
\path{AbelianHostFieldConditions.lean}, line 60, within its definition
at 59--78. The field \path{Fields.Ok.quad} at line 108 requires degree
two; \path{Fields.Core} still retains this requirement at line 157.
In $(\mathbf Z/7\mathbf Z)^\times$ the powers of $3$ are
$3,2,6,4,5,1$. Thus $\sigma_3$ generates the order-six Galois group.
Equality $\sigma_1(z)=\sigma_3(z)$ forces $z$ to be fixed by that group,
so $z\in\mathbf Q$; conversely every rational element is fixed by all
selected embeddings. An injective coefficient-field enlargement does not
change these equalities. Simultaneous conjugation or relabeling conjugates
the relative-automorphism subgroup, which remains the full group.
This proves the exact obstruction and its invariances.

For the general assertion, let $E/\mathbf Q$ be Galois CM of degree $2g$
and let the common field $L$ have degree two.
All $g$ selected embeddings restrict to the same embedding of $L$.
That embedding has exactly $[E:L]=g$ extensions, so the type consists
of all of them.
If $L$ were real, the $g$ complex-conjugate embeddings would have the same
restriction, giving $2g$ distinct extensions instead, a contradiction.
Hence $L$ is imaginary quadratic and the type is induced from it.
Conversely an induced type is a coset of $\operatorname{Gal}(E/L)$;
its relative automorphisms fill that subgroup, whose fixed field is $L$.

The contrasting type $\{1,2,4\}$ has an order-three relative subgroup.
Its fixed element $z=\zeta_7+\zeta_7^2+\zeta_7^4$ satisfies
$\bar z=-1-z$ and $z\bar z=2$, hence $z^2+z+2=0$.
Its common field is therefore $\mathbf Q(\sqrt{-7})$.
That is a different CM type, not a relabeling of the actual Jacobian type.
This proves (d).
\end{proof}

Two further source distinctions prevent the conclusion from overstating
what has been constructed.
First, \path{SpanPowerMembership/ThetaFamilyAdmissibleThree.lean}
does construct an abstract \path{thetaG} at 71--153, proves its six
statements at 215--274 and states \path{thetaG_exists} at 277--280.
Its spaces are a common auxiliary Hilbert space, its group actions are
identity and its periods are inner products.
This is a realization of the abstract interface, but does not identify
the actual mixed-CM theta germs with that realization's kernels.
Second, \path{SpanPowerMembership/NZColC.lean}, \path{periodN}, 160--167,
accepts an arbitrary bilinear form on the terminal class functions and
uses identity evaluation/lowering maps.
Its existence likewise supplies no geometric period identification.
The source's initial column nontriviality results
\path{Columns/NonDeg.lean}, \path{MramL_ne_bot}, 522--525, and
\path{finrank_MramV_pos_of_sizes}, 900--902, use $2fk<g_0-1$,
which stays unchanged when only $m$ decreases.
They are not nonvanishing statements for every terminal iterated image.

Together, DP44 and DP52 remove a rank, parameter, abstract induction and
endpoint-period existence obstruction to rebuilding a shorter even chain.
The remaining construction must identify one actual mixed-type theta
family with these periods, its algebraic theta germs and selected local
kernel operators; supply its six theta-correspondence identities; and
prove its local/global quantitative certificates.
The failed quadratic field wrapper is explicitly bypassed at the
endpoint layer, rather than treated as an intrinsic period obstruction.
The net deficit in DP44 must still be evaluated for that actual rebuilt
family. Neither ordinary primitive-path rates nor equal-direction averages
may be substituted for its selected gains, and no minimality of the old
lift count follows from this proposition.

\subsection{Prime redistribution and globally fixed derivative directions}

\begin{theorem}[DP53: the cyclotomic redistribution preserves the leading bare-jet budget]\label{DP53}
Let $J=\operatorname{Jac}(y^2=x^7-1)$ with the good integral model of DP22, and let
\[
 \Phi=\{1,2,3\},\qquad\Psi=\{1,2,4\}
       \subset(\Z/7)^\times
\]
be respectively its actual covariant CM type and the comparison type induced from $\Q(\sqrt{-7})$.
For odd $p\ne7$, label the three actual invariant directions by the global integral differential basis
$\omega_i=x^{i-1}dx/y$, $i=1,2,3$.
Put
\[
 c_2(p)=\frac{p}{p^2-1},\qquad
 c_3(p)=\frac1{p-1}-\frac1{p^3-1},\qquad
 D(p)=(p-1)(p^3+1),
\]
\[
 T(p)=\frac{p(p^2+p+1)}{D(p)}.
\]
The sharp \emph{bare integral directional} slopes are
\[
\begin{array}{c|ccc|c}
p\bmod7&c_1^\Phi(p)&c_2^\Phi(p)&c_3^\Phi(p)&B_\Phi(p)\\\hline
1&0&0&0&0\\
2&p/(p^3-1)&p^2/(p^3-1)&c_3(p)&2c_3(p)\\
4&p^2/(p^3-1)&p/(p^3-1)&c_3(p)&2c_3(p)\\
3&p/D(p)&p^3/D(p)&p^2/D(p)&T(p)\\
5&p^3/D(p)&p/D(p)&p^2/D(p)&T(p)\\
6&c_2(p)&c_2(p)&c_2(p)&3c_2(p)
\end{array}
\]
Every selected direction of the induced type has slope
\[
 c^\Psi(p)=
 \begin{cases}0,&p\equiv1,2,4\pmod7,\\
 c_2(p),&p\equiv3,5,6\pmod7.\end{cases}
\]

For arbitrary nonnegative integers $n_1,n_2,n_3$, held fixed as the prime varies, let $N=n_1+n_2+n_3$. Uniformly in these orders and $X\ge2$,
\begin{align}
 \sum_{\substack{p\le X\\p\ne2,7}}\log p
       \sum_{i=1}^3 n_i c_i^\Phi(p)
       &=\frac N2\log X+O(N),\label{eq:cmprimeactual}\\
 \sum_{\substack{p\le X\\p\ne2,7}}\log p
       \sum_{i=1}^3 n_i c^\Psi(p)
       &=\frac N2\log X+O(N).\label{eq:cmprimeinduced}
\end{align}
Their difference is $O(N)$. The same $O(N)$ difference bound holds for the full integer periodic-floor lower certificates of DP43, not merely for their linear slopes.
These are statements about bare integral derivatives and their lower certificates, not upper bounds on actual theta-selected gain.
\end{theorem}

\begin{proof}
\emph{Actual integral directions at every good prime.}
The two-chart model of DP22 and its relative Jacobian are defined over $\Z[1/14]$. The displayed differentials form a basis on every good fibre, hence a primitive integral relative basis. Their dual Lie basis and the corresponding invariant directions $D_i$ are therefore rational and integral at every prime under consideration.
The covariant cyclotomic automorphism has cotangent and tangent characters $1,2,3$; the ordinary transpose retains each scalar eigenvalue.
After the unramified extension containing $\zeta_7$, the distinct seventh-root characters split the tangent lattice into rank-one direct summands. Thus a local Teichm\"uller tangent generator of a fixed character differs from this global generator only by a unit. No prime-dependent direction is being selected.

Let $f=\operatorname{ord}_7(p)$. The CM height/type theorem gives height $f$ to each Frobenius orbit and dimension equal to its intersection with $\Phi$; this is the standard local formula already used in DP22, \cite[\S5, (10*) and Theorem 3, pp.~106--108]{Tate352}.
An orbit with positive selected count is connected: a nonzero \'etale quotient would carry the unramified degree-$f$ field action, so its rational Tate dimension would be a positive multiple of $f$. It would have full height and force dimension zero, a contradiction. For a partially selected orbit the dual is connected by the same argument.
The integral idempotents and the connected formal-group equivalence identify actual formal factors, not merely isogeny factors.

For each partially selected orbit, apply the integral argument of DP45 and DP49. Over the $p$-nilpotent quotients of the unramified base, its display has rank $f$; the Witt character idempotents are free summands, semilinear Frobenius permutes them cyclically, and its linearization becomes invertible after inverting $p$. The summands therefore all have rank one. The exact Hodge quotient forces $Q_i=IP_i$ on selected characters and $Q_i=P_i$ on the others. The normal-decomposition matrix is monomial and invertible, so all arrows are units. Compatible generators lift through the Witt reductions. This gives the actual integral display required by DP40 at every such prime.

For the actual type, the selected full orbits occur only for $p\equiv1$, when $f=1$. Each is a height-one formal group with sharp slope zero. To check exactness directly, kill logarithm terms of degrees $2,\ldots,p-1$ by an integral strict coordinate. If $b_p$ is the next coefficient, height one makes $pb_p$ a unit, since the degree-$p$ coefficient of $[p]$ is $(p-p^p)b_p$.
The integral translation coefficient is $D^p/p!+b_pD$.
Hence $D^p$ reduces to a unit multiple of $D$, and every iterated order $p^r$ takes a unit value on a normalized coordinate. Integrality gives the opposite slope bound zero.
For the induced comparison type, the split factors are ordinary of multiplicative type. After a faithfully unramified strictly Henselian extension their \'etale Cartier duals are constant and their formal groups are multiplicative. The restricted operation is invertible; a CM eigendirection has unit $p^f$-power, again giving sharp slope zero. In the other classes the type alternates, so every selected gap is two and DP40 gives $c_2(p)$ in each direction.

\emph{The globally labelled spectrum.}
For $a=p\bmod7$ and $i\in\Phi$, the exact cyclic formula of DP40--DP43 is
\begin{equation}\label{eq:cmprimecycle}
 c_i^\Phi(p)
 =\sum_{r\ge1}\mathbf1_{a^ri\notin\Phi}p^{-r}
 =\frac{\sum_{r=1}^f\mathbf1_{a^ri\notin\Phi}p^{f-r}}{p^f-1}.
\end{equation}
It also gives zero on the selected ordinary orbits.
For $a=2$, use $(1,2,4)$ and $(3,6,5)$; for $a=4$, use their inverse orders. These give the two height-three rows of the table.
For $a=3$, the order $(1,3,2,6,4,5)$ has its first three positions selected; for $a=5$, use $(2,3,1,5,4,6)$. DP45's ordered triple $(p,p^2,p^3)/D(p)$ therefore becomes respectively $(p,p^3,p^2)/D(p)$ and $(p^3,p,p^2)/D(p)$ in the fixed global label order.
For $a=6$, the orbits $(1,6),(2,5),(3,4)$ each contain one selected index. This proves the whole table. The induced type is constant on its orbits for $a\in\Psi$ and alternating for $a\notin\Psi$, proving its formula as well.

As an exact common-variable identity for $p>1$,
\begin{equation}\label{eq:cmprimeformalaverage}
 \frac{4c_3(p)+2T(p)+3c_2(p)}6-\frac{9c_2(p)}6
       =\frac{4p^2(p^2+1)}{3(p^6-1)}.
\end{equation}
Indeed, over denominator $p^6-1$, the numerators of $c_3,T,c_2$ are
$p(p+1)(p^3+1)$, $p(p^2+p+1)^2$, and $p(p^4+p^2+1)$.
The numerator of $4c_3+2T-6c_2$ is $8p^2(p^2+1)$.
This formal averaging identity does not determine the sign of the actual prime-sum constant.

\emph{The unconditional prime sum.}
Set
\[
 \alpha_i(a)=\mathbf1_{ai\notin\Phi},\qquad
 \alpha^\Psi(a)=\mathbf1_{a\in\{3,5,6\}}.
\]
The cyclic series gives the uniform tail estimate
\[
 0\le c_i^\Phi(p)-\alpha_i(a)/p
       \le\sum_{r\ge2}p^{-r}=\frac1{p(p-1)},
\]
and the same estimate holds for the induced type.
For each fixed $i$, multiplication by $i$ permutes the six nonzero residues, exactly three of which land outside $\Phi$. Therefore
\[
 \sum_{a=1}^6\alpha_i(a)=3
       =\sum_{a=1}^6\alpha^\Psi(a).
\]
Use the established fixed-modulus prime-sum estimate
\begin{equation}\label{eq:cmprimeAP}
 \sum_{\substack{p\le X\\p\equiv a\ (7)}}\frac{\log p}{p}
       =\frac16\log X+O(1),\qquad a=1,\ldots,6.
\end{equation}
An exact direct reference is Rudnick, \emph{Primes in Arithmetic Progressions: Fixed Modulus}, May 26, 2015, \S2, p.~2, equation~(1); its preceding paragraph gives the unconditional quantitative fixed-modulus PNT.\footnote{Official author PDF: \url{https://www.math.tau.ac.il/~rudnick/courses/sieves2015/fixed\%20progressions.pdf}. The original page and fixed-modulus hypotheses were checked. The known PNT is cited, not reproved.}
The tail errors weighted by $\log p$ converge absolutely, by comparison with $\sum_{m\ge2}(\log m)/m^2$. Thus for each globally fixed direction
\[
 \sum_{\substack{p\le X\\p\ne2,7}}c_i^\Phi(p)\log p
       =\frac12\log X+O(1),\qquad
 \sum_{\substack{p\le X\\p\ne2,7}}c^\Psi(p)\log p
       =\frac12\log X+O(1).
\]
The finitely many bad primes affect only the bounded terms.
The three labelwise errors are uniformly bounded for all $X\ge2$.
Multiplying by $n_i$ and adding proves
\eqref{eq:cmprimeactual}--\eqref{eq:cmprimeinduced}, with constants independent of the orders. These orders may vary with an outer parameter; they must be fixed while the prime in the sum varies.

\emph{The integer periodic-floor certificates.}
Write
\[
 E_{i,p}^\Phi(n)
       =\sum_{r\ge1}\mathbf1_{a^ri\notin\Phi}
          \floor{n/p^r}.
\]
The induced certificate $E_p^\Psi(n)$ is zero in the split classes and
$\sum_{r\text{ odd}}\floor{n/p^r}$ in the inert classes.
These are the actual coefficientwise lower certificates of DP43; pointwise exactness is not asserted.
We prove, uniformly in $X$,
\begin{equation}\label{eq:cmprimefloors}
 \left|\sum_{\substack{p\le X\\p\ne2,7}}\log p
       \sum_{i=1}^3\bigl(E_{i,p}^\Phi(n_i)-E_p^\Psi(n_i)\bigr)\right|
       =O(N).
\end{equation}
Fix $i$ and put $n=n_i$. If $n\le1$, all floors vanish.
Otherwise put $Y=\min(X,n)$ and $\delta_i(a)=\alpha_i(a)-\alpha^\Psi(a)$.
The first-power contribution is
\begin{align*}
 &\sum_{\substack{p\le Y\\p\ne2,7}}
       \delta_i(a)\floor{n/p}\log p\\
 &\qquad=n\sum_{\substack{p\le Y\\p\ne2,7}}
       \delta_i(a)\frac{\log p}{p}
       +O\left(\sum_{p\le Y}\log p\right)=O(n).
\end{align*}
The first sum is bounded by \eqref{eq:cmprimeAP} because $\sum_a\delta_i(a)=0$; the last sum is $O(Y)=O(n)$ by the same fixed-modulus PNT summed over the six classes.
The absolute higher-power difference is at most
\[
 n\sum_p\log p\sum_{r\ge2}p^{-r}
       =n\sum_p\frac{\log p}{p(p-1)}=O(n).
\]
Adding over the three labels proves \eqref{eq:cmprimefloors}.
This checks rounding and the entire periodic pattern, not only its first real slope.
\end{proof}

For equal direction orders the actual first coefficients across the six classes are $(0,2,1,2,1,3)$, whereas the induced type gives $(0,0,3,0,3,3)$. Their sums are both nine. Some prime classes gain while others lose; the signed first-order terms require prime equidistribution, not an absolute same-prime summability argument.

The combinatorial reason has a precise general form. If a finite group $\Gamma$ acts transitively on $\Omega$ with $|\Omega|=2g$ and $|A|=g$, then
\[
 \frac1{|\Gamma|}\sum_{\sigma\in\Gamma}
       \#\{i\in A:\sigma i\notin A\}=\frac g2.
\]
For each fixed $i$, half the equally frequent targets lie outside $A$; fixed nonnegative weights $n_i$ give average $\tfrac12\sum_i n_i$ by the same count.
In the present abelian cyclotomic case, Frobenius is the actual residue element and \eqref{eq:cmprimeAP} supplies this distribution. For a nonabelian embedding action the displayed function need not be constant on conjugacy classes, and sampling primes of a reflex or other base field can change the distribution. The group average alone is not asserted to prove a general prime-average theorem for every CM variety.

Thus the concrete Jacobian does redistribute bare derivative divisibility into prime classes where the induced type has none, but for arbitrary orders assigned to its fixed global rational directions the linear and periodic-floor budgets change by only $O(N)$. Their leading logarithmic coefficient is unchanged. This does not identify them with normalized theta-selected values, reverse the source's one-way \path{LinkG} lower certificate, or permit prime-dependent reassignment of prescribed directions. A finite $O(N)$ improvement can still matter in a finite inequality. Special sections, selected harmonic weights, the new endpoint period and the actual normalization/cost terms remain additional data.


\subsection{Global sections realizing each finite sharp theta jet}

\begin{theorem}[DP54: finite-jet globalization at square tensor levels]\label{DP54}
Use the actual polarized Jacobian of DP49 over
$R=W(\mathbf F_{11^3})$, with fraction field $K$, and put
\[
 \mathcal L=(1,w_0)^*\mathcal P_\lambda,
 \qquad w_0=-(a+a^2+a^4).
\]
This trace-one theta line is relatively ample.
In the invariant product canonical gauge of DP47, fix an integral formal frame and write the normalized germ of a section $s$ of $\mathcal L^{\otimes\kappa}$ as
\[
 \Theta_s(z)=f_s(z)U(z)^\kappa B(z)^{-\kappa},
 \qquad U\in1+(z)R[[z]],\qquad B=\exp b.
\]
Here $B$ is the actual primitive arbitrary-unit Poincar\'e splitting of DP46.
Let $E$ be the inverse normalized group logarithm and put $\beta=b\circ E$.

For every $n\ge0$ there is $\kappa_0(n)$ such that for every integer
$\kappa\ge\kappa_0(n)$ some global integral section
$s_\kappa\in H^0(J,\mathcal L^{\otimes\kappa})$ satisfies
\begin{equation}\label{eq:globalfinitejet}
 \Theta_{s_\kappa}(E(Z))
 \equiv\exp(-\kappa\beta(Z))\pmod{(Z)^{n+1}}.
\end{equation}
Its algebraic coefficient is a unit at zero. The same section used in any fixed finite number of copies has the corresponding product pure-kernel jet through total degree $n$.

Fix finitely many linear functionals on finite jets of those copies whose pure-kernel values at exponent $-1$ are nonzero. For some integer $B_0$, every sufficiently large positive integer
\[
 M\equiv1\pmod{11^{B_0}},\qquad \kappa=M^2,
\]
admits such a section for which every selected value has exactly the valuation of its pure-kernel value at exponent $-1$.
For a homogeneous degree-$n$ differential functional, multiplication by the source factor $M^{-n}$ preserves that valuation.

Consequently the sharp ordinary primitive-line sequences of DP46, and the first-negative-row two-column harmonic sequence of DP51, are realized in valuation by global sections of varying square tensor powers of this actual theta line. Their rates are respectively $1/1220$ or $1/120$ for the two primitive-line classes, and $1/1220$ for the first-row harmonic family.
The sections may also be chosen over $\Q(\zeta_7)$, integral at the selected prime above $11$, preserving these finite valuations.
The section and tensor power may vary with the derivative order; no tensor-size, height, cost, or prescribed translated-family bound is asserted.
\end{theorem}

\begin{proof}
\emph{Relative ampleness with the correct sign.}
On the geometric closed fibre let $\mathcal T$ be an ample theta line inducing the canonical principal polarization. DP49 proves
\[
 \phi_{\mathcal L}=\phi_{\mathcal T}=\lambda.
\]
Thus $\mathcal N=\mathcal L\otimes\mathcal T^{-1}$ has
$\phi_{\mathcal N}=0$ and belongs to $\operatorname{Pic}^0$.
The known Picard facts needed here are Milne, \emph{Abelian Varieties} (1986), Proposition~9.2, pp.~117--118, the following definition of $\operatorname{Pic}^0$ on p.~118, and Proposition~10.1, p.~119: an ample line over an algebraically closed field has a surjective map $\phi_{\mathcal T}$ onto $\operatorname{Pic}^0$.
Choose $c$ such that
\[
 \mathcal N\cong t_c^*\mathcal T\otimes\mathcal T^{-1}.
\]
Then $\mathcal L\cong t_c^*\mathcal T$ is ample on the geometric closed fibre.
Ampleness of the fixed line descends along the ground-field extension, by Stacks Project Lemma~35.25.6, and closed-fibre ampleness spreads to an open neighbourhood by Lemma~30.21.4.
Every open neighbourhood of the closed point of $\operatorname{Spec}R$ is the whole spectrum. Therefore $\mathcal L$ is relatively ample, and is ample on $J$ because the base is affine. This uses equality with the positive principal polarization, not merely that $\phi_{\mathcal L}$ is an isogeny.\footnote{The exact primary locators are J.~S.~Milne, ``Abelian Varieties,'' in \emph{Arithmetic Geometry} (1986), pp.~117--119, official author copy \url{https://www.jmilne.org/math/articles/1986b.pdf}; Stacks Project Lemma~35.25.6 in \url{https://stacks.math.columbia.edu/tag/02LB}, and Lemma~30.21.4, \url{https://stacks.math.columbia.edu/tag/0D2N}.}

\emph{Integral finite-jet extension.}
Let $I$ be the ideal of the zero section and $Z_n=V(I^{n+1})$.
Smoothness and the formal parameters give
\[
 \mathcal O(Z_n)
 =R[[z_1,z_2,z_3]]/(z_1,z_2,z_3)^{n+1}.
\]
This is finite free over $R$. The formal frame restricts to a genuine frame on this finite thickening; no algebraization of the entire infinite frame is required.
The ideal sheaf $I^{n+1}$ is coherent, so Serre vanishing gives
\[
 H^1(J,I^{n+1}\otimes\mathcal L^{\otimes\kappa})=0
 \qquad(\kappa\ge\kappa_0(n)).
\]
The exact known form is Stacks Project Lemma~30.16.1(4), or its relative version Lemma~30.16.2.\footnote{\url{https://stacks.math.columbia.edu/tag/0B5T}, part~(4), and \url{https://stacks.math.columbia.edu/tag/02O1}.}
The restriction sequence therefore yields the integral surjection
\[
 H^0(J,\mathcal L^{\otimes\kappa})
 \twoheadrightarrow
 H^0(Z_n,\mathcal L^{\otimes\kappa}|_{Z_n}).
\]
Prescribe $f_s\equiv U^{-\kappa}\pmod{I^{n+1}}$ in the restricted frame and lift it. This target is integral and has value one at zero. Thus $s$ is nonzero and
$f_sU^\kappa\equiv1\pmod{I^{n+1}}$.
Substitution by $E$ preserves total order because it has zero constant term and identity tangent map. Multiplication by $\exp(-\kappa\beta)$ also preserves the ideal of terms of degree at least $n+1$. This proves \eqref{eq:globalfinitejet}.
In a product of finitely many copies, every nonconstant prefactor error already has total degree at least $n+1$, so the same section realizes the product jet.

\emph{Polynomial dependence on the tensor exponent.}
For a fixed finite jet functional $\ell$ of order at most $n$ on $r$ copies, put
\[
 F_\ell(T)=\ell\left(
 \exp\left[-T\sum_{b=1}^r\beta(Z^{(b)})\right]\right).
\]
This lies in $K[T]$: the logarithmic sum has zero constant term, so only its powers of order at most $n$ can contribute to the chosen jet.
Let $m_\ell$ be the minimum valuation of the nonzero coefficients of this polynomial. For an integer $T\equiv1\pmod{11^B}$,
\[
 v_{11}\bigl(F_\ell(T)-F_\ell(1)\bigr)\ge B+m_\ell,
\]
since $T^a-1$ is divisible by $11^B$ for every positive integer $a$.
Choose $B$ with
\[
 B+m_\ell>v_{11}(F_\ell(1)).
\]
Then $v_{11}(F_\ell(T))=v_{11}(F_\ell(1))$.
For finitely many selected nonzero values take the maximum of the finitely many required bounds.
There are arbitrarily large positive integers $M\equiv1\pmod{11^B}$; choose any with $M^2\ge\kappa_0(n)$ and set $T=M^2$.
The exact finite-jet lift then proves all the selected valuation equalities.
The source's definition \texttt{thExpr}, \texttt{IsThRule.lean}:66--70, uses tensor exponent $M^2$ and factor $M^{-\deg P}$. Here the latter is a unit, so it does not change the valuation.

\emph{Sharp sequences.}
For a primitive logarithmic line of the actual pair in DP49, DP46 gives
\[
 \liminf_{n\longrightarrow\infty}
 \frac1n v_{11}\bigl(n![t^n]\exp[-\beta(ut,vt)]\bigr)
 =c_k,
\]
where $c_2=1/1220$ if the first dual component is a unit and $c_1=1/120$ if the first is divisible by $11$ and the second is a unit.
Choose a sequence of nonzero values attaining this liminf and apply the finite-jet argument separately at each selected degree.
It produces global sections $s_{M_\nu^2}$ with exactly the same selected valuations, including the unit factor $M_\nu^{-n_\nu}$.

For the first dual row the two-column integral polynomials
\[
 P_{a,b}=\sum_{i=0}^{\min(a,b)}(-1)^i
 \binom ai\binom bi
 X_1^{a-i}Y_{1,1}^{b-i}X_2^iY_{1,2}^i
\]
are the source-shaped harmonic family of DP51.
Its proof applies to the actual arbitrary-unit display: DP46 retains the same unique initial monomial with unit coefficient, the gap to the other terms, and the primitive-path threshold used in that proof.
It supplies nonzero values of rate $1/1220$ as $a+b$ tends to infinity. Use two copies of the same finite-jet lift with $n=a+b$, and apply the same polynomial congruence argument in the tensor exponent. The resulting product of global normalized sections realizes each selected valuation.
This gives a sequence of sections and tensor powers, not the asymptotic sequence of a fixed section or a fixed power.

\emph{Number-field sections.}
Let $E_0=\Q(\zeta_7)$, choose the prime $\mathfrak p\mid11$ underlying the given embedding, and let $R_0=\mathcal O_{E_0,\mathfrak p}$. Its completion is $R$.
The curve, its Jacobian, $a$, $w_0$, the polarization, and $\mathcal L$ are defined over $R_0$ by the good model and formulas of DP49.
For a fixed tensor power already chosen, the module
\[
 V_0=H^0(J_{R_0},\mathcal L^{\otimes\kappa})
\]
is finite and torsion-free over the DVR $R_0$, hence free. Proper coherent cohomology gives finiteness; flatness of $J_{R_0}$ and invertibility of the sheaf give injectivity of multiplication by $11$, hence torsion-freeness.
Flat cohomology base change gives
\[
 V_0\otimes_{R_0}R
 =H^0(J_R,\mathcal L^{\otimes\kappa}).
\]
The exact locator is Stacks Project Lemma~30.5.2(2), Tag~02KH.\footnote{\url{https://stacks.math.columbia.edu/tag/02KH}.}
Density of $R_0$ in its completion makes $V_0$ dense in the right-hand side in a finite free basis. Each selected finite jet evaluation is a finite linear expression in the section coordinates; for a product of copies it is a finite polynomial expression. It is therefore continuous. A sufficiently close section in $V_0$ preserves every prescribed nonzero valuation and its unit value at zero.
This supplies a number-field section integral at $\mathfrak p$, with the selected valuations preserved. It need not preserve the exact formal coefficients, and it gives no height bound.
\end{proof}

The theorem removes the objection that the sharp prefactor-$1$ examples exist only as formal germs: each relevant finite jet is realized by global algebraic theta sections at sufficiently large square tensor levels. What remains is quantitative and specific to the prescribed family: bounds for the tensor level, the section heights and norms, and the realization inside the required translated family with its endpoint and nonvanishing inputs. No such bound, family membership, or shorter chain follows from this existence argument.


\subsection{Detection by the actual endpoint period}

\begin{proposition}[DP55: endpoint detection and a nonzero principal-level period]\label{DP55}
Assume the weak conjugation hypotheses and specified endpoint weights and
level of DP52(c). Let $S,S'$ be the actual point and row data, and
$\mathcal E=\operatorname{ptEv}$ their point evaluation object. Put
\[
 \mathcal F=\operatorname{span}_{\mathbf C}
    \{\mathcal E.\operatorname{emb}(\phi):
                    \phi\in S.\operatorname{EqFun}\}.
\]
This is precisely the source's $\mathcal E.\operatorname{Fi}$.
\begin{enumerate}[label=(\alph*),leftmargin=*]
\item The source row representing map
$J:\mathbf C^{\dim_K S'.\operatorname{EqFun}}\to\mathcal F$
is conjugate-linear and surjective under these weaker hypotheses.
The actual weakened endpoint period satisfies
\[
 \operatorname{per}(F,z)=\langle Jz,F\rangle_{\mathrm{cl}},
\]
where the classwise form is positive definite on $\mathcal F$.
Thus every nonzero $F\in\mathcal F$ has a row vector $z$ with
\[
 Jz=F,\qquad
 \operatorname{per}(F,z)=\langle F,F\rangle_{\mathrm{cl}}>0.
\]
This requires neither a free class action nor a change of the given level.
\item Add an auxiliary place $P$ above an odd rational prime $\ell$
unramified in the CM field $E$, outside the old finite endpoint support.
Set its level subgroup equal to $\{1\}$, with trivial character, retaining
the other level data.
The decorated-lattice action is then free. For every dominant endpoint
coefficient weight, the point and row spaces are nonzero and the actual
endpoint period has a pair with value $1$.
This applies to the actual type $\{1,2,3\}$ of $\mathbf Q(\zeta_7)$
at every admissible rebuilt even lift count.
\item A vector belongs to $\mathcal F$ exactly when its value at every
class lies in the coefficient-evaluation span and is invariant under the
class stabilizer. The source's adelic identification additionally requires
the explicit inputs $\operatorname{DecoPtSurj}$ and
$\operatorname{IsLevOf}$; that theorem has no
$\operatorname{Fields.Ok}$ argument.
\end{enumerate}
The proposition supplies no quantitative bound on the chosen row
coordinates or the cost of the auxiliary level, and does not construct
a nonzero image of the prescribed iterated theta kernel.
\end{proposition}

\begin{proof}
The source paths in this proof begin with
\path{Formalism/LocalToGlobal/Columns/}.
The actual point and row representations are respectively the round
and square coefficient representations, with inverse finite-level
characters. They have the same decorated-lattice set and identified
acting groups.
Their perfect bilinear coefficient pairing is
$B=\operatorname{pairBK}_\lambda$.
Its exact definition is the restriction of
\[
 B_\lambda(f,g)=
 \frac{\sum_\alpha f_\alpha g_\alpha\,\alpha!}
 {\prod_\tau\sum_\alpha
       (\operatorname{extremal}_{\lambda_\tau})_\alpha^2\,\alpha!}.
\]
The actual theorems \path{pairingBInvariant_holds} and
\path{pairingBNondegenerate_holds} in \path{CoefModB.lean},
188--203, give
\[
 B(\rho_{\mathrm{pt}}(g)x,\rho_{\mathrm{row}}(g)y)=B(x,y)
\]
and nondegeneracy. The inverse finite characters cancel.
Thus \path{TiedPer.lean}, \path{perDua}, 232--239, constructs the
required coefficient duality with no common-quadratic-subfield hypothesis.
Here \path{PairBNondeg} at \path{CoefModB.lean}, 88--90, is only the
proposition; the cited theorem at 193--203 is its proof.

Let $H$ be the point coefficient hermitian form.
For $\phi\in S.\operatorname{EqFun}$, perfectness gives a unique
row value $\psi(d)$ at each decorated lattice such that
\[
 B(x,\psi(d))=H(\phi(d),x)\quad\text{for all }x.
\]
This assignment is row-equivariant, since
\[
 \begin{aligned}
 B(x,\psi(gd))
 &=H(\rho_{\mathrm{pt}}(g)\phi(d),x)\\
 &=H(\phi(d),\rho_{\mathrm{pt}}(g)^{-1}x)\\
 &=B(\rho_{\mathrm{pt}}(g)^{-1}x,\psi(d))\\
 &=B(x,\rho_{\mathrm{row}}(g)\psi(d)).
 \end{aligned}
\]
Uniqueness in the second variable proves the assertion.
This is the complete argument of \path{GSD.Dua.partner} at
\path{TiedPer.lean}, 104--117; it uses DP52's explicit point-goodness
invariance.

Choose a $K$-basis $(b_a)$ of row equivariant functions and write
$\psi=\sum_a c_ab_a$.
For chosen point-class representatives $d_c$, the representing map is
\[
 (Jz)(c)=\sum_a\overline{z_a}\,
  \mathcal E.\operatorname{sharp}
  \bigl(x\longmapsto\sigma_0(B(x,b_a(d_c)))\bigr).
\]
The coefficient Gram matrix is positive definite, hence invertible:
this is exactly the matrix identified in the proof
\path{RecordInst.ptEv_gram_unit}, 47--58.
The operation \path{sharp} represents the displayed functional in the
first, conjugate-linear slot of the coefficient inner product.
Putting $z_a=\sigma_0(c_a)$ gives
$Jz=\mathcal E.\operatorname{emb}(\phi)$.
Indeed, the pairings of the two sides with every evaluated coefficient
vector agree, both sides lie in the coefficient-evaluation span, and
the Gram form is nonsingular there.

The displayed formula makes $J$ additive and conjugate-linear.
If $F=\sum_i t_i\mathcal E.\operatorname{emb}(\phi_i)$ and
$Jz_i=\mathcal E.\operatorname{emb}(\phi_i)$, then
$J(\sum_i\overline{t_i}z_i)=F$.
Conversely every $Jz$ is coefficient-valued and stabilizer invariant by
row equivariance and pairing invariance; the averaging argument below
puts it in $\mathcal F$.
This proves exactly the stated surjectivity.
The corresponding full source proofs are
\path{jOf_mem_Fi}, 133--148,
\path{jOf_partner}, 150--162, and \path{jOf_onto}, 172--189,
in \path{TiedPer.lean}.
Its \path{per_jMC_onto}, 250--252, uses
$h_X:X.\operatorname{Ok}$ only through \path{ptOk}.
Replacing that argument by DP52's proved point-goodness certificate
preserves this entire construction.

The classwise form is
\[
 \langle F_1,F_2\rangle_{\mathrm{cl}}
 =\sum_c\frac{\langle F_1(c),F_2(c)\rangle_{\mathrm{coef}}}
                {|\operatorname{Stab}(d_c)|}.
\]
Every denominator is a positive finite integer.
Each coefficient form is positive definite on its evaluation span,
so the sum is positive for $F_1=F_2=F\ne0$.
The finite-sum period identity is
\path{Record.periodC_eq_ipCl}, 1580--1590, or
\path{RecordInst.periodC_eq_ipCl'}, 60--64, with DP52's weakened
goodness input.
Taking $Jz=F$ proves (a).

For the membership assertion, every evaluated equivariant function is
coefficient-valued and stabilizer invariant, and so is its complex span.
Conversely, express each invariant value $F(c)$ in the evaluated
$K$-coefficient basis, retaining its complex scalar coefficients.
Average each $K$-basis vector over the finite stabilizer of $d_c$,
dividing by the group order.
Their evaluated complex combination remains $F(c)$ by invariance.
Each averaged vector is a $K$-coefficient vector fixed by the stabilizer,
so it extends equivariantly to the orbit and by zero to other orbits.
Summing these finitely many $K$-valued orbit functions with the original
complex coefficients gives $F\in\mathcal F$.
This is the full argument in \path{TiedValFi.lean}, 16--133.

The same file's \path{fi_iff_FiC}, 239--255, identifies this space with
adelic automorphic coefficient values under an actual CM datum,
the goodness certificate, \path{DecoPtSurj} and \path{IsLevOf}.
The last condition, at 151--154, identifies the level subgroup with
\path{levU} and its character with \path{levTyp}.
The theorem constructs an adelic function from an equivariant function
in one direction; in the other, the level transformation law gives
stabilizer invariance and the preceding averaging proof applies.
There is no \path{Fields.Ok} argument in this general theorem.
This proves (c), with precisely its stated remaining inputs.

It remains to prove the principal-level assertion.
First, if $R$ is the ring of integers of a finite unramified extension
of $\mathbf Q_\ell$, with $\ell$ odd, the group
$1+\ell M_n(R)$ has no nonidentity torsion.
Otherwise a power would have prime order $q$ and could be written
$A=I+\ell^aC$, where $a\ge1$ and some entry of $C$ is a unit.
For $q\ne\ell$,
\[
 A^q\equiv I+q\ell^aC\pmod{\ell^{a+1}M_n(R)}\ne I.
\]
For $q=\ell$,
\[
 A^\ell\equiv I+\ell^{a+1}C
                  \pmod{\ell^{a+2}M_n(R)}\ne I.
\]
Indeed, intermediate binomial terms have valuation at least
$1+2a\ge a+2$, while the last has $a\ell\ge a+2$.
Both conclusions contradict prime order.

Now choose the auxiliary place in (b); only finitely many rational primes
are excluded by ramification or the old support.
If $(g,h)$ stabilizes a decorated lattice $d$ with lattice $N$, then
$gN=N$.
The standard hermitian form is positive at every complex embedding.
Its rational trace norm is positive definite on the underlying
$\mathbf Q$-space.
An isometry sends each member of a $\mathbf Z$-basis of $N$ into the
finite set of lattice points with its prescribed norm.
Thus the lattice-isometry group is finite and $g$ has finite order.

The exact decoration action in \path{StdLatUnimodular.lean}, 108--111,
is $\operatorname{be}_P(x)\mapsto
h_P\operatorname{be}_P(g^{-1}x)$.
At the new place $h_P=1$.
Stabilization and the \path{be_add}, \path{be_ker} fields of
\path{Deco.Ok}, 81--94, give
\[
 (g^{-1}-1)N\subseteq PN.
\]
After completion at $P$, the lattice is free over the unramified local
ring, and $P$ is generated by $\ell$.
The matrix of $g^{-1}$ therefore belongs to the torsion-free group just
considered. Hence $g=1$ locally, and so globally.
At every other level place, stabilization now gives
$h_w\operatorname{be}_w(x)=\operatorname{be}_w(x)$.
Surjectivity of \path{be_w}, another \path{Deco.Ok} field, forces
$h_w=1$. The full stabilizer is trivial.
The row and point actions have the same underlying action, so both
are free.

The class set is nonempty, by the actual standard-lattice construction
\path{DecoSet.base} at \path{StdLatUnimodular.lean}, line 650.
The endpoint rank is positive.
The coefficient module is nonzero: its tensor of highest-weight
extremal polynomials is a product of nonzero polynomials in disjoint
variable blocks. The full source membership and nonvanishing proofs
are \path{CoefColLemmas.lean}, 320--342.
For a free action, arbitrary values at class representatives extend
uniquely by $\phi(gd_c)=\rho(g)\phi(d_c)$;
this is \path{GSetLemmas.resFree}, 221--225.

Choose a class $c_0$, a nonzero coefficient $x$, and by perfectness
$y$ with $B(x,y)=1$.
Assign $x$ to the point value at $d_{c_0}$ and zero at other classes.
Assign $y$ to the row value at that same representative and zero elsewhere.
Expand the resulting row function in its chosen basis and apply
$\sigma_0$ to its coordinates, obtaining $z$.
The exact identity \path{Record.perCl_emb}, 331--337, has just one
contribution, with stabilizer order one:
\[
 \operatorname{per}(\mathcal E.\operatorname{emb}\phi,z)
                    =\sigma_0(B(x,y))=1.
\]
This proves (b) for the actual period.
DP52 verifies all weak assumptions for the genus-three type;
any odd auxiliary prime other than $7$, outside its finite support,
is unramified and may be used.
\end{proof}

The endpoint period now detects every nonzero vector in the exact source
coefficient space. An auxiliary principal level also gives a nonzero
vector and an actual period value $1$ for every dominant coefficient weight.
What remains is to identify a nonzero image of the intended iterated theta
family in this space, with its stated adelic, level and kernel ties, and to
control the row and level costs in the same global comparison.
A finite global theta-jet realization on a Jacobian does not alone supply
that identification. The proposition establishes neither a shorter
global lift chain nor minimality of the old chain.


\subsection{Linear jet interpolation and subexponential tensor levels}

\begin{theorem}[DP56: quantitative tensor levels for sharp finite jets]\label{DP56}
Let $X$ be proper and smooth over a complete DVR $R$, let $e$ be a section with ideal $I$, and let $\mathcal L$ be relatively ample.
There are fixed integers $r\ge1$ and $b\ge0$ such that, for every $n\ge0$ and $\kappa\ge rn+b$,
\begin{equation}\label{eq:linearjetthreshold}
 H^0(X,\mathcal L^{\otimes\kappa})
 \twoheadrightarrow
 H^0\bigl(V(I^{n+1}),\mathcal L^{\otimes\kappa}|_{V(I^{n+1})}\bigr).
\end{equation}

For the actual theta line of DP54 at $p=11$, the sharp ordinary primitive-line sequences of DP46 and the first-row two-column harmonic sequence of DP51 can be realized by global integral sections of powers $\mathcal L^{\otimes\kappa_\nu}$ with
\[
 \kappa_\nu=M_\nu^2,\qquad p\nmid M_\nu,
 \qquad \log M_\nu=o(n_\nu),
 \qquad \log\kappa_\nu=o(n_\nu),
\]
where $n_\nu$ is the selected derivative degree.
The normalized selected values, including $M_\nu^{-n_\nu}$, have exactly the valuations of the corresponding pure primitive-kernel values at exponent $-1$.
The same tensor powers can be retained for number-field sections integral at the chosen prime, preserving those finite valuations.
These bounds concern tensor levels, not the archimedean sizes of the sections or their membership in a prescribed translated family.
\end{theorem}

\begin{proof}
\emph{A fixed integral projective coordinate system.}
Choose $r\ge1$ and an immersion
\[
 \iota:X\longrightarrow\mathbf P^N_R,
 \qquad \mathcal L^{\otimes r}\cong\iota^*\mathcal O(1).
\]
This is the known ample-power result, Stacks Project Lemma~29.40.4, Tag~01VT; properness makes the immersion closed.
Let $h_0,\ldots,h_N$ be the coordinate sections.
After choosing a generator of $e^*\mathcal L^{\otimes r}$, some $h_0(e)$ is a unit; rescale it to one.
Replace $h_i$ by $h_i-h_i(e)h_0$ for $i>0$.
The affine ratios $h_i/h_0$ now vanish along $e$ and generate its conormal module, because the chart is closed in affine space.
This conormal module is free of rank $g=\operatorname{rank}_R(e^*\Omega^1_{X/R})$.
Choose $g$ ratios whose cotangent matrix has a unit minor, by choosing a basis modulo the maximal ideal and using Nakayama's lemma.
Relabel their sections $h_1,\ldots,h_g$ and put $t_i=h_i/h_0$.
Their linear matrix is invertible over $R$, so formal inversion over $R$ identifies the formal neighbourhood with $R[[t_1,\ldots,t_g]]$.

In the frame $h_0^n$, any integral $n$-jet is a polynomial
\[
 Q(t)=\sum_{|\alpha|\le n}q_\alpha t^\alpha,
 \qquad q_\alpha\in R.
\]
It is the restriction of the global section
\[
 \sum_{|\alpha|\le n}q_\alpha
 h_0^{n-|\alpha|}h_1^{\alpha_1}\cdots h_g^{\alpha_g}
 \in H^0(X,\mathcal L^{\otimes rn}).
\]
All exponents are nonnegative and every coefficient is integral.

Choose $b$ so that $\mathcal L^{\otimes m}$ is globally generated for every $m\ge b$; this is Stacks Project Proposition~28.27.13(7), Tag~01Q3.
For $\kappa\ge rn+b$, a global section $\eta$ of $\mathcal L^{\otimes(\kappa-rn)}$ has unit value at $e$.
Express the desired jet in the frame $\eta h_0^n$ and multiply the displayed homogenization by $\eta$.
This proves \eqref{eq:linearjetthreshold}, including $n=0$.\footnote{The exact known projective facts used here are \url{https://stacks.math.columbia.edu/tag/01VT}, parts~(2) and~(5), and \url{https://stacks.math.columbia.edu/tag/01Q3}, part~(7).}
In particular the inverse unit correction $U^{-\kappa}$ in DP54 is an allowed target. The change from its specified integral frame to this projective frame is an integral unit, so the resulting normalized germ agrees with the pure exponent-$-\kappa$ kernel through degree $n$.

\emph{A coefficientwise estimate in the tensor parameter.}
In the mixed-characteristic application put $K=\operatorname{Frac}R$ and $s=1/(p-1)$, with $v_p(p)=1$.
Suppose a zero-constant series $\gamma(\mathbf Z)$ in $D$ additive variables satisfies
\begin{equation}\label{eq:tensorlogprofile}
 \inf_\alpha\bigl(v_p([\mathbf Z^\alpha]\gamma)+w|\alpha|\bigr)\ge s
\end{equation}
for a rational $w>0$, and put $c=s-w$.
For a nonzero integral homogeneous polynomial $P$ of degree $n$, set
\[
 F_P(T)=P(\partial)\exp(-T\gamma)(0)\in K[T],
 \qquad
 \sigma(P)=s\max_{\alpha\in\operatorname{supp}P}
                      \sum_{i=1}^D s_p(\alpha_i).
\]
Each coefficient of this polynomial satisfies
\begin{equation}\label{eq:tensorcoefficientbound}
 v_p([T^m]F_P)\ge cn-\sigma(P),
 \qquad
 \sigma(P)\le D(1+\lfloor\log_p n\rfloor)\quad(n\ge1).
\end{equation}

To prove the coefficient assertion, choose a finite extension containing $\varpi$ with $v_p(\varpi)=w$ and set
$\widetilde\gamma(\mathbf Z)=\gamma(\varpi\mathbf Z)$.
Every coefficient of $\widetilde\gamma$ has valuation at least $s$.
For every $m\ge1$, every coefficient of $\widetilde\gamma^m/m!$ has valuation at least
\[
 ms-v_p(m!)=s\,s_p(m)\ge0.
\]
Thus every coefficient in the additive variables of
$\exp(-T\widetilde\gamma)$ is a polynomial in $T$ with integral coefficients.
Each is a polynomial because $\widetilde\gamma$ has zero constant term.
It follows that every $\mathbf Z^\alpha T^m$ coefficient of $\exp(-T\gamma)$ has valuation at least $-w|\alpha|$.
Applying $P(\partial)$ multiplies this coefficient by the integral coefficient $P_\alpha$ and by $\alpha!$.
Legendre's identity gives
\[
 v_p(\alpha!)=s|\alpha|-s\sum_i s_p(\alpha_i),
\]
which proves the first inequality in \eqref{eq:tensorcoefficientbound}.
The digit estimate proves the second.
This is a coefficientwise polynomial bound; it has not been inferred from evaluations at integral $T$.

\emph{The actual logarithm profiles.}
For a primitive logarithmic line, the complete initial-form computation of DP46 gives
\[
 \inf_n\bigl(v_p([t^n]\beta(ut,vt))+w_kn\bigr)=s.
\]
Use $\gamma(t)=\beta(ut,vt)$, $D=1$ and $P(Z)=Z^n$.
For the first-row harmonic family, set
\[
 w_*=\frac{p^{h-1}}{(p-1)(p^{h-1}+1)},
 \qquad c_*=\frac1{(p-1)(p^{h-1}+1)}.
\]
At the $j=1$ endpoint profile of DP46, the scalar and first dual weights equal $w_*$; the later dual weights are
\[
 \omega_i=s-p^{i-1}(s-w_*)<w_*\quad(i>1).
\]
The contracted inverse logarithms are integral and the group-coordinate primitive logarithm has coefficient infimum $s$.
Giving every additive input the common weight $w_*$ supplies nonnegative extra factors in the later coordinates.
Composition with the integral inverse maps therefore proves \eqref{eq:tensorlogprofile} for the full primitive logarithm.
This is the logarithmic estimate in DP51's full-polydisc corollary, before exponentiation.
For two copies, their sum also satisfies it. Take $D=2h$, that sum as $\gamma$, and the first-row polynomial $P_{a,b}$ of DP51.
Arbitrary unit arrows retain the same coefficient valuations and the unit initial coefficient in these profile computations.

\emph{Strict integer congruences and a bounded choice of level.}
Fix one selected nonzero value of degree $n\ge1$ and put
\[
 V=v_p(F_P(1)),\qquad L=cn-\sigma(P),\qquad
 B=\max\{1,\lfloor V-L\rfloor+1\}.
\]
The coefficientwise bound implies $V\ge L$, so $B>V-L$ and $B\le V-L+1$.
For any integer $T\equiv1\pmod{p^B}$,
\[
 v_p(F_P(T)-F_P(1))\ge B+L>V,
\]
because $T^m-1$ is divisible by $p^B$ for each positive integer $m$.
Hence $v_p(F_P(T))=V$.
The strict floor-plus-one choice covers integral and nonintegral $V-L$ without an endpoint error.

Let $r,b$ be the fixed interpolation constants and define
\[
 m_0=\lceil\sqrt{rn+b}\rceil,\qquad
 q_0=\left\lceil\frac{m_0-1}{p^B}\right\rceil,
 \qquad M=1+q_0p^B.
\]
Here $n\ge1$ and $r\ge1$, so $m_0\ge1$ and $q_0\ge0$.
These definitions give
\[
 M\equiv1\pmod{p^B},\qquad
 M\ge m_0,\qquad
 M<m_0+p^B\le\sqrt{rn+b}+1+p^B.
\]
Set $\kappa=M^2$. Then $\kappa\ge rn+b$, so the integral interpolation realizes the pure exponent-$-\kappa$ jet, and $\kappa\equiv1\pmod{p^B}$ preserves its selected valuation $V$.
Since $p\nmid M$, the source factor $M^{-n}$ preserves it as well.

Along a sharp sequence, $V_\nu/n_\nu\to c$ and the fixed number of variables gives $\sigma(P_\nu)=O(\log n_\nu)$.
Consequently
\[
 0\le\frac{B_\nu}{n_\nu}
 \le\frac{V_\nu-cn_\nu+\sigma(P_\nu)+1}{n_\nu}
 \longrightarrow0.
\]
The right side is nonnegative by the coefficient bound, even when $V_\nu-cn_\nu$ is negative.
It follows that
\[
 0\le\frac{\log M_\nu}{n_\nu}
 \le\frac{\log(\sqrt{rn_\nu+b}+1+p^{B_\nu})}{n_\nu}
 \longrightarrow0.
\]
Multiplying by two proves $\log\kappa_\nu=o(n_\nu)$.

\emph{Number-field sections at unchanged powers.}
For each fixed selected $\kappa_\nu$, the flat-base-change and density argument of DP54 identifies the completed global-section module with the completion of the finite free lattice over $\mathcal O_{\Q(\zeta_7),\mathfrak p}$.
Approximate the exact integral interpolant sufficiently closely to preserve its selected nonzero finite valuation.
This does not change $\kappa_\nu$, so the tensor-level estimate remains valid for the resulting number-field section integral at $\mathfrak p$.
Exact equality of arbitrary $p$-adic jet coefficients over the number field is not asserted.
\end{proof}

The tensor level has thus been controlled along the sharp sequence: its logarithm is sublinear in the selected degree. The argument does not control the archimedean norm or height of the interpolated section, the cost of imposing the prescribed translated-family conditions, or the source's selected global row-degree sequence. Those quantities still require their own estimates before a shorter theta-lift chain can be concluded or excluded.


\subsection{An explicit CM repair and the actual adelic endpoint}

\begin{proposition}[DP57: unramified CM repair and endpoint realization]\label{DP57}
Let $E=\mathbf Q(\zeta _7)$, $F=E^+$, and let $J$ be the principally
polarized genus-three Jacobian of DP22 and DP49, of type $\{1,2,3\}$.
Put
\[
 E'=E(i)=\mathbf Q(\zeta _{28}),\qquad F'=(E')^+=F(\sqrt7).
\]
\begin{enumerate}[label=(\alph*),leftmargin=*]
\item The original $E$ fails the literal $\operatorname{IsCMUnram}$
condition used in the source's $\operatorname{CMData}$.
The pair $(F',E')$ satisfies $\operatorname{CMData}$ directly.
The extension $\mathcal O_{E'}/\mathcal O_{F'}$ is finite
\emph{étale}, and the field-level toroidal $\operatorname{Host}$
can also be supplied with $d=7$.
\item The product $A=J^2$ admits an integral $\mathcal O_{E'}$-action
whose Rosati involution for the product principal polarization is
canonical conjugation. Its CM type is
\[
 \Phi'=\{a\in(\mathbf Z/28)^\times:a\bmod7\in\{1,2,3\}\}
      =\{1,3,9,15,17,23\}.
\]
At $11$, its two $E'$-prime factors have heights $6,6$,
dimensions $4,2$, and slopes $2/3,1/3$.
The underlying product retains its height-three summands, but these
are not individual $E'$-stable factors.
\item For every positive Hermitian rank and dominant coefficient
weights, the generic source construction supplies an actual endpoint
level with $T=1$ and principal reduction level at
\[
 P=(29,\zeta _{28}-2),\qquad
 S_T=\varnothing,\quad S_G=\{P\}.
\]
For this fresh level, $\operatorname{DecoPtSurj}$ and
$\operatorname{IsLevOf}$ hold.
The source's finite coefficient-value space is exactly the restriction
of its adelic finite-weight space. It contains a nonzero vector with
an actual detecting partner period equal to $1$.
\item The type $\Phi'$ is induced from no imaginary quadratic
subfield. Thus this construction does not supply the old
$\operatorname{Datum.dom}$ or $\operatorname{CMObj.HasType}$
theta input for this geometric $A$.
The separate old wrapper condition $8\mid f$ also fails, since $f=6$;
neither restriction occurs in the generic endpoint construction in (c).
\end{enumerate}
The construction doubles the geometric dimension and changes the field
and level. Its product polarization has degree one.
No preservation of a prescribed earlier level, quantitative global-cost
bound, or membership in the prescribed iterated theta image is asserted.
\end{proposition}

\begin{proof}
The source's $\operatorname{CMData}$ is
\path{Formalism/LocalRingCardinality/DataSpaces.lean}, 14--19.
Its $\operatorname{IsCMUnram}$ field is explicitly
\[
 \begin{gathered}
 E\text{ totally complex, with conjugation compatible with embeddings},
 \\
 \forall Q\subset\mathcal O_E\text{ maximal},\
       \exists\theta\in\mathcal O_E:
                    \theta-\bar\theta\notin Q.
 \end{gathered}
\]
See \path{Formalism/HermitianLattices/StatementMass.lean}, 44--48.
In $\mathcal O_E=\mathbf Z[\zeta _7]$, the prime
$\mathfrak p=(1-\zeta _7)$ has residue field $\mathbf F_7$.
Conjugation is the identity modulo $\mathfrak p$, so every
$\theta-\bar\theta$ belongs to it. This proves the first failure.

For $E'$ use the explicit norm identity
\[
 N_{E/\mathbf Q}(\zeta _7-\zeta _7^{-1})
 =N(\zeta _7^{-1})N(\zeta _7-1)N(\zeta _7+1)
 =1\cdot7\cdot1=7.
\]
Thus $\theta=\zeta _7$ satisfies the required condition at every
prime of $E'$ not over $7$; at primes over $7$ use $\theta=i$,
whose difference from its conjugate is the unit $2i$.
The field is cyclotomic and totally complex, so the other
conjugation and real-fixed-field conditions hold.
For the exact cyclotomic integer-ring and prime-power assertions,
see J.\ S.\ Milne, \emph{Algebraic Number Theory}, version 3.08,
Proposition 6.2(b),(c), printed pp.~96--98,
and Theorem 6.4(a),(b), printed p.~99,
\url{https://www.jmilne.org/math/CourseNotes/ANT.pdf}.

Write
\[
 z=\zeta _7+\zeta _7^2+\zeta _7^4,\qquad
 z+\bar z=-1,\qquad z\bar z=2,\qquad
 \delta=1+2z,\quad\delta^2=-7.
\]
Then $i\delta$ is real with square $7$.
Since $[F:\mathbf Q]=3$, the field $F(\sqrt7)$ has degree six
and is the fixed field of $E'$.

The same witnesses prove literal unramifiedness of the integer rings.
At a finite prime $v$ of $F'$, let $R=(\mathcal O_{F'})_v$ and
choose $\theta=\zeta _7$ if $v\nmid7$, and $\theta=i$ otherwise.
It generates the quadratic extension, its trace and norm belong to $R$,
and
\[
 \operatorname{disc}(1,\theta)=(\theta-\bar\theta)^2\in R^\times.
\]
Hence $R[\theta]$ is finite étale over $R$ and is normal.
For the latter implication one may use Stacks Project,
Lemma 10.163.9, Tag 033C,
\url{https://stacks.math.columbia.edu/tag/033C}.
Its fraction field is $E'$, and every element integral over $R$
is integral over $R[\theta]$; normality therefore makes this order the
full integral closure. This proves étaleness at every finite prime.
In
\path{ReciprocityTheory/ReductiveGroups/ArithmeticGeometry/BirationalGeometry/ToroidalCompactification/TwoWeightEqOneAddRoot.lean},
14--26, take
\[
 d=7,\qquad rt=-\delta,\qquad w_0=-z,\qquad
 rt_0=\sigma_0(rt),
\]
with canonical conjugation. Then $rt^2=-7$ and $2w_0=1+rt$,
and the just-proved étaleness supplies $\operatorname{Host.unram}$.
This completes the field-level construction in (a).

For (b), let $E$ act diagonally on $A$ and let
\[
 I=\begin{pmatrix}0&-1\\1&0\end{pmatrix}
\]
act on the two factors. It commutes with $E$ and has square $-1$.
The unique quadratic subfield of $E$ is $\mathbf Q(\sqrt{-7})$,
so $i\notin E$ and this gives an injective action of $E(i)$.
With compatible roots,
\[
 \zeta _7=\zeta _{28}^4,\quad i=\zeta _{28}^7,\quad
 \zeta _{28}=\zeta _7^2i^{-1},\qquad
 \mathcal O_{E'}=\mathbf Z[\zeta _7,i].
\]
Both generators act integrally, so this is the full integer-ring
action without any isogeny.
Rosati for the product polarization sends $I$ to
$I^{\mathsf T}=-I$ and $a\in E$ to $\bar a$.
It is therefore canonical conjugation on $E'$.
Equivalently $-\delta I$ is self-adjoint and has square $7$.
The polarization remains principal; the degree $7^6$ arising from
the alternative presentation with polarization
$\operatorname{diag}(1,7)$ is not a necessary cost.

On each old tangent-character multiplicity space, $I$ has the
two distinct eigenvalues $i,-i$. Thus each old selected embedding
has both extensions selected, giving the stated $\Phi'$.
The Frobenius orbits at $11$, with their selected indicators, are
\[
 \begin{array}{c|c}
 (1,11,9,15,25,23)&(1,0,1,1,0,1)\\
 (3,5,27,17,19,13)&(1,0,0,1,0,0).
 \end{array}
\]
The dimensions are four and two; the underlying factors are the
squares of the two height-three factors of $J$ from DP22 and DP49.
This proves their stated heights and slopes.
Each local $E'$-field at $11$ is unramified of degree six.
After splitting its action on a rational Dieudonné module,
Frobenius cyclically identifies the six eigenspaces. They have
equal dimensions, so any nonzero factor stable under this field
has height divisible by six. This rules out treating the old
height-three product summands as individual $E'$-stable factors.

For (c), $2\in\mathbf F_{29}^\times$ has order $28$:
$2^{14}=-1$ and $2^4=16\ne1$.
Thus evaluation $\zeta _{28}\mapsto2$ gives the displayed prime
with residue field $\mathbf F_{29}$.
All primitive $28$-th roots are simple modulo $29$, so its
completion is $\mathbf Q_{29}$.
Conjugation replaces $2$ by $2^{-1}=15$, proving $P\ne\bar P$.
The definitions \path{IsConjPl} and \path{IsSplitSet} in
\path{Formalism/LocalToGlobal/Divisibility/DefGroupLemmas.lean},
23--27, therefore give
$\operatorname{IsSplitSet}(S_T\cup S_G)$.
Disjointness is automatic since $S_T$ is empty.

Hereafter the source paths beginning with \path{Columns/} are under
\path{Formalism/LocalToGlobal/}.
The theorem \path{Columns/InputsGenus.lean},
\path{unimodGenus_holds}, 17--19, supplies
$\operatorname{UnimodGenus}(hc,n)$ for the CM datum just constructed
and every $n>0$.
The theorem \path{reduction_surjective_bis} in the preceding
\path{DefGroupLemmas.lean}, 1475--1481, supplies
$\operatorname{RedSurj}$.
Its definition at 64--68 allows arbitrary invertible reductions
at an $\operatorname{IsSplitSet}$ for a framed definite unimodular lattice.
The theorem \path{Columns/AdelicLemmas.lean},
\path{decoPtSurj_of}, 1245, consumes exactly these two statements,
the split-set condition and disjointness.
It proves the actual $\operatorname{DecoPtSurj}$ input.

Take the finite reduction subgroup $H_P=\{1\}$ with trivial
character. In the generic constructor
\path{Columns/RecLevelStep.lean}, \path{stepOf}, 20--51,
set
\[
 T=1,\qquad \pi=1,\qquad f=6,\qquad S_T=\varnothing.
\]
Divisibility by $\pi$ is automatic, $\pi\bar\pi=T$ is $1=1$,
and $1$ is in no prime. The conditions at the empty set and at
primes dividing $T$ are vacuous.
The standard lattice is self-dual, so its required relative index
is $1=T^{2f}$.
These are precisely the inputs already used by
\path{Columns/TiedXp.lean}, \path{xp_basePi_one}, 21--23,
and \path{tieLevP}, 35--41; the generic constructor accepts the
present arbitrary singleton $S_G$ without an
$\operatorname{Fields.Ok}$ argument.
Its fields are literally
\[
 L.U=\operatorname{levU}(hc,n,\varnothing,S_G,\Lambda.\operatorname{grp}),
 \qquad
 L.\chi=\operatorname{levTyp}(\sigma_0\circ\Lambda.\operatorname{chr}),
\]
at \path{RecLevelStep.lean}, 46 and 50.
These are exactly the two fields of
$\operatorname{IsLevOf}$ at \path{Columns/TiedValFi.lean}, 151--154.
Both equalities follow by definition, just as in that file's
\path{tieLev_isLevOf}, 265--266.
The resulting level is compact and open, principal modulo $29$
at the selected split place and standard hyperspecial elsewhere.

Choose $K=E'$ and the embeddings $\Phi'$ as the coefficient data.
They satisfy the conjugation compatibilities of DP52.
Its weak-goodness proof supplies
$\operatorname{formGSD.Ok}$, with no
$\operatorname{Fields.Ok.quad}$ requirement.
All hypotheses of \path{Columns/TiedValFi.lean},
\path{fi_iff_FiC}, 239--255, are now supplied:
goodness, $\operatorname{DecoPtSurj}$ and
$\operatorname{IsLevOf}$.
That theorem identifies $\operatorname{Fi}$ with the restrictions
of actual adelic finite-weight forms.
Its forward direction uses \path{adOf}, 179--181, and
\path{adOf_mem}, 185--201; its reverse direction uses
the precise stabilizer-invariance theorem at 209--237.

At the chosen prime the principal congruence group is
$1+29M_n(\mathbf Z_{29})$, which is torsion-free by DP55.
Every decorated positive-definite lattice has finite stabilizer,
so the principal decoration makes that stabilizer trivial.
The standard decorated lattice and the nonzero extremal coefficient
give a nonzero finite value supported on a free class.
It is the restriction of an actual adelic form by the preceding
identification. DP55's perfect partner pairing and period formula
then produce a partner with period $1$.
This proves (c) at the fresh level constructed here, rather than
asserting preservation of some earlier fixed level.

Finally, $3,15\in\Phi'$ and
\[
 \langle3\rangle=\{1,3,9,27,25,19\},\qquad
 15\notin\langle3\rangle,\qquad15^2=1.
\]
They generate the full twelve-element Galois group.
If $\Phi'$ restricted to one embedding of an imaginary quadratic
subfield, these automorphisms, compared with the identity in
$\Phi'$, would fix that subfield. It would therefore be contained
in $\mathbf Q$, a contradiction. The common restriction field
$K_q$ is consequently $\mathbf Q$.
The source's uniform two-type condition is
\path{Development/AbelianVarietyTheory/Definitions/CMLattices/IsDiffOf.lean},
\path{HasType}, 64--67, with
$\operatorname{CMObj.typeP,typeM,typeB}$ at 115--122.
Its matching trace condition is $\operatorname{Datum.dom}$ at
\path{TwoWeightEqOneAddRoot.lean}, 43--49.
The elliptic Serre tensor tie $\operatorname{IsTen}$ at
\path{IsDiffOf.lean}, 73--77, is not furnished by an adelic
endpoint identification.
The separate $8\mid f$ requirement is
\path{Columns/AbelianHostFieldConditions.lean}, 110--120,
and fails for $f=6$.
These facts establish (d).
\end{proof}

The unramified CM datum, the field-level toroidal host, the actual
adelic endpoint and a nonzero detecting period have all been constructed.
The remaining mixed-type input is the particular theta family and its
coefficient, path and nonvanishing identities, with the corresponding
global costs. The doubled geometry and its old height-three product
factors do not furnish those identities or a shorter lift chain by
themselves.


\subsection{Explicit sharp degrees and linear square tensor levels}

\begin{theorem}[DP58: explicit unit-parameter witnesses]\label{DP58}
Use the actual arbitrary-unit pair of DP46, with $p$ odd and $h\ge3$, and its primitive logarithmic splitting $\beta$ in additive coordinates and a normalized invariant integral frame.
For a primitive mixed line $(x,\mathbf y)=(ut,\mathbf vt)$, with $u\in R^\times$ and $\mathbf v$ primitive, let $j$ be the first unit dual component and put
\[
 k=h-j,\qquad N=p^k,\qquad d=N+1,\qquad
 s=\frac1{p-1},\qquad
 w=\frac{sN}{N+1},\qquad c=\frac{s}{N+1}.
\]
For every unit $\lambda\in R^\times$ write
\[
 \exp\bigl(\lambda\beta(ut,\mathbf vt)\bigr)
 =\sum_n A_n(\lambda)t^n.
\]
At every degree $n_r=(p^k+1)p^r$, $r\ge0$, one has
\begin{align}
 v_p(A_{n_r}(\lambda))
 &=-\frac{p^{k+r}-1}{p-1},\label{eq:explicitkernelcoefficient}\\
 v_p(n_r!A_{n_r}(\lambda))
 &=\frac{p^r-1}{p-1}=cn_r-s.\label{eq:explicitkerneljet}
\end{align}
These valuations are independent of the unit $\lambda$.

For an eigenaxis $j$, put
$F_\lambda(x,y)=\exp(\lambda\beta(x,e_jy))$ and use the two-column polynomial $P_{a,b}$ of DP51.
Its value on $F_\lambda(x_1,y_1)F_\lambda(x_2,y_2)$ at zero satisfies
\begin{equation}\label{eq:explicitharmonicjet}
 v_p(H_{p^r,p^{k+r}}(\lambda))
 =\frac{p^r-1}{p-1}.
\end{equation}
For $j=1$ these are the source-shaped first-negative-row $\mathrm{HaO}$ polynomials; for other $j$ the assertion is only about the stated summed contraction unless a permitted row identification is supplied.

On the actual Jacobian of DP49 at $11$, every finite prefix of these explicit witnesses can be realized by global sections at square tensor powers satisfying
\[
 \kappa=M^2=An+O(\sqrt n),\qquad 11\nmid M,
\]
where $n$ is the largest selected degree and $A$ is a fixed interpolation constant from DP56.
The selected valuations include the source normalization $M^{-n}$, and the same tensor powers may be retained for number-field sections integral at the selected prime.
No section-height or prescribed-family conclusion is included here.
\end{theorem}

\begin{proof}
\emph{The full initial form of a critical exponential.}
Let $K'$ be a complete discretely valued extension with $v(p)=1$, containing $\mu$ with $v(\mu)=s$.
Suppose a zero-constant formal series $z$ has coefficients of valuation at least $s$ and its entire lowest-valuation form is
\[
 \overline{\mu^{-1}z}=u\mathbf Z^\alpha,
 \qquad u\ne0,\quad |\alpha|>0.
\]
Set $Q=\mu^{-1}z$.
The scalar in the $m$th term of
\[
 \mu^{-1}(\exp z-1)
 =\sum_{m\ge1}\frac{\mu^{m-1}}{m!}Q^m
\]
has valuation
\[
 (m-1)s-v_p(m!)=s(s_p(m)-1)\ge0.
\]
It is a unit precisely for $m=p^r$.
Reduction therefore gives the exact formal identity
\begin{equation}\label{eq:criticalexponentialinitial}
 \overline{\mu^{-1}(\exp z-1)}
 =\sum_{r\ge0}\epsilon_r u^{p^r}\mathbf Z^{p^r\alpha},
 \qquad
 \epsilon_r=\overline{\frac{\mu^{p^r-1}}{(p^r)!}}\ne0.
\end{equation}
Each coefficient before reduction is a finite sum because $Q$ has zero constant term.
The exponents on the right are distinct and their degrees tend to infinity.
Thus no convergence of an infinite sum in one fixed degree is being assumed, and no displayed leading coefficients cancel.
The hypothesis concerns the entire initial form; a single first ordinary-degree coefficient would not suffice.

\emph{Ordinary primitive lines.}
At the critical weight, the full initial-form calculation of DP46 gives a finite ramified scale $\varpi$ with $v(\varpi)=w$ and
\[
 S=\frac{N-1}{p-1},\qquad
 \mu=\frac{\varpi^{N+1}}{p^S},\qquad v(\mu)=s,
\]
such that
\[
 \overline{\mu^{-1}\beta(u\varpi T,\mathbf v\varpi T)}
 =\xi T^{N+1},\qquad \xi\ne0.
\]
This is the full initial form after the actual inverse logarithms; its scalar is the unit arrow product times
$(-1)^{k-1}\bar u\bar v_j^{p^k}$.
The invariant integral frame corrections are strictly higher by the invariant-frame part of DP46.
Apply \eqref{eq:criticalexponentialinitial} to $z=\lambda\beta(u\varpi T,\mathbf v\varpi T)$.
Every coefficient at degree $(N+1)p^r$ in the scaled exponential has valuation exactly $s$.
Hence
\[
 v_p(A_{n_r})=s-wn_r
 =-\frac{p^{k+r}-1}{p-1},
\]
which proves \eqref{eq:explicitkernelcoefficient}.
Since $n_r=p^{k+r}+p^r$ has digit sum two,
\[
 v_p(n_r!)=s(n_r-2),
\]
and addition gives \eqref{eq:explicitkerneljet}.
Every unit $\lambda$ multiplies the initial scalar by a nonzero residue, so the valuations are uniform over all such parameters.

\emph{The harmonic diagonal factor.}
With independent inputs, the same complete calculation gives
\[
 \overline{\mu^{-1}\lambda\beta(\varpi X,e_j\varpi Y)}
 =\xi_\lambda XY^N,\qquad \xi_\lambda\ne0.
\]
This is the two-variable initial-form calculation in DP51: the reduced scalar inverse is $X$, and the reduced last dual inverse is a unit times $Y^{p^{h-j-1}}$.
For $\widetilde F(X,Y)=F_\lambda(\varpi X,\varpi Y)$, the exponential identity yields
\[
 \overline{\mu^{-1}(\widetilde F-1)}
 =\sum_{r\ge0}\epsilon_r\xi_\lambda^{p^r}
       X^{p^r}Y^{Np^r}.
\]
Write $f_{a,b}=[x^ay^b]F_\lambda$ and set
\[
 D(Z)=\sum_{i\ge0}f_{i,i}Z^i,
 \qquad C(x,y)=F_\lambda(x,y)D(-xy).
\]
As $N>1$, the displayed initial form has no positive diagonal monomial.
Thus $v_p(\varpi^{2i}f_{i,i})>s$ for every $i\ge1$.
The values lie in one discrete value group, so the gap is uniform.
Consequently $\mu^{-1}(D(-\varpi^2XY)-1)$ is integral with zero reduction.
Multiplication by the integral series $\widetilde F$ preserves this zero reduction, and therefore
\[
 \overline{\mu^{-1}(C(\varpi X,\varpi Y)-1)}
 =\sum_{r\ge0}\epsilon_r\xi_\lambda^{p^r}
       X^{p^r}Y^{Np^r}.
\]
In particular $v_p(C_{p^r,Np^r})=s-w(N+1)p^r$.

The exact harmonic coefficient identity is
\begin{align*}
 H_{a,b}(\lambda)
 &=\sum_i(-1)^i\binom ai\binom bi
 (a-i)!(b-i)!(i!)^2 f_{a-i,b-i}f_{i,i}\\
 &=a!b!\,[x^ay^b]C(x,y),
\end{align*}
because the factorial factor in every summand equals $a!b!$.
This includes all the alternating terms.
At $(a,b)=(p^r,Np^r)$ both digit sums are one, so
\[
 v_p(a!b!)=s(a+b-2).
\]
Together with the coefficient of $C$ just computed, this proves \eqref{eq:explicitharmonicjet}.
The row-degree ratio is exactly $1:N$.

\emph{Linear square tensor powers.}
On the actual Jacobian at $11$, let $A\ge1,B_0\ge0$ be fixed constants such that every $n$-jet extends in every tensor power at least $An+B_0$, by DP56.
Put $m_0=\lceil\sqrt{An+B_0}\rceil$, choose $M=m_0$ unless $11\mid m_0$, and in that case choose $M=m_0+1$.
Then
\[
 11\nmid M,
 \qquad An+B_0\le M^2
 <An+B_0+4\sqrt{An+B_0}+4.
\]
At the chosen power $\kappa=M^2$, interpolate the algebraic prefactor to $U^{-\kappa}$ through degree $n$, as in DP54.
The normalized germ then has the pure kernel jet with unit exponent $\lambda=-\kappa$.
The exact formulas above hold for every unit exponent, so no further congruence on $M$ is necessary.
The source factor $M^{-n}$ is a unit.

The explicit ordinary degrees at $11$ are $12\cdot11^r$ and $122\cdot11^r$, each with derivative valuation $(11^r-1)/10$.
For the first-row harmonic family they are $(a,b)=(11^r,121\cdot11^r)$, with that same valuation.
Interpolating to the largest degree of any finite prefix realizes every earlier witness as well, with the same section and unit exponent.
For two columns use two copies of that section.
DP54's number-field approximation preserves all these finitely many nonzero values in valuation at the same chosen tensor power.
\end{proof}

This theorem concerns pure primitive kernels and global sections with deliberately prescribed finite prefactor jets. It does not assert these exact-degree values for every integral prefactor, does not realize infinitely many witnesses with one fixed section, and does not control the sections' archimedean norms or identify the prescribed translated family. The primitive projections and unit tensor parameter remain essential hypotheses; the nonunit regime of DP50 is separate.


\subsection{The integral span of permitted theta translates}

\begin{theorem}[DP59: permitted translates contain explicit sharp witnesses]\label{DP59}
Use the actual Jacobian and trace-one line of DP49 over \(R=W(\mathbf F_{11^3})\): \[
 \begin{gathered}
 A=\operatorname{Jac}(y^2=x^7-1),\qquad g=3,\qquad \delta^2=-7,\\
 \delta=1+2(a+a^2+a^4),\qquad w_0=(1-\delta)/2,\\
 L=(1,w_0)^*\mathcal P_\lambda.
 \end{gathered}
\] The principal polarization is \(\lambda\), and \(\delta^\dagger=-\delta\). Choose a generator \(s\) of the rank-one module \(H^0(A,L)\) obtained from a section over \(\mathbf Q(\zeta_7)\), primitive at the chosen prime. Its existence is proved in the arithmetic-realization part of the proof. The local span conclusions would also hold after multiplying it by any unit of \(R\). Set \[
 \alpha_0=1+2\delta=3-4w_0,\qquad
 \alpha=\alpha_0^e,\qquad M=29^e,\qquad \kappa=M^2,
 \qquad e\ge1.
 \tag{TS1}
\] Then \(\alpha\bar\alpha=M\), and the ideals \((\alpha)\) and \((\bar\alpha)\) are relatively prime. All these levels are odd and prime to 11. The ideal \(\mathfrak b=(\alpha)\) satisfies the source's arithmetic \texttt{Level} conditions involving \(M,\mathfrak b,w_0\).

After a finite unramified extension \(R'/R\), all torsion points needed below and their symmetric theta lifts are defined. Let \[
 H=A[\alpha^2](R'),\qquad D=A[M](R'),\qquad C=A[\alpha](R').
 \tag{TS2}
\] Write \(s_M=[M]^*s\) as a section of \(L^{M^2}\) using its canonical rigidified multiplication isomorphism. For every \(t'\in H\), let \[
 s_{t'}=U_{t'}s_M,
 \tag{TS3}
\] where \(U_{t'}\) is the symmetric theta lift of order dividing \(M^2\).

\textbf{DP59(a).} The integral span of these sections is exactly \[
 W_M=[\alpha]^*H^0(A_{R'},L^M)
       \ \subset H^0(A_{R'},L^{M^2}).
 \tag{TS4}
\] Its rank is \(M^g\), while the ambient module has rank \(M^{2g}\). Choosing one representative of each coset in \(H/C\) gives an \(R'\)-basis of \(W_M\), up to multiplying individual representatives by units. In particular the allowed half-torsion orbit does not span the entire ambient section module when \(M>1\).

\textbf{DP59(b).} There are fixed integers \(r\ge1,b\ge0\) such that, whenever \(M\ge rn+b\), the module \(W_M\) surjects onto the integral jets of order \(n\) of \(L^{M^2}\) at zero. For each explicit ordinary or first-row two-column harmonic degree in DP58, some allowed translate, or some pair of allowed translates, has exactly the DP58 normalized valuation \[
 V=\frac{11^q-1}{10}.
 \tag{TS5}
\] Here the ordinary degrees are \(n=(11^k+1)11^q\), with \(k=1\) or \(2\); the harmonic pair is \((a,b)=(11^q,121\cdot11^q)\), of total degree \(n=122\cdot11^q\). The differential operator is the one specified in DP58, and the factor \(M^{-n}\) is included.

The level may be chosen among the powers of 29 with \[
 M=O(n),\qquad \kappa=M^2=O(n^2).
 \tag{TS6}
\] The sections in (TS3), their torsion points and their theta lifts are algebraic over finite number-field extensions and integral at the chosen prime over 11. The completion of that extension at this prime is unramified.

\textbf{DP59(c).} Each section (TS3) satisfies the generic source relation \texttt{IsVtOver} with \(N=M\), base section \(s\), allowed root \(t'\), point \(t=[M]t'\), and condition ``annihilated by \(\mathfrak b^2\).'' The canonical multiplication isomorphism satisfies \texttt{IsNormMul}. This does not construct a point of a source \texttt{Moduli\ Dw\ l} with the old prescribed \texttt{Datum.dom}/\texttt{HasType}, nor identify these sections with the nonzero iterated theta image. Those remain distinct interfaces.
\end{theorem}

\paragraph{Exact inputs and source interfaces.}
The classical representation input is Mumford, \emph{On the Equations Defining Abelian Varieties I}, Inventiones 1 (1966), §1: the basic cohomology facts (II), printed p.289; Theorem 1, pp.293--294, on the nondegenerate theta commutator; Proposition 3, pp.295--296, on the weight-one representation; and Theorem 2, p.297, that \(H^0(A,L)\) is irreducible for an ample line of separable type. These known results are cited, not asserted as new. The integral lattice, restricted-orbit and jet conclusions below are proved here. Primary copy: \url{https://pazuki.perso.math.cnrs.fr/index_fichiers/Mumford66.pdf}. The actual original pages 289 and 297 were visually read; the neighboring statement/proof text was read as well.

The precise recovered definitions are:

\begin{itemize}
\item
  \path{Formalism/LocalToGlobal/ProjectiveLattices/IsThRule.lean}: \texttt{IsVtOver},32--42; \texttt{thExpr},66--70; \texttt{ThValAt},72--79.
\item
  \path{ReciprocityTheory/ReductiveGroups/ArithmeticGeometry/ExtensionDescent/ThetaPt.lean}: \texttt{SymLift},44--51; \texttt{IsTwistPull},73--81; \texttt{IsActOf},83--84.
\item
  The same directory, \path{ThetaVV.lean}: \texttt{IsNormMul},103--108.
\item
  \path{ReciprocityTheory/ReductiveGroups/ArithmeticGeometry/BirationalGeometry/ToroidalCompactification/TwoWeightEqOneAddRoot.lean}: \texttt{Datum.dom},41--51; \texttt{Level},53--62; \texttt{SAB},72--83.
\end{itemize}

\texttt{SymLift} demands the specified point, the stated finite power equal to the identity, and compatibility with a symmetry of the line. \texttt{IsVtOver} permits roots killed by the specified ideal, not all points of \(A[M^2]\). That distinction is essential to (TS4).

\begin{proof}
\emph{Symmetry, multiplication, and the split principal level.}
The normalized Poincaré line is a biextension. Applying multiplication by \(M\) in both arguments gives the rigidified isomorphism \[
 [M]^*L\simeq L^{M^2}.
 \tag{TS7}
\] Applying minus one in both arguments similarly gives \([-1]^*L\simeq L\). These isomorphisms preserve the rigidification at the identity, by the unit compatibility of the biextension. In particular the actual trace-one line is symmetric; no arbitrary principal-line symmetry assumption or unaccounted Picard-zero twist is used.

For a CM endomorphism \(u\) commuting with \(w_0\), Poincaré functoriality and the definition of the Rosati adjoint give \[
 (u,w_0u)^*\mathcal P_\lambda
 \simeq (1,u^\dagger w_0u)^*\mathcal P_\lambda.
\] Since \(\alpha^\dagger\alpha=M\), this gives, rigidified at zero, \[
 [\alpha]^*L\simeq L^M,
 \qquad [\bar\alpha]^*L\simeq L^M,
 \qquad [\alpha]^*L^M\simeq L^{M^2}.
 \tag{TS8}
\] Composing these isomorphisms agrees with (TS7), again by biextension functoriality.

The norm identity follows directly from \(\delta^2=-7\): \((1+2\delta)(1-2\delta)=29\). The ideals generated by \(\alpha_0\) and \(\bar\alpha_0\) are relatively prime: a common divisor would divide both their sum2 and their product29; alternatively \(1=15(\alpha_0+\bar\alpha_0)-\alpha_0\bar\alpha_0\). Their powers remain relatively prime. The generator \(\alpha_0=3-4w_0\) lies in \(\mathbf Z[w_0]\), so \((\alpha)\) is generated by elements of that subring. This proves exactly \texttt{fromK}; the norm and coprimality identities prove \texttt{prod} and \texttt{coprime}. Also \(M\ge29\) and \(M\) is odd; the auxiliary positive \texttt{A} of the \texttt{Level} record may be chosen1. No claim about the different \texttt{Datum} fields follows from these arithmetic checks.

For completeness, the degree of \([\alpha]\) is \(M^g\). Over the complex numbers the integral homology lattice has rank one over the CM field \(\mathbf Q(\zeta_7)\) after tensoring with \(\mathbf Q\); multiplication by \(\alpha\) has rational determinant \(\operatorname{Norm}_{E/\mathbf Q}(\alpha)=M^3\). Its lattice index is the degree of the corresponding complex torus covering, so the degree is \(M^g\). These kernels at the chosen good prime are finite étale: they are direct factors, by the coprime-ideal identities, of torsion groups killed by powers of29. Their degrees therefore agree on both fibres.

All the finite étale torsion becomes constant after a finite unramified extension of the complete DVR. The finite theta isomorphisms used below can be defined after further such extensions. Each fibre of the theta group over a rational torsion point is a multiplicative-group torsor over a local ring, hence is trivial because an invertible module over a local ring is free. Thus it has an integral lift. All scalar comparisons between integral lifts are units.

\emph{Symmetric finite-order lifts.}
Let \(N=M^2\), which is odd. Fix the normalized symmetry of \(L^N\), and let \(\iota\) be its induced automorphism of the theta group. It fixes central scalars and induces \(t\mapsto-t\) on the torsion quotient. An integral lift \(U\) over an \(N\)-torsion point has \(U^N\) central and \(\iota(U)U\) central. Write the latter scalar as \(c\in R'^\times\).

After adjoining a square root of the unit \(c^{-1}\), replace \(U\) by \(aU\) with \(a^2c=1\). This extension is finite étale because 2 is a unit. The new lift satisfies \(\iota(U)=U^{-1}\). Its central power \(u=U^N\) therefore satisfies \(u=u^{-1}\), so \(u=1\) or \(-1\). If needed replace \(U\) by \(-U\); since \(N\) is odd, its \(N\)-th power becomes1 and the symmetry condition persists.

These are exactly the power and symmetry identities in \texttt{SymLift}. Uniqueness follows too: two such lifts differ by a central unit \(a\) with \(a^2=a^N=1\), hence \(a=1\). For an isotropic subgroup of odd exponent, these lifts form a subgroup: its lifts commute, their product satisfies the same power and symmetry identities, and uniqueness identifies it with the lift of the sum. This argument also shows that replacing any chosen integral lifts by the symmetric ones only multiplies the translated sections by units.

\emph{The exact restricted orbit and its integral basis.}
Put \(N=M^2\). The theta commutator of \(L^N\) is a perfect alternating pairing on \(A[N]\). It is balanced for CM adjoints: \[
 e_N(u x,y)=e_N(x,u^\dagger y).
 \tag{TS9}
\] This identity is Poincaré functoriality applied to the two restrictions defining the commutator; it uses the same adjoint identity as (TS8).

The subgroup \(H=A[\alpha^2]\) is isotropic. Indeed, choose \(c,d\) with \(c\alpha^2+d\bar\alpha^2=1\). For \(x,y\in H\), write \(y=d\bar\alpha^2y\) and use (TS9) to get \(e_N(x,y)=e_N(\bar d\alpha^2x,y)=1\). Its order is \(M^{2g}\), half the total rank in the square-root sense, so it is maximal isotropic.

The subgroup \(D=A[M]\) is also maximal isotropic. More concretely, \(L^N=[M]^*L\) has canonical descent data for the étale isogeny \([M]\); its deck transformations lift to commuting theta transformations fixing \(s_M\). They preserve symmetry, so they are the symmetric lifts just constructed. The orders of \(D\) and \(H\) are both \(M^{2g}\). The coprime-ideal identity gives \[
 H\cap D=A[\alpha]=C,\qquad |C|=M^g,
 \qquad |H/C|=M^g.
 \tag{TS10}
\]

Every translated section in (TS3) is invariant under the canonical deck group \(C\) of \([\alpha]\). In fact \(C\subset D\) fixes \(s_M\), and \(C\subset H\) commutes with the lift of \(t'\). By faithfully flat descent, its invariant sections are precisely \([\alpha]^*H^0(A,L^M)\), using (TS8). The canonical deck lifts on \(C\) obtained from \([\alpha]^*L^M\) agree with the restrictions of the \([M]\)-deck lifts: the rigidified identifications (TS8) respect normalized symmetry, and both lifts have odd finite order and satisfy that symmetry, so the uniqueness of symmetric finite-order lifts proved above applies. Thus the orbit span is contained in \(W_M\).

We show that coset representatives give an integral basis, not only a basis after inverting11. First reduce to the residue field, extending it finitely if needed to contain all torsion and characters. The reduction of \(s\) is nonzero, and so is that of \(s_M\), since \([M]\) is faithfully flat. Here the base-change assertion follows directly from the classical vanishing in Mumford §1(II): \(H^1(A_k,L_k^j)=0\) for \(j>0\). Proper coherent cohomology makes \(H^1(A_R,L^j)\) finite. The exact sequence for multiplication by11 embeds its quotient modulo11 into the vanishing special-fibre group; Nakayama therefore gives \(H^1(A_R,L^j)=0\). The same exact sequence gives \[
 H^0(A_R,L^j)/11\cong H^0(A_k,L_k^j).
 \tag{TS11}
\] The left module is finite and torsion-free, hence free. Its rank is \(j^g\), by the same classical cohomology formula for a principal line. In particular \(H^0(A_R,L)\) has a primitive generator as used in the statement.

The vector \(s_M\) has trivial character under the deck group \(D\). Its translate by \(t'\in H\) has character \(d\mapsto e_N(d,t')\), up to the harmless common inverse convention for the action. Two such characters agree precisely when the difference of their points lies in \(D^\perp\cap H=D\cap H=C\), because \(D\) is maximal isotropic. Therefore the \(M^g\) coset representatives give nonzero vectors in distinct character spaces. They are linearly independent over the residue field: a finite abelian group of order prime to11 has the mutually orthogonal character projectors obtained by averaging its characters, so distinct characters cannot cancel.

The target \(W_M\), identified with \(H^0(A,L^M)\), is free of rank \(M^g\), and its reduction is the corresponding invariant space by (TS11) and the prime-to11 averaging projector. Thus the \(M^g\) columns just found form a basis modulo11. Their determinant in any integral basis is a unit. They form an \(R'\)-basis of \(W_M\). This proves (TS4), including its rank and its strict difference from the ambient space. It also proves that two representatives in the same coset differ only by an invertible scalar in the orbit, so no larger span is hidden in the unused points.

\emph{Integral jets inside the smaller linear system.}
Apply DP56(a) to the fixed ample line \(L\): for \(M\ge rn+b\), its global sections of power \(M\) realize every integral jet of order \(n\). These same fixed constants work after the flat extension \(R'/R\).

The isogeny \([\alpha]\) is an isomorphism of the formal group at zero. Its differential is invertible over \(R'\), since its finite kernel has order prime to11; the integral formal inverse follows by recursive formal inversion with that unit linear matrix. The pullback isomorphism (TS8) identifies integral frames by an integral unit. Consequently pullback by \([\alpha]\) induces an isomorphism between the two finite jet modules at zero. Combining that isomorphism with the surjection for \(L^M\) proves the jet surjectivity of \(W_M\) asserted in (b).

In particular choose \(F\in W_M\) whose algebraic prefactor in the frame of DP54 is \(U^{-M^2}\) through degree \(n\). Its normalized germ then agrees with the pure primitive kernel of exponent \(-M^2\) through that degree. DP58 applies because \(-M^2\) is a unit. The selected derivative of \(F\), or harmonic derivative of two copies of \(F\), has exactly (TS5).

Choose the least \(e\ge1\) such that \(29^e\ge rn+b\). For \(rn+b>29\) this gives \(rn+b\le M<29(rn+b)\); the finitely many smaller degrees only change the constant. Hence (TS6) follows. This choice obeys all the arithmetic level conditions verified in the first part of the proof.

\emph{An individual allowed translate attains the sharp valuation.}
Write the interpolating section as its integral basis expansion \[
 F=\sum_{t'\in H/C}c_{t'}s_{t'},\qquad c_{t'}\in R'.
 \tag{TS12}
\] For an ordinary derivative, evaluate (TS12). For a two-copy harmonic derivative, expand the product into the finite sum indexed by two cosets, with coefficients \(c_{t'}c_{u'}\). If every selected value of the individual translates or tuples had valuation strictly greater than \(V\), the sum would too. Thus at least one has valuation at most \(V\).

The matching lower bound is important. At the critical weight \(w\), every normalized germ with integral algebraic prefactor has integral scaled coefficients, by the actual profile argument of DP46 and the full two-copy common-weight argument of DP51. The normalization factor \(M^{-n}\) is a unit. In the ordinary case the selected degree has digit sum two, so Legendre's formula gives the lower bound \[
 cn-2s=V-s,\qquad s=1/10.
 \tag{TS13}
\] For the harmonic polynomial \(P_{a,b}\), each supported differential monomial has its factorial multiplied by its polynomial coefficient equal, up to sign, to the same \(a!b!\): \[
 \binom ai\binom bi(a-i)!(b-i)!(i!)^2=a!b!.
\] Since \(a\) and \(b\) are powers of11, the same lower bound (TS13) follows for every pair of integral prefactors. It is not necessary to bound the separate factorial digit losses of every monomial crudely.

All these values lie in the fraction field of the unramified DVR \(R'\), whose normalized valuation group is \(\mathbf Z\). The number \(V=(11^q-1)/10\) is an integer and \(V-s>V-1\). Therefore (TS13) implies the integer lower bound \(v\ge V\). The tuple with \(v\le V\) consequently has \(v=V\). This proves the claimed individual translated-section witnesses, rather than merely a value of an arbitrary linear combination.

The valuation-group point cannot be discarded if one allows arbitrary ramified changes of lattice. All extensions used for the torsion, unit theta isomorphisms and odd-order symmetric lifts above were finite unramified extensions at11.

\emph{Exact generic source membership and arithmetic realization.}
For \(t'\in A[\alpha^2]\), put \(t=[M]t'\). Then \(t\in A[\alpha]\), since \(M=\alpha\bar\alpha\), and \(t'\) is killed by \(\mathfrak b^2=(\alpha^2)\). Use the canonical multiplication isomorphism (TS7) for \(\alpha_0\) in \texttt{IsVtOver}, its pullback to the base for \(\alpha\) in that definition, \(s_1=s_M\), and the symmetric lift \(U_{t'}\). The equality \([M]t'=t\), the annihilation condition, the multiplication pullback and the theta action are exactly its remaining clauses. Rigidification of (TS7) is exactly \texttt{IsNormMul}. This verifies the generic section relation clause by clause.

The curve, its CM endomorphisms and the line are defined over \(\mathbf Q(\zeta_7)\). A generator of its one-dimensional section space can be chosen over that number field and scaled to be primitive at the chosen prime: the finite free section lattice and flat base change of DP54 give this. The torsion points are algebraic. Their finite étale good models at11 give unramified completions. Integral isomorphisms in their theta-group fibres and their symmetric normalizations are defined over finite algebraic extensions; the unit roots used to normalize them are étale at11. One may choose an algebraic theta isomorphism primitive at the selected place, so these same integral lifts are realized algebraically. The witness tuples just proved therefore consist of global algebraic translated sections, not only complete-local formal sections.

There is no assertion here that the selected tuple is the one fixed in advance by a global lifting argument. Nor does an allowed tuple instantiate the old source \texttt{Datum.dom} or its induced-quadratic \texttt{HasType}. The result closes a narrower question: the exact generic permitted translate family itself already contains sharp ordinary and first-row harmonic witnesses on this higher-height Jacobian, at polynomially bounded admissible arithmetic levels. A uniform improved local bound for all these allowed translates is therefore impossible at those rates. Particular selected sections, prescribed row-degree sequences, global theta-image identification and the full net lift-cost comparison remain separate.

\end{proof}


\subsection{Algebraic interpolation with linear logarithmic section size}

\begin{theorem}[DP60: bounded representatives and section norms]\label{DP60}
Let $E=\mathbf Q(\zeta_7)$, let $\mathfrak p\mid11$, and put
\[
 A_0=\mathcal O_E[1/14],\qquad
 R_0=\mathcal O_{E,\mathfrak p},\qquad
 R=\widehat R_0=W(\mathbf F_{11^3}).
\]
Use the good Jacobian model $J/A_0$, the trace-one ample line $\mathcal L$, and the invariant product normalization of DP49 and DP54:
\[
 \Theta_s=f_sU^\kappa B^{-\kappa},
 \qquad U\in1+(z)R[[z]],\qquad B=\exp b.
\]
Fix continuous Hermitian metrics on $\mathcal L$ at the complex embeddings and use their tensor powers.
There are constants $a\ge1$, $b_0\ge0$, $C_1,C_2,C_3$, independent of $n$ and $\kappa$, such that for every $n\ge1$ and every $\kappa\ge an+b_0$ with $11\nmid\kappa$, there is a section
\[
 0\ne s_{n,\kappa}\in H^0(J_{A_0},\mathcal L^{\otimes\kappa})
\]
satisfying
\begin{align}
 f_{s_{n,\kappa}}U^\kappa
 &\equiv1\pmod{11R[[z]]+(z)^{n+1}},
 \label{eq:boundedrepresentativejet}\\
 \log\max\{1,\|s_{n,\kappa}\|_{\sigma,\sup}\}
 &\le C_1\kappa+C_2\log(n+1)+C_3
 \label{eq:linearsectionsize}
\end{align}
for every complex embedding $\sigma$.
Every explicit ordinary witness of DP58 of degree at most $n$ retains its exact pure-kernel valuation on this section.
The same holds for the explicit first-row two-column harmonic witnesses, using two copies of the section.

In particular one can choose
\[
 \kappa=M^2=an+O(\sqrt n),\qquad 11\nmid M,
\]
so both the tensor exponent and the logarithmic section sup-size are $O(n)$, while the exact selected valuations, including the unit factor $M^{-m}$ at each selected degree $m\le n$, are preserved.
The sections are integral at every finite place away from $2$ and $7$ in this fixed model.
For fixed tensor-compatible adelic metrics induced by the model outside these bad places, their adelic logarithmic sup-size is also $O(\kappa+\log(n+1))$.
\end{theorem}

\begin{proof}
\emph{Constant precision preserves the explicit jets.}
First work at any odd $p$ with the actual pair of DP46.
Write $s_*=1/(p-1)$ and, for a primitive line of intrinsic length $k$, put
\[
 N=p^k,\qquad w=\frac{s_*N}{N+1},\qquad
 c=\frac{s_*}{N+1}.
\]
At $m=(N+1)p^r$, the pure unit-parameter jet has valuation $cm-s_*$ by DP58.
If an integral prefactor is congruent to one modulo $p$ and order $n+1$, write
\[
 W=1+pV+H,\qquad V\in R[[z]],\quad H\in(z)^{n+1}.
\]
For $m\le n$, the term $H$ makes no contribution after inverse-logarithm substitution, since that substitution preserves total order.
DP46 makes the substituted integral function $V$ times the unit-parameter kernel integral at weight $w$.
As $m$ has digit sum two, the degree-$m$ error derivative therefore has valuation at least
\[
 1-wm+v_p(m!)=1+cm-2s_*>cm-s_*.
\]
The inequality follows from $s_*\le1/2$.
Thus reduction modulo $p$ is sufficient; the exact value is preserved for every unit exponent, including $-\kappa$.

For two copies, modulo total order $n+1$ one has
\[
 (1+pV_1)(1+pV_2)-1
 =p(V_1+V_2+pV_1V_2).
\]
The bracket is an integral function on the product formal group.
The full-polydisc argument of DP51 applies to any such function: the contracted inverse-logarithm coordinates and the product exponential are both integral.
At the first-row critical common weight
$w_*=s_*p^{h-1}/(p^{h-1}+1)$, with $c_*=s_*-w_*$, take
\[
 a_r=p^r,\qquad b_r=p^{h-1+r},\qquad m=a_r+b_r.
\]
For every supported differential monomial of $P_{a,b}$, its coefficient times its multi-factorial is, up to sign,
\[
 \binom ai\binom bi(a-i)!(b-i)!(i!)^2=a!b!.
\]
Consequently the full harmonic error has valuation at least
\[
 1-w_*m+v_p(a_r!b_r!)
 =1+c_*m-2s_*>c_*m-s_*.
\]
This is strictly above the exact pure value in DP58.
All product errors are covered, including nonseparable prefactors: the difference of products is $p$ times an integral product-coordinate function.

\emph{A finite set of global generators.}
Return to $p=11$.
Over $R_0$, choose the projective-coordinate sections from DP56,
\[
 h_0,h_1,h_2,h_3\in H^0(J_{R_0},\mathcal L^{\otimes a}),
\]
with $h_0(e)$ a unit, $h_i(e)=0$ for $i>0$, and $t_i=h_i/h_0$ integral formal parameters.
Choose $b$ so every power $\mathcal L^{\otimes m}$ with $m\ge b$ is globally generated.
For each residue class $\ell\bmod a$, choose one exponent $b_\ell\ge b$ in that class and one section
\[
 \eta_\ell\in H^0(J_{R_0},\mathcal L^{\otimes b_\ell})
\]
whose value at zero is a unit.
These choices use the same known ample-power and global-generation statements as DP56, namely Stacks Project Tags 01VT and 01Q3.
Put $b_0=\max_\ell b_\ell$.
For $\kappa\ge an+b_0$, choose $\ell\equiv\kappa\bmod a$ and set
\[
 q=\frac{\kappa-b_\ell}{a}\ge n.
\]
The section $\eta_\ell h_0^q$ is an integral local frame at zero.
Only finitely many auxiliary generators have been chosen.

Flat cohomology base change for $A_0\to R_0$ identifies these section modules over $R_0$ with the localizations of their modules over $A_0$.
The precise known statement is Stacks Project Lemma~30.5.2(2), Tag~02KH.\footnote{\url{https://stacks.math.columbia.edu/tag/02KH}.}
Clear the finitely many denominators using algebraic integers outside $\mathfrak p$.
One may do this integrally because a denominator in $A_0\setminus\mathfrak p$ has a numerator in $\mathcal O_E\setminus\mathfrak p$, and only powers of $14$ were inverted.
Use a common clearing scalar for all $h_i$, preserving their ratios.
The scalars are units over $R_0$, so the unit values, vanishing and formal-coordinate properties persist.
The resulting fixed generators are global sections over $A_0$; clearing can add divisibility at other primes but introduces no new denominator.

\emph{Bounded integral coefficient representatives.}
For every $C\ge1$ the map
\[
 \mathcal O_E/p^C\mathcal O_E\longrightarrow R/p^CR
\]
is surjective.
Indeed $p\mathcal O_E=\mathfrak p\overline{\mathfrak p}$, and the Chinese remainder theorem identifies the source with
\[
 \mathcal O_E/\mathfrak p^C\times
 \mathcal O_E/\overline{\mathfrak p}^{\,C}.
\]
Localization does not change the first quotient, since all denominators outside $\mathfrak p$ are units there, and completion preserves it as well.
Only $C=1$ is needed.
As $11$ has order three modulo $7$, the residue of $\zeta_7$ has degree three, and the fixed set
\[
 \mathcal S=\{c_0+c_1\zeta_7+c_2\zeta_7^2:
                 0\le c_i\le10\}
\]
represents $R/pR$ bijectively.
Every element is an algebraic integer and has absolute value at most $30$ at every complex embedding.

\emph{The actual global interpolant.}
Fix an integral formal frame $\tau$ of $\mathcal L$ and let
\[
 g_\kappa=\frac{\eta_\ell h_0^q}{\tau^\kappa}
 \in R[[z]]^\times.
\]
Express the integral unit $g_\kappa^{-1}U^{-\kappa}$ in the integral parameters $t_i$, truncate through degree $n$, and write the result as
\[
 \sum_{|\alpha|\le n}q^*_\alpha t^\alpha,
 \qquad q^*_\alpha\in R.
\]
Choose $q_\alpha\in\mathcal S$ congruent to $q^*_\alpha$ modulo $p$, and define
\begin{equation}\label{eq:boundedglobalinterpolant}
 s_{n,\kappa}
 =\eta_\ell\sum_{|\alpha|\le n}q_\alpha
      h_0^{q-|\alpha|}h_1^{\alpha_1}
      h_2^{\alpha_2}h_3^{\alpha_3}.
\end{equation}
All powers are nonnegative and each summand has tensor exponent $b_\ell+aq=\kappa$.
The generators are integral over $A_0$ and the coefficients lie in $\mathcal O_E$, so this is a global section over $A_0$.
Its formal coefficient multiplied by $U^\kappa$ is congruent to one modulo $p+(t)^{n+1}$, hence modulo $p+(z)^{n+1}$.
This proves \eqref{eq:boundedrepresentativejet} and makes the section nonzero, with unit value at the chosen prime.
The constant-precision argument proves all the selected jet claims, simultaneously for every witness of degree at most $n$.
The construction already lies over the number field; no uncontrolled final approximation in a completed section space is used.

\emph{Norms and all finite places.}
Choose $H\ge1$ bounding the sup norms of the finitely many $h_i$ and $\eta_\ell$ at all complex embeddings in their indicated tensor metrics.
There are $\binom{n+3}{3}$ terms in \eqref{eq:boundedglobalinterpolant}.
Tensor compatibility and the triangle inequality give the explicit estimate
\[
 \|s_{n,\kappa}\|_{\sigma,\sup}
 \le30\binom{n+3}{3}H^{q+1}.
\]
Since $q\le\kappa/a$, its logarithm is bounded by
\[
 \log30+\log\binom{n+3}{3}
       +(\kappa/a+1)\log H,
\]
proving \eqref{eq:linearsectionsize}.
At every model place outside $2,7$, integrality gives sup norm at most one.
For any fixed continuous metrics on the proper analytifications at the remaining finite places, bound the same finite generator set by constants $H_v\ge1$.
The coefficients are algebraic integers, so the ultrametric inequality gives
\[
 \|s_{n,\kappa}\|_{v,\sup}\le H_v^{q+1}.
\]
The normalized adelic sum of positive logarithmic sup norms is therefore $O(\kappa+\log(n+1))$.
There are no other finite-place contributions and no new good-prime denominators.

Finally choose $m_0=\lceil\sqrt{an+b_0}\rceil$ and let $M=m_0$ unless $11\mid m_0$, in which case set $M=m_0+1$.
Then $11\nmid M$ and $M^2=an+O(\sqrt n)$.
The preceding construction works at $\kappa=M^2$, and its normalized kernel exponent $-\kappa$ is a unit.
DP58 is uniform over all such units.
The factor $M^{-m}$ has valuation zero at every selected derivative degree.
This proves the linear bounds and exact-value assertions.
\end{proof}

Two fixed continuous Hermitian metrics on the line over a compact complex fibre differ by bounded positive factors; on its $\kappa$th tensor power the logarithmic discrepancy is $O(\kappa)$.
Thus the theorem also applies to any fixed geometric canonical theta metric.
In a fixed honest holomorphic frame on a relatively compact chart, dividing by the frame and multiplying by a fixed holomorphic gauge $\exp(\kappa h_L)$ costs only another $\exp(O(\kappa))$ factor.
This statement does not assign an archimedean norm to the $p$-adic display frame or identify it with a source complex gauge.
Level-dependent normalizations require their own comparison.

The result is a construction in the full algebraic section space, with exact finite local values and upper bounds for fixed norms.
It does not identify the prescribed torsion-translate span, force a prescribed degree sequence to lie on the explicit witness ray, or compare the constants with a global gain-minus-cost budget.
No conclusion about existence or impossibility of an entire shorter lift chain follows from this norm bound alone.


\subsection{The prescribed row component and its finite windows}

\begin{proposition}[DP61: selected row degrees and harmonic completion]\label{DP61}
Fix the source rank data $k,m,t$ and a nonfinal stage $1\le j<m$.
Put $P=k+j-1$, $Q=m-j$, $n_V=k+m-1+t$, and $H=(t+3)/2$.
For one embedding write the source HC parameters as $u,v,b,e$, and set
\[
 \sigma_i=\begin{cases}v_i&i<k-1,\\-b_{i-k+1}&i\ge k-1,\end{cases}
 \quad A_i=\sigma_i+e_{j-1}-(P-1-i)-(t+1)/2,
\]
\[
 C_c=b_{j+c}-e_{j-1}-(t+1)/2-c,\qquad B_n=C_{Q-1-n}.
\]
These are the actual positive and negative row partitions of the selected
source weight. In particular the negative order is reversed.
\begin{enumerate}[label=(\alph*),leftmargin=*]
\item Every active entry is a positive integer, both partitions are
antitone, and
\[
 A_0-A_1=v_0-v_1-1\ge0,
 \qquad B_0-B_1=b_{m-1}-b_{m-2}-1\ge t+2\quad(Q\ge2).
\]
Consequently the unchanged DP58 first-row two-column polynomial does not
lie in the full prescribed row component.
\item For any antitone nonnegative partitions $A,B$, if
\[
 0\le a\le A_0-A_1,\qquad 0\le b\le B_0-B_1,
\]
with a missing next entry interpreted as zero, and $n\ge P+Q+2$, the
first-row harmonic polynomial
\[
 P_{a,b}=\sum_{s=0}^{\min(a,b)}(-1)^s\binom as\binom bs
 X_{0,1}^{a-s}Y_{0,1}^{b-s}X_{0,2}^sY_{0,2}^s
\]
has a primitive integral determinant multiple in the full
$\operatorname{HaO}$ component of degrees $A,B$.
The actual source column count $n_V$ satisfies this numerical requirement.
\item If each column kernel satisfies $K_c(X,0)=K_c(0,Y)=1$, the
completion in (b) evaluates to zero on $\prod_cK_c$ whenever its
filler is nonconstant. For the partitions in (a) the filler is always
nonconstant.
\item For fixed source $\operatorname{SizeData}_2$ windows $Sc,\eta$,
\[
 A_0-A_1\le Sc-1.
\]
For $Q\ge2$,
\[
 B_0-B_1\le2\eta-1,
 \qquad B_0-B_1\le\eta-H-(Q-2)(t+3)-1;
\]
for $Q=1$, $B_0\le\eta-(t+1)/2$.
Thus only finitely many of the sufficient completions in (b), along
$a=\ell^s,b=\ell^{h-1+s}$, fit a fixed datum.
\end{enumerate}
The sufficient completion criterion is not a classification of every
harmonic polynomial in the component. These assertions do not identify
any mixed-type geometric object with the old source
$\operatorname{Datum.dom}$ or $\operatorname{CMObj.HasType}$.
\end{proposition}

\begin{proof}
The source HC conditions are
\[
 u_i>v_i>u_{i+1},\quad u_{k-1}>-b_0,\quad
 e_s-b_s\ge H,\quad b_{s+1}-e_s\ge H,
\]
with integral $v,b,e$ and half-integral $u$.
Hence successive $\sigma$ entries decrease by at least one and
$b_{s+1}-b_s\ge t+3$. Therefore $A$ is antitone and
\[
 A_{P-1}=e_{j-1}-b_{j-1}-(t+1)/2\ge1.
\]
Likewise $C_0\ge1$ and $C_{c+1}-C_c\ge t+2$, so reversing $C$
gives the positive antitone partition $B$ and the stated first gaps.
The rank assumptions give $P\ge5$ and $n_V\ge P+Q+2$.
The unchanged first-row polynomial has degree zero in later positive
rows, proving its literal exclusion.

For (b), set $A'=A-ae_0$, $B'=B-be_0$; the gap hypotheses make them
partitions. Reserve columns 1 and 2 for $P_{a,b}$ and disjoint further
column sets $C_+,C_-$ of sizes $P,Q$. For a partition $L$ of length $d$,
write
\[
 \mathcal D_L(Z)=\prod_{s=1}^{d}
 \det(Z_{i,c_k})_{0\le i,k<s}^{L_{s-1}-L_s},\qquad L_d=0.
\]
Use positive rows and $C_+$ for $\mathcal D_{A'}$, negative rows and
$C_-$ for $\mathcal D_{B'}$. Then
$\widetilde P=P_{a,b}\mathcal D_{A'}\mathcal D_{B'}$ has the required
row degrees by telescoping. Prefix-row determinants are invariant under
adding an earlier row to a later row: either the latter row is absent,
or determinant row addition applies. The first-row polynomial has the
same invariance. Every mixed contraction on filler columns vanishes
because the opposite row variables are absent. On columns 1 and 2,
only the first-row contraction can be nonzero, and it vanishes by
\[
 (a-s)(b-s)c_s+(s+1)^2c_{s+1}=0,
 \qquad c_s=(-1)^s\binom as\binom bs.
\]
These are exactly the identity-column-matrix contractions in the source
$\operatorname{HHa}$ definition. All coefficients are integral. The
product of determinant diagonal terms and the $s=0$ term of $P_{a,b}$
has coefficient one, so the result is nonzero and primitive. The same
construction works for all conjugate embedding pairs, since contractions
only couple an embedding with its conjugate.

For (c), every monomial differential operator of a nonconstant filler
has a positive-order derivative in some filler column and differentiates
only one side there. The corresponding axis-normalized kernel gives zero.
Every term of the product evaluation therefore vanishes. Later positive
row degrees remain positive after subtracting $ae_0$, so the source
filler is nonconstant.

Finally the source bounds $0\le v_i\le Sc$ and
$|b_i-e_s|\le\eta$ give the first and the coarse second estimates in (d).
For the sharper bound, HC separation gives
\[
 b_{m-2}-e_{j-1}\ge H+(Q-2)(t+3),
 \qquad b_{m-1}-e_{j-1}\le\eta.
\]
Subtracting and then subtracting one proves the claim. The $Q=1$ formula
is immediate from $B_0=b_{m-1}-e_{m-2}-(t+1)/2$.
\end{proof}

\paragraph{Exact source identification.}
The defining rank and HC fields are
\path{Formalism/LocalRingCardinality/Params.lean}, lines 6--69.
\path{DataKap.lean}, lines 13--34 and 60--116, gives the integral
formulas and positive last entries.
\path{IndefiniteReal/LinksDefs.lean}, lines 341--348, defines
\path{aPl/aMi};
\path{Development/AbelianVarietyTheory/Proofs/IwasawaDecomposition/AMiEqAMinusQ.lean},
lines 20--66, proves their identification and the negative reversal.
\path{IndefiniteReal/FWExist.lean}, lines 518--580, supplies
\texttt{KapIs}, \texttt{kapTop}, and its uniqueness.
\path{ProjectiveLattices/PartitionExtendByZeroAntitone.lean}, lines 67--69,
puts these partitions at the selected and conjugate embeddings via
\texttt{wtA}. The actual selected family is
\path{IsThRule.lean}, lines 84--91, and
\path{RecPJThAllNeBot.lean}, lines 31 and 80--87.
The polynomial conditions are \path{RowDeg.lean}, lines 21--48;
the fixed windows are \path{DataParams}\textsubscript{2}\path{.lean}, lines 148 and 160--164.
\path{Columns/AbelianHostFieldConditions.lean}, lines 443--460, fixes
$v,b,e$ and the first and last $u$ across words, so changing the entire
row datum is distinct from varying a word.
These are source reads, with original-author byte identity recorded;
no Lean execution is asserted.


\subsection{Integral theta lifts and exact canonical section norms}

\begin{theorem}[DP62: individual allowed translates]\label{DP62}
Use the good model of the Jacobian, its symmetric rigidified trace-one line $L$, and its principal polarization over
$S=\operatorname{Spec}\mathcal O_E[1/14]$, where $E=\mathbf Q(\zeta_7)$.
Fix a nonzero algebraic section $s\in H^0(A,L)$ primitive at the selected prime over $11$.
Let $M$ be odd, put $N=M^2$, and let
\[
 \Phi_M:[M]^*L\xrightarrow{\sim}L^N
\]
be the canonical rigidified multiplication isomorphism.
For a finite extension $F/E$ and $t\in A[N](F)$, the theta lift $U_t$ with
\[
 U_t^N=1,\qquad \iota(U_t)=U_t^{-1}
\]
exists uniquely over the field of definition of $t$ and extends integrally over $\mathcal O_F[1/14]$.
This includes places dividing $M$.
For any local lift $U$, it is given by
\begin{equation}\label{eq:rootfreethetalift}
 c=\iota(U)U,\quad d=U^N,\qquad
 U_t=c^{(N-1)/2}d^{-1}U.
\end{equation}
The formula glues without any root extraction.
Consequently every individual translate
\[
 s_t=U_t\bigl(\Phi_M([M]^*s)\bigr)
\]
is integral on the good model over $\mathcal O_F[1/14]$.

At every complex embedding, equip $L$ with the canonical smooth Hermitian metric whose curvature is translation-invariant and whose rigidification at zero is an isometry.
With the tensor metric on $L^N$, one has the exact equality
\begin{equation}\label{eq:individualthetanorm}
 \|s_t\|_{\sigma,\sup}
 =\|s\|_{\sigma|_E,\sup}.
\end{equation}
The same holds for $L^r$-norms, $1\le r<\infty$, with probability Haar measure.

For the permitted subgroup $H=A[\alpha^2]$ of DP59, any $r$-tuple of points, its canonical lifts and its translated sections are defined over a field of degree at most
\[
 [E(t_1,\ldots,t_r):E]\le M^{2gr}.
\]
Thus the ordinary and two-copy harmonic sharp witnesses in dimension three need degrees at most $M^6$ and $M^{12}$.
Their canonical section norms are uniformly bounded, and they have no new good-prime denominators, at the levels $M=29^e=O(n)$ and $M^2=O(n^2)$ of DP59.
No discriminant or complete source derivative/gauge-cost bound is asserted.
\end{theorem}

\begin{proof}
\emph{The finite-flat theta torsor.}
For the principally polarized line,
\[
 \phi_{L^N}=N\phi_L,\qquad K(L^N)=A[N].
\]
This kernel is finite locally free even at primes dividing $N$.
Indeed each fibre of multiplication by $N$ is finite locally free of degree $N^{2g}$ by the known field theorem, Stacks Project Lemma~39.9.8, Tag~0BFG.
The relative source is smooth, hence flat and locally of finite presentation over $S$; the target is locally of finite type.
The fibrewise criterion, Lemma~37.16.4, Tag~039E, gives relative flatness.
The map is proper with finite fibres and hence finite by Lemma~37.44.1, Tag~02LS.
Over this noetherian base, finite flatness gives local freeness, of the stated fibre rank.
The kernel is the base change along zero.\footnote{The exact known statements are \url{https://stacks.math.columbia.edu/tag/0BFG}, \url{https://stacks.math.columbia.edu/tag/039E}, and \url{https://stacks.math.columbia.edu/tag/02LS}.}

The relative theta group is the central extension
\[
 1\longrightarrow\mathbf G_m\longrightarrow
 \mathcal G(L^N)\longrightarrow A[N]\longrightarrow1.
\]
Its fibre is the torsor of translation isomorphisms of $L^N$.
The normalized symmetry induces an automorphism $\iota$ which fixes the central scalars and covers inversion on $A[N]$.
The relative finite-flat setting and this symmetry convention are the known theta-group statements in Polishchuk, \emph{Analogue of Weil representation for abelian schemes}, \S1, author preprint pp.~4--5, before Lemma~1.1.
Proposition~1.2 and Corollary~1.3 also establish the odd symmetric torsor's triviality; the following formula makes its normalization explicit.\footnote{\url{https://arxiv.org/pdf/alg-geom/9712021}. The stated page numbers refer to this directly read author version.}

\emph{Normalization over arbitrary test bases.}
For a point $t\in A[N](T)$, choose a lift $U$ locally in the fppf topology.
The elements $d=U^N$ and $c=\iota(U)U$ lie over zero, so are invertible central scalars.
Since $\iota(U)=cU^{-1}$, these two lifts commute, and
\[
 c^N=\iota(U)^NU^N=d^2.
\]
For $a=c^{(N-1)/2}d^{-1}$ this gives
\[
 a^2c=c^Nd^{-2}=1,\qquad
 a^Nd=c^{N(N-1)/2}d^{1-N}=1.
\]
Hence $aU$ has both required identities.
Replacing the local lift by $uU$ changes $c$ to $u^2c$, $d$ to $u^Nd$, and $a$ to $u^{-1}a$.
Thus $aU$ is unchanged and descends on every overlap.
Only multiplication and inversion of units occur; neither $N$ nor any root is inverted or adjoined.
The construction commutes with arbitrary base change.
Two normalized lifts differ by a unit with $u^2=u^N=1$; since $N$ is odd, an integer B\'ezout identity gives $u=1$, including over nonreduced bases.

This is a section of the underlying torsor, not a group splitting of the whole central extension.
On an isotropic subgroup the lifts commute; their products again satisfy the two identities, so uniqueness makes that restriction a homomorphism.
The normalized symmetry and the power identity are exactly the clauses of the recovered \texttt{ThetaPt.SymLift} definition, lines~44--51; the action on sections is \texttt{IsActOf}, lines~83--84.

\emph{Integrality of the individual algebraic sections.}
A torsion point over $F$ extends uniquely over every good DVR by properness of the abelian scheme.
Stacks Project Lemma~29.43.5, Tag~0BX7, extends the rational map to a neighborhood of each such point: the source is noetherian and its local ring is a valuation ring.
The valuative criterion is Lemma~29.43.1, Tag~0BX5.
The neighborhood extensions agree on overlaps by separatedness and the dense common generic fibre, so they glue over $\operatorname{Spec}\mathcal O_F[1/14]$.\footnote{\url{https://stacks.math.columbia.edu/tag/0BX7}; \url{https://stacks.math.columbia.edu/tag/0BX5}.}
Separatedness preserves $Nt=0$.
Applying \eqref{eq:rootfreethetalift} to the resulting integral theta torsor gives an integral lift on this same base.
In particular the field of the point already contains the canonical lift.

DP59's Poincar\'e construction makes the symmetry and $\Phi_M$ integral on the good model.
Pulling back the fixed integral section by $[M]$, applying $\Phi_M$, and then applying the integral theta action therefore gives $s_t$ integrally.
This includes the places above $29$ when $M=29^e$; their torsion kernels need not be \'etale.
For these DP59 levels the torsion at $11$ is finite \'etale, so the already unramified point field requires no further extension for the lift.
The integer valuation group used in DP59 is preserved.
No assertion that the fixed section is primitive at every other good prime is needed.
The existence of a fixed global $s$ primitive at $11$ follows by clearing the local generator's denominators outside that prime, as in DP60.

\emph{Canonical metrics.}
The known metric theorem is Moret-Bailly, \emph{Pinceaux de vari\'et\'es ab\'eliennes}, Ast\'erisque~129 (1985), Proposition~II.2.1, printed p.~48: a cubical line on a complex abelian variety has its unique smooth Hermitian metric compatible with the rigidification and with translation-invariant curvature.
The equivalent characterization is also stated in de Jong--Shokrieh, \emph{Faltings height and N\'eron--Tate height of a theta divisor}, Compositio Mathematica~158 (2022), \S12.2.\footnote{Original primary page: \url{https://www.numdam.org/item/AST_1985__129__1_0.pdf}, printed p.~48, PDF page~51. The original page and the later published statement were both directly read.}
We use this known theorem, rather than re-proving metric existence.

Let $\omega$ be the invariant curvature of $L$ at a fixed complex embedding.
Because the differential of $[M]$ is multiplication by $M$,
\[
 [M]^*\omega=M^2\omega=N\omega.
\]
The pullback and tensor metrics compared by $\Phi_M$ therefore have the same curvature.
Their logarithmic ratio is a smooth real pluriharmonic function on the compact connected torus and hence is constant.
The rigidified isomorphism is an isometry at zero, so that constant is zero.
Thus $\Phi_M$ is isometric and
\[
 \|\Phi_M([M]^*s)(x)\|=\|s([M]x)\|.
\]
Surjectivity of $[M]$ gives equality of the sup norms.

For a finite-order theta lift $U_t$, the comparison of the metric with its translate has zero curvature because $T_t^*(N\omega)=N\omega$.
Its positive norm ratio is therefore a constant $r$.
Composing the lift $N$ times gives ratio $r^N$, while $U_t^N=1$ gives ratio one.
Hence $r=1$: the theta lift is isometric.
Translation is bijective, proving \eqref{eq:individualthetanorm}.
With probability Haar measure, translation preserves measure and the pushforward by $[M]$ is again probability Haar measure; this proves the finite $L^r$ equalities.
Every embedding of $F$ restricts to one of the finitely many embeddings of the fixed field $E$, so all the constants are uniform in the level and the point.

\emph{The actual tuple field.}
The subgroup $H=A[\alpha^2]$ is defined over $E$, is \'etale in characteristic zero, and has degree $M^{2g}$ by DP59.
Consequently $H^r$ has degree $M^{2gr}$, bounding the residue-field degree of any tuple by that number.
The normalized lift is unique and functorial, so its coefficients and the translated section descend to the same field.
This gives the stated $M^6$ and $M^{12}$ bounds in dimension three.
For comparison, the full torsion splitting field has Galois group injecting into $\operatorname{GL}_{2g}(\mathbf Z/N\mathbf Z)$ and therefore degree at most $N^{4g^2}=M^{8g^2}$, though no full splitting field is needed for an individual tuple.
No discriminant estimate follows merely from either degree bound.
\end{proof}

The norm in DP62 is the canonical norm of the section itself.
Dividing by a source analytic frame, multiplying by its exponential gauge, and applying a specified differential polynomial can introduce additional level and degree factors.
Those operations, bad-prime norm comparisons, discriminants, prescribed row weights, the mixed-type moduli interface and membership in the iterated theta image require their own estimates.
The theorem removes an uncontrolled canonical section norm and new good-prime denominators for the individual allowed witnesses; it does not by itself compare the net cost of whole theta-lift chains.


\subsection{Scalable row partitions and an exact filler cost}

\begin{proposition}[DP63: fixed factorial fillers and permitted translates]\label{DP63}
Fix source rank data $k,m,t$ and $1\le j<m$, and put
$P=k+j-1$, $Q=m-j$, $n_V=P+Q+t$, $H=(t+3)/2$.
There are fixed antitone nonnegative integral partitions $A^*,B^*$ of
lengths $P,Q$ such that, for every $a\ge0$ and $b\ge t+2$, an actual
$\operatorname{HCParam}$ has selected partitions
\[
 A=A^*+ae_0,\qquad B=B^*+be_0.
\]
Taking the same data at every word supplies the literal algebraic
$\operatorname{Wts.Ok}$ conditions. Set
\[
 D=|A^*|+|B^*|,\qquad
 \mathcal F=\prod_i A_i^*!\prod_s B_s^*!.
\]
The DP61 determinant completion has an integral algebraic-coordinate
prefactor of degree $D$ for which its normalized evaluation is exactly
$\mathcal F H_{a,b}(\lambda)$, with $H_{a,b}$ as in DP58.

On the actual Jacobian at $11$, let
\[
 a_r=11^r,\quad b_r=121\,11^r,\quad T_r=122\,11^r+D.
\]
On a fixed product of copies of that Jacobian, there is an actual tuple
of the generic allowed translates of DP59, at
$M=29^e=O(T_r)$ and $\kappa=M^2=O(T_r^2)$, whose completed-row
polynomial evaluation, including $M^{-T_r}$, has valuation at most
\[
 V_r=(11^r-1)/10+v_{11}(\mathcal F).
\]
Hence no uniform bound $\gamma T_r-C$ with $\gamma>1/1220$ and fixed
$C$ holds for all these evaluations as the HC datum varies.

The geometric assertion uses explicitly assigned formal directions for
its two abstract row blocks. It does not identify those blocks with the
old source's conjugate CM-embedding blocks, instantiate its
$\operatorname{Datum.dom}/\operatorname{CMObj.HasType}$, or supply a
varying $\operatorname{SizeData}_2$ package or theta-image condition.
\end{proposition}

\begin{proof}
Use $\beta,\epsilon$ for the HC arrays normally called $b,e$.
Set
\[
 v_i=k-2-i\ (1\le i\le k-2),\qquad v_0=k-2+a,
\]
\[
 u_i=v_i+\tfrac12\ (0\le i\le k-2),\qquad u_{k-1}=-\tfrac12.
\]
Let $\beta_0=1$, put $\beta_{s+1}-\beta_s=2H$ for $0\le s<m-2$,
and define $\epsilon_s=\beta_s+H$ for every $0\le s<m$, independently
of the gap range. For the final gap take
\[
 \beta_{m-1}-\beta_{m-2}=
 \begin{cases}b+1&Q\ge2,\\b+2H-1&Q=1,\end{cases}
 \qquad\epsilon_{m-1}=\beta_{m-1}+H.
\]
The first range is empty when $m=2$. Every $v,\beta,\epsilon$ is
integral and every $u$ is half-integral. Consecutive $v$ gaps are one
except for the first gap $a+1$, so all strict interlacing conditions
hold. The last inequality is $-1/2>-\beta_0=-1$.
Every $\epsilon_s-\beta_s=H$ and every
$\beta_{s+1}-\epsilon_s\ge H$, including the final gap because
$b\ge t+2=2H-1$. These verify all HC fields.

At stage $j<m$, no positive row uses the variable final entries
$\beta_{m-1},\epsilon_{m-1}$; only $A_0$ varies, and
$A_0-A_1=a$. When $Q\ge2$, only $B_0$ varies and
$B_0-B_1=\beta_{m-1}-\beta_{m-2}-1=b$.
When $Q=1$,
\[
 B_0=(b+2H-1)-H-(H-1)=b.
\]
Thus the residual partitions are fixed, antitone and nonnegative.
All full entries are positive by DP61. For $\operatorname{Wts.Ok}$,
use this same tuple at every word and its common $\operatorname{RowPar}$;
its row witness and equality fields are immediate. No larger source
datum is inferred from this algebraic record.

For the exact filler evaluation, take the disjoint column sets
$C_+=\{c_i\}$, $C_-=\{d_s\}$ from DP61. In
$Q_0=\mathcal D_{A^*}(X)\mathcal D_{B^*}(Y)$ the monomial
\[
 Z^\nu=\prod_iX_{i,c_i}^{A_i^*}\prod_sY_{s,d_s}^{B_s^*}
\]
has coefficient one, because every determinant must use its diagonal
term to avoid an off-diagonal variable. Choose integral formal
coordinates $z$ with the prescribed tangent basis. In normalized
additive coordinates, $z=Z+O(\deg2)$. The integral monomial
\[
 f=\prod_i z_{X_i,c_i}^{A_i^*}\prod_s z_{Y_s,d_s}^{B_s^*}
\]
has initial ordinary homogeneous form $Z^\nu$ of degree $D$.
The filler-column kernels have constant one, so
\[
 Q_0(\partial)\left(f\prod_{c\in C_+\cup C_-}K_c\right)(0)
 =Q_0(\partial)Z^\nu(0)=\mathcal F.
\]
All higher prefactor or kernel terms have degree greater than $D$.
The remaining operator $P_{a,b}$ uses only columns 1 and 2, giving the
exact product $\mathcal F H_{a,b}(\lambda)$. This argument uses integral
algebraic coordinates, not an integrality assertion for logarithms.
DP58 then gives the displayed valuation at $a_r,b_r$ for every unit
$\lambda$.

Here is a precise actual geometric realization of the abstract row
labels. Put $B=J^{P+Q-1}$ with product polarization and trace-one line.
In each column assign $X_0=x,Y_0=y_1$ to copy0; assign $X_i=x$ to
copy $i$ for $1\le i<P$; and assign $Y_s=y_1$ to copy $P+s-1$ for
$1\le s<Q$. All unused directions are set to zero. The product kernel
restricts to the sharp kernel in copy0 and to normalized axes in all
other copies. There are $n_V$ independent column copies of $B$.
This is an actual invariant differential evaluation, while the abstract
row labels are not claimed to coincide with CM conjugate pairs:
DP49 gives characters $3$ and $2$ for $x$ and $y_1$.

The tensor products of DP59 factor basis sections are a finite integral
generating family for the allowed product-translate span. The points lie in $B[\alpha^2]$ and the symmetric lifts,
multiplication isomorphisms and rigidifications multiply, proving the
same generic $\operatorname{IsVtOver}$ relation. Product jets of total
order $T$ are generated integrally: factor each coordinate monomial
among the $J$ factors, realize each factor jet to order $T$ by DP59,
and multiply. Every error has total order at least $T+1$.
Consequently $M\ge cT+d$ suffices, with fixed constants for this product.

For each column prescribe, in the exact integral frame and compatible
gauge of DP59, the jet $U^{-\kappa}$ times its desired prefactor
monomial, or $U^{-\kappa}$ alone in the first two and unused columns.
It is integral and realizes the desired normalized kernel jet through
order $T_r$. Thus the full selected evaluation of these interpolating
sections has valuation exactly $V_r$. Expand each section integrally
in this allowed-translate generating family and use multilinearity over the $n_V$
columns. If every resulting allowed tuple had valuation greater than
$V_r$, so would their finite integral linear combination. Hence some
tuple has valuation at most $V_r$. No matching full-row lower bound or
exact equality for this tuple is asserted. The factor $M^{-T_r}$ is a
unit. Choosing the least sufficient power of29 gives the stated
linear bound on $M$ and quadratic bound on $\kappa$.
The tuple consists of algebraic allowed generating translates even if the
interpolation coefficients lie in the completed ring.
Finally $V_r/T_r\to1/1220$, proving the uniform-rate exclusion with
its stated varying-datum and chosen-direction scope.
\end{proof}


\subsection{The bare prime budget for every fixed cyclotomic type}

\begin{theorem}[DP64: fixed cyclotomic types and saturated integer budgets]\label{DP64}
Fix a CM subfield $E\subset\mathbf Q(\zeta_m)$, its embedding set
$\Omega$, a CM type $\Phi\subset\Omega$, and a finite prime set $S$
containing the divisors of $m$.
For $p\notin S$, let $\sigma_p\in\operatorname{Gal}(E/\mathbf Q)$
be the restriction of $\zeta_m\mapsto\zeta_m^p$.
For $i\in\Phi$ set
\[
 c_{i,p}=\sum_{r\ge1}\frac{\mathbf1_{\sigma_p^ri\notin\Phi}}{p^r},
 \qquad E_{i,p}(n)=\sum_{r\ge1}\mathbf1_{\sigma_p^ri\notin\Phi}
                     \left\lfloor\frac n{p^r}\right\rfloor.
\]
These are the DP40/43 bare directional slopes and integer lower
certificates for the associated integral cyclic displays; applying
them to geometry requires those displays and fixed global primitive
eigendirections.
For arbitrary nonnegative integers $n_i$, held fixed as $p$ varies,
put $N=\sum_i n_i$. Uniformly in the type, labels, orders and $X\ge2$,
\begin{align}
 \sum_{\substack{p\le X\\p\notin S}}\log p\sum_i n_i c_{i,p}
 &=\frac N2\log X+O_{m,S}(N),\label{eq:generalcmlinear}\\
 \sum_{\substack{p\le X\\p\notin S}}\log p\sum_i E_{i,p}(n_i)
 &=\frac12\sum_i n_i\log\bigl(\max\{1,\min(X,n_i)\}\bigr)
       +O_{m,S}(N).\label{eq:generalcmfloors}
\end{align}
Thus data in the same fixed cyclotomic field with the same order
multiset have budget difference $O_{m,S}(N)$.
If at most $d$ directions are used, even arbitrary distributions of
the same total order satisfy
\begin{equation}\label{eq:generalcmentropy}
 \sum_{\substack{p\le X\\p\notin S}}\log p\sum_i E_{i,p}(n_i)
 =\frac N2\log\bigl(\max\{1,\min(X,N)\}\bigr)
       +O_{m,S,d}(N).
\end{equation}
The constants are not asserted uniform in growing $m,S,d$.
The orders may vary with an outer parameter, but are fixed during
each prime sum.
\end{theorem}

\begin{proof}
The known analytic input is Rudnick, \emph{Primes in Arithmetic
Progressions: Fixed Modulus}, May 26, 2015, \S2, printed p.~2,
equation~(1): for fixed $m$ and $(a,m)=1$,
\[
 \sum_{\substack{p\le Y\\p\equiv a\pmod m}}\frac{\log p}{p}
 =\frac{\log Y}{\varphi(m)}+O_m(1).
\]
The quantitative fixed-modulus PNT immediately above it also gives
$\sum_{p\le Y}\log p=O(Y)$.
These are cited known results, independent of RH.\footnote{The exact
primary page was reopened and read:
\url{https://www.math.tau.ac.il/~rudnick/courses/sieves2015/fixed\%20progressions.pdf}.}

Put $b_i(a)=\mathbf1_{\sigma_a i\notin\Phi}$.
Restriction from $(\mathbf Z/m)^\times$ onto
$G=\operatorname{Gal}(E/\mathbf Q)$ is surjective with equal fibres,
and $G\to\Omega$, $\sigma\mapsto\sigma i$, is bijective.
Exactly half of $\Omega$ lies outside $\Phi$, so
\[
 \frac1{\varphi(m)}\sum_{a\in(\mathbf Z/m)^\times}b_i(a)=\frac12.
\]
Summing the analytic input in the finitely many classes, and deleting
$S$, gives
\begin{equation}\label{eq:generalcmfirst}
 \sum_{\substack{p\le Y\\p\notin S}}b_i(p)\frac{\log p}{p}
 =\frac12\log Y+O_{m,S}(1).
\end{equation}
The sum of the absolute class errors bounds the error for every
$0$--$1$ function $b_i$, giving the stated uniformity in types and
labels. Abelianness makes $\sigma_p$ a well-defined element specified
by the actual residue class; there is no conjugacy-class sampling
assumption hidden in this step.

The higher-power tail satisfies
\[
 0\le c_{i,p}-\frac{b_i(p)}p\le\frac1{p(p-1)}.
\]
The series $\sum_p\log p/[p(p-1)]$ converges by comparison with
$\sum_{n\ge2}\log n/[n(n-1)]$.
Multiply \eqref{eq:generalcmfirst} by $n_i$, add the tails and sum
over $i$ to prove \eqref{eq:generalcmlinear}.
Entirely selected orbits contribute zero, as their formulas require.

For the integer assertion fix $i$ and write $n=n_i$.
If $n=0$ or $1$, every floor and the asserted main term vanish.
For $n\ge2$, put $Y=\min(X,n)$.
Replacing each first-power floor by $n/p$ makes an error bounded
by $\sum_{p\le Y}\log p=O(Y)=O(n)$; therefore
\[
 \sum_{\substack{p\le Y\\p\notin S}}
 b_i(p)\left\lfloor\frac np\right\rfloor\log p
 =\frac n2\log Y+O_{m,S}(n).
\]
The full higher-power contribution is nonnegative and at most
\[
 n\sum_p\log p\sum_{r\ge2}p^{-r}=O(n).
\]
Adding proves \eqref{eq:generalcmfloors}.
Its main term saturates at $X=n$ for each direction, matching the
finite floor sum; no indefinitely growing fixed-order main term is
asserted. Subtracting the formula for a common order multiset proves
the first difference statement.

Finally assume $N>0$, let $I=\{i:n_i>0\}$, put $w_i=n_i/N$,
and set $L=\min(X,N)\ge1$.
For every $i\in I$,
\[
 0\le\log L-\log\min(X,n_i)\le\log(N/n_i).
\]
Indeed, if $X\le n_i$ the difference is zero; otherwise it is
$\log(L/n_i)$ with $L\le N$.
Multiplying by $n_i$ and adding yields
\[
 0\le N\log L-\sum_{i\in I}n_i\log\min(X,n_i)
 \le N\sum_{i\in I}w_i\log(1/w_i)\le N\log d.
\]
For the last inequality, concavity of $\log$ gives
$\sum_iw_i\log(1/w_i)\le\log\sum_iw_i/w_i=\log|I|$.
When $N=0$ every term is zero.
Together with \eqref{eq:generalcmfloors} this proves
\eqref{eq:generalcmentropy}.
\end{proof}

These comparisons concern bare lower certificates and do not upper-bound
selected theta values or the source's actual gain. Changes of order $N$
can still matter after leading terms cancel. Changing the conductor
requires control of constants not supplied here.

Prime-dependent choice is a distinct possibility. For
$\Phi=\{1,2,3\}\subset(\mathbf Z/7)^\times$,
$\max_{i\in\Phi}\mathbf1_{ai\notin\Phi}$ equals zero for $a=1$
and one for each other class: for $a=2,3,4,5,6$ the sets $a\Phi$ are
respectively $\{2,4,6\}$, $\{2,3,6\}$, $\{1,4,5\}$,
$\{1,3,5\}$ and $\{4,5,6\}$.
Its class mean is therefore $5/6$, exceeding $1/2$.
This finite calculation is not a global realization of varying local
directions with negligible cost. That construction and its norm and
selection constraints require a separate argument. Nonabelian and
reflex-field prime sampling are likewise outside the stated hypotheses.


\subsection{Actual embedding-labelled generalized harmonic witnesses}

\begin{proposition}[DP65: all CM rows with bounded filler cost]\label{DP65}
Let $E=\mathbf Q(\zeta_7)$, $\Phi=\{1,2,3\}$, and let
$\overline J$ denote the Jacobian $J$ with conjugated $E$-action and
unchanged underlying polarization. Fix source rank data and $1\le j<m$;
put $P=k+j-1$, $Q=m-j$, $n_V=P+Q+t$, $H=(t+3)/2$, $d=t+2$.
On the actual principally polarized $E$-abelian variety
\[
 C=J^P\times\overline J^{\,Q},
\]
embedding $\tau_a$ has multiplicity $P$ for $a\in\Phi$ and $Q$ for
$a\in-\Phi$. Use the source row-invariance, integral-coefficient and
row-degree equations, with all conjugate-pair contractions
\[
 \Delta_{a,i,s}=\sum_c
 \partial_{Z_{a,i,c}}\partial_{Z_{7-a,s,c}}\qquad(a\in\Phi).
\]
Thus the variables have their actual CM labels; the type predicate is
$\Phi$, in place of the old quadratic $\operatorname{InPsi}$.

Put $c_j=1+d(j-1)$ and define fixed positive partitions
\[
 A_i^*=\begin{cases}c_j&i<k-1,\\1+d(P-1-i)&k-1\le i<P,\end{cases}
 \qquad B_s^*=1+d(Q-1-s).
\]
For $a_r=11^r$, $b_r=121\,11^r$, there are source HC parameters,
word-independent at each embedding, with partitions
\[
 A^{(3)}=A^*+a_re_0,\quad A^{(2)}=A^*+b_re_0,\quad
 A^{(1)}=A^*,\qquad B^{(a)}=B^*\quad(a\in\Phi).
\]
Set
\[
 D_0=3(|A^*|+|B^*|),\quad
 \mathcal F_0=\left(\prod_iA_i^*!\prod_sB_s^*!\right)^3,
 \quad T_r=122\,11^r+D_0.
\]
There is a primitive polynomial in the full actual embedding-labelled
generalized harmonic component with these partitions, and an integral
prefactor normalized theta jet whose evaluation has valuation exactly
\[
 V_r=(11^r-1)/10+v_{11}(\mathcal F_0).
\]
There is also an actual generic permitted-translate tuple with
$M=29^e=O(T_r)$, $\kappa=M^2=O(T_r^2)$ and normalized valuation at most
$V_r$, including $M^{-T_r}$. Consequently no uniform bound
$\gamma T_r-C_0$ with $\gamma>1/1220$ and fixed $C_0$ holds for all
these evaluations as the HC family varies.

The product variety, column count and filler cost are fixed. The result
does not instantiate the old quadratic-induced
$\operatorname{Datum.dom}/\operatorname{CMObj.HasType}$, a growing
$\operatorname{SizeData}_2$ package, or a prescribed theta image.
\end{proposition}

\begin{proof}
DP49 gives the actual tangent characters $x:\tau_3$, $y_1:\tau_2$,
$y_2:\tau_1$. Conjugating the $E$-action reverses these characters and
preserves Rosati conjugation. Hence the stated multiplicities on $C$
are exact. The trace-one line on a conjugated factor is integrally
isomorphic to the original one: adjunction and symmetry give
\[
 (1,\bar w_0)^*\mathcal P_\lambda
 \simeq(w_0,1)^*\mathcal P_\lambda
 \simeq(1,w_0)^*\mathcal P_\lambda,
\]
respecting rigidifications. Thus the product has the same principal ample
trace-one lines and primitive product base section. Integral eigenspaces
exist because the order-seven projectors are integral at11 and the
characters split over the fixed unramified coefficient ring.

At an embedding $a\in\Phi$, assign a nonnegative $\eta_a$ and take
\[
 v_i^{(a)}=k-2-i\ (1\le i\le k-2),\quad
 v_0^{(a)}=k-2+\eta_a,
\]
\[
 u_i^{(a)}=v_i^{(a)}+\tfrac12\ (0\le i\le k-2),\quad
 u_{k-1}^{(a)}=-\tfrac12,
 \quad\beta_s=1+2Hs,\quad\epsilon_s=\beta_s+H\ (0\le s<m).
\]
The consecutive $v$ gaps are one except for $\eta_a+1$ at the first
row; $0>-1/2>-1=-\beta_0$ checks the last interlacing.
Both HC separation bounds equal $H$, and all integral/half-integral
conditions hold. Substitution in DP61's formulas gives exactly $A^*$
and $B^*$ above, with only the first positive entry increased by
$\eta_a$. Taking $\eta_3=a_r$, $\eta_2=b_r$, $\eta_1=0$ proves the
weight assertion. The common-row and fixed-endpoint-$u$ equations across
words hold by taking identical tuples for every word. The representative
set here is $\Phi$; no old quadratic \texttt{Fields.Tau} is inferred.

In column1 use the sharp monomial
\[
 P_r^0=Z_{3,0,1}^{a_r}Z_{2,0,1}^{b_r}.
\]
It has no variable at a negative embedding, so every conjugate-pair
contraction kills it. It is invariant under later-row additions.
Choose disjoint further column sets $C_+,C_-$ of sizes $P,Q$;
$n_V\ge P+Q+1$. With the prefix determinants from DP61 put
\[
 Q_0=\prod_{a\in\Phi}\mathcal D_{A^*}(Z_a;C_+)
                            \mathcal D_{B^*}(Z_{7-a};C_-),
 \qquad P_r=P_r^0Q_0.
\]
Prefix-row invariance and telescoping degrees give every prescribed row.
On $C_+$ there are no negative variables; on $C_-$ there are no positive
variables; on column1 there are only the two positive variables above.
Thus all contractions, including product-rule cross terms, vanish.
Different embedding pairs can reuse the same column sets because only
conjugate embeddings contract. The determinant diagonal monomial has
coefficient one, proving primitivity and nonvanishing.

The matching monomial in integral algebraic coordinates on filler
columns has initial additive-coordinate degree $D_0$. Applying
$Q_0(\partial)$ to that prefactor times the filler kernels gives exactly
$\mathcal F_0$, by the coefficient-one and lowest-degree argument of
DP63. This uses all three positive and all three negative embedding
blocks, so no unfilled actual row is hidden.

In the first ordinary $J$ factor of column1, the two selected directions
are precisely $x,y_1$. The two-variable initial-form computation in the proof of DP58 gives
\[
 v_{11}([x^{a_r}y^{b_r}]F_\lambda)=s-w(a_r+b_r),
 \qquad s=\tfrac1{10},\quad w=\tfrac{121}{1220}.
\]
There is no harmonic diagonal correction: the mixed monomial is already
harmonic for the actual conjugate labels. Both exponents are powers of11,
so $v_{11}(a_r!b_r!)=s(a_r+b_r-2)$ and the derivative valuation is
$(s-w)(a_r+b_r)-s=(11^r-1)/10$. The trace-one parameter is the unit
$-\kappa$; unit choices of primitive eigenbasis do not alter valuations.
Multiplying by the filler gives the exact displayed value.

For permitted translates, apply DP59 to each ordinary factor and its
conjugate version to each conjugate-action factor. Exchanging
$\alpha$ and $\bar\alpha$ preserves all norm, coprimality and prime-to11
properties. The allowed span on a conjugate factor is
$[\bar\alpha]^*H^0(J,L^M)$, whose roots are killed by
$\bar\alpha^2$, exactly the action of the ideal $(\alpha^2)$ on that
conjugated object. Product symmetric lifts and rigidifications give the
same generic \texttt{IsVtOver} relation on $C$.
Products of factor basis sections give a finite integral generating family
for the allowed product-translate span. Their factor jets realize every
integral total-order jet; the number of factors is fixed, so $M\ge cT+d$ suffices.
Prescribe the pure first-column kernel and the diagonal filler jets,
including the inverse integral frame units of DP59, through order $T_r$.
The interpolating sections have the exact value above. Expanding all
$n_V$ sections in these integral generating families and applying multilinearity
shows that one allowed tuple has valuation at most that value.
The normalization $M^{-T_r}$ is a unit; the least sufficient power29
gives $M=O(T_r)$ and $\kappa=O(T_r^2)$. The extracted tuple is algebraic
even when interpolation coefficients are complete-local.
Finally $V_r/T_r\to1/1220$, proving the uniform-rate assertion with
its stated varying-HC and generalized-type scope.
\end{proof}


\section{Cote 54, Cartier duality, and the crystalline setting}

\subsection{The exact question in the supplied text}

Cote 54, scan 3, writes $B=A^\vee$, identifies $t_B=H^1(A,\mathcal O_A)$ and $t_A^\vee=H^0(A,\Omega_A^1)$, and introduces
\begin{equation}\label{eq:SA}
 S_A:\quad 0\longrightarrow t_A^\vee\longrightarrow H^1_\dR(A)
 \longrightarrow t_B\longrightarrow0.
\end{equation}
It relates this to universal vector extensions. Scan 5 asks to compare the extensions and dualities with Cartier duality and places $A[F]$ inside $A[p]$. Scan 6 is an interleaved flattening argument, ``Proposition A'', not an intermediate step in the duality calculation. Scan 7 returns to the Frobenius--Verschiebung diagram, writes the corresponding finite-group extension, and asks whether the restricted-Lie extension can be recovered from $A[p]$. Its box says that $t_A^\vee$ has zero $p$-power.

\begin{proposition}[DP10: the zero operation and what Cartier duality recovers]\label{DP10}
The boxed assertion concerns the vector-group kernel of a universal extension and holds for every abelian variety in characteristic $p$. It does not assert that $\Lie A$ has zero restricted $p$-power. Over a perfect field the Cartier sequence is
\begin{equation}\label{eq:cartier}
 0\longrightarrow A[F]\longrightarrow A[p]
 \longrightarrow (B[F])^D\longrightarrow0.
\end{equation}
The filtered $F,V$-module $H^1_\dR(A)$, including its Hodge extension, is recoverable from $A[p]$ by Oda's theorem. In the usual restricted-Lie interpretation $x^{[p]}=F(x)$, the Hodge subspace has zero $p$-operation.
\end{proposition}
\begin{proof}
The invariant derivations of $\Ga^g$ are constant linear combinations of the coordinate derivatives. In characteristic $p$, their $p$th powers are zero: the summands commute, and each $p$fold derivative of a monomial has a factor divisible by $p$. Maps of smooth groups respect the restricted operation, so every vector kernel has this property. In contrast on $\Gm$, $(x\,d/dx)^p=x\,d/dx$, as is seen on every monomial. An ordinary elliptic curve has height-one formal group and nonzero restricted operation by DP6, while its universal extension still has the same additive vector kernel. This proves the distinction without assuming supersingularity.

Since $[p]=V_AF_A$, restriction of $F_A$ gives an fppf-surjection $A[p]\to A^{(p)}[V_A]$ with kernel $A[F]$. Cartier duality for isogeny kernels identifies $A^{(p)}[V_A]$ with $(B[F])^D$ \cite[Theorem 1.1 and Corollary 1.3, pp.~67--70]{Oda}, giving \eqref{eq:cartier}. The same perfect pairing makes $A[F]$ and $B[F]$ annihilators of one another, which is the orthogonality asserted in scan 5.

Write $M=M(A[p])$ for the contravariant Dieudonn\'e module. Oda \cite[Corollary 5.11, p.~131]{Oda} identifies it canonically with $H^1_\dR(A)$, respecting $F,V$ and identifying $VM=\ker F$ with $H^0(A,\Omega_A^1)$. The quotient is $H^1(A,\mathcal O_A)$. Thus the Hodge extension and its Frobenius are functorial in $A[p]$. Since $FV=0$ on $M$, $F$ is zero on this subspace.

To verify that this is the actual restricted operation, use the classical realization of the universal extension of $B$ by algebraically trivial line bundles on $A$, rigidified at the identity and equipped with integrable connection \cite[Chapter I, 2.6; Propositions 2.6.7, 3.2.3, 4.2.1; (4.6.3)]{MM}. A tangent vector has transition functions $1+\epsilon f_{ij}$ and connection forms $\epsilon\omega_i$, representing the de Rham \v Cech cocycle $(f_{ij},\omega_i)$. Tensor $p$ copies with independent square-zero parameters $\epsilon_1,\ldots,\epsilon_p$. Its transitions are
\[
 \prod_{a=1}^p(1+\epsilon_af_{ij})=\sum_{r=0}^p\sigma_rf_{ij}^r,
\]
and its connection forms are $\sigma_1\omega_i$. The symbolic restricted-power construction of SGA 3 VII$_A$, 6.2--6.3, sends $\sigma_1,\ldots,\sigma_{p-1}$ to zero and $\sigma_p$ to $\epsilon$. The resulting cocycle is $(f_{ij}^p,0)$, precisely Oda's de Rham Frobenius \cite[Definition 5.3]{Oda}. Thus the identification respects the restricted operation, not merely a chosen semilinear endomorphism. Finally an abelian restricted Lie algebra is an abelian Lie algebra with a $p$-semilinear operation; the corresponding height-one group is recovered by SGA 3, VII$_A$, 7.2--7.4 and 8.1 \cite{SGA3}. Oda explicitly identifies his comparison as answering Grothendieck's conjecture.
\end{proof}

This answer was already available in 1969. It is not made true by RH, and it is not a reconstruction of a mixed-characteristic deformation or a theta trivialization from a zero restricted operation. The local bound uses the integral $p$-series; its higher-height distinctions survive after the first restricted operation has become zero.

Cartier duality interchanges $F$ and $V$ \cite[Theorem 3.19, p.~89]{Oda}. For realizations, the Poincar\'e pairing takes values in $W(k)(-1)$ or $\Q_\ell(-1)$, so the corresponding duality is $H^1(B)\simeq H^1(A)^\vee(-1)$. The plain underlying de Rham vector-space duality in scan 3 is not a transcription error requiring an added twist.

\subsection{Cote 4 and the letter to Tate}

Cote 4, scan 3, (8.1.1), asks for derived crystalline--de Rham base change, and (8.1.2) retains, for a smooth proper $X/k$, the integral universal-coefficient sequence
\[
 0\to H^i_\crys(X/W)\otimes_Wk\to H^i_\dR(X/k)
 \to\operatorname{Tor}_1^W(H^{i+1}_\crys(X/W),k)\to0.
\]
The occurrence of $i+1$ in the Tor term is essential. The relation belongs to the established crystalline formalism, not to a new consequence of derivative divisibility; see Berthelot \cite[V.3.6]{Berthelot} and the exact historical discussion in Illusie \cite[\S3.2, pp.~9--11]{Illusie}. The integral degree-one abelian comparison continues Oda's theorem through the Dieudonn\'e crystal. Illusie \cite[\S4.2(iv), pp.~16--17]{Illusie} explains the universal vector extension and the required crystalline modification of the original formulation.

Scans 5--6 of cote 4, 8.4.2 and Lemma 8.4.3, use the Frobenius-bijective stable part; scans 7--10 retain torsion corrections and pass to \emph{\'etale} $\F_p$ local systems. Scan 12, 8.6, asks about irreducibility and full monodromy. DP1 controls an integral local derivative operator. It does not compute the monodromy of those local systems, prove the assumed integral weak Lefschetz statements, or establish a motivic level claim for vanishing cohomology. The Tate letter is not a separate attachment here; the located historical reference is its discussion in Illusie \S3.2 and \S4.2(iv), including the letter's \S2.6. We do not invent an archive scan for it.

\subsection{Placement in SGA 4, SGA 5, and the weight formalism}

Let $a:A\to S$ be an abelian scheme over a locally noetherian base, with $\ell$ invertible on $S$. Applied to the constant sheaves $\Z/\ell^r\Z$, proper and smooth base change identify the geometric stalks of $R^1a_*\Z/\ell^r\Z$ with fibre cohomology and make these sheaves locally constant: SGA 4, XII, Theorem 5.1 and Corollary 5.2; XV, Theorem 2.1; XVI, Theorem 2.1 \cite{SGA4}. The coefficient-system formalism then gives the lisse sheaf $R^1a_*\Q_\ell$ of rank $2\dim(A/S)$; see SGA 5, VI, Definition 1.2.1 and \S1.4.2, Definition 1.4.3, pp.~265--266 \cite{SGA5}. We use constant, hence locally constant, coefficients; the omitted local-constancy hypothesis in the original printing of SGA 4 XVI Corollary 2.2 is not needed to justify this application.

Over a finite-field base, its closed fibres are pure of weight one: every geometric Frobenius eigenvalue on $H^1_{\mathrm{\acute et}}(A_{\bar s},\Q_\ell)$ has complex absolute value $|\kappa(s)|^{1/2}$. This is Weil I, Theorem (1.6), p.~276, for smooth projective fibres, and Weil II, Corollary (3.3.9), p.~207, for smooth proper fibres \cite{WeilI,WeilII}. Both ordinary and supersingular abelian varieties satisfy that statement. The $\ell$-adic and crystalline realizations of $h^1(A)$ therefore retain different additional invariants: complex weights do not determine the integral $p$-series, its height, or the period lattice used in an automorphic construction. The local results above concern precisely those integral $p$-adic data. No new weight or purity theorem is asserted.

\section{Binomial structures: the precise affirmative statement}

\begin{theorem}[DP11: identity Frobenius lifts and binomial rings]\label{DP11}
Let $S$ be a torsion-free commutative unital ring. The following conditions are equivalent:
\begin{enumerate}[label=(\roman*)]
\item $\binom an=a(a-1)\cdots(a-n+1)/n!$ belongs to $S\subset S\otimes\Q$ for every $a$ and every $n\ge0$;
\item $a^p-a\in pS$ for every prime $p$ and every $a\in S$, so $\id_S$ is a Frobenius lift at every prime;
\item $S$ has a special $\lambda$-structure with all Adams operations equal to the identity.
\end{enumerate}
When it exists the structure is unique and $\lambda_t(a)=(1+t)^a$. Every ring map between such rings preserves these operations.
\end{theorem}
\begin{proof}
The equivalence (i)$\Leftrightarrow$(ii) is Elliott's criterion \cite[Theorem 4.1]{Elliott}, recorded explicitly in \cite[Remark 1.3]{Horel}. Since $\Spec S$ is flat over $\Z$, commuting Frobenius lifts are equivalent to a special $\lambda$-structure; apply Borger's supplied introduction and \S1.1 to the commuting identity lifts \cite{Borger}. This gives (ii)$\Leftrightarrow$(iii).

For the stated normalization, work in $S\otimes\Q$. Newton's identity gives
\[
 \frac{d}{dt}\log\lambda_t(a)
 =\sum_{m\ge1}(-1)^{m-1}\psi^m(a)t^{m-1}
 =\frac{a}{1+t}.
\]
The constant term is one, so $\lambda_t(a)=(1+t)^a$, proving uniqueness. A ring map preserves these rational polynomials after tensoring with $\Q$; injectivity of the torsion-free target into its rationalization proves the claimed integral compatibility. This proves only the explicit normalized and functorial form of the cited equivalence.
\end{proof}

This criterion is known: Elliott \cite[Theorem 4.1 and Proposition 8.3]{Elliott}, also Horel \cite[Definition 1.1 and Remark 1.3]{Horel}. Borger's supplied source, introduction and \S1.1, places commuting Frobenius lifts on $\Z$-flat schemes in the $\lambda$-ring framework \cite{Borger}. Torsion-freeness is necessary for the formulation above. For instance $\F_p$ satisfies the residue congruences, but is not a binomial ring in this torsion-free sense. Identity lifts at all primes on a flat affine scheme are a special structure, not a structure on every scheme over $\Z$.

\begin{proposition}[DP12: finite differences and an associated-graded bridge]\label{DP12}
Let $\Int(\Z)=\{P\in\Q[x]:P(\Z)\subset\Z\}$, with its degree filtration, and put $b_n(x)=\binom xn$. Then
\[
 \Int(\Z)_{\le n}=\bigoplus_{j=0}^n\Z b_j,
 \qquad\gr\Int(\Z)\simeq\Gamma_\Z(u)
\]
as graded Hopf algebras. The finite difference $\Delta P(x)=P(x+1)-P(x)$ is integral and satisfies $\Delta b_n=b_{n-1}$ for $n\ge1$, with $\Delta b_0=0$, whereas ordinary differentiation does not preserve $\Int(\Z)$. In a nonzero binomial ring the identity lift gives the $p$-derivation
\[
 \delta_p(a)=\frac{a-a^p}{p},
\]
which is not an ordinary derivation.
\end{proposition}
\begin{proof}
The binomial basis is the standard numerical-polynomial basis: Elliott \cite[Lemma 2.2]{Elliott}, explicitly recalled in Horel's proof of Proposition 2.2, author-PDF pp.~4--5 \cite{Horel}. Pascal's identity gives $\Delta b_n=b_{n-1}$ for $n\ge1$ and $\Delta b_0=0$. We verify the associated-graded Hopf form and the two calculus distinctions needed here.

The leading coefficient of $b_ib_j$ is $1/(i!j!)$, so its coefficient on $b_{i+j}$ is $\binom{i+j}i$, giving the divided-power multiplication in the associated graded. Vandermonde's identity
\[
 b_n(x+y)=\sum_{i+j=n}b_i(x)b_j(y)
\]
gives the corresponding coproduct; evaluation at zero is the counit, and substitution $x\mapsto-x$ induces the antipode $u^{[n]}\mapsto(-1)^nu^{[n]}$. This proves the Hopf statement. But $b_2'(x)=x-1/2$ fails to be integer-valued at zero, proving the assertion about ordinary differentiation.

DP11 makes $\delta_p$ integral and unique. Expansion gives
\begin{align*}
 \delta_p(a+b)&=\delta_p(a)+\delta_p(b)+\frac{a^p+b^p-(a+b)^p}{p},\\
 \delta_p(ab)&=a^p\delta_p(b)+b^p\delta_p(a)+p\delta_p(a)\delta_p(b).
\end{align*}
These are the $p$-derivation identities. In particular $\delta_p(p)=1-p^{p-1}\ne0$, whereas every ordinary derivation kills integers. This proves the distinction.
\end{proof}

This gives a precise sense in which operations on numerical functions provide an integral calculus. It does not make $d/dx$ integral, and $\Delta$ satisfies a translated product rule rather than the ordinary Leibniz rule. The associated graded in DP12 is essential; the filtered ring itself is not being identified with a divided-power polynomial ring.

\begin{theorem}[DP13: the precise homotopy consequence]\label{DP13}
For a simplicial set $X$, define the cosimplicial ring $A_X^n=\operatorname{Map}(X_n,\Z)$. It has canonical degreewise binomial operations and identity Adams operations. If $X$ is connected, nilpotent and of finite type, then
\[
 X\xrightarrow{\sim}\RMap_{\mathrm{cBRing}}(A_X,\Z).
\]
For such $X,Y$ the resulting functor is fully faithful, with derived mapping spaces
\[
 \operatorname{Map}(X,Y)\simeq\RMap_{\mathrm{cBRing}}(A_Y,A_X).
\]
\end{theorem}
\begin{proof}
Each degree is a product of copies of $\Z$. Binomial operations are defined pointwise, and pullback of functions along the simplicial maps preserves them. Additive normalization is the usual integral simplicial cochain complex. Horel's model structure and reconstruction theorem \cite[Theorems 3.7 and 5.4]{Horel} apply to this cosimplicial binomial ring and give the unit equivalence. The mapping-space equivalence follows from the derived adjunction and the unit equivalence for $Y$. This is an application of the cited reconstruction theorem, not a new proof of that theorem.
\end{proof}

The cosimplicial and derived qualifiers cannot be omitted. The ordinary cup-product cochain complex is not being asserted to be a strictly commutative binomial ring. For connected finite-type $X$ without the nilpotence hypothesis, Horel's Theorem 6.1 identifies the unit with Bousfield--Kan $\Z$-completion, not an automatic equivalence.

\begin{proposition}[DP14: the proposed identifications that fail]\label{DP14}
The formal-disc coordinate ring $\Z_p[[t]]$ does not have the identity as a Frobenius lift at $p$. The binomial homotopy construction does not give integral ordinary derivatives by identifying its coefficient ring with this formal-disc ring.
\end{proposition}
\begin{proof}
The coefficient of $t$ in $t^p-t$ is $-1$, so $t^p-t\notin p\Z_p[[t]]$. This violates DP11(ii). Independently, DP12 gives an integer-valued polynomial whose ordinary derivative is not integer-valued. These counterexamples refute the claimed identifications. They do not preclude a richer comparison functor with additional structure, but no such functor has been constructed here.
\end{proof}

Suggestion 9 is therefore affirmative in the exact sense of DP11--DP13: identity Frobenius lifts under flatness provide an established integral homotopy model. Its further identification with the multiplicity mechanism is not a consequence of that theorem. The precise proposed B7 comparison statement was not supplied. Targeted public recovery does, however, establish exact historical locators \cite{PublicArchive}: cote 153, physical PDF p.~3 (archive/manuscript p.~2), defines the binomial polynomials, displays their product and composition expansions, explicitly annotates integral product coefficients, and writes their $\Z$-span in $\Q[T]$. Cote 136, PDF pp.~4--6 (archive pp.~3--5), moves from rational Sullivan models and loss of torsion through integration denominators to extra PD-base data. On PDF p.~6 (author's circled p.~4), the failure of $1-\sum X_i$ to lie in the PD ideal leads to a coefficient ring $S$, an element $t$ in its PD ideal, and the relation $\sum X_i=t$. These are precise motivations for the distinctions in DP5 and DP14, not a theorem transferring binomial homotopy reconstruction to theta multiplicity.

A second targeted facsimile pass closes the remaining historical page locators. Cote 133, PDF p.~6 (archive and handwritten p.~5), calls $H^0(S,\Z)$ the binomial ring, describes $WG(S)$ as a pre-$\lambda$-ring augmented to it, and explicitly asks whether the universal identities for $\lambda^i(xy)$ and $\lambda^i(\lambda^j x)$ hold. PDF pp.~27--29 (archive pp.~26--28) describe sum, product and composition of finite-set functors and then $\lambda$- and symmetric-group representation operations. These are exact composition-content locators; no claim that the literal word ``plethysm'' occurs there is required.

Cote 137, PDF pp.~38--39 (archive and handwritten pp.~37--38), proposes the cochain functor $F_X(M)=\operatorname{Hom}(X,M)$ with tensor and $\Gamma$-operation maps, and asks how far its multiplicative divided-power structures reconstruct the originating homotopy type. This is a historical question, not a proved reconstruction theorem. DP13 supplies a precise modern theorem in its stated cosimplicial binomial category; identifying that category with every unstated manuscript axiom or with an unspecified B7 bridge is not asserted.

The official catalogue dates cote 153 from 1982; the four cotes should not all be dated to 1976. The earlier bounded pass read twelve specified public pages; the second pass directly read six further pages, with exact page labels, image hashes and source-reading events in \path{research/public_archive_q5e_second_pass.md} and its pixel metadata. The earlier account in \path{research/public_archive_q5e_sources.md} remains a truthful record of its earlier scope. A typescript headed ``Hirsch Method'' on another inspected cote 137 leaf is not silently attributed to Grothendieck.

\section{Cotes 16 and 47: a different divided-power calculation}

\begin{proposition}[DP15: exterior divided powers do not imply integral Lefschetz]\label{DP15}
Let $M=\Z^{2g}$ with symplectic basis $e_i,f_i$ and $\xi=\sum_i e_i\wedge f_i$. Then $\xi^m/m!\in\bigwedge^{2m}M$ is integral. Nevertheless, for $g=3$, multiplication by $\xi$ from the span of the three $e_i\wedge f_i$ to the span of their pairwise products has cokernel $\Z/2$.
\end{proposition}
\begin{proof}
Put $p_i=e_i\wedge f_i$. The degree-two elements commute and satisfy $p_i^2=0$. Expanding gives
\[
 \frac{\xi^m}{m!}=\sum_{|S|=m}\prod_{i\in S}p_i,
\]
which is integral. For $g=3$, order the target basis by the pair complementary to each single index. The matrix is $J-I$, with zero diagonal and ones elsewhere. Its determinant is $2$; a $1\times1$ minor and a $2\times2$ minor are units. Its Smith invariants are therefore $1,1,2$, proving the cokernel assertion. This is a direct block of the full map $\xi\wedge-:\bigwedge^2M\to\bigwedge^4M$: no other exterior basis monomial can contribute to a target consisting of two complete pairs. Hence the full integral Lefschetz map also fails to be an isomorphism.
\end{proof}

This is the concrete integral exterior-algebra setting of cote 16, scan 14, \S7 and the ring $\Phi_0$, and the smallest test in the supplied H8 \cite{Hunches}. It does not prove H8(iv)'s index statement for an integral order, and no relation of its Lefschetz strings to the step-8 theta chain is supplied by DP1. The familiar occurrence of factorials is not a comparison theorem between those two objects.

\begin{proposition}[DP16: forced coefficient corrections in cote 47]\label{DP16}
For $\chi\in\Z\setminus\{0\}$, conditional on the displayed total-Chern generating formula in cote 47, scan 9, interpreted degree by degree over $\Q$,
\[
 c_{\rm tot}=\left(1-\frac{D'}\chi\right)^\chi,
\]
the degree-$n$ coefficient is
\begin{equation}\label{eq:chern}
 (-1)^n\frac{\chi(\chi-1)\cdots(\chi-n+1)}{n!\chi^n}D'^n.
\end{equation}
The recorded quadratic, cubic and general expansions each lack one power of $\chi$ in the denominator. This does not establish the generating formula in integral Chow or $K$-theory.
\end{proposition}
\begin{proof}
Expand by the binomial polynomial: $\binom\chi n(-D'/\chi)^n$ is exactly \eqref{eq:chern}. At degrees two and three the denominators, before cancelling the numerator's first $\chi$, are $2\chi^2$ and $3!\chi^3$, respectively. The recorded denominators are $2\chi$ and $3!\chi^2$. The general recorded denominator is $n!\chi^{n-1}$ with the initial $\chi$ still in the numerator. Each mismatch follows by direct algebra from the displayed generating formula.
\end{proof}

The final integrality expression in the same scan, with denominator $k!\chi^{k-2}$ after cancelling the initial $\chi$, is \emph{not} a forced error: it is, up to sign, $\chi$ times the actual $k$th Chern coefficient, so can be a deliberately weaker valid integrality consequence. It is excluded from the forced-correction rows. Distinguishing this from the genuinely inconsistent expansion is part of the proof check.

\subsection{Transcription candidates and exact scope}

There are five review rows in \texttt{TRANSCRIPTION\_NOTES.tsv}. Their line numbers refer to the supplied \LaTeX\ files. The official public PDFs \cite{PublicArchive} were subsequently recovered and the target pages independently inspected. Supplied scans 54/005, 54/007 and 47/009 match physical PDF pages 5, 7 and 9, carrying archive/manuscript labels 4, 6 and 8 respectively. Matching the page content, rather than assuming a cover-page offset, establishes this correspondence.

The autograph itself visibly has $\Omega_A^p$ in the cohomology on $B$, the upper-right diagram vertex $V$, and the Chern denominators $2\chi$ and $3!\chi^2$ with the displayed initial $\chi$ retained in their numerators. These four supplied readings are therefore faithful to mathematical slips in the source. The general exponent is small and faint, compatible with $n-1$; its visual reading remains qualified. The mathematical emendations are still required: the sheaf on $B$ is $\Omega_B^p$, the codomain of $V:A^{(p)}\to A$ is $A$, and DP16 gives the Chern coefficients. Preserve the diplomatic readings and append these corrections as editorial notes; do not silently classify them as transcription errors. The boxed $t_A^\vee$ is also correctly transcribed; its proposed supersingularity interpretation was the error. No new confirmed transcription error is asserted by these five rows.

\part{Unsettled applications, exact gates, and limits}

\section{What is not proved in this package}

There is no unlabelled draft proof in the checked part. The following are open applications or unavailable source identifications, not established findings.

\subsection{U1: source-identification gate closed}

The initial attachment-only gate was to locate the theta normalization and period-to-jet interface. That gate is now closed by Section~\ref{sec:recovered}, the exact declared comparisons, and byte verification against the author's archive. The answer to item 2 is precise: the scalar leading coefficient is universal height two; its theta application additionally uses normalization and geometric comparison data, and the later multivariate construction has the listed kernel, minor, path and denominator hypotheses. A whole existing theorem is not reduced to a bare height condition. This conclusion uses the accepted source theorems and does not claim an independent reconstruction or compilation of Alp\"oge's complete argument.

\subsection{U2: whether a different higher-height construction shortens the chain}

The literal unchanged-logarithm and balanced-product proposals are refuted by DP21 and DP23. The question about a \emph{different} construction remains open. DP25 identifies the first geometric requirement: a theta correspondence for the needed mixed CM type, replacing the source's induced-type/Serre-tensor comparison and specifying its source and terminal representations. DP28--DP31 close the generalized two-scalar analytic test, with the full $c_h$ slope under its stated hypotheses.

For the actual dual pair, the local model, primitive class and quantitative ordinary paths are now explicit. DP32--DP34 give the dual and its sharp derivative estimates; DP35--DP38 construct the primitive frame and its first poles. DP40 and DP43 extend the sharp directional and optimal affine bounds to all cyclic types, and prove that the equally weighted type-plus-dual sum is largest for alternating type. DP41 computes the entire primitive logarithm and proves the exact multivariate exponential threshold, robust under every nonzero integral prefactor. DP42 determines every mixed logarithmic line with primitive projections in the stated unramified lattices: its best factorial slope is $\rho_2$, independent of $h$, and longer restricted-$p$ chains give less. These settled families cannot be left as an unspecified local-kernel gap.

DP46 removes the unit-one restriction for all singleton/complement cyclic models and their invariant integral frames. DP47 identifies the actual trace-one diagonal on a prime-to-$p$ polarized conjugate-idempotent block with one primitive Poincar\'e factor, up to an integral unit and a quadratic logarithmic gauge; DP48 proves that such a quadratic term cannot improve the ordinary primitive-line slope. DP44 also closes the generic smaller-rank construction issue: every smaller even count at least two admits compatible rank/HC/host data and the actual Hermitian Spaces constructor. DP52 now transports the exact complex parameter conditions except for the explicit remaining bound $1801\le k+M'-1+t$, proves the six-input nonzero-period induction at every admissible even count, and constructs the actual mixed-type endpoint period under weaker conjugation assumptions. Its next nonautomatic input is identification with the prescribed iterated theta image and the selected quantitative ties.

DP49 realizes the actual polarized theta-germ identity on the genus-three Jacobian, with a nonzero restriction in every primitive residue class. DP50 determines the exact nonunit-tensor kernel and proves that the source $M^{-n}$ level normalization outweighs this ordinary-path improvement. DP51 supplies the sharp universal bound for all integral homogeneous differential polynomials, including the exact digit-sum loss, and proves sharpness already on the specified two-column harmonic module. The remaining local task is to identify the prescribed geometric sections and row-degree sequence, actual connection lattice and corrections, selected minor module and geometric path. If they lie in DP42's class, that theorem supplies the exact local rate; a proposed improvement must then use a genuinely different selection or path mechanism. DP39 proves why full-series nonintegrality and summed harmonicity alone decide neither the prescribed selected value nor its quantitative bound. The source supplies integral conjugated connections and gauge independence. The quantitative universal polynomial bound is now settled; any improvement must use the actual geometric section or the prescribed selected weights and must establish controlled nonvanishing for those data. Finally the construction must recompute the period main term and the norm, rank, level, polarization and covolume costs, and prove a net inequality shortening the integer lift count.

DP54 and DP56 close finite global-jet existence with a fixed linear interpolation threshold. DP58 gives explicit sharp degrees uniform in every unit tensor parameter, and DP60 realizes them over the number field with constant $p$-adic precision, no new good-prime denominators and logarithmic section size $O(n)$. DP59 identifies the exact smaller integral span of the generic permitted torsion translates and proves that individual permitted translates or pairs already attain the sharp values at polynomial levels. DP62 proves their canonical section norms are exactly the fixed base norm, extends their normalized theta lifts integrally across every good level prime, and bounds the fields of the actual tuples. These results do not identify the source's level-dependent analytic gauge or its selected automorphic image.

DP55 and DP57 close intrinsic endpoint detection and its actual adelic realization: over the explicit unramified CM host $\mathbf Q(\zeta_{28})$, a fresh split principal level yields an actual period of value one. The old induced-quadratic type wrapper still fails. DP61 identifies the precise positive and reversed negative row partitions. The unchanged two-row witness is excluded from that component; a primitive determinant completion exists under explicit gap inequalities, but its pure-kernel evaluation is zero. DP63 supplies an integral prefactor with fixed factorial cost in a scalable HC family. DP65 then realizes every row with its actual CM embedding label on $J^P\times\overline J^{\,Q}$: the sharp mixed monomial lies in two distinct positive embeddings, so all conjugate-pair contractions vanish, and the permitted-translate span produces a tuple with rate at most $1/1220$. The remaining question is compatibility with one common unramified moduli/host family, its actual theta image, the growing quantitative windows and all source costs.

DP64 also removes a broader bare-prime alternative: within any fixed cyclotomic modulus, every CM type and every distribution of a fixed total order among finitely many globally fixed directions has the same leading integer budget, including rounding and saturation, up to $O(N)$. A direction that changes with the prime can have a larger formal first-coefficient average, but requires a global realization and its arithmetic cost. Neither this certificate comparison nor the norm conclusions upper-bound all actual selected theta gains.

The following attempted routes have definite stopping points:
\begin{enumerate}[leftmargin=*,label=\arabic*.]
\item \textbf{Substitute $c_h$ into the existing slope input.} The functional and its lattice are now located. Its entries involve both $\gamma$ and $\rho$, fixed logarithms, and exponential comparison identities. DP21 refutes retaining the original scalar logarithm; a one-coefficient change does not realize the typed input.
\item \textbf{Increase the CM dimension.} DP8 gives the complementary dimension forced by the CM type, and DP22 constructs an actual height-three example. Neither computes the cost of a new global construction.
\item \textbf{Use total height.} This route is refuted by $\widehat\Gm^{\,r}$ in DP8.
\item \textbf{Use the ratio of local gains.} DP9 bounds the ratio by $1+1/p$, but it does not divide the length of a theta chain by that ratio.
\item \textbf{Change a frame or rescale a section.} DP4 and the unit/scalar examples track the valuation changes while generic-fibre multiplicity need not change. A nonunit also changes the tangent lattice.
\item \textbf{Use Lefschetz or Arthur-$\mathrm{SL}_2$ terminology.} The supplied suggestions give an analogy, not a functor identifying integral filtrations. DP15 exposes a separate integral obstruction; no comparison theorem is supplied.
\item \textbf{Use the improvement after summing primes.} DP17 proves that pure replacement changes the fixed ramp certificate by only $O(\sum A_t)$. It preserves a leading logarithmic coefficient; it does not exclude all lower-order savings or upper-bound actual geometric gain. DP53 treats the actual cyclotomic prime redistribution: for arbitrary derivative orders fixed across primes, the linear bare slope budgets share leading term $(N/2)\log X$ and differ by $O(N)$. The difference of the integer periodic-floor lower certificates is also $O(N)$, uniformly in the cutoff. It therefore includes gains at new prime classes, without replacing an actual selected theta value by that lower certificate.
\item \textbf{Keep the old error estimate with fewer lifts.} DP18 exposes the unchanged $2x/m+4m$ balance. A bounded coefficient improvement cannot change that allowance's exponent; a new error estimate is a different possibility.
\item \textbf{Balance a higher-height line with its Cartier dual, or restrict scalars.} DP23 proves that these concrete operations do not produce the old alternating rank-one blocks. The example DP22 also has the wrong orbit and slopes for that unchanged interface. DP25 identifies its earlier failure at the actual CM-type and surjective-comparison fields.
\item \textbf{Enlarge the CM field while preserving the induced type.} DP25 proves that this keeps the inert-prime factors at slope $1/2$; the explicit height-three Jacobian fails the existing type and inert-prime inputs.
\item \textbf{Save a finite number of lifts, truncate, or combine kernels.} DP26 proves a parity-compatible numerical saving can occur in a bare certificate. DP27 and the cited chain definitions show why truncation or formal kernel regrouping does not preserve the old endpoint or reduce its arithmetic cost. A different geometric correspondence remains to be built.
\item \textbf{Change the logarithm and prove the generalized scalar kernel estimate.} DP28 constructs a full-hypothesis counterexample to the direct intermediate estimate. DP29 recovers a valid bound and an exact nonzero example. DP31 then proves the full scalar slope using integral biextension trivialization. This route is closed under its stated scalar hypotheses; it does not realize the genuinely dual multivariate kernel.
\item \textbf{Assign the same estimate to the Cartier-dual directions.} DP30 refutes that order-$p$ assertion for an actual dual. DP32--DP34 then determine the actual positive asymptotic slopes and optimal affine constants; there is no residual uncertainty about those displayed directions. Their sum is $c_h$, but an equally weighted mean does not determine the geometric cost.
\item \textbf{Use the actual dual's unequal weights in the kernel argument.} DP36 restores the symmetry quantization and proves the stated sufficient bound. DP37 realizes its primitive local class. DP41 then proves the full multivariate inflation threshold exactly, and DP42 replaces the sufficient last-axis residual $\rho_h$ by the sharp $\rho_2$. It classifies all primitive mixed lines, while leaving their identification with a prescribed global selection explicit.
\item \textbf{Optimize the whole cyclic CM type.} DP40 proves that the equally weighted sum across a cyclic type and its dual is at most $hc_2$, with equality exactly for alternating type. Individual directions can exceed $c_2$, so unequal selections need a separate estimate. DP43 gives their exact codimension affine loss and periodic floor sums; it does not identify the global cost with a directional sum.
\item \textbf{Clear the primitive exponential, or discard it by harmonicity.} DP35 proves that bounded denominators force an integral splitting and DP38 computes the first primitive poles and unit jets. DP39(a) nevertheless shows that unbounded full-series denominators can be invisible to selected values. DP39(b) then refutes the opposite shortcut: the exact summed harmonic algebraic conditions alone do not recover the bare integral-function divisibility. DP41 further rules out clearing the full primitive exponential by any nonzero bounded integral prefactor. DP51 now determines the universal selected-polynomial rate and its sharp harmonic family. It leaves additional divisibility of the prescribed geometric sections and degree sequence as explicit inputs.
\end{enumerate}
The recovered source has the period/jet comparison and explicit global gain and error budgets; there is no free height parameter in those displayed interfaces. The present report proves specific substitutions false and a numerical replacement insufficient for the stated asymptotic changes. Step 8 nevertheless asserts an achieved multiplicity for a particular chain, not an upper bound on every possible chain. Therefore claiming ``no shorter proof exists'' would be as unsupported as claiming that one has been constructed.

\subsection{Other referenced programmes}

The local calculations do not settle H8(iv), the integral Chow/$K$-theory question of cote 47, the monodromy questions of cote 4 scan 12, or unspecified questions in cotes 133, 153 and 136--137. These are recorded as unaffected implications, not silently added to the proof claims. No conditional result in the cited literature is made unconditional by Alp\"oge's RH/GRH \emph{through the arguments in this report}; the local and homotopy results used here are already independent of that premise. No result is declared conditional on Masullo in the absence of a named supplied statement.

\section{Framework and audit map}

{\small
\begin{longtable}{@{}>{\raggedright\arraybackslash}p{.17\textwidth}>{\raggedright\arraybackslash}p{.37\textwidth}>{\raggedright\arraybackslash}p{.39\textwidth}@{}}
\toprule Finding & Mathematical setting & Exact text/framework locator\\\midrule
\endfirsthead
\toprule Finding & Mathematical setting & Exact text/framework locator\\\midrule
\endhead
DP1--DP4 & Integral formal completion over $W(k)$; invariant distributions and their lattices & Fable suggestion 2; Katz Corollary 1.8; EGA IV$_4$ 16.11.2; SGA 3 VII$_A$ 1--3.\\
DP5--DP6 & Divided powers, restricted Lie operation and height-two elliptic formal groups & Cote 54 scans 3, 5, 7; SGA 3 VII$_A$ 6.2--6.4 and 7.2--7.4; note p.~3 step 2.\\
DP7--DP9 & Integral CM directions and local $p$-divisible factors & Tate expos\'e 352 \S5, (10*) and Theorem 3; note p.~4 steps 8--10, 14, 18.\\
DP10 & Universal extensions, Cartier duality and de Rham/Dieudonn\'e modules & Cote 54 scans 3, 5--7; Oda Corollaries 1.3 and 5.11, Theorem 3.19; Mazur--Messing locations in the proof.\\
Crystalline discussion & Integral derived base change, PD-site descent, $F$-crystals & Cote 4 scan 3 (8.1.1)--(8.1.2); Berthelot V.3.6; Stacks 07JH; Illusie \S3.2 and \S4.2(iv).\\
DP11--DP14 & Flat schemes over $\Z$, $\lambda$-rings, numerical functions and binomial affine homotopy & Borger introduction and \S1.1; Elliott Theorem 4.1; Horel Theorems 3.7, 5.4, 6.1; recovered cote 153 PDF p.~3/archive p.~2 and cote 136 PDF pp.~4--6/archive pp.~3--5.\\
DP15--DP16 & Exterior Lefschetz operators; rational Chern expansions and integral-cycle questions & Cote 16 scan 14 \S7; H8; cote 47 scan 9.\\
DP17--DP18 & Prime-weighted arithmetic degrees and the explicit global size allowance & \path{LocalRingCardinality/DataParams}\textsubscript{2}\path{.lean}, \path{vonMangoldtClassSum}; \path{SizeWindowBrackets/GrowthBoundGivesRiemannHypothesis.lean}, \path{SizeWindows}; note p.~4 steps 9--10, 14.\\
DP19--DP20 & Integral unramified CM action, ramified deformation, and the formal derivative lattice & Lubin--Tate 1966 Proposition 1.1, pp.~50--51; Tate expos\'e 352 \S5; caveat to the proposed height-only extension.\\
DP21, DP23--DP24 & Fixed logarithm, Newton slopes, and controlled nonvanishing & \path{GaussJetBounds/Statement.lean}, \path{oneLog}; \path{CartierLocalBounds.lean}, \path{vLe}, \path{vGe}, \path{ThmP}; Tate expos\'e 352 \S5.\\
DP25--DP27 & Induced CM types, finite ramp counts and geometric chain endpoints & \path{CMLattices/IsDiffOf.lean}, \path{CMObj}; \path{LocalRingCardinality/Params.lean}, \path{Ranks}; \path{DataSpaces.lean}, \path{DataKernels.lean}; Tate expos\'e 352 \S5.\\
DP28--DP29, DP31 & Scalar formal kernels, integral biextensions and the full derivative slope & Hazewinkel Lemma 2.3; Zink \S4.4, Proposition 121, (243), Theorem 122; recovered \path{weightQuantization} and \path{LemmaS}.\\
DP30, DP32--DP34 & Actual Cartier dual, restricted Lie algebra, integral monomial-cycle logarithm and sharp affine directional bounds & Hazewinkel Lemma 2.3, (2.8), \S5; Tate 1967 \S2.2 diagram (3), \S2.3 Proposition 3; cote 54 scans 5--7.\\
DP35, DP37--DP38 & Primitive formal Poincar\'e biextensions, invariant frames and exact first poles & Zink Proposition 106, Lemma 111, (196), (211), (218), Proposition 121 and Theorem 122; Tate 1967 \S4 Corollary 2, p.~180; cote 54 scans 5--7.\\
DP39 & Harmonic differential selection and conjugated integral connections & Recovered \path{RowDeg.lap}, \path{HaO}, \path{ExpConjConnectionIntegral} and \path{CanonGaugeIndependenceAll}; precise declarations in Section~\ref{sec:connection}.\\
DP40, DP43 & Cyclic integral displays, optimal distribution lattices and CM-type comparison & Zink normal decomposition p.~129, Theorem 5 p.~131, Definition 19/(22), Theorem 81/(145)--(156); SGA 3 VII$_A$ 1--3.\\
DP41--DP42 & Exact primitive logarithm, exponential denominator radius and restricted-$p$ paths & Zink (147)--(156), Proposition 106, (196), (210)--(211), (218); cote 54 scans 5--7; the precise new estimates are proved here.\\
DP36 & Unequal weighted valuations on a genuine formal Cartier-dual pair & DP32's integral model; recovered \path{weightQuantization} and \path{LemmaS} motivate the explicitly stated generalized hypotheses.\\
DP44 & Smaller even rank data, Hermitian spaces, endpoint periods and exact gain/cost sums & \path{Ranks}, \path{HCParam}, \path{Host}; \path{hostSpacesOfNorm} 137--142; \path{spacesOfRanks} 325--386; \path{NZRowE}, \path{ExtraNZIn}, \path{LinkG_at}, \path{LinkDIV_at}.\\
DP45 & Actual integral CM Jacobian eigenspaces and Hodge filtration & DP22's FGA inputs; Tate 1967 \S1.4 and \S2.2--2.3; Zink Theorems 5 and 9 and Definition 19/(22).\\
DP46--DP48 & Unit-arrow primitive pairings, polarized trace-one pullbacks, and quadratic gauge radii & Zink (216)--(220), Proposition 121 and Theorem 122; source \path{IsCanonGauge}, \path{hL}, \path{CanonGaugeSubIsBiadd}; cote 54 scans 5--7.\\
DP52 & Smaller complex parameters, nonzero-period induction and actual mixed-type endpoint periods & \path{DataTheta.TPar}, \path{CayleySignature.Hyp3/AdmS}; \path{Fields.Kq/Ok}; \path{formGSD_ok} 83--119; \path{periodC} 1568--1590.\\
DP49 & Actual polarized Jacobian theta germs and nonzero primitive residue-class paths & DP22's actual model; Zink Theorems 5 and 9; trace-one source condition; Stacks Lemma 10.51.4, Tag 00IP.\\
DP50 & Nonunit tensor powers and the source degree normalization & Variable-ghost primitive pairing; \path{IsThRule.thExpr}, lines 66--70, with tensor $M^2$ and factor $M^{-n}$.\\
DP51 & Sharp integral polynomial and harmonic distribution bounds & \path{RowDeg.lap}, \path{harmK}, \path{RowInv}, \path{RowDeg}, \path{HaO}, lines 21--48; \path{harmA}, lines 86--87; DP42 and DP46.\\
DP53 & CM prime redistribution and integer periodic-floor certificates & Tate expos\'e 352 \S5; actual cyclic distributions of DP43; Rudnick 2015 \S2, p.~2, equation (1).\\
DP22 & A CM Jacobian and its integral formal direct factors & FGA expos\'e 232 Theorem 3.1; expos\'e 236 Corollaries 1.7, 2.3, Proposition 2.10(ii); Tate 1967 \S2.2 Proposition 1 and \S2.3 Example (b).\\
DP54, DP56, DP58, DP60 & Integral finite jets, ample line powers and arithmetic section norms & Milne abelian varieties Propositions9.2/10.1; Stacks0D2N,0B5T,02KH,01VT,01Q3; exact interpolation and residue proofs here.\\
DP55, DP57 & Exact endpoint Gram pairing and unramified adelic host & \path{TiedPer}79--252; \path{TiedValFi}239--255; \path{RecLevelStep}20--51; \path{InputsGenus}17--19; Stacks033C.\\
DP59, DP62 & Integral theta group, torsion-translate lattice and canonical metric & Mumford1966§1(II), Theorem1, Proposition3; Polishchuk§1 Proposition1.2/Corollary1.3; Moret-BaillyII.2.1; Stacks0BX7.\\
DP61 & Exact row partitions, determinant harmonicity and finite windows & \path{DataKap}13--116; \path{AMiEqAMinusQ}20--66; \path{FWExist}518--580; \path{RowDeg}21--48; source windows148,160--164.\\
DP63, DP65 & Full-row local theta jets, exact factorial fillers and actual CM embedding labels & DP49 character table; DP58 initial-form calculation; DP59 integral translate span; DP61 HC and harmonic source interfaces.\\
DP64 & Fixed cyclotomic prime distributions and saturated integer certificates & DP40/43 cyclic distribution lattices; Rudnick2015§2, p.~2, equation(1); all new floor and entropy deductions proved here.\\
\bottomrule
\end{longtable}
}

Every current proof in Part I has a recorded complete reading. The event file distinguishes an author's whole-proof reading, an independent re-proof, finite exact-arithmetic checks, source statements used and rechecked, and background carried without a new proof. It records model, run, actor, date and actual outcome. The runtime did not expose a numeric effort setting; the record says so instead of inventing one. No Lean formalization of the findings in this paper was produced.

The final package also includes the dated ledger, a question-by-question status table, implications, the transcription candidates, and research notes documenting failed approaches and review corrections. The original attachments suffice to recover the supplied inputs; cited external sources are identified by URLs and locators rather than redistributed as entire third-party papers.

\begin{thebibliography}{99}
\bibitem{Alpoge} L.~Alp\"oge, \emph{Remark}, supplied four-page PDF; p.~3 timeline step 2 and p.~4 steps 8--10, 12--15, 18. \url{https://alpo.ge/remark.pdf}. Supplied final Lean declarations: \texttt{ALPOGE\_Main.lean}, importing \texttt{Statements} and \path{Supplements.Kernels.MainTheoremsAll}; definitions: \texttt{ALPOGE\_Statements.lean}. Step-to-file evidence: supplied \texttt{ALPOGE\_LEAN\_MAP\_20261120.md}, \S4 rows 2 and 8. Cited as directed by the prompt; not independently recompiled here.
\bibitem{AlpogeSource} L.~Alp\"oge, accompanying Lean source archive, \url{https://alpo.ge/lean.tar.gz.part01} through \path{part14}; 239,275,554 compressed bytes. The expanded targeted inventory contains 1029 modules, recovered either at mirror commit \path{7cbf7a0c32fca1badf05931879a37c3f8a2b03b4} or directly from the author's archive. The recorded sweep compared 526 local copies byte for byte with the original archive, all matching; 503 further modules were extracted directly from that archive with exact member and byte hashes, as recorded in the successive manifests. Exact declaration locators and actual reading scopes are in Sections~\ref{sec:recovered}--\ref{sec:connection} and the included source/review manifests. Most recovered modules were not read in full. Mirror root: \url{https://github.com/CocoSgt/riemann-hypothesis-lean}. Source identity verification is not a Lean compilation or proof audit.
\bibitem{Fable} Fable programme suggestions, run 20261120, \texttt{FABLE\_PROGRAM\_SUGGESTIONS\_20261120.md}, suggestions 2 and 9, with its explained companion. Supplied, explicitly unchecked proposals.
\bibitem{Hunches} \texttt{WARGAME\_MATH\_HUNCHES\_20261112.md}, run 20261112, H8, especially its exterior divided-power statement, smallest test, and clause (iv). Supplied.
\bibitem{Archive} A.~Grothendieck, archive transcriptions supplied in this task: cote 54 scans 3, 5--7; cote 4 scans 2--12, especially (8.1.1)--(8.1.2), 8.4.2--8.4.3, (8.5.1), 8.5.9, 8.6; cote 16 scans 14--18, especially scan 14 \S7; cote 47 scan 9. Locators refer to the filenames and mathematical labels of the supplied transcriptions.
\bibitem{PublicArchive} Universit\'e de Montpellier, \emph{Archives Grothendieck}, official public catalogue and digitized manuscripts. \url{https://grothendieck.umontpellier.fr/archives-grothendieck/}. Cote 54: \url{https://grothendieck.umontpellier.fr/54.pdf}, physical PDF pp.~5 and 7, archive/manuscript labels 4 and 6. Cote 47: \url{https://grothendieck.umontpellier.fr/47.pdf}, physical PDF p.~9, label 8. Cote 153: \url{https://grothendieck.umontpellier.fr/153.pdf}, PDF pp.~2--3, archive labels 1--2. Cote 136: \url{https://grothendieck.umontpellier.fr/136.pdf}, PDF pp.~4--6, archive labels 3--5. Cote 133: \url{https://grothendieck.umontpellier.fr/133.pdf}, PDF pp.~6, 27--29, archive/manuscript pp.~5, 26--28. Cote 137: \url{https://grothendieck.umontpellier.fr/137.pdf}, PDF pp.~38--39, archive/manuscript pp.~37--38. Public recovery metadata, hashes, exact page mapping and bounded visual-check provenance are included in the research notes.
\bibitem{Katz81} N.~M.~Katz, \emph{Divisibilities, congruences, and Cartier duality}, J. Fac. Sci. Univ. Tokyo Sect. IA Math. \textbf{28} (1981), no.~3, 667--678 (catalogued publication 1982). Theorem 1.6 p.~668; Corollary 1.8 p.~669; Main Theorem 3.11(I) p.~672; \S IV (4.1)--(4.14), pp.~675--677. \url{https://web.math.princeton.edu/~nmk/old/divcongCartier.pdf}.
\bibitem{Katz77} N.~M.~Katz, \emph{Formal groups and $p$-adic interpolation}, Ast\'erisque \textbf{41--42} (1977), 55--65. \S IV main theorem pp.~60--61; \S V, Lemma 1 and Corollary 2, pp.~62--63, with the trace-zero hypothesis. \url{https://www.numdam.org/article/AST_1977__41-42__55_0.pdf}.
\bibitem{BK} K.~Bannai and S.~Kobayashi, \emph{Integral structures on $p$-adic Fourier theory}, Ann. Inst. Fourier \textbf{66} (2016), no.~2, 521--550. \S2 p.~527; \S3 definition p.~528; Propositions 3.3 and 3.5 pp.~529, 531. \url{https://www.numdam.org/article/AIF_2016__66_2_521_0.pdf}.
\bibitem{LT} J.~Lubin and J.~Tate, \emph{Formal complex multiplication in local fields}, Ann. of Math. (2) \textbf{81} (1965), 380--387, initial construction and Lemma 1. \url{https://math.bu.edu/people/jsweinst/AWS/Files/LubinTateFormalCMinLocalFields.pdf}.
\bibitem{LT66} J.~Lubin and J.~Tate, \emph{Formal moduli for one-parameter formal Lie groups}, Bull. Soc. Math. France \textbf{94} (1966), 49--59. Proposition 1.1, pp.~50--51, and the normalization immediately preceding it; Theorem 3.1, p.~56. Only Proposition 1.1 is needed for DP20. \url{https://www.numdam.org/item/10.24033/bsmf.1633.pdf}.
\bibitem{Oda} T.~Oda, \emph{The first de Rham cohomology group and Dieudonn\'e modules}, Ann. Sci. \'Ecole Norm. Sup. (4) \textbf{2} (1969), 63--135. Theorem 1.1 and Corollary 1.3; Theorem 3.19 p.~89; Definition 5.3; Corollary 5.11 p.~131. \url{https://www.numdam.org/article/ASENS_1969_4_2_1_63_0.pdf}.
\bibitem{SGA3} M.~Demazure and A.~Grothendieck, \emph{Sch\'emas en groupes}, SGA 3, Expos\'e VII$_A$ (P.~Gabriel), \S\S1--3; 6.2--6.4; 7.2--7.4 and 8.1. Original-preserving transcription: \url{https://webusers.imj-prg.fr/~patrick.polo/SGA3/origExp7A.pdf}. In editions where a reference says ``7.4.1'', distinguish the theorem 7.4 from its proof subsection.
\bibitem{EGA} A.~Grothendieck and J.~Dieudonn\'e, \emph{\'El\'ements de g\'eom\'etrie alg\'ebrique IV, quatri\`eme partie}, Publ. Math. IHES \textbf{32} (1967), Theorem 16.11.2, printed p.~54, formulas (16.11.2.1)--(16.11.2.2). \url{https://www.numdam.org/item/PMIHES_1967__32__5_0.pdf}.
\bibitem{SGA4} M.~Artin, A.~Grothendieck, J.-L.~Verdier et al., \emph{Th\'eorie des topos et cohomologie \'etale des sch\'emas}, SGA 4: XII, Theorem 5.1 and Corollary 5.2; XV, Theorem 2.1; XVI, Theorem 2.1, original pp.~213--214. Primary-text re-editions: \url{https://www.normalesup.org/~forgogozo/SGA4/12/12.pdf}, \url{https://www.normalesup.org/~forgogozo/SGA4/15/15.pdf}, \url{https://www.normalesup.org/~forgogozo/SGA4/16/16.pdf}.
\bibitem{SGA5} A.~Grothendieck et al., \emph{Cohomologie $\ell$-adique et fonctions $L$}, SGA 5, Lecture Notes in Mathematics 589, Springer, 1977. Expos\'e VI, Definition 1.2.1, p.~254; \S1.4.2 and Definition 1.4.3, pp.~265--266, for lisse integral systems and constructible $\Q_\ell$ sheaves.
\bibitem{WeilI} P.~Deligne, \emph{La conjecture de Weil. I}, Publ. Math. IHES \textbf{43} (1974), 273--307. Theorem (1.6), p.~276; Frobenius conventions in (1.15), p.~279. \url{https://www.numdam.org/item/PMIHES_1974__43__273_0.pdf}.
\bibitem{WeilII} P.~Deligne, \emph{La conjecture de Weil. II}, Publ. Math. IHES \textbf{52} (1980), 137--252. Corollary (3.3.9), p.~207. \url{https://publications.ias.edu/sites/default/files/Number40.pdf}.
\bibitem{Berthelot} P.~Berthelot, \emph{Cohomologie cristalline des sch\'emas de caract\'eristique $p>0$}, Lecture Notes in Mathematics 407, Springer, 1974, V.3.6. Exact historical attribution checked through Illusie, \S3.2.
\bibitem{MM} B.~Mazur and W.~Messing, \emph{Universal Extensions and One Dimensional Crystalline Cohomology}, Lecture Notes in Mathematics 370, Springer, 1974. Chapter I, 2.6; Propositions 2.6.7, 3.2.3, 4.2.1; (4.6.3), for universal extensions and the line-bundle-with-connection realization. The combined statement and these locations are also recorded by G.~Oh and K.~Shimizu, \emph{Moduli stacks of crystals and isocrystals}, Theorem 9.6 and Remark 9.7, pp.~94--95, in the inspected preprint. \url{https://arxiv.org/pdf/2504.14801}.
\bibitem{Illusie} L.~Illusie, \emph{Grothendieck at Pisa: crystals and Barsotti--Tate groups}, author manuscript of the 23 April 2013 talk, \S2.1, \S3.2 (pp.~9--11), \S4.2(iv) (pp.~16--17). \url{https://www.imo.universite-paris-saclay.fr/~luc.illusie/Illusie-Pisa5.pdf}.
\bibitem{Stacks} The Stacks Project, \emph{Crystalline Cohomology}, Section 60.17, Proposition 60.17.4, tags 07J7 and 07JH. \url{https://stacks.math.columbia.edu/tag/07JH}.
\bibitem{Tate352} J.~Tate, \emph{Classes d'isog\'enie des vari\'et\'es ab\'eliennes sur un corps fini}, S\'eminaire Bourbaki, expos\'e 352 (November 1968), collected volume 1968--1969, pp.~95--110. \S5, definition p.~106, formula (10*) and Theorem 3 p.~107, proof through p.~108. \url{https://www.numdam.org/item/SB_1968-1969__11__95_0.pdf}.
\bibitem{Tate67} J.~Tate, \emph{$p$-Divisible Groups}, in \emph{Proceedings of a Conference on Local Fields}, Springer, 1967, pp.~158--183. \S1.2--1.3, pp.~159--160; \S2.2, Proposition 1, printed p.~162, diagram (3) and proof pp.~163--164; \S2.3, Proposition 3 and Example (b), p.~166; \S4, Corollary 2 following Theorem 3, p.~180, with its natural exact sequence. \url{https://www.math.purdue.edu/~tongliu/teaching/598/p-divisible.pdf}.
\bibitem{Hazewinkel} M.~Hazewinkel, \emph{On formal groups. The functional equation lemma and some of its applications}, Ast\'erisque \textbf{63} (1979), 73--82. Lemma 2.3, (2.2)--(2.8), pp.~74--75; \S5, (5.1)--(5.4), pp.~76--77. \url{https://www.numdam.org/item/AST_1979__63__73_0.pdf}.
\bibitem{Katz79} N.~M. Katz, \emph{Slope filtration of $F$-crystals}, Ast\'erisque \textbf{63} (1979), 113--163. \S1.3, pp.~119--120, especially the unnumbered semisimplicity theorem in the first paragraph of p.~120. \url{https://www.numdam.org/item/AST_1979__63__113_0.pdf}.
\bibitem{Zink} T.~Zink, \emph{The display of a formal $p$-divisible group}, Ast\'erisque \textbf{278} (2002), 127--248. Theorem 9, p.~132; Definition 19 and (22)--(24), pp.~143--144; Theorem 103, p.~221; Proposition 106, p.~226; Lemma 111, p.~230; \S4.3, (216)--(218), pp.~238--239; \S4.4, Proposition 121, pp.~241--246, injection (243), p.~246, and Theorem 122, p.~247. \url{https://www.numdam.org/article/AST_2002__278__127_0.pdf}. The author's version \url{https://www.math.uni-bielefeld.de/~zink/display.pdf} numbers the injection (244), p.~118; the published numbering is used here.
\bibitem{FGAPic} A.~Grothendieck, \emph{Technique de descente et th\'eor\`emes d'existence en g\'eom\'etrie alg\'ebrique}, FGA, V: \emph{Les sch\'emas de Picard: th\'eor\`emes d'existence}, S\'eminaire Bourbaki expos\'e 232, Theorem 3.1, pp.~148--149; VI: \emph{Les sch\'emas de Picard: propri\'et\'es g\'en\'erales}, expos\'e 236, Corollary 1.7, p.~229, Corollary 2.3, p.~232, Proposition 2.10(ii), p.~236. \url{https://www.numdam.org/article/SB_1961-1962__7__143_0.pdf}; \url{https://www.numdam.org/article/SB_1961-1962__7__221_0.pdf}.
\bibitem{Borger} J.~Borger, \emph{$\Lambda$-rings and the field with one element}, supplied \LaTeX\ source \texttt{lrfoe13.tex}, introduction and \S1.1 ``General $\Lambda$-spaces'', with \S2.1 for the point. No integral homotopy reconstruction theorem is attributed to this source.
\bibitem{Elliott} J.~Elliott, \emph{Binomial rings, integer-valued polynomials, and $\lambda$-rings}, J. Pure Appl. Algebra \textbf{207} (2006), 165--185. Lemma 2.2; Theorem 4.1; Proposition 8.3. \url{https://doi.org/10.1016/j.jpaa.2005.09.003}. The basis and Fermat-criterion locators are confirmed in Horel's proof of Proposition 2.2 and Remark 1.3, respectively; the original full proofs are cited rather than newly audited.
\bibitem{Horel} G.~Horel, \emph{Binomial rings and homotopy theory}, J. reine angew. Math. \textbf{813} (2024), 283--305. Definition 1.1, Remark 1.3; Proposition 2.2 and its proof, author-PDF pp.~4--5; Theorems 3.7, 5.4, 6.1; Proposition 8.1 and Notation 8.4. Author PDF: \url{https://geoffroy.horel.org/binomial.pdf}; Theorem 5.4 is on author-PDF p.~13. \url{https://doi.org/10.1515/crelle-2024-0039}.
\end{thebibliography}
\end{document}
